\documentclass[11pt,a4paper]{article}
\usepackage{fontspec}
\setmainfont{Latin Modern Roman}[SmallCapsFont={Latin Modern Roman Caps}]
\usepackage{polyglossia}
\setdefaultlanguage{french}
\newcommand{\og}{« }
\newcommand{\fg}{ »}
\usepackage{amsmath,amssymb,amsthm}
\IfFileExists{dsfont.sty}{\usepackage{dsfont}}{}
\usepackage{booktabs}
\usepackage{array}
\usepackage{longtable}
\usepackage[margin=2.3cm]{geometry}
\usepackage{enumitem}
\usepackage[most]{tcolorbox}
\usepackage{microtype}
\usepackage[colorlinks=true,linkcolor=blue!60!black,urlcolor=blue!60!black,citecolor=blue!60!black]{hyperref}


\usepackage{tikz}
\usetikzlibrary{positioning,arrows,shapes.geometric,calc,fit}
\tikzset{
  agent/.style={draw,rounded corners=3pt,thick,fill=blue!6,minimum width=2.6cm,minimum height=0.9cm,align=center,font=\small},
  actif/.style={draw,rounded corners=2pt,fill=orange!10,minimum width=2.2cm,minimum height=0.8cm,align=center,font=\small},
  pol/.style={draw,rounded corners=2pt,fill=green!10,minimum width=2.2cm,minimum height=0.8cm,align=center,font=\small},
  var/.style={draw,ellipse,fill=gray!10,align=center,font=\small,inner sep=2pt},
  crise/.style={draw,rounded corners=2pt,fill=red!8,minimum width=2.4cm,minimum height=0.8cm,align=center,font=\small},
  flux/.style={->,>=latex,thick},
  fluxm/.style={->,>=latex,thick,blue!60!black},
  fluxf/.style={->,>=latex,thick,orange!80!black},
  fluxp/.style={->,>=latex,thick,green!50!black,dashed},
  lab/.style={font=\scriptsize,fill=white,inner sep=1pt,align=center},
}


\newtcolorbox{lecture}[1][]{breakable,enhanced,colback=yellow!4,colframe=yellow!50!black,fonttitle=\bfseries\small,title={Lecture --- #1},fontupper=\small,left=4pt,right=4pt,top=3pt,bottom=3pt}
\newcommand{\variables}{\par\noindent\textbf{Variables.} }
\newcommand{\sens}{\par\noindent\textbf{Ce que dit l'équation.} }
\newcommand{\hyp}{\par\noindent\textbf{Hypothèses.} }
\newcommand{\limites}{\par\noindent\textbf{Limites.} }
\newcommand{\up}{$\uparrow$}
\newcommand{\dn}{$\downarrow$}

\setlength{\parskip}{4pt}
\setlist{nosep}
\newcommand{\E}{\mathbb{E}}
\newcommand{\R}{\mathbb{R}}
\newcommand{\pos}[1]{{\left[#1\right]^{+}}}
\newcommand{\clip}[3]{\operatorname{clip}\!\left(#1;\,#2,\,#3\right)}
\IfFileExists{dsfont.sty}{\newcommand{\ind}{\mathds{1}}}{\newcommand{\ind}{\mathbf{1}}}   % repli si dsfont est absent de l'environnement de compilation
\DeclareMathOperator{\sgn}{sgn}

\newtcolorbox{vic}[1][]{colback=orange!6,colframe=orange!60!black,fonttitle=\bfseries,title={Ce que l'on retient de Victoria~3 / E\&F --- #1}}
\newtcolorbox{design}[1][]{colback=blue!5,colframe=blue!50!black,fonttitle=\bfseries,title={Choix de conception --- #1}}
\newtcolorbox{lit}[1][]{colback=green!5,colframe=green!40!black,fonttitle=\bfseries,title={Ancrage théorique --- #1}}
\newtcolorbox{joueur}[1][]{colback=gray!8,colframe=black!60,fonttitle=\bfseries,title={Leviers du joueur --- #1}}

\title{\textbf{Nations \& Marchés}\\[4pt]
\large Simulation macroéconomique jouable : mécaniques et équations\\[2pt]
\normalsize Document de conception --- livrable 1 (version 1.5)}
\author{}
\date{Septembre 2026}

\begin{document}
\maketitle

\begin{abstract}
\noindent Ce document décrit la mécanique et le système d'équations d'un jeu de simulation économique multi-pays. Chaque joueur incarne un État : il contrôle la politique monétaire (taux directeur, création monétaire, régime de change) et les finances publiques (impôts, dépenses, dette), et peut faire glisser son pays sur un continuum institutionnel allant de l'économie de marché libérale à l'économie planifiée. Le secteur privé (ménages, entreprises, banques) est autonome et réagit aux signaux de prix et aux décisions publiques. Le modèle emprunte au jeu \emph{Victoria~3} et au mod \emph{Economic \& Financial Mod (E\&F) V4} leur logique de marché par ordres d'achat/vente, de biens intermédiaires, de besoins hiérarchisés des populations, de banques centrales émettrices et de régimes monétaires ; il s'appuie sur la littérature macroéconomique (Ramsey--Cass--Koopmans, IS--LM et Mundell--Fleming, courbe de Phillips néo-keynésienne, règle de Taylor, dynamique de la dette, modèles stock-flux cohérents, économie de la pénurie de Kornai) pour en fixer les équations. Le document spécifie également l'hétérogénéité des pays (terre, ressources, population, structure duale), un module d'actifs (obligations, actions, immobilier, devises) et vérifie que les grandes crises historiques (1929, Japon 1990, 2008, crises de change, hyperinflation) ainsi que les deux issues possibles d'un assouplissement quantitatif émergent des équations. La v0.6 ajoute le régime de souveraineté monétaire comme dimension de l'hétérogénéité des pays, un taux apparent à mémoire, un bilan vivant de banque centrale et sept configurations contemporaines à reproduire. Aucun code n'est fourni : l'objet est la spécification.
\end{abstract}

\tableofcontents
\newpage
%=====================================================================
\section*{Ce qui change en v1.5 : relevé d'incohérences du cours}
\addcontentsline{toc}{section}{Ce qui change en v1.5 : relevé d'incohérences du cours}
%=====================================================================
La préparation du cours magistral (table de correspondance des 65 équations, chapitre témoin) a produit un relevé de treize incohérences du \emph{document} — des tableaux qui ne reproduisaient plus les équations, des symboles à plusieurs sens, des termes cités sans définition. La règle d'arbitrage adoptée est celle que le cours applique : les équations priment sur les tableaux, les blocs de révision priment sur la table de calibration initiale. Chaque point a été vérifié sur le prototype avant correction ; deux ont conduit à modifier le moteur (répression financière et plancher conditionnel), un a été renvoyé en chantier après que le prototype a montré que la correction envisagée n'aurait aucun effet (accélérateur financier). Trente-sept tests passent (un ajouté).

\begin{center}\small
\begin{longtable}{p{0.5cm}p{6.6cm}p{2.1cm}p{5.8cm}}
\toprule
\textbf{\#} & \textbf{Modification} & \textbf{Section} & \textbf{Ce que la revue ou le prototype a montré}\\
\midrule
\endfirsthead
\toprule
\textbf{\#} & \textbf{Modification} & \textbf{Section} & \textbf{Ce que la revue ou le prototype a montré}\\
\midrule
\endhead
1 & \textbf{Grille des dix configurations recalculée} avec \eqref{eq:spread} telle qu'écrite (dette de marché, $\varkappa^{for}$, $\psi_{QE}$ sur la prime de terme) & \S\ref{sec:configs} & La colonne « valeurs calculées » utilisait encore la forme antérieure (convexité sur $b$ total pondérée par $1-\varkappa^{dom}$, rabais $s_7\omega$). Cinq configurations changent : 109 $\to$ 169, 13 $\to$ 76, 168 $\to$ 166, 99 $\to$ 145, 79 $\to$ 88 pb. L'écart « union vertueuse » du \S\ref{sec:chantiers} est résolu\\
2 & \textbf{Plancher conditionnel} $\underline\varkappa(1-\chi^K\varrho^B)$ dans \eqref{eq:spread} ; répression $\varrho^B$ et $s_9$ ajoutés au prototype & \S\ref{sec:spread} & Le plancher 0,3 faisait payer 218 pb à une dette détenue à 95\,\% par des créanciers captifs ($\approx30$ observés). Test H6 : le contrôle des capitaux seul est sans effet, la répression déplace le coût de la dette du taux (5,4 contre 8,8\,\%) vers l'inflation (7,7 contre 5,2\,\%). T9 insensible au plancher : la convexité ne porte pas la dynamique\\
3 & \textbf{Table de calibration nettoyée} : six lignes périmées supprimées ($\iota$, $\delta_j$, $\rho$, $u^n$, $\theta_{1\text{--}3}$, $s_7$), liste des paramètres retirés complétée, $u^n_0$ et $\underline u^n$ séparés & \S\ref{sec:calib}, glossaire & Un lecteur prenant la table dans l'ordre calculait un investissement faux d'un facteur deux\\
4 & \textbf{(35) précisée} : $\hat r^\ast_0=\rho+\sigma g_0$ & \S\ref{sec:bc} & $g_0$ donne les 2,5\,\% affichés ; $g$ réalisé aurait donné 2,9\\
5 & \textbf{$\varpi$ dédoublé} : $\varpi$ coefficient de vitesse d'indexation des prix (0,5 par tick sur $\pi^e/52$), $\varpi_w$ degré d'indexation des salaires (borné par 1) & \eqref{eq:price}, \eqref{eq:wage}, \S\ref{sec:qe} & Trois échelles pour un symbole : 0,02/mois, par tick, « $\to1$ »\\
6 & \textbf{$q^\ast=1+g/(\iota\Gamma\Lambda\Lambda^{Bk})$} ; indicateur joueur $q/q^\ast$ & \S\ref{sec:invest}, \S\ref{sec:configs} & $q=1$ n'est pas le seuil de rentabilité d'une économie qui croît ; $q^\ast=1{,}11$ est ce que le prototype affiche\\
7 & \textbf{Écart de QE assumé} ($-8$ pb pour 20\,\% du PIB) & \S\ref{sec:qe} & Conforme aux équations ; l'écart à la littérature est un effet d'anticipation de trajectoire hors modèle\\
8 & \textbf{Accélérateur financier} : constat d'inactivité, correction par seuil écartée, forme continue renvoyée en chantier & \S\ref{sec:invest}, \S\ref{sec:chantiers} & Levier 0,32 en économie fermée, 0,08 avec actifs ; balayage de $\bar\ell$ de 0,6 à 0,35 strictement sans effet sur A1, A2, A4, T1\\
9 & \textbf{Termes non spécifiés} : $\eta_{agg}$ retiré ; $\eta_H$ (0,01/an) et $q^{file}$ ($\zeta_q=0{,}1$) définis, marqués « non simulés » & \S\ref{sec:hetero}, \S\ref{sec:etat}, \eqref{eq:sol} & Aucun des trois n'était dans le prototype ; règle : retrait sans fondement documenté, spécification sans implémentation quand le gain n'est pas mesurable avant l'étape~5\\
10 & \textbf{(41) et (43)} encadrées « choix de conception, pas résultat théorique » & \S\ref{sec:reforme}, \S\ref{sec:collapse} & Aucune référence ne les adosse ; le dire vaut mieux que fabriquer un fondement\\
11 & \textbf{Deux états de référence, rôles déclarés} : actifs $=$ jeu, fermée $=$ calibration ; $I/Y$ recalculé à 0,17, $K/Y$ 2,85, levier 0,32 & \S\ref{sec:ssref} & $I/Y=(\delta+g)K/Y$, non 0,18 ; l'écart d'un quart sur $K/Y$ reste le chantier du partage de la valeur ajoutée\\
12 & \textbf{Colonne « Prix » de l'économie planifiée dédoublée} selon $\varrho^p$ & \S\ref{sec:systemes} & La même contrainte budgétaire lâche produit la pénurie à prix administrés et la déflation exportée à prix libres\\
\bottomrule
\end{longtable}
\end{center}

%=====================================================================
\section*{Ce qui change en v1.4 : revue des chantiers ouverts}
\addcontentsline{toc}{section}{Ce qui change en v1.4 : revue des chantiers ouverts}
%=====================================================================
Les vingt-cinq chantiers ouverts par les versions 0.9 à 1.3 ont été repris un à un : chacun a été vérifié sur le prototype, résolu quand c'était possible, et son état est consigné dans la table de la \S\ref{sec:chantiers}, qui remplace les listes empilées des versions précédentes. Dix-sept sont clos (dont six par des travaux antérieurs qui n'avaient pas été reportés), trois partiellement, cinq restent ouverts. Trente-six tests passent. Les modifications du modèle sont les suivantes.

\begin{center}\small
\begin{longtable}{p{0.5cm}p{6.6cm}p{2.1cm}p{5.8cm}}
\toprule
\textbf{\#} & \textbf{Modification} & \textbf{Section} & \textbf{Ce que la revue a montré}\\
\midrule
\endfirsthead
\toprule
\textbf{\#} & \textbf{Modification} & \textbf{Section} & \textbf{Ce que la revue a montré}\\
\midrule
\endhead
1 & \textbf{Ratio de liquidité bancaire} (option $\phi_{liq}$) : la banque détient des titres publics au moins égaux à une part des dépôts & \S\ref{sec:banques}, \S\ref{sec:qe} & Avec $\phi_{liq}=0{,}3$, le QE s'achète à la banque : monnaie inchangée ($-0{,}4\,\%$), souverain $-8$ pb, actions $+0{,}8\,\%$, inflation $+0{,}1$ pt — le cas « réserves oisives » du \S\ref{sec:qe} est reproduit ; sans, achat aux ménages : monnaie $+15\,\%$, inflation $+0{,}9$ pt\\
2 & \textbf{Fécondité et mortalité endogènes} (option) : fécondité $\propto SoL^{-\eta_f}$, mortalité des âgés $\propto SoL^{-\eta_m}$ & \S\ref{sec:hetero} & Transition démographique : part des âgés 33\,\% contre 26\,\% à soixante ans ; natalité de référence portée à 1,4\,\% des actifs pour une population stable\\
3 & \textbf{Bornes de jeu} sur $\chi_Z$, $\chi_E$, LTV, marge, $\kappa^{CAR}$, $\pi^\ast$ & \S\ref{sec:calib} & Au-delà, cycles d'actifs de grande amplitude (configuration « 1929 »)\\
4 & \textbf{Actions} : poids du rappel fondamental porté à 2 & \S\ref{sec:actifs} & Cours moyen 0,77 au lieu de 0,70 ; le reste de l'écart tient au rapport dividendes / valeur comptable\\
\bottomrule
\end{longtable}
\end{center}

%=====================================================================
\section*{Ce qui change en v1.3 : étape 4 et effondrement}
\addcontentsline{toc}{section}{Ce qui change en v1.3 : étape 4 et effondrement}
%=====================================================================
L'étape~4 de la feuille de route (terre et climat, ressource épuisable, cohortes, migration rurale, archétype agraire) a été programmée et soumise à cinq tests, et les cinq chantiers ouverts par l'audit de la v1.2 ont été traités (\S\ref{sec:proto4}). L'état de mort économique, resté ouvert en v1.2, est résolu par une \emph{machine à états} : l'économie bascule dans un régime de survie stylisé et en sort par une réforme monétaire réussie. Trente-six tests passent.

\begin{center}\small
\begin{longtable}{p{0.5cm}p{6.6cm}p{2.1cm}p{5.8cm}}
\toprule
\textbf{\#} & \textbf{Modification} & \textbf{Section} & \textbf{Ce que le prototype a montré}\\
\midrule
\endfirsthead
\toprule
\textbf{\#} & \textbf{Modification} & \textbf{Section} & \textbf{Ce que le prototype a montré}\\
\midrule
\endhead
1 & \textbf{Machine à états « économie effondrée »} \eqref{eq:collapse} : entrée (anticipations au plafond six mois, ou deux réformes en cinq ans), régime de survie (production à 60\,\% du potentiel, unité de compte stable, ni crédit ni investissement net), sortie par réforme avec conversion monétaire & \S\ref{sec:collapse} & La mort économique était un régime dégénéré du système de prix (cibles de production astronomiques, ventes nulles) ; aucune correction de mécanisme ne le résolvait, et l'une d'elles rendait la trappe à taux zéro explosive\\
2 & \textbf{Règle budgétaire} \eqref{eq:fiscrule} : achats, investissement public et minima corrigés de l'écart à une cible de dette & \S\ref{sec:etat} & Sans elle, la dette de l'état de référence fond vers zéro et prive la banque centrale de titres ; un bug d'unité (dette rapportée à 52 PIB annuels) la rendait explosive\\
3 & \textbf{Ressource épuisable} : plateau tant que la moitié du stock subsiste, puis coût croissant, plancher de substituts (40\,\%) ; stock de référence 150 ans & \S\ref{sec:hetero} & Trente ans de réserves et un coût croissant dès la première tonne tuaient l'économie (énergie intermédiaire Leontief) ; le plancher des substituts fait la différence entre transition ($-18\,\%$) et effondrement\\
4 & \textbf{Terre, climat, cohortes, migration rurale} : formes opérationnelles et calibration & \S\ref{sec:hetero}, \S\ref{sec:proto4} & Lewis émerge (réservoir rural épuisé en quarante ans, PIB $+37\,\%$) ; le vieillissement triple la dette ; une mauvaise récolte transmet le climat à l'inflation par la strate basse\\
5 & \textbf{Plancher d'actualisation} des valeurs fondamentales (3\,\%) & \S\ref{sec:actifs} & La valeur fondamentale des actions divergeait à taux zéro et transformait la trappe en hyperinflation\\
6 & \textbf{Liquidation unique par épisode de faillite} & \S\ref{sec:prod} & En hyperinflation, une liquidation de 5\,\% répétée chaque semaine de réorganisation annihilait le capital\\
\bottomrule
\end{longtable}
\end{center}
Décompte : deux équations nouvelles (\eqref{eq:collapse}, \eqref{eq:fiscrule}), la ressource réécrite, quatre blocs d'hétérogénéité rendus opérationnels, une sous-section de résultats. Un chantier reste ouvert : l'économie sortie de l'effondrement dérive vers une inflation chronique (23\,\% à vingt ans) faute d'ancre, ce qui appelle une seconde stabilisation ou un ancrage externe (régime E).

%=====================================================================
\section*{Ce qui change en v1.2 : audit du prototype}
\addcontentsline{toc}{section}{Ce qui change en v1.2 : audit du prototype}
%=====================================================================
Avant l'étape~4, le prototype a été audité : état stationnaire sur soixante ans, quarante valeurs de leviers modifiées \emph{en cours de partie}, huit politiques extrêmes maintenues vingt ans (\S\ref{sec:audit}). L'audit a trouvé cinq défauts, dont deux rendaient l'état stationnaire avec actifs non viable au-delà de quarante-cinq ans, et une faille méthodologique. Aucune méta-stratégie ne subsiste après correction. Les équations touchées sont celles ci-dessous ; le reste du document est inchangé.

\begin{center}\small
\begin{longtable}{p{0.5cm}p{6.6cm}p{2.1cm}p{5.8cm}}
\toprule
\textbf{\#} & \textbf{Modification} & \textbf{Section} & \textbf{Ce que l'audit avait révélé}\\
\midrule
\endfirsthead
\toprule
\textbf{\#} & \textbf{Modification} & \textbf{Section} & \textbf{Ce que l'audit avait révélé}\\
\midrule
\endhead
1 & \textbf{Seigneuriage} \eqref{eq:ecb} : versé sur le \emph{profit de la période}, et seulement si les fonds propres statutaires sont reconstitués (la règle était déjà écrite ainsi ; le prototype distribuait 10\,\%/semaine du \emph{stock} de fonds propres) & \S\ref{sec:bc} & Avec actifs, la banque centrale démarrait avec 0,7 PIB de fonds propres apparents (créance de refinancement sur la banque) et les versait à l'État : $+50\,\%$ de dépôts la première année, dérive inflationniste, effondrement à l'an 45\\
2 & \textbf{Réserves} \eqref{eq:ecb} : seules les réserves adossées à des actifs porteurs d'intérêt sont rémunérées ; opérations structurelles (le portefeuille $B^{CB}$ suit les réserves) & \S\ref{sec:bc} & Rémunérer des réserves sans contrepartie est une perte structurelle (2,5\,\% du PIB par an), c'est-à-dire une monétisation cachée\\
3 & \textbf{Prix d'actifs} : anticipations de plus-value \emph{réelles} et valeur fondamentale au coût d'usage réel net de la croissance & \S\ref{sec:actifs} & Double comptage de l'inflation ; valeur fondamentale négative dès $\pi^\ast=8\,\%$\\
4 & \textbf{Placement de la dette} : la banque absorbe jusqu'à 40\,\% de son bilan ; au-delà, demande d'obligations des ménages sensible au rendement ($\kappa^{hB}=20$ par point d'écart souverain-dépôts) ; ensuite rationnement des dépenses & \S\ref{sec:etat} & Toute expansion budgétaire de plus de cinq points de PIB finissait en hyperinflation par financement bancaire illimité ; sans demande élastique, un plafond seul provoquait des dépressions par rationnement\\
5 & \textbf{Chômage naturel} \eqref{eq:un} : plancher frictionnel $u^n\ge3\,\%$ & \S\ref{sec:prod} & Une surchauffe permanente ramenait $u^n$ à 1,8\,\% et effaçait l'accélération de l'inflation\\
6 & \textbf{Niveau de vie} \eqref{eq:sol} : inclut la consommation publique effective par tête (élasticité 0,15) & \S\ref{sec:menages} & « Supprimer les impôts et laisser rationner les dépenses » gagnait 11\,\% de niveau de vie ; l'indice ignorait les services publics\\
7 & Section~\ref{sec:audit} : méthode, résultats du balayage, effets excessifs restants, faille méthodologique & \S\ref{sec:audit} & Un balayage effectué \emph{avant} la calibration comparait des économies initiales différentes ($+62\,\%$ de niveau de vie apparents)\\
\bottomrule
\end{longtable}
\end{center}
Décompte : trois équations modifiées (\eqref{eq:ecb}, \eqref{eq:un}, \eqref{eq:sol}), une règle de placement de la dette, une section d'audit. Les 26 tests des étapes 1 à 3 passent après corrections.

%=====================================================================
\section*{Ce qui change en v1.1 : actifs, capital bancaire, corrections}
\addcontentsline{toc}{section}{Ce qui change en v1.1 : actifs, capital bancaire, corrections}
%=====================================================================
L'étape~3 de la feuille de route (immobilier, actions, crédit hypothécaire et sur marge, effet richesse, ruées, assouplissement quantitatif à coupon figé) a été programmée et soumise à cinq tests, tous passés (\S\ref{sec:proto3}). Elle apporte le résultat annoncé depuis la v0.9 : \textbf{la prime de risque recyclée en richesse ferme l'écart entre taux réel effectif et taux naturel théorique}. Elle a aussi révélé une fragilité (la stabilisation post-hyperinflation dépend de la capacité bancaire résiduelle) dont le remède — l'émission d'actions bancaires — était le mécanisme manquant du bloc bancaire. Trois ajustements de modélisation ont par ailleurs été testés et retenus.

\begin{center}\small
\begin{longtable}{p{0.5cm}p{6.6cm}p{2.1cm}p{5.8cm}}
\toprule
\textbf{\#} & \textbf{Modification} & \textbf{Section} & \textbf{Ce que le prototype a montré}\\
\midrule
\endfirsthead
\toprule
\textbf{\#} & \textbf{Modification} & \textbf{Section} & \textbf{Ce que le prototype a montré}\\
\midrule
\endhead
1 & \textbf{Émission d'actions bancaires} \eqref{eq:bankeq} : quand $E^{Bk}<1{,}1\,\kappa^{CAR}L$, la banque lève des fonds propres auprès des ménages (jusqu'à 2\,\% de leurs dépôts par semaine) ; recapitalisation publique dès la sous-capitalisation et non seulement l'insolvabilité & \S\ref{sec:banques} & Sans elle, le capital bancaire ne croît que par rétention de profits et la contrainte de Bâle devient une falaise : une économie sortant d'une hyperinflation avec 3\,\% de fonds propres rechutait (chômage 34\,\%) parce que la reprise du crédit butait sur le rationnement\\
2 & \textbf{Dividendes bancaires selon la couverture} \eqref{eq:bankeq} : taux de distribution du capital excédentaire multiplié par $\clip{CAR/(1{,}3\kappa^{CAR})}{1}{3}$ & \S\ref{sec:banques} & Une demande de crédit faible laisse le capital s'accumuler ; sans ce terme, la banque thésaurise un capital sans emploi\\
3 & \textbf{Prime souveraine hors banque centrale} \eqref{eq:spread} : la convexité porte sur $(B-B^{CB})/Y$, pondérée par la part des non-résidents dans la dette \emph{hors banque centrale} (plancher 0,3) ; le rabais $-s_7B^{CB}/B$ disparaît et le canal de portefeuille du QE agit sur la seule prime de terme, $s_0(1-\psi_{QE}B^{CB}/B)$ & \S\ref{sec:spread} & Le double comptage était réel : la dette détenue par la banque centrale était à la fois comptée dans le risque de refinancement et rabattue par $s_7$. Effet mesuré : $-8$ pb sur le souverain lors d'un QE, dette et taux inchangés en régime normal\\
4 & \textbf{Prix mondial des biens} (option) : pour l'énergie, prix mondial = moyenne des prix nationaux en numéraire pondérée par la production, vers lequel chaque prix national est rappelé ($\lambda_{lop}=0{,}05$/semaine) & \S\ref{sec:marches} & Un prix mondial \emph{imposé} (fixé par l'excès de demande mondial) diverge numériquement ; le rappel par arbitrage est stable et rend le choc pétrolier plus net (inflation $1{,}6\to2{,}6\,\%$ chez l'importateur)\\
5 & \textbf{Actifs} (étape 3) : immobilier, actions, crédit hypothécaire et sur marge, effet richesse, ruées, QE à coupon figé ; bornes de crédit et de construction & \S\ref{sec:actifs}, \S\ref{sec:proto3} & Le taux réel d'équilibre passe de 1,1 à 3,0\,\% et rejoint $\rho+\sigma g$ ; sept versions du bloc immobilier ont divergé avant d'être stables\\
\bottomrule
\end{longtable}
\end{center}
Décompte : deux équations modifiées (\eqref{eq:spread}, compte de résultat de la banque devenu \eqref{eq:bankeq}), une option de prix mondial, une sous-section de résultats. Les équations d'actifs de la \S\ref{sec:actifs} sont conservées telles quelles ; la \S\ref{sec:proto3} donne leur forme opérationnelle et leur calibration.

%=====================================================================
\section*{Ce qui change en v1.0 : économie ouverte vérifiée}
\addcontentsline{toc}{section}{Ce qui change en v1.0 : économie ouverte vérifiée}
%=====================================================================
La v1.0 clôt les deux premières étapes de la feuille de route. L'étape~1 (économie fermée) a reçu la séquence de correction de la transmission monétaire annoncée en v0.9 ; l'étape~2 (monde à $N$ pays, cinq régimes de souveraineté, contrôle des capitaux, dollarisation, attaques de première et seconde génération, énergie échangeable) a été programmée et soumise à treize tests, tous passés. Les équations de l'économie ouverte sont réécrites là où le prototype les a contredites ; le reste du document est inchangé.

\begin{center}\small
\begin{longtable}{p{0.5cm}p{6.6cm}p{2.1cm}p{5.8cm}}
\toprule
\textbf{\#} & \textbf{Modification} & \textbf{Section} & \textbf{Ce que le prototype avait révélé}\\
\midrule
\endfirsthead
\toprule
\textbf{\#} & \textbf{Modification} & \textbf{Section} & \textbf{Ce que le prototype avait révélé}\\
\midrule
\endhead
1 & \textbf{Investissement} \eqref{eq:invest} : réponse à $q$ saturante autour du $q$ de long terme ; rétention de main-d'œuvre dans les biens d'équipement ; dividendes sur profit lissé & \S\ref{sec:prod} & Une hausse de 1 pt maintenue un an coûtait $-5{,}9\,\%$ de PIB ; $-2{,}9\,\%$ après correction, sans cycle. Les retards de transmission (taux long, ajustement partiel) et l'accélérateur ont été essayés et rejetés : cycles de 3 à 12 ans\\
2 & \textbf{Taux de change} \eqref{eq:fx} : le change est un prix d'actif rappelé vers une cible (PPA corrigée du différentiel de taux réel), plus la pression des flux nets & \S\ref{sec:change} & La parité des taux écrite en flux ne produisait pas la surréaction de Dornbusch (appréciation de 2\,\% en 19 mois) ; la forme en niveau la produit (4,4\,\% en 12 mois, retour à 1,5\,\%)\\
3 & \textbf{Change fixe} : taux directeur importé ; attaque quand les réserves atteignent le seuil \emph{et} que les fondamentaux sont incohérents ; \textbf{clause d'abandon} (seconde génération) en option & \S\ref{sec:change} & Avec sa règle de Taylor, le pays ancré défendait la parité par les taux (attaque à l'an 17 au lieu de 12) ; une probabilité de dévaluation fondée sur les seules réserves déclenchait des attaques auto-réalisatrices dès l'an 3 sans faute de politique\\
4 & \textbf{Armington} \eqref{eq:armington} : parts ajustées graduellement (contrats) & \S\ref{sec:marches} & À élasticité 6 et parts instantanées, les parts oscillaient de 0,45 à 0,83 en trois semaines et le change explosait\\
5 & \textbf{Règlement international} : sur montants payés, comptes de change des banques centrales, position extérieure du numéraire, intérêts aux détenteurs étrangers & \S\ref{sec:sfc}, \S\ref{sec:bop} & Le règlement sur quantités livrées faisait dériver les réserves mondiales de 5\,\% en huit ans ; deux signes de flux de capitaux inversés, invisibles aux identités locales, révélés par la conservation mondiale\\
6 & \textbf{Dollarisation} \eqref{eq:sfx} : part désirée fonction de l'inflation réalisée, hystérèse 0,9995 par semaine & \S\ref{sec:souv} & Une part dollarisée de 18\,\% après réforme n'est retombée qu'à 13\,\% en treize ans d'inflation à 2\,\%\\
7 & Section~\ref{sec:proto} : résultats de l'étape~2 (treize tests), enseignements, chantiers & \S\ref{sec:proto} & ---\\
\bottomrule
\end{longtable}
\end{center}
Décompte : quatre équations modifiées (\eqref{eq:invest}, \eqref{eq:fx}, \eqref{eq:armington}, \eqref{eq:sfx}), une règle d'abandon ajoutée en option, aucune équation nouvelle.

%=====================================================================
\section*{Ce qui change en v0.9 : retour du prototype}
\addcontentsline{toc}{section}{Ce qui change en v0.9 : retour du prototype}
%=====================================================================
La v0.9 est la première version écrite \emph{après} un prototype exécutable. L'étape~1 de la feuille de route (économie fermée : quatre biens, trois strates, une banque, une banque centrale, un État, pas hebdomadaire) a été programmée et soumise à quinze tests (\S\ref{sec:proto}). Le prototype a confirmé l'architecture (cohérence stock-flux exacte, émergence de l'inflation, de l'hyperinflation, des cycles du crédit, des crises de demande) et a révélé une douzaine de défauts que le raisonnement à la main n'avait pas vus. Chaque correction ci-dessous est une modification d'une équation existante ou l'ajout d'une règle de bord ; aucune n'ajoute de sous-système. La section~\ref{sec:proto} donne l'état stationnaire obtenu et les propriétés identifiées.

\begin{center}\small
\begin{longtable}{p{0.5cm}p{6.6cm}p{2.1cm}p{5.8cm}}
\toprule
\textbf{\#} & \textbf{Modification} & \textbf{Section} & \textbf{Ce que le prototype avait révélé}\\
\midrule
\endfirsthead
\toprule
\textbf{\#} & \textbf{Modification} & \textbf{Section} & \textbf{Ce que le prototype avait révélé}\\
\midrule
\endhead
1 & \textbf{Prix} \eqref{eq:price} : rappel \emph{symétrique} vers le coût unitaire majoré d'une marge normale ; indexation sur $\pi^e$ pondérée par $\pos{1+z}$ ; flexibilité croissante avec $\pi^e$ ; borne hebdomadaire & \S\ref{sec:marches} & Sans rappel par le haut, la part salariale dérivait de 0,65 à 0,39 en trente ans (régime d'accumulation sans salaires) ; sans indexation dépendante de l'état, l'inflation continuait après l'arrêt de la monnaie\\
2 & \textbf{Salaires} \eqref{eq:wage} : indexation sur la productivité \emph{observée} lissée ($\gamma_A=1$) ; terme de surenchère en pénurie de main-d'œuvre (Phillips convexe) ; indexation sur l'inflation réalisée pondérée par $(1-c)$, affaiblie par le chômage de masse ; baisse forcée si le salaire rend l'emploi non rentable ; plafond payable & \S\ref{sec:prod} & Pas d'inflation de surchauffe à 1,8\,\% de chômage ; hyperinflation impossible sans indexation ; salaires nominaux croissant sans monnaie après une réforme\\
3 & \textbf{Investissement} \eqref{eq:invest} : prime de risque $\rho_E$ dans le coût du capital ; investissement net conditionné au taux d'utilisation ; plafond $2\delta$ ; marge normale calibrée sur $r^\ast+\rho_E$ & \S\ref{sec:prod} & Sans prime de risque, $K/Y$ tombe à 2,3 et $I/Y$ à 13\,\% ; sans plafond, cycles\\
4 & \textbf{Dividendes} \eqref{eq:div} prioritaires ($\psi$ du profit net), l'investissement se finance ensuite sur trésorerie puis crédit & \S\ref{sec:prod} & Les entreprises thésaurisaient (dividendes 0,7\,\% du PIB) et autofinançaient un investissement de 30\,\% du PIB\\
5 & \textbf{Règle de Taylor} \eqref{eq:taylor} : taux naturel \emph{estimé} par apprentissage ; inflation récente mesurée en log borné et lissée ; anticipations bornées par l'inflation observée ; lissage $\rho_i$ ; plafond & \S\ref{sec:bc} & Ancrée sur $\rho+\sigma g$, la règle maintenait 7\,\% de chômage et 1\,\% d'inflation ; la mesure annualisée du mois amplifiait le bruit (±20\,\% $\to$ 790\,\%/an) et bloquait le taux à 100\,\% après une réforme\\
6 & \textbf{Levier de réforme monétaire} : redénomination automatique à $P>10^6$, arrêt des avances, remise à zéro de $\pi^e$, crédibilité $=0{,}8\times$ plafond, taux ramené au neutre & \S\ref{sec:inflation} & Sans ancre nouvelle, aucune stabilisation ; débordement numérique en année 12,7\\
7 & \textbf{Contraintes de caisse} : masse salariale et achats d'intrants bornés par la trésorerie (relâchés par la vitesse de circulation) ; grand livre sans découvert & \S\ref{sec:prod}, \S\ref{sec:sfc} & Les découverts absorbés par la banque étaient une création monétaire illimitée\\
8 & \textbf{État} : taux d'exécution des dépenses = trésorerie disponible / dépenses prévues ; intérêts impayés capitalisés ; financement direct sur l'épargne des ménages si la banque est insolvable & \S\ref{sec:etat} & Transferts indexés financés sans limite en hyperinflation\\
9 & \textbf{Banque} : recapitalisation publique automatique (au plus annuelle) ; insolvable, elle ne prête ni ne rémunère & \S\ref{sec:banques} & Une banque insolvable non recapitalisée bloquait le crédit à vie\\
10 & \textbf{Faillite} : liquidation de 5\,\% du capital (et non 25), réorganisation d'un trimestre & \S\ref{sec:prod} & Des faillites répétées annihilaient le capital en un an\\
11 & \textbf{Bornes numériques} : taux directeur $\le100\,\%$, prime $\le30$ pts, prix $\le+19\,\%$/semaine, salaires $\le\times2$/mois, $\pi^e\le1\,200\,\%$, vitesse saturée à $\pi^e=100\,\%$ & \S\ref{sec:calib} & Divergences numériques en hyperinflation\\
12 & \textbf{Capital public} accumulé et actif dans la PGF (neutre au ratio initial) ; \textbf{répartition des salaires} 55/30/15 entre strates ; \textbf{calibration} révisée (\S\ref{sec:calib}) & \S\ref{sec:etat} & Investissement public sans effet ; revenu par tête plus élevé dans la strate basse que dans la haute\\
\bottomrule
\end{longtable}
\end{center}
Décompte : sept équations modifiées (\eqref{eq:price}, \eqref{eq:wage}, \eqref{eq:invest}, \eqref{eq:div}, \eqref{eq:taylor}, \eqref{eq:tfp} par $K^G$, \eqref{eq:debtdyn} par la contrainte de trésorerie), une nouvelle (réforme monétaire), douze règles de bord, et une section de résultats.

%=====================================================================
\section*{Ce qui change en v0.8 : recentrage}
\addcontentsline{toc}{section}{Ce qui change en v0.8 : recentrage}
%=====================================================================
\paragraph{Historique.} La v0.5 a posé le modèle complet (production, ménages, marchés, monnaie, État, systèmes économiques, crises) avec ses encadrés de lecture et son analyse d'impact. La v0.6 a ajouté le régime de souveraineté monétaire, le taux apparent à mémoire, le bilan vivant de banque centrale, le saut de prime et la défense d'un plafond de change. La v0.7 a corrigé trois calibrations, ajouté le régime E et la grille de dix configurations chiffrées.

\paragraph{Principe de la v0.8.} \emph{Un jeu, pas un modèle de banque centrale.} Tout terme, paramètre ou curseur qui ne change pas le comportement d'au moins une des dix configurations de la \S\ref{sec:configs} est retiré. Les écarts entre valeurs calculées et valeurs observées sont acceptés tant qu'ils ne produisent pas une trajectoire \emph{qualitativement} différente (un pays sûr qui saute, une dette stable qui explose, une inflation qui n'apparaît pas). Aucun ajout.

\begin{center}\small
\begin{longtable}{p{0.5cm}p{6.4cm}p{2.2cm}p{5.8cm}}
\toprule
\textbf{\#} & \textbf{Retrait ou fusion} & \textbf{Section} & \textbf{Pourquoi c'est sans perte}\\
\midrule
\endfirsthead
\toprule
\textbf{\#} & \textbf{Retrait ou fusion} & \textbf{Section} & \textbf{Pourquoi c'est sans perte}\\
\midrule
\endhead
1 & Prime souveraine \eqref{eq:spread} : suppression de la prime d'illiquidité $s_5\chi^K$ ; prime de terme $s_T(\bar T-\bar T_0)$ fusionnée dans $s_0\,\bar T/\bar T_0$ ; $s_{10}$ fixé à 1 & \S\ref{sec:spread} & $s_5$ ne jouait que pour le pays planifié, où la répression $-s_9$ l'annulait ; le coût de la maturité reste (proportionnel à $\bar T$) avec un paramètre de moins. Douze termes deviennent neuf\\
2 & Crédibilité \eqref{eq:cred} : suppression de la dominance budgétaire $\nu_5\Psi$, de l'undershoot $\nu_6$ et de la surprise $\lambda_2$ & \S\ref{sec:inflation} & $\Psi$ aurait fait perdre sa crédibilité à la configuration 4 (banque centrale à 50\,\% de la dette), à rebours de l'observation ; $\nu_2|\pi-\pi^\ast|$ est déjà symétrique ; la surprise n'a d'effet dans aucune configuration\\
3 & Bilan de banque centrale \eqref{eq:ecb} : suppression du fonds pour risques généraux $F^{CB},\vartheta$, du \emph{tiering} $\varrho^{tier}$, de la réévaluation de l'or et de la duration $D^{CB}$ (coupon figé conservé) & \S\ref{sec:bc} & Report à nouveau des pertes et suspension du dividende donnent le même amortisseur budgétaire avec zéro paramètre ; le tiering est un débat de banquier central, pas un choix de joueur\\
4 & Suppression de la fraction utilisable des réserves $\mu^{FX}$ et de la décote pays $\delta_i$ & \S\ref{sec:change}, \S\ref{sec:hetero} & Introduits pour un producteur sous sanctions, qui n'est dans aucune des dix configurations ; renvoyés aux pistes écartées\\
5 & Curseurs institutionnels : retrait de $\mu^{FX}$ et du tiering du tableau ; $\bar T$ conservé & \S\ref{sec:systemes} & Idem\\
6 & Calibration : dix lignes retirées ; glossaire, analyse d'impact et tests allégés d'autant & \S\ref{sec:calib} & ---\\
\bottomrule
\end{longtable}
\end{center}
Décompte : quatre équations \emph{allégées} (\eqref{eq:spread}, \eqref{eq:cred}, \eqref{eq:ecb}, \eqref{eq:fonds}), aucune ajoutée ; onze paramètres retirés ($s_5$, $s_T$, $s_{10}$, $\nu_5$, $\bar\varphi$, $\nu_6$, $\lambda_2$, $\vartheta$, $\varrho^{tier}$, $D^{CB}$, $\mu^{FX}$, $\delta_i$). Le document perd deux pages (72 contre 74).

%=====================================================================
\section{Objectifs, principes et sources d'inspiration}
%=====================================================================

\subsection{Ce que le jeu doit permettre}
\begin{itemize}
\item \textbf{Un joueur = un État.} Le joueur ne construit pas les usines (sauf en mode planifié) : il fixe le cadre. Ses leviers sont monétaires (taux, monnaie, change), budgétaires (recettes, dépenses, dette) et institutionnels (propriété, prix, commerce, capitaux).
\item \textbf{Un secteur privé autonome.} Entreprises, ménages et banques prennent leurs décisions à partir de règles comportementales dérivées de l'optimisation (demande de travail, investissement selon le $q$ de Tobin, consommation selon l'équation d'Euler ou une règle keynésienne selon la strate, portefeuille selon les rendements relatifs).
\item \textbf{Plusieurs approches viables.} Une économie libérale et une économie planifiée doivent toutes deux être jouables, avec des forces et des faiblesses différentes (croissance et cycles d'un côté ; plein emploi, égalité, pénuries et inefficience allocative de l'autre) plutôt qu'une hiérarchie imposée par le concepteur.
\item \textbf{Émergence plutôt que script.} Inflation, chômage, cycles, crises de change, crises bancaires et défauts souverains doivent résulter des interactions entre agents et des décisions des joueurs, pas d'événements tirés au hasard.
\item \textbf{Cohérence comptable stricte.} Tout flux monétaire quitte un bilan pour entrer dans un autre (approche \emph{stock-flux cohérente}, Godley--Lavoie 2007). Aucune monnaie n'est \og brûlée\fg{} ni créée hors des mécanismes bancaires explicites.
\end{itemize}

\subsection{Comment lire ce document}
Chaque équation est suivie d'un encadré \emph{Lecture} en quatre parties : la liste des \textbf{variables} qu'elle utilise, \textbf{ce qu'elle dit} en langage courant, les \textbf{hypothèses} qu'elle suppose et ses \textbf{limites}. Les équations qui portent un nom dans la littérature (courbe de Phillips, règle de Taylor, $q$ de Tobin\ldots) le conservent à titre d'information ; il n'est pas nécessaire de connaître ces travaux pour suivre. La section~\ref{sec:rappels} rappelle les quelques notions de base utilisées partout ; l'annexe~\ref{sec:glossaire} rassemble toute la notation ; la section~\ref{sec:impacts} décrit, variable par variable, ce qui se passe quand on l'augmente ou la diminue, en distinguant les effets directs des effets de second tour.

\paragraph{Conventions.} Un indice $t$ désigne le temps (le tick, ou le mois selon le contexte), $j$ un secteur ou un bien, $h$ une strate de ménages, $i$ et $k$ des pays. Une barre ($\bar x$) désigne une valeur de référence ou une cible, un chapeau ($\hat x$) un écart à cette référence ou une valeur visée, un astérisque ($x^\ast$) une valeur étrangère ou d'équilibre, $\Delta x = x_{t+1}-x_t$ une variation, $\E_t x$ une anticipation formée en $t$, $\pos{x}=\max(x,0)$ la partie positive, $\ind_{\{\cdot\}}$ une indicatrice (1 si la condition est vraie, 0 sinon) et $\operatorname{clip}(x;a,b)$ la valeur $x$ ramenée dans l'intervalle $[a,b]$. Les taux d'intérêt et d'inflation sont en base annuelle.

\subsection{Ce que nous apprennent Victoria 3 et le mod E\&F}

\begin{vic}[l'économie de Victoria 3]
\begin{itemize}
\item \textbf{Marché par ordres.} Chaque bien a un prix de base ; le prix de marché est fixé par le rapport entre ordres d'achat (consommation des pops, intrants des bâtiments, exportations) et ordres de vente (production, importations), borné à $\pm 75\,\%$ du prix de base. Tous les ordres sont servis quelle que soit la disjonction offre/demande : le déséquilibre passe uniquement par le prix.
\item \textbf{Chaînes de valeur.} Les bâtiments consomment des intrants en proportions fixes (Leontief) et produisent selon un \emph{throughput} ; une pénurie d'intrant (achats $> 2\times$ ventes) réduit la production et se propage en aval.
\item \textbf{Pops et besoins.} Les populations sont classées par strate et profession ; leurs besoins sont hiérarchisés (nourriture, chauffage, vêtements, luxe) et leur consommation augmente par paliers avec la richesse ; le niveau de vie (\emph{Standard of Living}) en découle et pilote démographie et radicalisation.
\item \textbf{Réserves de trésorerie et crédit.} Les bâtiments accumulent des réserves ; la somme des réserves des bâtiments (plus une base et une fraction du PIB) constitue le plafond de crédit de l'État ; l'excédent budgétaire alimente une réserve d'or plafonnée. Le taux d'intérêt souverain dépend du rang et des technologies. Au-delà du plafond : défaut.
\item \textbf{Monnaie nominale sans inflation.} Une livre de 1836 vaut une livre de 1936 ; il n'y a pas de politique monétaire dans le jeu de base --- c'est précisément la lacune que comble E\&F et que notre jeu place au centre.
\end{itemize}
\end{vic}

\begin{vic}[le mod Economic \& Financial (E\&F) V4]
\begin{itemize}
\item \textbf{Banques centrales émettrices.} Chaque grande puissance possède une banque centrale unique qui \emph{produit} une devise (livre, dollar, franc, mark\ldots) comme un bien, à partir de papier, de services et d'or ; les pops et l'économie ont un \emph{besoin} de monnaie et de produits financiers. Un déficit d'émission crée une \og pénurie de monnaie\fg{} (renchérissement de la devise, contraction de l'activité) ; un excédent crée de l'inflation.
\item \textbf{Régimes monétaires.} Méthodes de production de la banque centrale : étalon-or, étalon-argent, absence d'étalon (économie isolée), monnaie de réserve (fin de partie), auxquels s'ajoutent dévaluation et réévaluation (plus de monnaie pour moins d'or, au prix d'inflation et de prestige, ou l'inverse). La monnaie fiduciaire y est présentée comme un outil de crise.
\item \textbf{Forex, bourse, balance courante.} Achat/vente de devises étrangères, réserves de change des banques centrales, marché boursier, banques commerciales, bons du Trésor, suivi de la balance commerciale et du compte courant.
\end{itemize}
\end{vic}

\begin{design}[ce que l'on garde, ce que l'on change]
\begin{itemize}
\item \textbf{On garde} : le marché par ordres avec ajustement borné des prix (mais rendu dynamique, donc capable d'inflation), les intrants Leontief avec propagation des pénuries, la hiérarchie des besoins (formalisée par un système de demande de Stone--Geary), les réserves de trésorerie des entreprises comme mécanisme de faillite, l'idée E\&F d'une banque centrale au bilan explicite, de régimes de change et de monnaie non convertible.
\item \textbf{On change} : la monnaie n'est plus un \og bien produit\fg{} mais un passif du système bancaire, créé par le crédit et par la banque centrale ; il existe une inflation, des taux d'intérêt réels et nominaux, des anticipations ; l'État n'est plus le constructeur universel mais un régulateur qui peut, s'il le choisit, devenir planificateur.
\item \textbf{On ajoute} : chômage et courbe de Phillips, cycle du crédit, prime de risque souveraine et défaut, contrôle des capitaux, planification input--output avec contrainte budgétaire lâche et pénalités d'inefficience.
\end{itemize}
\end{design}

\subsection{Ancrages théoriques mobilisés}
\begin{center}\small
\begin{tabular}{p{4.6cm}p{10.4cm}}
\toprule
\textbf{Bloc du jeu} & \textbf{Références}\\
\midrule
Croissance et taux naturel & Solow (1956) ; Ramsey (1928), Cass (1965), Koopmans (1965) : $r^\ast = \rho + \sigma g$ \\
Consommation & Équation d'Euler (Ramsey) pour les ménages non contraints ; fonction de consommation keynésienne pour les contraints ; système de dépenses linéaire de Stone (1954), Geary (1950)\\
Investissement & Tobin (1969), Hayashi (1982) ; accélérateur financier de Bernanke--Gertler--Gilchrist (1999)\\
Demande globale & IS--LM (Hicks 1937) ; IS dynamique néo-keynésienne (Galí 2008) ; Mundell--Fleming en économie ouverte\\
Inflation & Courbe de Phillips (Phillips 1958, Friedman 1968) ; version hybride (Galí--Gertler 1999) ; théorie quantitative comme ancre de long terme ; hyperinflation de Cagan (1956)\\
Politique monétaire & Règle de Taylor (1993) ; Woodford (2003) ; crédibilité et ancrage des anticipations\\
Change & Parité des taux d'intérêt, PPA, surréaction de Dornbusch (1976) ; attaques spéculatives (Krugman 1979)\\
Finances publiques & Dynamique de la dette (Blanchard 1990, Domar 1944) ; équivalence ricardienne \emph{vs} contrainte de liquidité ; défaut souverain (Eaton--Gersovitz 1981, Reinhart--Rogoff 2009)\\
Banques et crédit & Monnaie endogène (Moore 1988), cohérence stock-flux (Godley--Lavoie 2007), instabilité financière (Minsky 1986)\\
Commerce & Armington (1969) ; élasticités de substitution ; tarifs\\
Planification & Tableau input--output de Leontief ; Kantorovitch ; Lange (1936) ; Kornai (1980, 1992) : contrainte budgétaire lâche et économie de la pénurie ; Hayek (1945) : problème de l'information\\
\bottomrule
\end{tabular}
\end{center}

%=====================================================================
\section{Rappels : dix notions pour lire la suite}
\label{sec:rappels}
%=====================================================================
\begin{description}[style=nextline,leftmargin=1.2em]
\item[PIB, production potentielle, écart de production.] Le PIB $Y$ est la valeur de tout ce qui est produit. La production \emph{potentielle} $Y^{pot}$ est ce que l'économie produirait à plein emploi de ses moyens sans surchauffe. L'écart de production $\hat y=\ln(Y/Y^{pot})$ est positif en surchauffe (l'inflation accélère) et négatif en récession (le chômage monte). Le PIB nominal est $PY$ : quantités multipliées par les prix.
\item[Inflation, prix relatifs, indice des prix.] L'inflation $\pi$ est la hausse de l'indice général des prix $P$. Un prix qui monte plus vite que les autres change les \emph{prix relatifs} et réoriente la production ; quand tous les prix montent ensemble, c'est de l'inflation pure, qui ne change rien aux quantités mais érode la monnaie et les dettes.
\item[Taux nominal et taux réel.] Un taux d'intérêt nominal $i$ est ce qui est écrit sur le contrat ; le taux réel $r = i - \pi^e$ est ce qu'il rapporte une fois l'inflation anticipée déduite (relation de Fisher). C'est le taux réel qui compte pour investir et épargner : un taux nominal de 10\,\% avec 12\,\% d'inflation est un cadeau à l'emprunteur.
\item[Anticipations.] Les décisions dépendent de ce que les agents \emph{croient} du futur. Des anticipations \emph{adaptatives} extrapolent le passé ; des anticipations \emph{ancrées} croient la cible annoncée par une banque centrale crédible. Le passage de l'une à l'autre est au cœur des régimes d'inflation.
\item[Monnaie : base et masse.] La \emph{base monétaire} $H$ est ce que la banque centrale crée directement (billets et réserves des banques). La \emph{masse monétaire} $M$ est ce que le public utilise (billets et dépôts). Ce sont les prêts des banques qui créent les dépôts : la masse monétaire est « endogène », la banque centrale l'influence par les taux et la réglementation, elle ne la fixe pas.
\item[Bilan, actif, passif, fonds propres.] Tout agent possède des actifs (ce qu'il détient) et des passifs (ce qu'il doit) ; la différence est sa richesse nette ou, pour une entreprise ou une banque, ses \emph{fonds propres}. Un dépôt est un actif pour le ménage et un passif pour la banque ; une obligation est un actif pour l'épargnant et un passif pour l'État. Le modèle vérifie qu'à tout instant, la somme des richesses nettes de tous les agents d'un pays plus sa position vis-à-vis de l'étranger vaut le capital physique du pays : rien ne se perd.
\item[Levier et collatéral.] Le levier est le rapport entre la dette et les actifs (ou les fonds propres) : un agent à fort levier gagne beaucoup si ses actifs montent et perd tout s'ils baissent un peu. Le collatéral est l'actif mis en garantie d'un prêt ; sa valeur détermine combien on peut emprunter, d'où un cercle : prix des actifs \up{} $\Rightarrow$ crédit \up{} $\Rightarrow$ prix des actifs \up.
\item[Solde primaire et soutenabilité.] Le solde primaire est le budget hors intérêts. La dette est soutenable si le ratio dette/PIB ne diverge pas ; cela dépend de la comparaison entre le taux d'intérêt réel et la croissance, et du solde primaire.
\item[Balance commerciale, compte courant, change.] Le compte courant est ce que le pays gagne sur l'étranger (exportations et revenus) moins ce qu'il lui paie. Un déficit se finance par de l'endettement extérieur ou par une baisse des réserves de change. Le taux de change $e$ est le prix de la monnaie étrangère : quand $e$ monte, la monnaie nationale se \emph{déprécie}, les importations coûtent plus cher et les exportations deviennent plus compétitives.
\item[Toutes choses égales par ailleurs.] Les analyses de la section~\ref{sec:impacts} font varier une variable en gelant les autres, puis relâchent progressivement les autres pour décrire les effets de second tour. Dans le jeu, tout bouge en même temps ; la décomposition est un outil de raisonnement, pas une prédiction.
\end{description}

%=====================================================================
\section{Architecture générale}
%=====================================================================

\subsection{Entités et indices}
\begin{itemize}
\item Pays : $i \in \{1,\dots,N\}$ (un joueur ou une IA par pays). Sauf mention contraire, les équations sont écrites pour un pays et l'indice $i$ est omis.
\item Biens (secteurs) : $j \in \mathcal{J}$. Liste proposée pour la version~1 : \emph{alimentation}, \emph{énergie et matières premières}, \emph{biens intermédiaires} (acier, chimie), \emph{biens de consommation}, \emph{biens d'équipement}, \emph{construction}, \emph{services}, plus un secteur de \emph{ressources} (pétrole/or) pour les pays dotés (\S\ref{sec:hetero}). L'agriculture utilise en outre un facteur fixe, la terre $T$. Chaque secteur $j$ est modélisé comme une entreprise représentative (agrégat de firmes), à la manière des \og bâtiments\fg{} de Victoria~3 qui représentent un secteur et non une usine.
\item Strates de ménages : $h \in \{B, M, H\}$ (basse : ouvriers et paysans ; moyenne : employés qualifiés, artisans, petits propriétaires ; haute : détenteurs de capital). Effectifs $N_h$, part de la propriété du capital privé $\omega_h$ avec $\sum_h \omega_h = 1$.
\item Secteur bancaire (une banque représentative par pays), banque centrale, État, reste du monde.
\end{itemize}

\subsection{Temps}
Deux échelles :
\begin{itemize}
\item le \textbf{tick} $t$ (une semaine) pour les marchés des biens, du travail, le crédit et les paiements ;
\item le \textbf{mois} (4 ticks) pour la mise à jour des salaires, des prix administrés, de l'indice des prix, des anticipations, de la démographie, de la productivité et pour les décisions de politique économique ; certaines décisions structurelles (lois, régime de change, nationalisations) ont un délai d'application de plusieurs mois.
\end{itemize}
Les taux d'intérêt sont exprimés en base annuelle et convertis au tick par $i_t^{(w)} = (1+i)^{1/52}-1$.

\subsection{Bilans et cohérence stock-flux}
Chaque agent possède un bilan ; la matrice des flux est fermée. Notations principales :
\begin{center}\small
\begin{tabular}{p{2.6cm}p{6.6cm}p{5.6cm}}
\toprule
\textbf{Agent} & \textbf{Actifs} & \textbf{Passifs}\\
\midrule
Ménages $h$ & dépôts $D_h$, obligations $B_h$, actions $p^E E_h$, immobilier $p^Z Z_h$ & crédits conso. et hypothécaires $L_h$, prêts sur marge $L^{marg}_h$\\
Entreprises $j$ & trésorerie $M_j$, capital $p_K K_j$, stocks $S^{inv}_j$ & dette bancaire $L_j$, fonds propres\\
Banque & réserves à la BC $Res$, crédits $L$, obligations $B^{Bk}$ & dépôts $D$, refinancement BC $L^{CB}$, fonds propres $E^{Bk}$\\
Banque centrale & or $G$, devises $R^{FX}$, obligations $B^{CB}$, refinancement $L^{CB}$ & billets $C$, réserves des banques $Res$\\
État & dépôt au Trésor $M^G$, participations $\theta_j$ & dette $B = B_h + B^{Bk} + B^{CB} + B^{\ast}$\\
Reste du monde & créances sur le pays & dettes envers le pays\\
\bottomrule
\end{tabular}
\end{center}
La masse monétaire au sens large est $M = C + D$ ; la base monétaire est $H = C + Res$.

\subsection{Ordre de résolution d'un tick}
\begin{enumerate}
\item \textbf{Plans.} Les entreprises fixent production visée, demande de travail et commandes d'intrants ; les ménages fixent leur budget de consommation ; l'État verse ses dépenses programmées.
\item \textbf{Marché du travail.} Embauches et licenciements (ajustement partiel) ; les salaires sont révisés une fois par mois.
\item \textbf{Marchés des biens.} Agrégation des ordres d'achat (ménages, entreprises, État, exportations) et de vente (production + stocks + importations) ; mise à jour des prix ; rationnement si nécessaire ; flux de commerce international.
\item \textbf{Paiements.} Salaires, achats, impôts et cotisations, dividendes ; mise à jour des trésoreries et des dépôts.
\item \textbf{Finance.} Demandes de crédit, offre bancaire, taux ; émission de dette publique et adjudication ; opérations de la banque centrale ; marché des changes et réserves.
\item \textbf{Capital.} Investissement, dépréciation, faillites et liquidations.
\item \textbf{État.} Solde budgétaire, dette, intérêts, éventuel défaut.
\item \textbf{Fin de mois.} Indice des prix, inflation, anticipations, crédibilité, productivité, démographie, niveau de vie, soutien politique, décisions des joueurs et des IA.
\end{enumerate}


\begin{figure}[htbp]
\centering
\resizebox{0.95\textwidth}{!}{%
\begin{tikzpicture}[x=1cm,y=1cm]
\node[agent] (men)  at (0,7)     {Ménages\\ $B,M,H$};
\node[agent] (ent)  at (10,7)    {Entreprises\\ secteurs $j$};
\node[agent] (bk)   at (0,2)     {Banque\\ commerciale};
\node[agent] (etat) at (10,2)    {État\\ (Trésor)};
\node[agent] (bc)   at (5,-3)     {Banque centrale};
\node[agent,fill=gray!12] (rdm) at (16.4,4.5) {Reste du monde};

% ménages <-> entreprises
\draw[flux] (ent.north) to[out=110,in=70] node[lab,pos=0.5]{salaires $WL$, dividendes $(1-\theta)Div$} (men.north);
\draw[flux] (men.east) -- node[lab,pos=0.5,above]{consommation $\sum_j p_j x_{hj}$} (ent.west);

% ménages <-> banque
\draw[fluxf] ([xshift=-9mm]men.south) -- node[lab,pos=0.5,left]{dépôts $D_h$} ([xshift=-9mm]bk.north);
\draw[fluxf] ([xshift=9mm]bk.north) -- node[lab,pos=0.5,right]{crédits $L_h$} ([xshift=9mm]men.south);

% banque -> entreprises (diagonale droite)
\draw[fluxf] (bk.north east) -- node[lab,pos=0.16,above,sloped]{crédits $L_j$, intérêts $i^L$} (ent.south west);

% ménages -> état (diagonale droite)
\draw[flux] (men.south east) -- node[lab,pos=0.84,above,sloped]{impôts, transferts $Tr$,\\ obligations $B_h$} (etat.north west);

% entreprises <-> état
\draw[flux] ([xshift=11mm]ent.south) -- node[lab,pos=0.5,right]{$\tau_\Pi$, TVA, $\theta_j Div_j$} ([xshift=11mm]etat.north);
\draw[flux] ([xshift=-4mm]etat.north) -- node[lab,pos=0.5,left]{achats $G_j$, $Sub_j$} ([xshift=-4mm]ent.south);

% banque <-> état
\draw[fluxf] ([yshift=3mm]bk.east) -- node[lab,pos=0.62,above]{obligations $B^{Bk}$} ([yshift=3mm]etat.west);
\draw[fluxf] ([yshift=-3mm]etat.west) -- node[lab,pos=0.38,below]{intérêts $i^B$} ([yshift=-3mm]bk.east);

% commerce extérieur
\draw[flux] (ent.north east) to[out=30,in=120] node[lab,pos=0.5,above,sloped]{exportations $EX$} (rdm.north west);
\draw[flux] (rdm.west) to[out=170,in=-20] node[lab,pos=0.5,below,sloped]{importations $IM$} (ent.east);
\draw[fluxf] (rdm.south west) to[out=210,in=30] node[lab,pos=0.5,below,sloped]{dette externe $B^\ast$, capitaux $KA$} (etat.east);

% banque centrale
\draw[fluxp] (bc.west) to[out=180,in=-90] node[lab,pos=0.45,left]{refinancement $L^{CB}$ à $i^{CB}$} (bk.south);
\draw[fluxp] (bc.east) to[out=0,in=-90] node[lab,pos=0.55,right]{QE $\Delta B^{CB}$, avances $\Delta A^G$,\\ seigneuriage $\Sigma$} (etat.south);
\draw[fluxp] (bc.south east) to[out=-30,in=-100] node[lab,pos=0.6,below]{réserves $R^{FX}$, or $G$} (rdm.south);
\end{tikzpicture}}
\caption{Flux entre agents d'un pays. Noir : flux réels et fiscaux ; orange : flux financiers (crédit, titres) ; vert pointillé : opérations de banque centrale. Chaque flèche est un débit sur un bilan et un crédit sur un autre (cohérence stock-flux). Les marchés d'actifs sont détaillés en figure~\ref{fig:fin}.}
\label{fig:flux}
\end{figure}

%=====================================================================
\subsection{Règles de caisse et bornes numériques (v0.9)}
\label{sec:sfc}
Trois règles de bord garantissent que la cohérence stock-flux tient dans tous les régimes, hyperinflation comprise, et que la monnaie ne se crée que par les quatre portes prévues (crédit bancaire, absorption de dette publique par une banque solvable, avances de la banque centrale, recapitalisation publique) :
\begin{itemize}
\item \textbf{Aucun découvert} : ménages, entreprises et État ne peuvent transférer plus que leur dépôt ; un paiement impossible est réduit (arriéré) et non financé par la banque.
\item \textbf{Contraintes de caisse hebdomadaires} : la masse salariale d'un secteur est bornée par $v(\pi^e)\,(M_j+\tfrac12 p_jQ_j)$ et ses achats d'intrants par le solde après salaires, où $v(\pi^e)=\exp(a\min(\pi^e,1))$ est le facteur de vitesse (en forte inflation la monnaie tourne plus vite et la contrainte se relâche, jusqu'à sept fois). La contrainte de caisse est ce qui, en hyperinflation, transmet la pénurie de monnaie réelle à l'emploi, et ce qui, après une réforme, empêche les salaires nominaux de croître sans monnaie.
\item \textbf{Bornes} : prix $\times[0{,}7\,;1{,}19]$ par semaine, salaires $\times[0{,}9\,;2]$ par mois (0,5 si non payables), $\pi^e\le12$, vitesse saturée à $\pi^e=1$, taux directeur $\le1$, prime souveraine $\le0{,}3$. Elles n'agissent qu'en régime extrême et rendent l'hyperinflation finie sans en changer la nature (la redénomination \eqref{eq:reform} prend le relais).
\end{itemize}
\paragraph{Ajout de la v1.0 : règlement international.} Chaque banque centrale tient un \emph{compte de change} (dépôt en monnaie nationale à la banque) : les importateurs y paient la valeur rendue de leurs achats (droits de douane reversés à l'État), les exportateurs en sont payés. Les banques centrales se règlent en réserves libellées dans la monnaie de référence ; la position du pays émetteur de cette monnaie peut être négative (passif extérieur : privilège de l'émetteur). Trois règles rendent le système exactement cohérent : (i) l'exportateur reçoit \emph{ce que l'importateur a effectivement payé}, converti au change du jour de paiement, et non la valeur des quantités livrées (avec des prix qui bougent entre la commande et la livraison et des paiements bornés par la trésorerie, le règlement sur quantités faisait dériver les réserves mondiales de 5\,\% en huit ans) ; (ii) les flux de capitaux (achats d'obligations étrangères par les ménages, rapatriements) passent par les comptes de change et se règlent en réserves dans les deux sens ; (iii) les intérêts versés aux détenteurs étrangers d'obligations sont payés par l'État débiteur et reçus par les ménages créanciers, réglés de même. Deux tests gardent le système : les identités de bilan de chaque banque centrale (réserves valorisées au change courant, plus-values de change au bilan), et la \emph{conservation des réserves mondiales} au montant en transit près (une semaine d'importations) ; c'est ce second test, global, qui a révélé deux signes inversés dans les flux de capitaux que les identités locales ne voyaient pas.

%=====================================================================
\section{Secteur productif}
\label{sec:prod}
%=====================================================================

\subsection{Technologie}
Chaque secteur $j$ combine capital $K_j$, travail $L_j$ et intrants $X_{kj}$ (quantité du bien $k$ utilisée par $j$). La capacité est une Cobb--Douglas, les intrants sont en proportions fixes (Leontief) comme les méthodes de production de Victoria~3 :
\begin{equation}
Y_j^{cap} = A_j\, (K_j)^{\alpha_j} (H\, L_j)^{1-\alpha_j}, \qquad
Y_j = Y_j^{cap} \cdot \min\!\Big(1,\ \min_{k \in \mathcal{K}_j} \frac{X_{kj}^{obt}}{a_{kj}\, Y_j^{cap}}\Big),
\label{eq:prod}
\end{equation}
\begin{lecture}[production d'un secteur]
\variables $Y_j$ : quantité produite par le secteur $j$ pendant le tick ; $Y_j^{cap}$ : capacité, c'est-à-dire ce que le secteur pourrait produire s'il disposait de tous ses intrants ; $A_j$ : productivité globale des facteurs (efficacité avec laquelle machines et travail sont combinés : organisation, technologie) ; $K_j$ : stock de capital (machines, bâtiments) ; $L_j$ : nombre de travailleurs ; $H$ : capital humain moyen (qualification, santé), qui multiplie l'efficacité de chaque travailleur ; $\alpha_j\in(0,1)$ : part du capital dans la production ; $X^{obt}_{kj}$ : quantité du bien $k$ effectivement reçue comme intrant ; $a_{kj}$ : quantité de $k$ nécessaire par unité de $j$ (par exemple 1,5 tonne de minerai par tonne d'acier).
\sens Deux règles superposées. La première (Cobb--Douglas, 1928) dit qu'avec plus de machines ou plus de travailleurs on produit plus, mais avec des \emph{rendements décroissants} : doubler seulement le travail sans changer les machines ne double pas la production. Doubler les deux la double exactement (rendements d'échelle constants). La seconde (Leontief) dit qu'il faut des intrants en proportions fixes : sans minerai, pas d'acier, quelle que soit la taille de l'usine. Le facteur $\min(1,\cdot)$ mesure le taux de fonctionnement effectif : si le secteur n'a reçu que 60\,\% du minerai dont il a besoin, il ne produit que 60\,\% de sa capacité.
\hyp Les entreprises d'un secteur sont agrégées en une seule entreprise représentative ; la substitution entre capital et travail est possible mais pas la substitution entre intrants ; pas de délai entre embauche et production.
\limites Les vraies entreprises peuvent, avec le temps, changer de recette (moins d'énergie par unité produite quand elle devient chère). Le modèle traite ce changement uniquement par la modification des coefficients $a_{kj}$ au fil des technologies, pas par une réaction aux prix.
\end{lecture}
où $a_{kj}$ est le coefficient technique (unités de $k$ par unité de $j$), $X^{obt}_{kj}$ la quantité effectivement obtenue après rationnement (\S\ref{sec:marches}), $H$ le capital humain moyen et $A_j$ la productivité globale des facteurs. Le facteur multiplicatif joue le rôle du \emph{throughput} de Victoria~3 : une pénurie d'intrant réduit la production sans réduire l'emploi immédiatement, ce qui dégrade la rentabilité et propage la pénurie vers l'aval.

\paragraph{Commande d'intrants.} $X^{dem}_{kj} = a_{kj}\, \hat Y_j + \lambda_S (s^\ast a_{kj}\hat Y_j - S^{inv}_{kj})$, où $\hat Y_j$ est la production visée et $S^{inv}$ le stock d'intrants (cible $s^\ast$ semaines de production).

\subsection{Production visée et demande de travail}
\label{sec:stocks}
La production visée suit la demande anticipée (moyenne mobile des ventes) corrigée du niveau des stocks de produits finis :
\begin{equation}
\hat Y_{j,t} = \bar Q_{j,t}\,(1 + g^e) + \lambda_{inv}\,(s^\ast_j \bar Q_{j,t} - S^{inv}_{j,t}),
\qquad \bar Q_{j,t} = (1-\mu)\bar Q_{j,t-1} + \mu\, Q_{j,t-1}.
\end{equation}
\begin{lecture}[production visée et gestion des stocks]
\variables $\hat Y_{j,t}$ : production que le secteur \emph{souhaite} réaliser au tick $t$ ; $\bar Q_{j,t}$ : ventes moyennes récentes (moyenne mobile) ; $Q_{j,t-1}$ : ventes du tick précédent ; $\mu\in(0,1)$ : vitesse d'oubli de la moyenne mobile ; $g^e$ : croissance anticipée de la demande ; $S^{inv}_{j,t}$ : stock de produits finis invendus ; $s^\ast_j$ : niveau de stock souhaité, en semaines de ventes ; $\lambda_{inv}$ : vitesse à laquelle le secteur corrige un écart de stock.
\sens L'entreprise ne connaît pas la demande future ; elle extrapole ses ventes récentes, puis ajuste : si ses stocks sont trop hauts, elle produit moins que ses ventes pour les écouler ; s'ils sont trop bas, elle produit plus. C'est le mécanisme classique du \emph{cycle des stocks} : une baisse de la demande fait chuter la production plus que proportionnellement, le temps que les stocks se résorbent.
\hyp Anticipations adaptatives (on regarde le passé), pas de connaissance du futur ; tous les secteurs ont la même règle.
\limites La règle ne fait pas la différence entre une baisse temporaire et une baisse durable de la demande ; c'est volontaire (cela crée des cycles réalistes), mais cela rend les secteurs myopes.
\end{lecture}
La demande de travail découle de l'inversion de \eqref{eq:prod} et de la condition de premier ordre $W_j = (1-\alpha_j)\, p_j Y_j / L_j$ ; l'ajustement est partiel (coûts d'ajustement, délais d'embauche) :
\begin{equation}
L^\ast_j = \min\!\Big\{\Big(\tfrac{\hat Y_j}{A_j K_j^{\alpha_j} H^{1-\alpha_j}}\Big)^{\frac{1}{1-\alpha_j}},\ 
\frac{(1-\alpha_j)\,p_j \hat Y_j}{W_j (1+\tau_{S})}\Big\}, \qquad
L_{j,t+1} = L_{j,t} + \lambda_L\,(L^\ast_j - L_{j,t}),
\label{eq:labour}
\end{equation}
\begin{lecture}[demande de travail]
\variables $L^\ast_j$ : effectif que le secteur souhaite employer ; $W_j$ : salaire nominal hebdomadaire ; $\tau_S$ : taux de cotisations sociales payées par l'employeur (le coût du travail est $W_j(1+\tau_S)$) ; $p_j$ : prix de vente du bien $j$ ; $\lambda_L\in(0,1)$ : fraction de l'écart entre effectif souhaité et effectif actuel comblée chaque semaine.
\sens Le premier terme du $\min$ répond à la question « combien de travailleurs pour produire $\hat Y_j$ avec le capital que j'ai ? » (on inverse la fonction de production). Le second répond à « jusqu'où l'embauche reste rentable ? » : un travailleur supplémentaire rapporte sa productivité marginale $(1-\alpha_j)p_jY_j/L_j$ et coûte $W_j(1+\tau_S)$ ; au-delà du point d'égalité, embaucher détruit du profit. Le secteur retient le plus petit des deux : il n'embauche jamais à perte, même si la demande existe. L'ajustement partiel ($\lambda_L$) représente les délais et coûts d'embauche et de licenciement.
\hyp Concurrence suffisante pour que le salaire égale la productivité marginale ; pas de contrat de long terme.
\limites Le modèle ne distingue pas les qualifications à l'intérieur d'un secteur ; en réalité une entreprise licencie rarement dès le premier trimestre de pertes (rétention de main-d'œuvre), ce que $\lambda_L$ petit approxime.
\end{lecture}
avec $\tau_S$ le taux de cotisations sociales employeur. Le second terme du $\min$ borne l'emploi par la rentabilité : à salaire trop élevé relativement au prix de vente, le secteur licencie même si la demande existe.

\subsection{Compte d'exploitation, trésorerie et dividendes}
\begin{align}
\Pi_j &= p_j Q_j - W_j(1+\tau_S) L_j - \sum_k p_k X^{obt}_{kj} - \delta_j\, p_K K_j - i^L_j L_j - \tau_\Pi \pos{\Pi^{avant\ IS}_j}, \label{eq:profit}\\
M_{j,t+1} &= M_{j,t} + \Pi_j - Div_j - p_K I_j + \Delta L_j + Sub_j, \label{eq:cash}\\
Div_j &= \min\!\big(\psi_j \pos{\Pi_j},\ \pos{M_j - R_j}\big) + \tfrac14\pos{M_j - R_j - \bar m\,p_jQ_j}, \label{eq:div}\\
&\quad Div_j \to \text{ménages } (\omega_h (1-\theta_j)) \text{ et État } (\theta_j), \nonumber
\end{align}
\begin{lecture}[compte d'exploitation, trésorerie, dividendes]
\variables $\Pi_j$ : profit net du tick ; $Q_j$ : quantité vendue ; $\delta_j$ : taux de dépréciation du capital (usure) ; $p_K$ : prix d'une unité de capital (biens d'équipement) ; $i^L_j$ : taux d'intérêt payé sur la dette bancaire $L_j$ ; $\tau_\Pi$ : taux de l'impôt sur les sociétés, appliqué seulement si le profit avant impôt est positif ($\pos{\cdot}$ signifie « partie positive ») ; $M_j$ : trésorerie ; $Div_j$ : dividendes versés ; $I_j$ : investissement (achat de nouveau capital) ; $\Delta L_j$ : nouveaux emprunts nets ; $Sub_j$ : subventions reçues de l'État ; $\psi_j$ : part du profit distribuée ; $\bar m$ : trésorerie de sécurité souhaitée, en semaines de chiffre d'affaires ; $\theta_j$ : part du capital détenue par l'État ; $\omega_h$ : part du capital privé détenue par la strate $h$.
\sens Le profit est ce qui reste des ventes après salaires, intrants, usure du capital, intérêts et impôts. La trésorerie augmente du profit et des emprunts, diminue des dividendes et de l'investissement. Les dividendes ne sont versés que si le profit est positif \emph{et} si la trésorerie de sécurité est reconstituée : c'est le comportement prudent d'une entreprise, et c'est ce qui fait qu'en récession les dividendes s'effondrent avant les faillites. Une trésorerie négative sans crédit disponible déclenche la faillite.
\hyp Les dividendes sont une fraction fixe du profit ; pas d'émission d'actions nouvelles pour se financer (le financement externe est bancaire).
\limites Les entreprises réelles ont aussi accès au marché obligataire et peuvent racheter leurs actions ; ce sont des raffinements possibles mais non nécessaires pour les mécanismes visés.
\paragraph{Ajouts de la v0.9.} Les dividendes sont \emph{prioritaires} : la fraction $\psi_j$ du profit net est versée dès que la trésorerie le permet (après remboursement $R_j$ d'un quart de la trésorerie excédentaire), puis un quart de l'excédent au-delà de la cible $\bar m$ est distribué ; l'investissement se finance ensuite sur la trésorerie restante puis par le crédit. Ordre inverse dans la v0.8 : l'investissement absorbait la trésorerie et les dividendes tombaient à 0,7\,\% du PIB, les entreprises autofinançant un investissement de 30\,\% du PIB. Avec $\psi=0{,}9$ : dividendes 9\,\% du PIB, crédits aux entreprises 1,0 PIB, levier 0,34. La faillite est également adoucie : 5\,\% du capital liquidé (le reste est repris par le même secteur) et un trimestre sans investir ; à 25\,\% et un mois, des faillites répétées annihilaient le capital d'un secteur en un an, plancher irréversible que le test de stabilisation a révélé.
\end{lecture}
$Q_j$ est la quantité vendue, $\theta_j \in [0,1]$ la part du capital détenue par l'État (\S\ref{sec:systemes}), $\bar m$ une cible de trésorerie (en semaines de chiffre d'affaires, analogue des \emph{cash reserves} de Victoria~3), $Sub_j$ les subventions publiques. Une entreprise dont la trésorerie devient négative demande un crédit ; si le crédit est refusé (\S\ref{sec:banques}), elle est en \textbf{faillite} : une fraction $\phi_F=5\,\%$ de son capital est liquidée à prix décoté (le reste est repris), \emph{une seule fois par épisode} (v1.3 : la liquidation répétée chaque semaine de réorganisation annihilait le capital en hyperinflation), le secteur n'investit pas pendant un trimestre, sa dette est passée en pertes par la banque, ses salariés sont licenciés.

\subsection{Investissement : $q$ de Tobin et accélérateur financier}
\label{sec:invest}
Le rendement attendu du capital est comparé au coût du capital (taux réel prêteur, prime de risque, dépréciation), et l'investissement net est conditionné au taux d'utilisation des capacités :
\begin{equation}
\begin{split}
q_j &= \frac{\E[\Pi^{brut}_j] / (p_K K_j)}{r^L_j + \rho_E + \delta_j}, \qquad r^L_j = i^L_j - \pi^e, \qquad \bar q_{j,t+1} = \bar q_{j,t} + \lambda_q\,(q_{j,t} - \bar q_{j,t}),\\
I_j &= K_j\Big(\delta_j + \min\!\Big(\iota\,(\bar q_j - 1) + \iota\, s_q \tanh\!\frac{q_j-\bar q_j}{s_q},\ 2\delta_j\Big)^{\!+}\cdot \Lambda_j\,\Lambda^{Bk}\,\Gamma_j\Big),\\
\Gamma_j &= \clip{\frac{\hat Y_j/Y^{cap}_j - \underline u}{\bar u - \underline u}}{0}{1},
\end{split}
\label{eq:invest}
\end{equation}
\begin{lecture}[investissement : le $q$ de Tobin]
\variables $q_j$ : rapport entre le rendement du capital et son coût ; $\E[\Pi^{brut}_j]$ : profit \emph{avant} dépréciation et intérêts, anticipé (moyenne mobile) ; $r^L_j = i^L_j - \pi^e$ : taux d'intérêt \emph{réel} du crédit, c'est-à-dire le taux nominal corrigé de l'inflation anticipée $\pi^e$ ; $\iota$ : sensibilité de l'investissement à $q$ ; $\Lambda_j\in[0,1]$ : facteur de rationnement du crédit ; $\ell_j = L_j/(p_K K_j)$ : levier, dette rapportée à la valeur du capital ; $\bar\ell$ : levier au-delà duquel les banques deviennent réticentes ; $\eta_\ell$ : sévérité de ce rationnement.
\sens Le $q$ de Tobin (Tobin 1969) compare ce que rapporte une unité de capital installée à ce qu'elle coûte à financer et à entretenir. Si $q>1$, un euro investi dans le secteur vaut plus qu'un euro placé au taux du crédit : le secteur fait de l'investissement net, d'autant plus que $q$ dépasse~1. Si $q<1$, il ne remplace même pas tout le capital usé. Mais $q=1$ n'est pas le point de repos d'une économie qui croît : voir la révision de la v1.5 ci-dessous. Le facteur $\Lambda_j$ dit qu'une entreprise déjà très endettée ne peut pas emprunter autant qu'elle le voudrait, même si $q$ est élevé : c'est l'\emph{accélérateur financier} (Bernanke, Gertler et Gilchrist 1999) : quand les bilans se dégradent, l'investissement chute plus que ne le justifie la rentabilité.
\hyp Le rendement futur est extrapolé du rendement récent ; le coût du capital est le taux bancaire (pas de financement par actions). Le taux réel est calculé avec les anticipations d'inflation, ce qui est essentiel : une inflation élevée avec des taux nominaux élevés peut laisser le taux réel bas et stimuler l'investissement.
\limites Version simplifiée sans coûts d'ajustement explicites (Hayashi 1982) ; l'incertitude n'apparaît que par la prime de risque bancaire.
\paragraph{Ajouts de la v0.9.} $\rho_E$ : prime de risque exigée sur le capital productif (2\,\% ; écart entre le rendement du capital observé, 8 à 12\,\%, et le taux sûr) ; elle entre aussi dans la marge normale de \eqref{eq:price}. $\Gamma_j$ : taux d'utilisation des capacités (production visée sur capacité) ramené entre 0 et 1 sur $[\underline u,\bar u]=[0{,}80 ; 0{,}95]$ : un secteur aux capacités inutilisées ne fait pas d'investissement net, quel que soit $q$. Le plafond $2\delta_j$ borne l'investissement brut à trois fois le remplacement. Ce que le prototype a montré : sans $\rho_E$, la condition $q=1$ donne $K/Y=\alpha/(r+\delta)\approx4{,}8$ que l'économie n'atteint jamais parce que la marge normale calibrée sur le seul taux sûr comprime les profits ($K/Y$ observé 2,3, $I/Y$ 13\,\%) ; avec $\rho_E=2\,\%$, $I/Y=18\,\%$ et $K/Y=2{,}9$. L'arbitrage à connaître : $\rho_E$ plus élevé relève $K/Y$ mais abaisse les salaires réels, donc la demande, donc le taux naturel effectif ($\rho_E=5\,\%$ : trappe à liquidité) ; il ne pourra monter vers sa valeur observée qu'avec les marchés d'actifs de l'étape~3, qui recyclent la prime en richesse et en consommation. Les variantes testées (achats publics 9,5\,\%, effet richesse doublé, baisse d'impôts) sont documentées en \S\ref{sec:proto} pour l'économie ouverte.
\paragraph{Ajouts de la v1.0.} $\bar q_j$ : $q$ de long terme du secteur (moyenne mobile d'un an, $\lambda_q=1/52$), qui porte le \emph{niveau} de l'investissement net ; $s_q=0{,}03$ : largeur de la réponse aux écarts de court terme $q-\bar q$, saturante (tangente hyperbolique). Une hausse temporaire du taux fait baisser $q$ d'environ 12\,\% par point ($1/(r+\rho_E+\delta)$) ; sans saturation, l'investissement net répondait intégralement et une hausse de 1 pt maintenue un an coûtait $-5{,}9\,\%$ de PIB (contre $-0{,}5$ à $-1{,}5$ observé) ; avec, $-2{,}9\,\%$, l'état stationnaire étant inchangé puisque le niveau suit $\bar q$. Deux compléments : la vitesse d'ajustement de l'emploi est divisée par deux dans le secteur des biens d'équipement (rétention de main-d'œuvre, $\Delta u$ au pic $4{,}6\to3{,}4$ pt) et les dividendes portent sur le profit lissé (\eqref{eq:div}, $\lambda=0{,}1$). Ce que le prototype a rejeté, et pourquoi : (i) la saturation autour de $q=1$ (et non de $\bar q$) plafonne l'investissement net quand $\bar q\approx1{,}1$ et fait tomber l'économie à la borne zéro ; (ii) un taux long lissé (moyenne des taux courts anticipés) réduit l'amplitude mais crée un cycle de 8 à 12 ans d'écart-type 4 à 5 pts de chômage, même avec une règle de Taylor prospective, parce que le retard dans la boucle taux $\to$ investissement $\to$ inflation $\to$ règle fait surréagir la banque centrale ; (iii) l'accélérateur (investissement net proportionnel à la croissance anticipée des commandes) crée les cycles de Samuelson--Hicks (écart-type 3 à 7 pts). Le facteur 2 résiduel par rapport aux estimations empiriques tient pour moitié au protocole (un choc maintenu un an, non un choc transitoire) et reste le chantier ouvert de la calibration.
\paragraph{Révision de la v1.5 : $q^\ast>1$ sur le sentier de croissance.} Sur un sentier de croissance équilibrée, $I/K=\delta+g$ ; la règle d'investissement impose donc $\iota\,\Gamma\,\Lambda\,\Lambda^{Bk}(\bar q-1)=g$, soit
\[
q^\ast = 1+\frac{g}{\iota\,\Gamma\,\Lambda\,\Lambda^{Bk}} .
\]
Avec $g=1{,}9\,\%$, $\iota=0{,}20$, $\Lambda=\Lambda^{Bk}=1$ et $\Gamma=0{,}86$ (taux d'utilisation 0,93), $q^\ast=1{,}11$ : c'est la valeur que le prototype affiche à l'état stationnaire, et elle ne signale aucune rentabilité anormale. L'écart $q^\ast-1$ mesure la croissance, pas un excès de rendement : présenter $q=1$ comme seuil de rentabilité, comme le faisait le texte ci-dessus jusqu'à la v1.4, donnerait au joueur un signal d'investissement permanent. L'indicateur affiché au joueur est désormais $q/q^\ast$ (\S\ref{sec:configs}), $q^\ast$ étant recalculé avec la croissance tendancielle $g_0$ et le taux d'utilisation courant. Vérification annexe : la borne $\min(\cdot,2\delta_j)$ ajoutée au terme $\delta_j$ plafonne bien l'investissement brut à $3\delta_jK_j$.
\paragraph{Révision de la v1.5 : l'accélérateur financier est inactif à la calibration de référence.} Avec $\bar\ell=0{,}6$ et $\eta_\ell=2$, $\Lambda_j$ reste saturé à 1 tant que $\ell_j<0{,}6$. En économie fermée, $\ell=0{,}32$ : il faudrait une chute de 47\,\% du prix du capital pour que $\Lambda<1$. Dans l'économie à actifs, qui est la référence de jeu (\S\ref{sec:proto3}), le levier mesuré est de $0{,}08$ (sectoriel 0,03 à 0,18) : les entreprises se financent par actions et les prêts aux entreprises tombent à 17\,\% du PIB contre 89\,\% en économie fermée. Un balayage de $\bar\ell$ de 0,6 à 0,35 ne change strictement rien à A1, A2, A4 ni T1. Le canal du collatéral, présenté ci-dessus comme central, ne joue donc jamais, et abaisser le seuil ne l'activerait pas. La correction retenue n'est pas un seuil plus bas mais une forme continue centrée sur le levier stationnaire $\ell^\ast$ (prime de financement externe croissante en $\ell/\ell^\ast$, à la Bernanke--Gertler--Gilchrist, avec aggravation non linéaire au-delà de $2\ell^\ast$) ; elle est renvoyée au chantier du partage de la valeur ajoutée (\S\ref{sec:chantiers}), parce que la cause première — le financement par actions qui évince le crédit — est la même.
\end{lecture}
Le facteur $\Lambda_j \in [0,1]$ est un rationnement du crédit décroissant avec le levier $\ell_j = L_j/(p_K K_j)$ : $\Lambda_j = \clip{1 - \eta_\ell(\ell_j - \bar\ell)}{0}{1}$. L'investissement est financé par la trésorerie puis par le crédit bancaire ; il s'exprime en achats de biens d'équipement et de construction (donc en ordres d'achat sur ces marchés), ce qui reproduit la pression que le secteur de la construction exerce sur les prix industriels dans Victoria~3. Accumulation : $K_{j,t+1} = (1-\delta_j) K_{j,t} + I_{j,t}$ (les livraisons sont différées de $T_K$ mois).

\begin{lit}[pourquoi Tobin plutôt qu'un simple accélérateur]
Le $q$ de Tobin relie directement l'investissement au taux d'intérêt réel fixé (indirectement) par le joueur via la banque centrale : c'est le canal de transmission monétaire du côté de l'offre. Il fait aussi apparaître naturellement le \emph{crowding out} lorsque la dette publique fait monter les taux, et l'accélérateur financier lorsque la dégradation des bilans réduit $\Lambda_j$ (Bernanke--Gertler--Gilchrist).
\end{lit}

\subsection{Productivité et progrès technique}
\begin{equation}
\frac{A_{j,t+1}}{A_{j,t}} = 1 + g_0 + \eta_R \frac{RD_t}{Y_t} + \eta_G \Big(\frac{K^G_t}{Y_t}\Big)^{\zeta} + \chi\Big(\frac{A^{front}_{j,t}}{A_{j,t}} - 1\Big) - \varepsilon^{plan}_{j,t} + \varepsilon^A_{j,t},
\label{eq:tfp}
\end{equation}
\begin{lecture}[croissance de la productivité]
\variables $A_{j,t}$ : productivité globale des facteurs du secteur $j$ ; $g_0$ : progrès technique tendanciel exogène ; $RD_t/Y_t$ : dépenses de recherche rapportées au PIB, avec rendement $\eta_R$ ; $K^G_t/Y_t$ : capital public (routes, réseaux, ports) rapporté au PIB, avec rendement $\eta_G$ et élasticité $\zeta<1$ (rendements décroissants) ; $A^{front}_{j,t}$ : meilleure productivité mondiale dans le secteur, et $\chi$ : vitesse de rattrapage ; $\varepsilon^{plan}_{j,t}$ : pénalité liée à la planification ; $\varepsilon^A_{j,t}$ : choc aléatoire.
\sens La productivité, moteur de la croissance à long terme (Solow 1956), progresse pour cinq raisons : une tendance de fond, la recherche, les infrastructures, l'imitation de ceux qui font mieux (un pays loin de la frontière rattrape vite, puis ralentit en s'approchant : c'est la \emph{convergence conditionnelle}) et le hasard. La planification retire une fraction de cette croissance.
\hyp La frontière technologique est accessible à tous (pas de brevets ni de secret) ; le capital public bénéficie à tous les secteurs uniformément.
\limites Le lien recherche--productivité est en réalité long et incertain (10 à 20 ans) ; le modèle l'applique avec un simple retard. Le rattrapage suppose un capital humain suffisant, que la formule ne conditionne pas explicitement (on pourrait multiplier $\chi$ par $H/H^{front}$).
\end{lecture}
avec $RD$ les dépenses de R\&D privées et publiques, $K^G$ le capital public (infrastructures, qui jouent le rôle du \emph{market access} de Victoria~3), $A^{front}$ la frontière technologique mondiale (max sur les pays ; terme de rattrapage), $\varepsilon^{plan}$ la pénalité d'inefficience liée au degré de planification (\S\ref{sec:systemes}) et $\varepsilon^A \sim AR(1)$ un choc. Le capital humain $H$ croît avec les dépenses d'éducation par tête (\S\ref{sec:etat}).

\subsection{Salaires et courbe de Phillips}
Le salaire nominal de chaque secteur est révisé mensuellement :
\begin{equation}
\begin{split}
\frac{W_{j,t+1}}{W_{j,t}} &= 1 + \omega_u\,\varpi_w\,\pi^{ref}_t + \phi_u\,(u^n - u_t) + \phi_x\min(x^L_t,1) + \phi_v\, v_{j,t} + \gamma_A\, \hat g^{prod}_t,
\qquad W_{j,t+1} \ge W^{min}_t,\\
\pi^{ref}_t &= c_t\,\pi^e_t + (1-c_t)\max(\pi^e_t,\pi_t), \qquad \omega_u = \clip{1-2\pos{u_t-u^n}}{0}{1},\\
x^L_t &= \pos{\frac{\sum_j L^\ast_j}{(1-u^n)L^S_t}-1},
\end{split}
\label{eq:wage}
\end{equation}
\begin{lecture}[salaires : la courbe de Phillips]
\variables $W_{j,t}$ : salaire nominal du secteur $j$ ; $\pi^e_t$ : inflation anticipée par les agents ; $u_t$ : taux de chômage ; $u^n$ : taux de chômage « naturel » ou d'équilibre, celui qui ne pousse les salaires ni à la hausse ni à la baisse, rendu endogène par \eqref{eq:un} ; $\phi_u$ : sensibilité des salaires au chômage ; $v_{j,t}$ : tension d'embauche du secteur (postes voulus non pourvus, en proportion de l'effectif) et $\phi_v$ sa sensibilité ; $g^A_j$ : croissance de la productivité du secteur, dont une part $\gamma_A$ est reversée aux salariés ; $W^{min}_t$ : salaire minimum légal.
\sens Les salaires nominaux augmentent (i) pour compenser l'inflation que salariés et employeurs \emph{anticipent}, (ii) plus vite quand le chômage est faible (les employeurs se disputent la main-d'œuvre), moins vite ou même baissent quand il est élevé, (iii) plus vite dans les secteurs qui manquent de bras, (iv) avec la productivité. C'est la \emph{courbe de Phillips augmentée des anticipations} (Phillips 1958, Friedman 1968). Son message central : on ne peut maintenir durablement le chômage sous $u^n$ qu'au prix d'une inflation toujours croissante, car les anticipations finissent par rattraper l'inflation réalisée.
\hyp Révision mensuelle des salaires (rigidité nominale) ; anticipations communes à tous les agents ; $u^n$ fixe.
\limites La pente $\phi_u$ semble s'être aplatie dans les économies avancées depuis les années 1990 (mondialisation, ancrage des anticipations) ; le modèle la traite comme un paramètre à calibrer. La formule ne capture pas les rigidités à la baisse des salaires (les gens acceptent mal une baisse nominale), que l'on peut ajouter par une borne $W_{t+1}\ge W_t$ en régime de faible inflation.
\paragraph{Ajouts de la v0.9.} $\pi^{ref}$ : inflation de référence pour l'indexation, qui est l'anticipation $\pi^e$ quand la banque centrale est crédible et l'inflation \emph{réalisée} quand elle ne l'est plus (c'est l'indexation de fait des hyperinflations) ; $\omega_u$ : l'indexation s'affaiblit avec le chômage de masse (nulle au-delà de $u^n+50$ pts) ; $x^L$ : excès de la demande de travail sur l'offre disponible au plein emploi, terme de \emph{surenchère} qui rend la courbe de Phillips convexe ($\phi_x=0{,}3$) ; $\hat g^{prod}$ : croissance lissée de la productivité apparente du travail (production réelle par emploi), et non plus la tendance $g_0$, avec $\gamma_A=1$ : c'est ce qui ancre la part salariale près de $1-\alpha$ et permet à l'inflation stationnaire d'égaler la cible à $u=u^n$ (avec $\gamma_A<1$, l'inflation stationnaire vaut $\pi^\ast-(1-\gamma_A)g$). Deux règles de bord : le salaire est plafonné à ce que la trésorerie du secteur peut payer (\S\ref{sec:sfc}), et il baisse de 20\,\% par mois quand l'emploi rentable tombe sous la moitié de l'emploi technique (le salaire rend l'emploi non rentable). Ce que le prototype a montré : sans surenchère, une demande monétisée de 20\,\% du PIB donnait 1\,\% de hausse des salaires à 1,8\,\% de chômage ; sans $\pi^{ref}$, elle convergeait vers 35\,\% d'inflation stationnaire au lieu de diverger ; sans $\omega_u$, les salaires nominaux continuaient de doubler chaque mois après l'arrêt de la monnaie.
\end{lecture}
où $u$ est le taux de chômage, $u^n$ le taux \og naturel\fg, $v_j = (L^\ast_j - L_j)/L_j$ la tension d'embauche sectorielle, $g^A_j$ la croissance de la productivité (partage des gains) et $W^{min}$ le salaire minimum légal. En mode planifié, le salaire est administré (\S\ref{sec:systemes}). L'équation \eqref{eq:wage} est la courbe de Phillips augmentée des anticipations : elle est le lieu où les anticipations d'inflation $\pi^e$ (\S\ref{sec:inflation}) entrent dans les coûts, et donc dans les prix.

%=====================================================================
\section{Ménages}
\label{sec:menages}
%=====================================================================

\subsection{Revenu disponible}
\begin{equation}
Y^d_h = (1-\tau_{W,h}) \sum_j W_j L_{jh} + (1-\tau_{K}) \big[\omega_h \textstyle\sum_j (1-\theta_j) Div_j + i^D D_h + i^B B_h\big] + Tr_h - T^{poll} N_h - i^{L}_h L_h ,
\label{eq:yd}
\end{equation}
\begin{lecture}[revenu disponible d'une strate]
\variables $Y^d_h$ : revenu disponible (après impôts et transferts) de la strate $h$ ; $L_{jh}$ : nombre de membres de la strate travaillant dans le secteur $j$ ; $\tau_{W,h}$ : taux d'imposition du travail pour la strate (progressivité : plus élevé pour la strate haute) ; $\tau_K$ : taux d'imposition des revenus du capital ; $\omega_h\sum_j(1-\theta_j)Div_j$ : dividendes reçus (part $\omega_h$ du capital privé, c'est-à-dire hors part publique $\theta_j$) ; $i^D D_h$ : intérêts sur dépôts ; $i^B B_h$ : intérêts sur obligations d'État ; $Tr_h$ : transferts publics (chômage, retraites, minima) et transferts privés reçus de l'étranger (remises des migrants $Rem_h$, en devises, insensibles au taux directeur ; v0.7) ; $T^{poll}$ : capitation (impôt fixe par tête) ; $i^L_h L_h$ : intérêts payés sur les crédits à la consommation et hypothécaires.
\sens C'est la comptabilité du revenu : salaires, revenus du capital, transferts, moins impôts et intérêts. Elle rend visible qui gagne et qui perd à chaque décision : une hausse de $\tau_K$ frappe surtout la strate $H$ (qui détient l'essentiel de $\omega_h$), une hausse des taux d'intérêt enrichit les épargnants et appauvrit les emprunteurs.
\hyp Trois strates homogènes ; les intérêts reçus sur la dette publique vont bien aux ménages (c'est ce qui fait que la dette publique est aussi un actif privé).
\limites La répartition des dividendes est fixée par $\omega_h$ ; la mobilité entre strates est traitée ailleurs de façon simplifiée.
\end{lecture}
avec $\tau_{W,h}$ le taux d'imposition du travail (progressif : $\tau_{W,B} \le \tau_{W,M} \le \tau_{W,H}$), $\tau_K$ l'imposition des revenus du capital, $Tr_h$ les transferts (allocations chômage, retraites, minima sociaux), $T^{poll}$ une éventuelle capitation. Les salariés sont répartis entre strates selon leur qualification ; les chômeurs perçoivent $Tr^{u} = \varrho\, \bar W$ si un régime d'assurance chômage existe.

\subsection{Dépense de consommation : deux régimes comportementaux}
\label{sec:conso}
La strate basse est contrainte par la liquidité et suit une règle keynésienne ; la strate haute lisse sa consommation selon l'équation d'Euler de Ramsey ; la strate moyenne est un mélange :
\begin{align}
E_B &= c_B\, Y^d_B + c^V_B V_B, \qquad c_h = c^0_h\big(1 - \psi_h\,(1 - cov_t)\big), \quad cov_t = \clip{\frac{Tr_t + G^{san}_t}{\overline{cov}\,P_tY_t}}{0}{1}, \label{eq:cB}\\
E_H &: \quad \frac{E_{H,t+1}}{E_{H,t}} = \Big(\beta\,\frac{1+i^D_t}{1+\pi^e_t}\Big)^{1/\sigma} \ \text{(cible)},\qquad
E_{H,t} \to \text{ajustement partiel vers } \kappa_E\,(Y^{perm}_H + r\, V_H), \label{eq:cH}\\
E_M &= \lambda_M E_M^{(B)} + (1-\lambda_M) E_M^{(H)} . \label{eq:cM}
\end{align}
\begin{lecture}[consommation : deux règles selon la strate]
\variables $E_h$ : dépense de consommation totale de la strate ; $c_h$ : propension marginale à consommer (part d'un euro de revenu supplémentaire qui est dépensée), égale à une valeur de base $c^0_h$ réduite par le motif de précaution $\psi_h$ quand la couverture sociale $cov$ (transferts, retraites et santé publics rapportés à un niveau de référence $\overline{cov}$ du PIB) est incomplète (v0.6) ; $c^V_h$ : propension à consommer la richesse (part de la richesse $V_h$ dépensée chaque année) ; $\beta$ : facteur d'escompte (poids donné au futur ; $\beta=0{,}98$ signifie qu'un euro dans un an vaut 0,98 euro aujourd'hui) ; $i^D$ : taux des dépôts ; $\sigma$ : aversion au déplacement de la consommation dans le temps ; $Y^{perm}_H$ : revenu permanent (revenu moyen attendu sur la durée) ; $r$ : taux réel ; $\kappa_E$ : fraction du revenu permanent et de la richesse consommée ; $\lambda_M$ : poids du comportement « contraint » dans la strate moyenne.
\sens La strate basse (\eqref{eq:cB}) dépense presque tout ce qu'elle gagne : elle n'a pas d'épargne pour lisser et pas d'accès au crédit ; un euro de transfert supplémentaire devient presque un euro de consommation (multiplicateur keynésien élevé). La strate haute (\eqref{eq:cH}) raisonne sur la durée : l'équation d'Euler de Ramsey dit qu'elle fait croître sa consommation plus vite quand le taux d'intérêt réel est élevé (épargner rapporte, donc on consomme moins aujourd'hui et plus demain) et plus lentement quand il est bas. Son niveau de consommation est ancré sur son revenu permanent plutôt que sur son revenu courant : un bonus exceptionnel est largement épargné. La strate moyenne est un mélange. Le terme de couverture dit qu'un ménage sans assurance maladie ni retraite épargne par précaution : c'est ainsi qu'une économie peut avoir une consommation à 39\,\% du PIB. Un transfert agit alors par deux canaux, le revenu et la réduction du motif de précaution.
\hyp La distinction entre ménages contraints et non contraints (Galí, López-Salido et Vallés 2007) ; anticipations de revenu adaptatives.
\limites Un pays à couverture sociale maximale dont le taux d'épargne monte quand même relève d'un mécanisme distinct et mal identifié empiriquement (incertitude sur la fiscalité future, retraite privée, richesse anticipée) que le modèle ne représente pas (v0.7). L'équation d'Euler suppose un consommateur rationnel et patient, ce que l'économie comportementale nuance (myopie, comptabilité mentale) ; on la garde parce qu'elle est le canal par lequel le taux d'intérêt agit sur la consommation, indispensable pour que la politique monétaire ait un effet.
\end{lecture}
Ici $V_h$ est la richesse nette de la strate, $Y^{perm}$ le revenu permanent (moyenne mobile exponentielle), $\sigma$ l'inverse de l'élasticité de substitution intertemporelle et $\beta$ le facteur d'escompte. L'équation \eqref{eq:cH} est le canal par lequel une hausse du taux réel réduit la consommation des ménages non contraints (courbe IS) ; \eqref{eq:cB} donne un multiplicateur budgétaire élevé sur les transferts aux ménages à bas revenus. Le reste, $S_h = Y^d_h - E_h$, est épargné.

\begin{lit}[Ramsey dans un jeu]
Le modèle de Ramsey ne peut pas être résolu par programmation dynamique à chaque tick pour chaque strate ; on en retient (i) l'équation d'Euler comme \emph{cible} de croissance de la consommation vers laquelle l'agent s'ajuste, (ii) sa condition d'état stationnaire $r^\ast = \rho + \sigma g$, utilisée comme taux d'intérêt naturel dans la règle de Taylor de l'IA et comme ancre du taux long. La coexistence de ménages ricardiens et de ménages contraints (Galí, López-Salido, Vallés 2007) est ce qui rend la politique budgétaire non neutre.
\end{lit}

\subsection{Système de demande de Stone--Geary (besoins hiérarchisés)}
La dépense $E_h$ est répartie entre biens selon le système linéaire de dépenses (LES) :
\begin{equation}
x_{hj} = \gamma_{hj} + \frac{\beta_{hj}}{p_j}\Big(E_h - \sum_k p_k \gamma_{hk}\Big), \qquad \sum_j \beta_{hj} = 1,
\label{eq:les}
\end{equation}
\begin{lecture}[répartition de la dépense : système de Stone--Geary]
\variables $x_{hj}$ : quantité du bien $j$ achetée par un ménage de la strate $h$ ; $\gamma_{hj}\ge0$ : quantité de subsistance, c'est-à-dire achetée en priorité quoi qu'il arrive (nourriture de base, chauffage) ; $\beta_{hj}$ : part du revenu « libre » (au-dessus de la subsistance) consacrée au bien $j$, les $\beta$ sommant à 1 ; $E_h - \sum_k p_k\gamma_{hk}$ : revenu libre, ce qui reste une fois la subsistance payée.
\sens Le ménage paie d'abord ses besoins de base, puis répartit le reste en proportions fixes (Stone 1954). Conséquences : quand on est pauvre, la nourriture absorbe presque tout le budget (loi d'Engel) ; quand le revenu augmente, la part des services et biens de confort croît ; quand le prix de la nourriture monte, le ménage réduit tout le reste pour la payer. C'est exactement la hiérarchie des besoins de Victoria~3 rendue continue.
\hyp Les besoins de subsistance sont fixes ; pas de substitution entre biens de base (si le pain double de prix, on n'achète pas deux fois plus de riz).
\limites La substitution entre biens est réelle à moyen terme ; un système plus riche (AIDS de Deaton et Muellbauer 1980) la permettrait, au prix de la simplicité.
\end{lecture}
où $\gamma_{hj} \ge 0$ est la consommation de subsistance (incompressible) du bien $j$ et $\beta_{hj}$ la propension marginale à dépenser le \og supernuméraire\fg{} dans $j$. Ce système est l'équivalent continu de la hiérarchie des besoins de Victoria~3 : quand $E_h$ est proche de $\sum_k p_k\gamma_{hk}$, la consommation est presque entièrement de subsistance (alimentation, énergie) ; quand $E_h$ croît, la part des biens de consommation et des services augmente. Si $E_h < \sum_k p_k \gamma_{hk}$, les quantités de subsistance sont réduites au prorata et le niveau de vie chute brutalement (famine si l'alimentation passe sous un seuil).

\subsection{Niveau de vie, indice des prix, inégalités}
\begin{align}
SoL_h &= \exp\Big(\sum_j \beta_{hj} \ln\big(x^{obt}_{hj} - \gamma_{hj} + \epsilon\big)\Big)\cdot(1 - \xi\, u_h)\cdot(1-\zeta_q\, q^{file}_h)\cdot\Big(\frac{g^{obt}_t}{g^{obt}_0}\Big)^{\omega_G}, \label{eq:sol}\\
P_t &= \sum_j w_j\, p_{j,t}, \qquad \pi_t = \frac{P_t}{P_{t-12}} - 1 \ \text{(glissement annuel)}, \label{eq:cpi}\\
\text{Gini} &= f\big(\{N_h, Y^d_h\}\big) \quad\text{(formule usuelle sur trois classes)}. 
\end{align}
\begin{lecture}[niveau de vie, indice des prix, inégalités]
\variables $SoL_h$ : niveau de vie (\emph{Standard of Living}) de la strate ; $x^{obt}_{hj}$ : quantité effectivement obtenue (après rationnement éventuel) ; $\epsilon$ : petite constante évitant $\ln 0$ ; $u_h$ : taux de chômage de la strate, et $\xi$ la perte de bien-être par point de chômage ; $q^{file}_h = \sum_j w_{hj}\,\pos{D_j/S_j-1}\,\ind_{\{\varrho^p_j=1\}}$ : indice de file d'attente, excès de demande pondéré par le panier de la strate sur les seuls biens à prix administrés (Kornai 1980), et $\zeta_q=0{,}1$ son poids (v1.5 : spécifié, non simulé avant l'étape~5, le prototype n'ayant pas de prix administrés) ; $P_t$ : indice des prix à la consommation ; $w_j$ : poids du bien $j$ dans le panier (Laspeyres : poids de l'année de base) ; $\pi_t$ : inflation en glissement annuel.
\sens Le niveau de vie est l'utilité de Stone--Geary : il mesure ce que le ménage consomme \emph{au-delà} de sa subsistance, en pondérant chaque bien par son importance ; il est très sensible quand on est proche de la subsistance et peu quand on est aisé (logarithme). Il est réduit par le chômage (au-delà de la seule perte de revenu) et par les files d'attente. L'indice des prix est la moyenne pondérée des prix ; l'inflation est sa variation sur douze mois.
\hyp Les poids du panier sont fixes (pas d'effet de substitution dans l'indice).
\limites Un indice de Laspeyres surestime l'inflation vécue quand les consommateurs se détournent des biens devenus chers ; l'écart est faible à l'horizon du jeu. Le coefficient de Gini est calculé sur trois classes, ce qui est grossier mais suffisant pour le pilotage.
\paragraph{Précision de la v1.2.} $g^{obt}_t$ : consommation publique \emph{effective} par tête (achats et investissement publics réellement exécutés, après rationnement $\varrho^G$ et après pénurie), rapportée à son niveau initial, avec une élasticité $\omega_G=0{,}15$. Sans ce terme, l'audit trouvait une méta-stratégie : supprimer les impôts et laisser la contrainte de trésorerie rationner les dépenses publiques gagnait 11\,\% de niveau de vie, les services publics disparus ne comptant pas ; avec, le gain tombe à 0--2,5\,\% pour 6\,\% d'inflation et une dette de 1,1 PIB.
\end{lecture}
$x^{obt}$ est la quantité effectivement obtenue après rationnement ; $q^{file}_h$ un indice de temps passé en file d'attente en régime de prix administrés (\S\ref{sec:systemes}) ; $w_j$ les poids de Laspeyres. Le niveau de vie moyen $\overline{SoL} = \sum_h N_h SoL_h / N$ et son taux de croissance sont les principaux indicateurs de score.

\subsection{Portefeuille et richesse}
L'épargne des strates $M$ et $H$ est répartie entre dépôts, obligations d'État et actions selon un modèle de portefeuille à la Tobin :
\begin{equation}
s^{B}_h = a_0 + a_1\,(i^B - i^D) - a_2\,(r^E - i^D) - a_3\,\rho^{def},\quad
s^{E}_h = b_0 + b_1\,(r^E - i^D) - b_2\,(i^B - i^D),\quad s^D_h = 1 - s^B_h - s^E_h,
\label{eq:portfolio}
\end{equation}
\begin{lecture}[choix de portefeuille]
\variables $s^B_h, s^E_h, s^D_h$ : parts de l'épargne de la strate placées respectivement en obligations d'État, en actions et en dépôts ; $i^B$ : taux des obligations ; $i^D$ : taux des dépôts ; $r^E$ : rendement attendu des actions (dividendes plus plus-value attendue) ; $\rho^{def}$ : probabilité de défaut de l'État perçue par les épargnants ; $a_k, b_k$ : sensibilités ; $V^E$ : valeur de marché de l'ensemble du capital privé coté ; $r^L+\rho^E-g^e$ : taux d'actualisation des dividendes (taux réel, plus prime de risque des actions, moins croissance attendue).
\sens Les épargnants comparent les rendements : si les obligations rapportent plus que les dépôts, ils en achètent davantage ; si les actions promettent plus, ils délaissent les obligations. Personne ne met tout au même endroit (diversification), d'où des parts qui bougent graduellement (Tobin 1969, Godley--Lavoie 2007). La valeur des actions est la somme actualisée des dividendes futurs : elle baisse quand les taux montent (les dividendes futurs valent moins aujourd'hui) et monte quand la croissance attendue augmente. La part en devises étrangères apparaît quand la monnaie nationale inspire la méfiance.
\hyp Les épargnants réagissent aux écarts de rendement avec des sensibilités constantes ; pas d'optimisation moyenne-variance explicite.
\limites Les parts peuvent en théorie sortir de $[0,1]$ ; on les borne. Les anticipations de plus-value sont décrites en \S\ref{sec:actifs}.
\end{lecture}
où $r^E$ est le rendement attendu des actions, $\rho^{def}$ la probabilité de défaut perçue de l'État. La valeur du capital coté suit un modèle de dividendes actualisés : $V^E = \sum_j (1-\theta_j)\, Div_j / (r^L + \rho^E - g^e)$ ; une hausse des taux déprime la valeur des actifs et, via $c^V_h$, la consommation (effet richesse). Une part désirée $s^{FX,\,d\acute{e}sir\acute{e}}_h = a_4\,\pi^e + a_5\,\pi^{dev}$ est détenue en devises étrangères lorsque la monnaie nationale est menacée (dollarisation), si $\chi^K < 1$ ; la part effective obéit à l'hystérèse de l'équation \eqref{eq:sfx}. La strate basse détient uniquement des dépôts et peut être endettée à la consommation. Les prix $p^E$ et $p^Z$ sont déterminés en \S\ref{sec:actifs}.

\subsection{Offre de travail, chômage, démographie}
\label{sec:demog}
\begin{equation}
L^S_t = \varsigma\, N_t, \qquad u_t = 1 - \frac{\sum_j L_{j,t} + L^G_t}{L^S_t}, \qquad
\frac{N_{h,t+1}}{N_{h,t}} = 1 + n_0 + n_1 \ln SoL_h - n_2 (\ln SoL_h)^2 + m_h ,
\end{equation}
\begin{lecture}[offre de travail, chômage, démographie]
\variables $L^S_t$ : population active ; $\varsigma$ : taux de participation (part de la population qui travaille ou cherche un emploi) ; $N_t$ : population ; $L^G_t$ : emploi public ; $u_t$ : taux de chômage ; $n_0, n_1, n_2$ : paramètres de la croissance démographique ; $m_h$ : solde migratoire de la strate.
\sens Le chômage est la part des actifs sans emploi. La population croît à un rythme qui augmente avec le niveau de vie quand celui-ci est bas (moins de mortalité), puis diminue quand il devient élevé (moins de naissances) : la courbe en cloche est la \emph{transition démographique}. Les migrations vont des pays à bas niveau de vie vers les pays ouverts à haut niveau de vie.
\hyp Taux de participation fixe à court terme (il évolue avec la structure d'âge, \S\ref{sec:hetero}) ; la transition démographique dépend du seul niveau de vie.
\limites L'éducation des femmes, l'urbanisation et les politiques familiales jouent aussi ; ils sont ici absorbés par $SoL$ et par la migration rurale--urbaine.
\paragraph{Forme opérationnelle (v1.3).} Trois cohortes (jeunes, actifs, âgés) : naissances $1{,}2\,\%$ des actifs par an, passage à l'activité à vingt ans, retraite après quarante-cinq ans d'activité, espérance de retraite dix-huit ans ; retraites à 50\,\% du salaire moyen, minima aux jeunes ; la participation est rapportée à la cohorte active. Test (quarante ans) : part des âgés 15 $\to$ 25\,\%, \textbf{dette 1,30 contre 0,38} sans vieillissement, chômage 2,3 contre 3,7\,\% (rareté du travail), croissance 1,0 contre 1,9\,\%, niveau de vie $-12{,}6\,\%$. Le vieillissement est d'abord un problème budgétaire, ensuite un problème de croissance.
\end{lecture}
avec $L^G$ l'emploi public et $m_h$ la migration nette (attirée par l'écart de $SoL$ avec les pays voisins ouverts). La forme concave en $SoL$ reproduit la transition démographique. Les ménages changent de strate par mobilité sociale (fonction de l'éducation) et par mobilité descendante (chômage de longue durée).

\paragraph{Chômage naturel endogène.} Le taux de chômage d'équilibre n'est pas fixe : un chômage élevé et durable dégrade les qualifications et détache une partie des chômeurs du marché (hystérèse), tandis qu'un marché du travail tendu les y ramène :
\begin{equation}
u^n_{t+1} = \max\!\Big(u^n_t + \eta_u\,(u_t - u^n_t) + \eta'_u\,(u^n_0 - u^n_t),\ \underline u^n\Big), \qquad \underline u^n = 3\,\%,
\label{eq:un}
\end{equation}
\begin{lecture}[hystérèse du chômage]
\variables $u^n$ : taux de chômage naturel qui entre dans la courbe de Phillips \eqref{eq:wage} ; $u$ : chômage observé ; $\eta_u$ : vitesse d'hystérèse ; $u^n_0$ : niveau structurel de long terme (institutions du marché du travail) ; $\eta'_u$ : vitesse de retour vers ce niveau.
\sens Chaque année passée avec un chômage supérieur au naturel relève un peu ce naturel (Blanchard et Summers 1986) : après une récession longue, la courbe de Phillips se déplace, et l'inflation redémarre à un chômage plus élevé qu'avant. Inversement, une économie maintenue en tension (point de Lewis, plein emploi durable) abaisse $u^n$. Une désinflation tardive coûte donc plus cher qu'une désinflation rapide : c'est l'argument, absent de la v0.5, en faveur de l'action précoce.
\hyp Hystérèse symétrique et linéaire ; retour lent vers un niveau institutionnel.
\limites Le niveau $u^n_0$ dépend des institutions (indemnisation, protection de l'emploi), que le modèle représente par $\varrho$ et le salaire minimum sans les relier explicitement à $u^n_0$.
\paragraph{Précision de la v1.2.} Plancher frictionnel $\underline u^n$ : l'hystérèse peut abaisser le chômage naturel, pas effacer le chômage de rotation (recherche d'emploi, appariement). Sans ce plancher, une surchauffe permanente ramenait $u^n$ au niveau du chômage observé (1,8\,\%), le terme $\phi_u(u^n-u)$ s'annulait et l'inflation cessait d'accélérer : la courbe de Phillips perdait sa verticalité de long terme.
\end{lecture}

%=====================================================================
\section{Marchés des biens et commerce international}
\label{sec:marches}
%=====================================================================

\subsection{Agrégation des ordres}
Pour chaque bien $j$ et chaque pays :
\begin{equation}
D_j = \underbrace{\sum_h N_h x_{hj}}_{\text{ménages}} + \underbrace{\sum_k X^{dem}_{jk}}_{\text{intrants}} + \underbrace{G_j}_{\text{État}} + \underbrace{I^{(j)}}_{\text{investissement}} + \underbrace{EX_j}_{\text{exportations}},
\qquad
S_j = Y_j + \lambda_S S^{inv}_j + IM_j .
\end{equation}
\begin{lecture}[ordres d'achat et de vente]
\variables $D_j$ : demande totale (ordres d'achat) adressée au bien $j$ ; $N_h x_{hj}$ : consommation des ménages ; $X^{dem}_{jk}$ : commandes d'intrants passées par les autres secteurs ; $G_j$ : achats publics ; $I^{(j)}$ : biens d'équipement et de construction achetés pour investir ; $EX_j$ : exportations ; $S_j$ : offre totale (ordres de vente) ; $\lambda_S S^{inv}_j$ : part des stocks mise en vente ; $IM_j$ : importations.
\sens La demande d'un bien vient de cinq sources et l'offre de trois. C'est la structure du marché de Victoria~3, où chaque acheteur et chaque vendeur passe un ordre au marché national plutôt que de traiter directement. La comparaison des deux totaux pilote le prix (équation suivante).
\hyp Un seul marché par pays et par bien (pas de marchés régionaux) ; les biens d'un même secteur sont homogènes.
\limites Pas de coûts de transport intérieurs ni de différenciation par la qualité.
\end{lecture}

\subsection{Formation des prix (mode marché)}
\begin{equation}
\begin{split}
z_j &= \frac{D_j - S_j}{D_j + S_j} \in [-1,1], \qquad f_t = 1 + 5\pos{\pi^e_t},\\
\frac{p_{j,t+1}}{p_{j,t}} &= \clip{1 + \tau_j + \pos{1+z_j}\,\varpi\,\frac{\pi^e_t}{52} + \mu_c\,\clip{\frac{c^{u}_{j,t}(1+m_j)}{p_{j,t}} - 1}{-\tfrac12}{\tfrac12}}{0{,}7}{1{,}19},\\
\tau_j &= \clip{\kappa_j f_t\, z_j}{-\bar\kappa f_t}{\bar\kappa f_t} \ \text{(tâtonnement)},\qquad m_j = \frac{(r^\ast+\rho_E)\,p_KK_j}{c^u_jY_j},
\end{split}
\label{eq:price}
\end{equation}
\begin{lecture}[formation des prix par tâtonnement]
\variables $z_j\in[-1,1]$ : indicateur de déséquilibre du marché ($z=0$ : offre égale demande ; $z>0$ : pénurie ; $z<0$ : excédent) ; $\kappa_j$ : vitesse de réaction du prix au déséquilibre ; $\bar\kappa$ : variation maximale du prix en une semaine ; $\varpi$ : degré d'indexation des prix sur l'inflation anticipée $\pi^e$ ; $c^u_{j,t}$ : coût unitaire de production ; $\mu_c$ : vitesse à laquelle un prix inférieur au coût est relevé ; $Q_j$ : quantité effectivement échangée.
\sens Le prix monte quand la demande dépasse l'offre et baisse dans le cas inverse : c'est le tâtonnement de Walras (1874), la loi de l'offre et de la demande rendue dynamique. Deux ajouts le rendent réaliste : les vendeurs répercutent l'inflation qu'ils anticipent (un commerçant qui s'attend à 10\,\% d'inflation relève ses étiquettes même sans pénurie), et personne ne vend durablement sous son coût. Contrairement à Victoria~3, où le prix est simplement $p^{base}(1+0{,}75z)$ et ne peut jamais dépasser 175\,\% de sa base, ici le prix n'a pas de plafond : une pénurie durable produit une inflation cumulative, et l'ensemble des prix peut dériver, ce qui est indispensable pour modéliser l'inflation. Quand l'offre est insuffisante, chaque acheteur n'obtient qu'une fraction de sa commande : c'est le rationnement, qui fait passer le déséquilibre par les quantités et non seulement par le prix.
\hyp Ajustement progressif (pas d'équilibre instantané) ; règle identique pour tous les biens sauf par $\kappa_j$.
\limites Le tâtonnement peut osciller si $\kappa_j$ est grand ; la borne $\bar\kappa$ et les stocks amortissent. Les prix réels sont souvent fixés par des entreprises en concurrence imparfaite (marge sur coût) plutôt que par un marché anonyme ; le terme en $\mu_c$ est une concession à cette réalité.
\paragraph{Ajouts de la v0.9.} $m_j$ : marge normale du secteur, telle que le prix couvre le coût unitaire et rémunère le capital au taux sûr plus la prime de risque ; le rappel $\mu_c$ est désormais \emph{symétrique} : un prix supérieur à la marge normale baisse, ce qui ancre le partage de la valeur ajoutée (part salariale $\approx1-\alpha$). $\pos{1+z}$ : un vendeur en excès d'offre n'indexe pas ses prix sur l'inflation anticipée. $f_t$ : les prix deviennent plus flexibles quand l'inflation anticipée est forte (les coûts de menu deviennent négligeables), ce qui permet à une contraction monétaire de faire baisser les prix pendant une stabilisation. La borne $[0{,}7\,;1{,}19]$ par semaine rend l'hyperinflation numériquement finie (environ $+100\,\%$ par mois au maximum). Ce que le prototype a montré : sans rappel symétrique, la part salariale dérive et l'investissement absorbe la moitié du PIB ; sans $\pos{1+z}$ et $f_t$, l'inflation se poursuit par pure indexation quand la monnaie ne croît plus.
\end{lecture}
où $c^u_j$ est le coût unitaire de production. Le premier terme est le tâtonnement walrasien ; c'est la version dynamique de la règle de Victoria~3 ($p = p^{base}(1+0{,}75 z)$), qui n'est que sa forme statique. Le terme $\varpi \pi^e$ (indexation partielle) et le rappel vers le coût unitaire (marge normale) rendent l'inflation persistante et cohérente avec les coûts : c'est ce qui permet une dynamique de prix sans ancre nominale fixe. Un terme de déséquilibre monétaire, défini en \S\ref{sec:qe} (éq.~\eqref{eq:monterm}), s'ajoute à moyen terme. La borne $\bar\kappa$ (par exemple 5\,\% par semaine) lisse les ajustements comme la borne $\pm75\,\%$ de Victoria~3, mais sans plafond absolu : une pénurie persistante finit par produire une inflation cumulative.

\paragraph{Rationnement.} La quantité échangée est $Q_j = \min(D_j, S_j)$. Si $D_j > S_j$, chaque acheteur obtient la fraction $S_j/D_j$ de sa commande (priorité possible aux intrants ou à l'État en mode planifié) ; si $S_j > D_j$, l'invendu va en stock, $S^{inv}_{j,t+1} = S^{inv}_{j,t} + (S_j - Q_j)$, ce qui pèse sur la production visée au tick suivant.

\begin{design}[rationner ou non ?]
Victoria~3 sert tous les ordres et fait passer le déséquilibre uniquement par le prix, ce qui évite les pénuries physiques mais casse la cohérence stock-flux (des biens apparaissent ou disparaissent). Nous rationnons : une pénurie se traduit à la fois par un prix plus élevé et par des quantités non obtenues (perte de niveau de vie, baisse du \emph{throughput} aval). Cela est indispensable pour que l'économie planifiée à prix administrés se comporte de manière réaliste (pénuries, files, marché noir) et pour que les stocks aient un rôle.
\end{design}

\subsection{Commerce international (Armington)}
Un acheteur du pays $i$ traite le bien $j$ produit dans le pays $k$ comme une variété imparfaitement substituable. Le prix rendu est
\begin{equation}
\tilde p_{ikj} = e_{ik}\, p_{kj}\,(1+\tau^M_{ij})(1+\tau^{tr}_{ik}),
\end{equation}
\begin{lecture}[prix rendu d'un bien importé]
\variables $\tilde p_{ikj}$ : prix payé par un acheteur du pays $i$ pour le bien $j$ produit dans le pays $k$ ; $e_{ik}$ : taux de change, nombre d'unités de monnaie du pays $i$ nécessaires pour acheter une unité de monnaie du pays $k$ (une hausse de $e_{ik}$ est une \emph{dépréciation} de la monnaie de $i$) ; $p_{kj}$ : prix du bien dans le pays producteur, en monnaie de $k$ ; $\tau^M_{ij}$ : tarif douanier appliqué par $i$ ; $\tau^{tr}_{ik}$ : coût de transport entre les deux pays.
\sens Le prix d'un bien importé combine son prix d'origine, le taux de change, la douane et le transport. Une dépréciation de la monnaie nationale rend toutes les importations plus chères ; un tarif fait de même pour les seules importations visées.
\hyp Transport et douane sont proportionnels au prix ; pas de délai de livraison.
\limites Les exportateurs réels ajustent leurs prix en monnaie de destination (\emph{pricing to market}), ce qui atténue la transmission du change ; le modèle suppose une transmission complète.
\end{lecture}
où $e_{ik}$ est le taux de change (unités de monnaie $i$ par unité de monnaie $k$), $\tau^M$ le tarif douanier fixé par le joueur et $\tau^{tr}$ le coût de transport. La part de la demande de $i$ adressée à la variété $k$ est de type CES :
\begin{equation}
\varsigma_{ikj} = \frac{\varphi_{ik}\,\tilde p_{ikj}^{\,1-\epsilon_j}}{\sum_{k'} \varphi_{ik'}\,\tilde p_{ik'j}^{\,1-\epsilon_j}}, \qquad
IM_{ikj} = \varsigma_{ikj}\, D_{ij}, \quad EX_{kj} = \sum_{i\neq k} IM_{ikj},
\label{eq:armington}
\end{equation}
\begin{lecture}[commerce international : le modèle d'Armington]
\variables $\varsigma_{ikj}$ : part de la demande du pays $i$ pour le bien $j$ qui s'adresse au producteur $k$ ; $\varphi_{ik}$ : préférence pour la variété $k$ (le biais domestique $\varphi_{ii}>\varphi_{ik}$ traduit que l'on achète d'abord local, par habitude, proximité ou réglementation) ; $\epsilon_j>1$ : élasticité de substitution, c'est-à-dire la facilité avec laquelle les acheteurs changent de fournisseur quand les prix relatifs bougent ; $IM_{ikj}$ : importations de $i$ en provenance de $k$ ; $EX_{kj}$ : exportations de $k$.
\sens Le même bien produit dans deux pays n'est pas parfaitement identique aux yeux des acheteurs (Armington 1969) : le pays le moins cher gagne des parts de marché, mais ne rafle pas tout. Plus $\epsilon_j$ est élevé, plus la sensibilité au prix est forte (matières premières : très substituables ; machines spécialisées : peu). C'est ce qui permet à un pays d'être en déficit commercial durable sans que son commerce s'effondre, et à une dévaluation de gagner des parts progressivement.
\hyp Préférences constantes ; ajustement graduel des parts (contrats).
\limites Pas de chaînes de valeur mondiales (un bien traverse une seule frontière) ; pas de commerce de services séparé.
\paragraph{Ajout de la v1.0 : ajustement graduel des parts.} Les parts effectives suivent les parts d'équilibre avec un retard, $\varsigma^{eff}_{t+1}=\varsigma^{eff}_t+\lambda_\varsigma(\varsigma_t-\varsigma^{eff}_t)$, $\lambda_\varsigma=0{,}1$ par semaine (contrats, habitudes, délais logistiques). Le prototype a montré qu'avec des parts instantanées et une élasticité de 6 (énergie), les parts oscillaient de 0,45 à 0,83 en trois semaines, les flux commerciaux de $\pm5\,\%$ du PIB par semaine, et le change divergeait : c'est l'hypothèse « ajustement graduel des parts » qui rend le commerce stable, et non l'élasticité elle-même (ramenée à 2 pour l'énergie). Le règlement international est décrit en \S\ref{sec:sfc}.
\paragraph{Ajout de la v1.1 : prix mondial (option).} Pour les biens homogènes — l'énergie et, à terme, les matières premières — l'hypothèse d'Armington est excessive : un baril est un baril. Le modèle peut alors basculer sur un \emph{prix mondial} $p^w_j$, défini comme la moyenne des prix nationaux exprimés dans la monnaie de référence, pondérée par les productions, vers lequel chaque prix national est rappelé par arbitrage : $p_{ij}\leftarrow(1-\lambda_{lop})p_{ij}+\lambda_{lop}\,e_i\,p^w_j$, avec $\lambda_{lop}=0{,}05$ par semaine. Ce que le prototype a montré : une formulation où le prix mondial est déterminé par l'excès de demande mondial \emph{et imposé} aux producteurs diverge numériquement (les quantités nationales ne peuvent pas s'ajuster à un prix qu'elles ne font pas) ; la formulation par consensus et rappel est stable. Effet : le choc pétrolier devient plus net chez l'importateur (inflation de 1,6 à 2,6\,\%, chômage de 4,9 à 6,1\,\%) et les prix relatifs de l'énergie convergent entre pays (0,85 contre 0,41 sans prix mondial). Un rappel rapide ($\lambda_{lop}=0{,}2$) synchronise les cycles nationaux et doit être évité. L'option est désactivée par défaut ; elle sera généralisée aux ressources de la \S\ref{sec:hetero}, dont la rente exige un prix mondial unique et une élasticité de demande faible.
\end{lecture}
avec un biais domestique $\varphi_{ii} > \varphi_{ik}$. Les parts ne s'ajustent que partiellement d'un mois sur l'autre (contrats, habitudes). Les tarifs sont une recette de l'État. Des quotas, embargos et blocs commerciaux (analogues aux \emph{power blocs} et zones douanières de Victoria~3) modifient $\varphi_{ik}$ et $\tau^M$.

\subsection{Balance des paiements}
\label{sec:bop}
\begin{align}
TB &= \sum_j \big(p_j EX_j - \textstyle\sum_k e_{ik} p_{kj} IM_{ikj}\big),\\
CA &= TB + i^\ast\,(\text{créances ext.}) - i^{B\ast}\, B^\ast - (\text{dividendes versés à l'étranger}),\\
KA &= \Delta(\text{entrées de capitaux}) - \Delta(\text{sorties}),\qquad \Delta R^{FX} = CA + KA \ \ (\text{change fixe}) .
\end{align}
\begin{lecture}[balance des paiements]
\variables $TB$ : balance commerciale (exportations moins importations, en monnaie nationale) ; $CA$ : compte courant, égal à la balance commerciale plus les remises des migrants $Rem$ (v0.7) et les revenus nets reçus de l'étranger (intérêts sur créances, moins intérêts sur la dette externe $B^\ast$ au taux $i^{B\ast}$, moins dividendes versés aux actionnaires étrangers) ; $KA$ : compte de capital, entrées nettes de capitaux (prêts et investissements étrangers moins placements des résidents à l'étranger) ; $R^{FX}$ : réserves de change de la banque centrale.
\sens La balance des paiements est une identité comptable : ce qu'un pays achète à l'étranger, il doit le payer soit par ce qu'il lui vend, soit en s'endettant (entrées de capitaux), soit en puisant dans ses réserves. Un déficit courant n'est donc ni bon ni mauvais en soi : il signifie que le pays investit ou consomme plus qu'il ne produit, financé par l'étranger. Il devient dangereux si les capitaux se retirent (\S\ref{sec:change}).
\hyp Comptabilité en monnaie nationale ; les revenus d'investissement sont simplifiés.
\limites Les erreurs et omissions, les transferts courants (remises des migrants) et les investissements directs ne sont pas distingués.
\end{lecture}
En change flottant sans intervention, $e$ s'ajuste pour que $CA + KA \approx 0$ (\S\ref{sec:change}).


%=====================================================================
\section{Hétérogénéité des pays : dotations, structure, ressources}
\label{sec:hetero}
%=====================================================================
Un pays est défini par ses \emph{dotations initiales} (terre, ressources, population et sa structure d'âge, capital, productivité, capital humain, institutions, régime de souveraineté monétaire) et par les équations ci-dessous qui font que ces dotations produisent des trajectoires différentes. Les archétypes (\S\ref{sec:archetypes}) ne sont que des jeux de conditions initiales.

\subsection{Terre et secteur agricole}
L'agriculture utilise un facteur fixe, la terre $T$ (qualité et surface) :
\begin{equation}
Y_{agr} = A_{agr}\, T^{\alpha_T} K_{agr}^{\alpha_K} (H L_{agr})^{1-\alpha_T-\alpha_K}, \qquad
\text{Rente}_{agr} = \alpha_T\, p_{agr} Y_{agr},
\label{eq:agr}
\end{equation}
\begin{lecture}[agriculture et terre]
\variables $Y_{agr}$ : production agricole ; $T$ : terre disponible (surface pondérée par sa fertilité), facteur fixe ; $\alpha_T$ : part de la terre ; $\alpha_K$ : part du capital agricole (machines, irrigation) ; $L_{agr}$ : main-d'œuvre agricole ; Rente : revenu qui revient aux propriétaires du sol.
\sens Ajouter des paysans sur une terre fixe augmente la récolte de moins en moins : c'est la loi des rendements décroissants de Ricardo (1817), à l'origine de la « trappe malthusienne » : quand la population croît plus vite que la terre, le salaire agricole tombe au minimum vital. La rente foncière est la part de la récolte qui échappe au travail et au capital pour aller aux propriétaires ; sa concentration est la source des inégalités des économies agraires et l'objet des réformes agraires.
\hyp Terre homogène ; qualité fixe (pas d'épuisement ni d'amélioration des sols, sauf via $A_{agr}$).
\limites L'irrigation et les engrais (capital et intrants) ont dans l'histoire repoussé la contrainte de la terre bien au-delà de ce que ce modèle simple suggère ; c'est le rôle de $K_{agr}$ et des intrants Leontief.
\paragraph{Forme opérationnelle (v1.3).} Le prototype porte la terre et le climat sur la productivité du secteur alimentaire : $A^{eff}_{agr}=A_{agr}\,\omega_t\,T^{\alpha_T}$, avec $\alpha_T=0{,}25$ et un aléa climatique annuel $\omega_t\sim\mathcal N(1,\sigma_\omega)$ borné à $[0{,}5\,;1{,}3]$. Test : une récolte à $-30\,\%$ porte l'inflation de 2,1 à 6,4\,\%, coûte 12,5\,\% de niveau de vie au creux (la strate basse, dont le panier est alimentaire, porte le choc) et 15\,\% de chômage au pic, avec retour en trois ans ; un archétype agraire ($\sigma_\omega=8\,\%$) a une volatilité de l'inflation cinquante fois celle de l'archétype industriel. C'est le canal alimentaire de la volatilité des pays pauvres, et la justification directe des stocks alimentaires publics comme levier.
\end{lecture}
la rente foncière allant aux propriétaires de la terre (strate $H$ ou $M$ selon la structure de propriété initiale ; l'État si $\theta_{agr}=1$).
Les rendements décroissants du travail sur une terre fixe produisent la pression malthusienne d'un pays agraire très peuplé : la productivité marginale du travail, donc le salaire, tombe vers le coût de subsistance $\sum_k p_k \gamma_{Bk}$. La réforme agraire (redistribution de $T$, donc de la rente) et la mécanisation ($K_{agr}$, intrants industriels) sont les issues classiques ; la première est un levier institutionnel du joueur (\S\ref{sec:systemes}), la seconde exige un secteur d'équipement ou des importations.

\subsection{Économie duale et transformation structurelle (Lewis, Harris--Todaro)}
La main-d'œuvre est répartie entre un secteur rural (agriculture, subsistance) et un secteur urbain (industrie, services). La migration suit l'écart de revenu espéré :
\begin{equation}
\frac{\Delta L_{urb}}{L_{rur}} = \lambda_{mig}\,\Big(\frac{W_{urb}(1-u_{urb})}{W_{rur}} - 1\Big)^{+}, \qquad
W_{rur} = \frac{(1-\alpha_T-\alpha_K)\, p_{agr} Y_{agr}}{L_{rur}} .
\label{eq:migr}
\end{equation}
\begin{lecture}[migration rurale--urbaine : Harris--Todaro et Lewis]
\variables $L_{rur}, L_{urb}$ : main-d'œuvre rurale et urbaine ; $W_{rur}$ : revenu par travailleur rural, égal à la productivité marginale du travail agricole ; $W_{urb}$ : salaire urbain ; $u_{urb}$ : chômage urbain, de sorte que $W_{urb}(1-u_{urb})$ est le salaire urbain \emph{espéré} (salaire multiplié par la probabilité de trouver un emploi) ; $\lambda_{mig}$ : vitesse de migration ; $\varkappa$ : prime que l'industrie doit payer au-dessus du revenu rural pour attirer des travailleurs ; $\bar\ell_{rur}$ : seuil de main-d'œuvre rurale sous lequel le surplus de travail est épuisé.
\sens Les gens quittent la campagne quand le salaire urbain espéré dépasse ce qu'ils gagnent aux champs (Harris et Todaro 1970) : le chômage urbain freine la migration sans l'arrêter. Tant que la campagne regorge de main-d'œuvre sous-employée, l'industrie peut embaucher sans hausse de salaire : la courbe de Phillips est plate (Lewis 1954). Quand le réservoir rural s'épuise, les salaires s'envolent : c'est le \emph{point de Lewis}, atteint par la Chine vers 2005--2010. L'urbanisation nourrit en outre la productivité par les économies d'agglomération.
\hyp Deux secteurs seulement ; migration réversible ; pas de coût de migration.
\limites La migration réelle dépend aussi des réseaux familiaux, du logement urbain et des politiques (hukou chinois) ; ces frictions sont absorbées par $\lambda_{mig}$.
\paragraph{Forme opérationnelle (v1.3).} Réservoir rural $L^{rur}$ hors marché du travail urbain, revenu rural $W^{rur}=w^{rur}\bar W\,(L^{rur}_0/L^{rur})^{0{,}3}$ (croissant avec l'exode : la terre par tête augmente), migration mensuelle $\Delta L^{rur} = -\frac{\mu}{12}\,\big[(1-u)\bar W/W^{rur} - 1 - c^{mig}\big]^+ L^{rur}$ avec $w^{rur}=0{,}5$, $\mu=0{,}05$ par an, $c^{mig}=0{,}2$ ; les migrants rejoignent l'offre de travail urbaine. Test (réservoir de 50\,\% des actifs, quarante ans) : ruraux 299 $\to$ 168 $\to$ 119, chômage urbain 4,7\,\% contre 3,9\,\% tant que le réservoir alimente l'offre, PIB $+37\,\%$ par absorption du réservoir, inflation inchangée : la trajectoire de Lewis émerge sans paramètre ad hoc.
\end{lecture}
Tant qu'il existe un \emph{surplus de travail} rural ($W_{rur}$ proche de la subsistance), le salaire urbain est ancré sur $W_{rur}(1+\varkappa)$ et la courbe de Phillips \eqref{eq:wage} est plate : l'industrie peut croître sans tension salariale (Lewis 1954). Cela se formalise en remplaçant $\phi_u$ par $\phi_u\cdot\ind_{\{L_{rur}/N\,<\,\bar\ell_{rur}\}}$. Une fois le surplus absorbé (\og point de Lewis\fg), la courbe de Phillips redevient pentue et les salaires décollent (Chine après 2005). L'urbanisation nourrit la productivité, mais le modèle ne lui donne pas de terme propre : le facteur d'agglomération $(L_{urb}/N)^{\eta_{agg}}$ mentionné jusqu'à la v1.4 est retiré (v1.5), faute de dérivation dans le cours et parce qu'il ferait double emploi avec le rattrapage $\chi$ de \eqref{eq:tfp} et avec la sortie du réservoir de Lewis, qui portent déjà ce gain (voir \S\ref{sec:ecartees}).

\subsection{Ressources naturelles : pétrole, or, matières premières}
Un secteur de ressources extrait un stock épuisable $\mathcal{R}$ :
\begin{equation}
\begin{split}
Y_{res} &= A^{eff}_{res} K_{res}^{\alpha} L_{res}^{1-\alpha}, \qquad A^{eff}_{res} = A_{res}\max\!\Big(\min\!\Big(\frac{\mathcal{R}_t/\mathcal{R}_0}{\varrho^{plat}},1\Big)^{\!1/2},\ \underline a\Big),\\
\mathcal{R}_{t+1} &= \mathcal{R}_t - Y_{res} + \text{Découvertes}_t\,(G^{explo}),
\end{split}
\end{equation}
\begin{lecture}[ressources naturelles]
\variables $Y_{res}$ : quantité extraite ; $\mathcal{R}_t$ : réserves prouvées restantes ; $\bar x$ : fraction maximale des réserves extractible par période (contrainte géologique) ; Découvertes : nouvelles réserves, fonction des dépenses d'exploration $G^{explo}$ ; $p^w_{res}$ : prix mondial, en monnaie de réserve ; $z^w$ : déséquilibre du marché mondial ; $\kappa^w$ : vitesse d'ajustement du prix mondial ; $\varepsilon^w$ : choc aléatoire persistant ; $c^u_{res}$ : coût unitaire d'extraction ; $\mathcal{S}_{res}$ : rente (recette au prix mondial moins coût) ; $\tau^{roy}$ : redevance prélevée par l'État.
\sens L'extraction est limitée à la fois par les équipements et par la géologie ; le stock s'épuise s'il n'est pas renouvelé par l'exploration. Le pétrole est vendu sur un marché mondial unique où tous les producteurs reçoivent le même prix, très volatil parce que l'offre et la demande sont peu sensibles au prix à court terme (une petite pénurie fait beaucoup monter le prix). La rente est la différence entre ce prix et le coût : elle est « tombée du ciel » et son partage entre propriétaires, État et étranger est un choix politique.
\hyp Bien homogène ; pas de stockage stratégique ; prix formé en monnaie de réserve.
\limites L'OPEP et les contrats de long terme lissent le prix réel ; le modèle permet le cartel comme décision de joueurs (quota), pas les contrats.
\paragraph{Réécriture de la v1.3.} La contrainte de stock n'est plus un plafond de flux mais un \emph{coût d'extraction} : la productivité du secteur reste à son plateau tant qu'il subsiste plus de $\varrho^{plat}=50\,\%$ du stock initial, puis décroît comme la racine du stock restant, avec un plancher $\underline a=40\,\%$ qui représente les substituts (renouvelables, importations, gisements marginaux). Le stock de référence vaut $\mathcal{R}_0=150$ années de production initiale. Ce que le prototype a montré : avec trente ans de réserves et un coût croissant dès la première tonne, l'énergie — intrant Leontief de tous les secteurs — se raréfiait et tuait la production entière ; avec 120 ans de réserves, la production reste sur son plateau quarante ans, le prix relatif de l'énergie double ensuite et le PIB perd 18\,\% à quatre-vingts ans sans effondrement ni inflation. Le plancher des substituts est le paramètre décisif : il fait la différence entre une transition coûteuse et un effondrement.
\end{lecture}
et se vend sur un \textbf{marché mondial unique} (bien homogène, $\epsilon \to \infty$) au prix $p^w_{res}$ exprimé dans la monnaie de réserve, formé par tâtonnement mondial : $p^w_{t+1}/p^w_t = 1 + \kappa^w z^w_t + \varepsilon^w_t$, avec $\varepsilon^w$ un choc AR(1) à forte variance (les prix des matières premières sont quatre à cinq fois plus volatils que ceux des biens manufacturés). La rente $\mathcal{S}_{res} = (e\, p^w_{res} - c^u_{res})\, Y_{res}$ est captée par les propriétaires et par l'État via une redevance $\tau^{roy}$, un impôt sur la rente ou la nationalisation ($\theta_{res}$).

\paragraph{Maladie hollandaise (mécanisme émergent).} Trois équations existantes suffisent : (i) les exportations de ressources gonflent $CA$, donc apprécient $e$ via \eqref{eq:fx} ; (ii) l'appréciation réduit la part Armington \eqref{eq:armington} des manufacturiers domestiques ; (iii) la productivité manufacturière dépend de l'apprentissage par la pratique :
\begin{equation}
\Delta \ln A_{man} \ \ni\ \eta_{LBD}\,\Big(\frac{Y_{man}}{Y} - \overline{s}_{man}\Big),
\label{eq:lbd}
\end{equation}
\begin{lecture}[apprentissage par la pratique et maladie hollandaise]
\variables $A_{man}$ : productivité du secteur manufacturier ; $Y_{man}/Y$ : poids de l'industrie dans le PIB ; $\overline{s}_{man}$ : poids de référence ; $\eta_{LBD}$ : intensité de l'apprentissage par la pratique (\emph{learning by doing}) ; $F_t$ : fonds souverain ; $\varsigma^F$ : part des redevances versée au fonds.
\sens Plus on produit, plus on apprend à produire : la productivité industrielle croît avec la taille de l'industrie (Arrow 1962). Conséquence : une désindustrialisation, même temporaire, laisse une trace permanente sur la productivité. C'est le cœur de la \emph{maladie hollandaise} (Corden et Neary 1982, Krugman 1987) : une manne pétrolière apprécie la monnaie, l'industrie devient non compétitive et se contracte, et quand le pétrole s'épuise ou baisse, le pays se retrouve sans industrie compétitive. Le fonds souverain, en plaçant la rente à l'étranger, neutralise l'appréciation.
\hyp Externalité d'apprentissage propre à l'industrie (pas aux services ni aux ressources) ; symétrie entre gain et perte.
\limites Le débat empirique sur la « malédiction des ressources » est ouvert (la Norvège s'en sort, le Venezuela non) ; la différence tient aux institutions, que le modèle représente par les choix du joueur (fonds, redevance, change).
\end{lecture}
\paragraph{Fonds souverain.} Le fonds souverain est investi en actifs risqués mondiaux (70\,\% d'actions en pratique) ; son rendement suit le marché mondial des actions plutôt qu'un taux sans risque :
\begin{equation}
\begin{split}
F_{t+1} &= (1+r^F_t)\,F_t + \varsigma^F\tau^{roy}\mathcal{S}_{res,t} - \text{Retraits}_t, \qquad
r^F_t = \bar r^F + \beta_F\,\Delta\ln p^{E,w}_t + \varepsilon^F_t,\\
F &= \underbrace{(1-\beta^{dom})F}_{\to\ KA,\ \eqref{eq:fx}} + \underbrace{\beta^{dom}F}_{\to\ D^E,\ \eqref{eq:pe}},
\end{split}
\label{eq:fonds}
\end{equation}
où $p^{E,w}$ est l'indice mondial des actions (moyenne des $p^E$ des pays pondérée par leur capitalisation). Une règle de prélèvement plafonné (par exemple retraits $\le 3\,\%$ de $F$ par an) est un curseur institutionnel (\S\ref{sec:systemes}).
\begin{lecture}[fonds souverain risqué]
\variables $F$ : encours du fonds ; $r^F$ : rendement annuel, de moyenne $\bar r^F$ ; $\beta_F$ : sensibilité au marché mondial des actions ($\approx0{,}7$) ; $\varepsilon^F$ : bruit ; $\varsigma^F\tau^{roy}\mathcal{S}_{res}$ : part de la rente versée au fonds ; Retraits : ponction budgétaire annuelle ; $\beta^{dom}\in[0,1]$ : part du fonds placée dans le pays (v0.7) : la part externe $(1-\beta^{dom})F$ est une sortie de capitaux qui neutralise l'appréciation du change ; la part domestique $\beta^{dom}F$ est une demande d'actions et de projets nationaux qui gonfle $p^E$, l'investissement et la demande intérieure.
\sens Le fonds ne supprime pas le risque de la rente, il le \emph{convertit} : à la place du prix du pétrole, le budget dépend du cours mondial des actions. Tant que le fonds est petit, c'est un gain de diversification ; quand il dépasse plusieurs fois le PIB, une année boursière à $-20\,\%$ coûte l'équivalent de plusieurs années de rente et sa volatilité devient le principal risque budgétaire du pays. Deux stérilisations distinctes : placer la rente à l'étranger neutralise le canal du \emph{change} (plus d'appréciation), mais la part de la rente réinjectée par le budget nourrit toujours la \emph{demande} : écart de production positif, salaires et inflation domestiques élevés. On peut stériliser l'un sans l'autre. Le placement (v0.7) est l'arbitrage central du rentier : à $\beta^{dom}=0$ le fonds stérilise le change et convertit la rente en risque boursier mondial ; à $\beta^{dom}=1$ il ne stérilise plus rien, injecte la rente dans la demande domestique (méga-projets, bulle d'actifs locale) et fait du pays un cas de maladie hollandaise avec un fonds. Deux rentiers à $\varsigma^F$ identique divergent entièrement selon ce seul paramètre.
\hyp Rendement du fonds exogène au pays (il est preneur de prix sur les marchés mondiaux) ; décote fixée par les décisions des autres joueurs (embargo) ou par la géographie.
\limites Un fonds très grand n'est plus preneur de prix ; on l'ignore. La décote ne distingue pas ses causes (logistique ou sanction), ce qui est voulu : le mécanisme économique est le même.
\end{lecture}

de sorte qu'une désindustrialisation temporaire a un coût permanent (Krugman 1987, Matsuyama 1992). La rente attire aussi la main-d'œuvre et le capital hors de l'industrie. Le joueur dispose des contre-mesures historiques : \textbf{fonds souverain} (équation \eqref{eq:fonds} ci-dessous), placé à l'étranger (donc neutralisant l'effet sur $CA$ et $e$), stérilisation par la banque centrale (accumulation de $R^{FX}$), ou subventions à l'industrie.

\paragraph{Or et étalon-or.} Sous étalon-or, la banque centrale achète l'or produit au prix fixe $1/\bar g$ : $\Delta G = Y_{or}$, donc $\Delta H = Y_{or}/\bar g$ (frappe de monnaie, comme le \emph{minting} de Victoria~3). Un grand producteur d'or sous étalon-or connaît une expansion monétaire et une inflation \og importée\fg{} par sa propre production (mécanisme prix--flux d'espèces de Hume) ; pour les autres pays, l'or est une ressource ordinaire vendue au prix mondial et un actif de réserve.

\subsection{Population : taille, structure d'âge, marché intérieur}
La population est divisée en trois cohortes (jeunes, actifs, âgés). La fécondité dépend du niveau de vie et de l'éducation ; la mortalité de $G^{san}$ et de $SoL$. En découlent :
\begin{itemize}
\item le taux de participation $\varsigma_t = N^{act}_t / N_t$ (dividende démographique quand la cohorte active gonfle ; charge des retraites $Tr^{ret} = \varrho^{ret}\bar W N^{ag}$ quand elle vieillit) ;
\item la taille du marché intérieur : le biais domestique Armington croît avec la population, $\varphi_{ii} = \varphi_0 N^{\eta_N}$, ce qui rend un grand pays moins dépendant du commerce et moins sensible au change ;
\item un poids sur les prix mondiaux : la demande d'un pays très peuplé déplace $p^w_{res}$ et les $z_j$ étrangers (la Chine dans le modèle est un preneur de prix qui n'en est pas un).
\end{itemize}

\subsection{Régime de souveraineté monétaire}
\label{sec:souv}
Les sections précédentes supposent qu'un pays émet la monnaie qui circule sur son territoire et en fixe librement le prix. C'est une dotation initiale au même titre que la terre ou les ressources, et elle n'est pas universelle : un pays peut avoir délégué sa monnaie à une union, l'avoir perdue au profit d'une devise étrangère qui circule, ou l'émettre sans pouvoir s'en servir parce que sa banque centrale est saturée. Le régime de souveraineté monétaire est donc une dimension explicite de l'hétérogénéité des pays. Il se réduit à une \emph{table de disponibilité des leviers} et à quelques désactivations d'équations.

\begin{center}\footnotesize
\begin{tabular}{p{2.0cm}p{2.4cm}p{2.6cm}p{2.2cm}p{2.4cm}p{2.5cm}}
\toprule
\textbf{Levier / équation} & \textbf{A. Émetteur souverain} & \textbf{B. Membre d'une union} & \textbf{C. Petit pays à monnaie propre} & \textbf{D. Partiellement dollarisé} & \textbf{E. Ancrage unilatéral rigide}\\
\midrule
Taux directeur $i^{CB}$ & libre \eqref{eq:taylor} & $i^{CB} = i^{U}$, fixé par l'union & libre, mais borne zéro vite atteinte & libre, n'agit que sur $(1-s^{FX})$ de l'économie \eqref{eq:reff} & $i^{CB}=i^{ancre}$, importé\\
Change $e$ & \eqref{eq:fx} & $e=1$ intra-union ; $e^{U}$ commun & \eqref{eq:fx} & \eqref{eq:fx}, avec $s^{FX}$ élevé & $\bar e$ défendu seul : contrainte $R^{FX}\ge\underline R = \varpi_E M$ (\S\ref{sec:change})\\
QE $\Delta B^{CB}$ & libre & décidé par l'union, clé de répartition & libre, bilan $E^{CB}$ mordant vite \eqref{eq:ecb} & limité au compartiment national & libre mais inutile (le taux est importé)\\
Avances $\Delta A^G$ & selon loi & interdit & selon loi & selon loi, inflationniste au carré ($M^{nat}$ seul) & impossible sans rompre l'ancrage\\
Contrôle des capitaux $\chi^K$ & libre & interdit intra-union & libre & libre (peu efficace) & libre\\
Prime souveraine \eqref{eq:spread} & $\bar b^{A}$ & $\bar b^{B}$ élevé si backstop, saut \eqref{eq:jump} sinon & $\bar b^{C}$ & $\bar b^{D}$ bas ; dette en devise & $\bar b^{E}$ ; prime de rupture $\pi^{dev}$ croissante quand $R^{FX}/M$ baisse ; saut \eqref{eq:jump}\\
Backstop & non (auto-assurance par $\Delta A^G$) & \textbf{conditionnel} au respect des règles budgétaires supranationales & non & non & non\\
Règles budgétaires & nationales & \textbf{supranationales} (déficit, dette) & nationales & nationales ; contrainte de change & nationales ; contrainte de réserves\\
Seigneuriage $\Sigma$ & \eqref{eq:ecb} & part de $\Sigma^{U}$ selon clé & \eqref{eq:ecb}, petit & réduit à $(1-s^{FX})$ & propre\\
Prêteur en dernier ressort & oui & union, conditionnel & oui, en monnaie nationale seulement & seulement pour $M^{nat}$ & limité par $R^{FX}$\\
\bottomrule
\end{tabular}
\end{center}

\paragraph{Régime B.} Le pays perd trois leviers (taux, change, avances) et gagne une assurance conditionnelle : tant qu'il respecte les règles budgétaires de l'union (déficit et trajectoire de dette), il bénéficie du backstop, qui abaisse sa prime ($-s_8$ dans \eqref{eq:spread}) et annule la probabilité de saut \eqref{eq:jump}. S'il les enfreint, le backstop est suspendu et il redevient un emprunteur en monnaie étrangère : c'est la décision de jeu centrale de ce régime. Le taux $i^{U}$ et le change $e^U$ sont fixés par l'union (une IA appliquant la règle de Taylor sur les agrégats de l'union, ou un joueur \og banque centrale de l'union\fg{} si l'on veut en faire un rôle). La \textbf{mutualisation} $\varsigma^{mut}\in[0,1]$ (v0.7) fixe la part des pertes de la banque centrale de l'union \eqref{eq:ecb} et des pertes de défaut \eqref{eq:default} partagée entre membres selon la clé de répartition ; le reste reste national. À $\varsigma^{mut}=0{,}2$ (valeur par défaut, proche des unions réelles), le backstop est une facilité de liquidité qui écarte le mauvais équilibre de \eqref{eq:jump} mais laisse le coût d'un défaut au pays ; à $\varsigma^{mut}=1$, c'est une assurance complète, et la règle budgétaire supranationale devient la seule protection des membres vertueux contre l'aléa moral.

\paragraph{Régime C.} Tous les leviers existent mais saturent vite : la borne zéro (\S\ref{sec:zlb}) est atteinte dès que le pays subit un afflux de capitaux refuge, et le bilan de la banque centrale (\S\ref{sec:bc}) devient rapidement grand devant le PIB, de sorte que la contrainte de bilan \eqref{eq:ecb} mord avant tout autre.

\paragraph{Régime D.} Une seconde monnaie circule comme moyen de paiement et unité de compte pour une part $s^{FX}$ de l'activité (\S\ref{sec:actifs}). Le taux directeur ne pilote que le compartiment national de la monnaie, $M^{nat}=(1-s^{FX})M$, et le taux réel qui entre dans les décisions d'investissement \eqref{eq:invest} et de consommation \eqref{eq:cH} est une moyenne des deux monnaies :
\begin{equation}
r^{eff}_t = (1-s^{FX}_t)\,\big(i^L_t - \pi^e_t\big) + s^{FX}_t\,\big(i^\ast_t + \rho^{pays}_t - \pi^{e\ast}_t\big), \qquad M^{nat}_t = (1-s^{FX}_t)\,M_t \ \text{dans \eqref{eq:monterm}}.
\label{eq:reff}
\end{equation}
\begin{lecture}[taux effectif en économie dollarisée]
\variables $r^{eff}$ : taux d'intérêt réel qui entre dans les décisions de l'économie ; $s^{FX}\in[0,1]$ : part de l'activité (contrats, prix, crédits) libellée en monnaie étrangère ; $i^L-\pi^e$ : taux réel en monnaie nationale ; $i^\ast+\rho^{pays}-\pi^{e\ast}$ : taux réel auquel les résidents empruntent en devise (taux étranger plus prime pays, déflaté de l'inflation étrangère) ; $M^{nat}$ : compartiment national de la masse monétaire.
\sens Dans une économie dollarisée à 80\,\%, relever le taux directeur ne renchérit que 20\,\% du crédit : la politique monétaire est diluée, et le seigneuriage avec elle. Le pays a de fait importé la politique monétaire de l'émetteur de la devise, sans en avoir le prêteur en dernier ressort. Le terme d'encaisses \eqref{eq:monterm} ne porte plus que sur la monnaie nationale : une émission de monnaie nationale est d'autant plus inflationniste que le compartiment est petit.
\hyp Les deux monnaies coexistent sans coût de conversion ; $s^{FX}$ évolue selon \eqref{eq:sfx}.
\limites En réalité les segments se recouvrent (salaires en monnaie nationale, loyers en dollars) et la dollarisation des dettes crée des désajustements de bilan lors d'une dépréciation, que le modèle capture via la dette externe $B^\ast$ mais pas pour les ménages.
\end{lecture}

\paragraph{Régime E (v0.7).} Un pays ancré unilatéralement (parité fixe depuis des décennies, \emph{currency board}) importe le taux directeur comme un membre d'union, mais sans union : pas de backstop, pas de partage de risque, pas de prêteur en dernier ressort au-delà de ses réserves. C'est la perte d'autonomie du régime B avec la solitude du régime C. Aucune équation nouvelle : c'est la contrainte de réserves de la défense d'un plancher (\S\ref{sec:change}) combinée au taux importé, avec un seuil de réserves exprimé en proportion de la masse monétaire, $\underline R=\varpi_E M$ ($\varpi_E=1$ pour un \emph{currency board} strict, 0,3 pour un ancrage ordinaire). Le joueur n'a que le budget, la réglementation et le fonds souverain ; la crise typique est la rupture d'ancrage par épuisement de $R^{FX}$ après une chute durable du prix de sa ressource.

\subsection{Grille dotations $\times$ souveraineté}
\label{sec:archetypes}
Les archétypes de départ et les régimes de souveraineté se recouvraient ; la v0.7 les fusionne en une seule grille. Une ligne est une dotation réelle (\S\ref{sec:hetero}), une colonne un régime monétaire (\S\ref{sec:souv}) ; une case est une configuration de la grille de validation (\S\ref{sec:configs}), avec sa dynamique attendue. Les cases vides sont possibles mais ne sont pas couvertes par la validation.
\begin{center}\footnotesize
\begin{tabular}{p{2.5cm}p{2.5cm}p{2.5cm}p{2.5cm}p{2.0cm}p{2.0cm}}
\toprule
\textbf{Dotation $\backslash$ Régime} & \textbf{A. Souverain} & \textbf{B. Union} & \textbf{C. Petit, monnaie propre} & \textbf{D. Dollarisé} & \textbf{E. Ancré}\\
\midrule
Industriel mûr ($A$ à la frontière, $H$ élevé, finance profonde) & \#4 dette monétisée en sortie de trappe ; \#3 émetteur de réserve (avec $-s_6$) & \#1 endetté ; \#2 vertueux ($\varrho^{fisc}$) & \#6 refuge à la borne zéro & --- & ---\\
Planifié en rattrapage post-Lewis ($\theta_j$, $\varrho^{plan}$ élevés, $\chi^K=1$) & \#5 sortie de bulle à contrainte lâche & --- & --- & --- & ---\\
Rentier ($Y_{res}/Y>0{,}3$) & \#7 fonds externe ($\beta^{dom}=0$) & --- & (\#7 selon taille) & --- & \#8 ancré, $\beta^{dom}$ élevé\\
Agraire ou émergent pré-Lewis ($L_{rur}/N$ élevé, $T/Y\approx0{,}1$) & --- & --- & \#9 faible capacité fiscale, dette en devises & (\#9 si $s^{FX}$ élevé) ; \#10 post-hyperinflation & ---\\
Producteur d'or, petit pays ouvert & étalon-or : frappe endogène (\S\ref{sec:hetero}) & --- & change dominant dans l'inflation & --- & \emph{currency board}\\
\bottomrule
\end{tabular}
\end{center}
Les dotations elles-mêmes sont définies par les variables du \S\ref{sec:hetero} : population et part rurale, terre, réserves de ressources, productivité relative à la frontière, capital humain, capacité fiscale $T/Y$, part domestique de la dette $\varkappa^{dom}$, et mémoire monétaire $M^{hist}$.

%=====================================================================
\section{Monnaie, banques et inflation}
\label{sec:monnaie}
%=====================================================================

\subsection{Banque centrale}
\label{sec:bc}
\begin{joueur}[politique monétaire]
\begin{itemize}
\item \textbf{Taux directeur} $i^{CB}$ (taux de refinancement des banques ; le taux de rémunération des réserves est $i^{CB}-\Delta$).
\item \textbf{Réserves obligatoires} $\varrho^{res} \in [0,1]$ sur les dépôts.
\item \textbf{Opérations d'open market / assouplissement quantitatif} : achats nets de dette publique $\Delta B^{CB}$ (et éventuellement d'actifs privés).
\item \textbf{Financement monétaire direct} du Trésor (\og planche à billets\fg) : avance $\Delta A^{G}$, interdite ou plafonnée selon la loi d'indépendance choisie.
\item \textbf{Régime de change} : flottant, géré, fixe (ancrage à une devise ou à l'or), monnaie non convertible ; parité $\bar e$ ; interventions $\Delta R^{FX}$.
\item \textbf{Contrôle des capitaux} $\chi^K \in [0,1]$ (0 : libre circulation ; 1 : fermeture).
\item \textbf{Cible d'inflation} annoncée $\pi^\ast$ et \textbf{statut} de la banque centrale (indépendante / dépendante), qui conditionnent la crédibilité.
\item Mode \og pilote automatique\fg{} : règle de Taylor \eqref{eq:taylor}.
\end{itemize}
\end{joueur}

Bilan : $G + e\,R^{FX} + B^{CB} + L^{CB} + A^G = C + Res + E^{CB}$. Une avance directe $\Delta A^G$ augmente la base monétaire $H$ d'autant. Le seigneuriage $\Sigma$ et les fonds propres $E^{CB}$ sont définis par l'équation \eqref{eq:ecb} ci-dessous.

\paragraph{Bilan vivant de la banque centrale.} Le taux de rémunération des réserves $i^{res}$ est distinct du taux directeur ($i^{res}=i^{CB}-\Delta^{res}$). Le portefeuille $B^{CB}$ porte un coupon moyen $\bar c^{CB}$ figé aux dates d'achat, distinct du taux courant ; ses fonds propres évoluent en deux termes et le dividende à l'État est le résultat positif après report des pertes passées :
\begin{equation}
\begin{split}
\Delta E^{CB}_t &= \underbrace{\bar c^{CB}_t B^{CB}_t + i^{CB}_t L^{CB}_t - i^{res}_t \min\!\big(Res_t,\ B^{CB}_t + A^G_t + L^{CB}_t\big)}_{\text{portage}}
 + \underbrace{\Delta e_t\, R^{FX}_t}_{\text{valorisation change}} - \Sigma_t,\\
\Sigma_t &= \pos{\Pi^{CB}_t - \pos{\bar E^{CB} - E^{CB}_t}},
\end{split}
\label{eq:ecb}
\end{equation}
où $\Pi^{CB}$ désigne le résultat avant distribution (portage plus valorisation) et $\bar E^{CB}$ le niveau statutaire de fonds propres : tant que les pertes passées n'ont pas été reconstituées, le dividende est nul (report à nouveau). Une vente de titres réalise la perte $(i^B-\bar c^{CB})$ sur le montant vendu. Contrainte de bilan : $E^{CB}_t/H_t \ge \underline\kappa^{CB}$, seuil politiquement tolérable (négatif : une banque centrale peut fonctionner avec des fonds propres négatifs tant que le public l'accepte) ; en dessous, la banque centrale abandonne l'opération qui creuse la perte (plafond de change, QE) ou obtient une recapitalisation qui accroît $B$.
\begin{lecture}[bilan et fonds propres de la banque centrale]
\variables $E^{CB}$ : fonds propres de la banque centrale ; $\bar c^{CB}$ : coupon moyen du portefeuille d'obligations, figé aux dates d'achat (moyenne pondérée des $i^B$ historiques) ; $i^{res}$ : taux payé sur les réserves des banques, $\Delta^{res}$ l'écart au taux directeur ; $L^{CB}$ : refinancement accordé aux banques ; $\Delta e\,R^{FX}$ : gain ou perte de change sur les réserves (une appréciation de la monnaie nationale, $\Delta e<0$, est une perte) ; $\Pi^{CB}$ : résultat avant distribution ; $\Sigma$ : dividende versé à l'État, nul tant que les pertes passées ne sont pas reconstituées ; $\bar E^{CB}$ : fonds propres statutaires ; $\underline\kappa^{CB}$ : seuil politiquement tolérable de $E^{CB}/H$.
\sens La banque centrale gagne l'écart entre ce que rapportent ses actifs (coupons figés, souvent bas s'ils ont été achetés pendant le QE) et ce qu'elle paie sur ses réserves (au taux courant). Quand le taux directeur monte, ses charges montent immédiatement, ses produits non : le seigneuriage s'effondre exactement au moment où la charge de la dette du Trésor augmente, corrélation perverse observée partout en 2022--2024. Les réserves de change et l'or la font gagner ou perdre au gré du change. Le fonds pour risques généraux lisse : les profits des années de taux bas sont provisionnés, consommés dans la phase haute ; quand il est épuisé, la perte est reportée, le dividende à l'État est suspendu jusqu'à reconstitution, sans recapitalisation. Vendre des titres avant échéance réalise la moins-value : le QE n'est ni gratuit ni réversible sans coût.
\hyp Les billets $C$ sont un passif gratuit : le portage négatif ne porte que sur la fraction du portefeuille financée par des réserves rémunérées, soit 70 à 80\,\% en pratique ; appliquer $(\bar c^{CB}-i^{res})B^{CB}$ à tout le portefeuille surestime la perte d'un facteur 2 à 3 (v0.7). Pas de transfert automatique du Trésor ; le seuil $\underline\kappa^{CB}$ est une variable politique, pas comptable.
\limites Les conventions comptables (valorisation au coût amorti ou au marché) diffèrent entre banques centrales ; le modèle retient le coût amorti avec pertes réalisées à la vente, qui est la convention la plus courante.
\paragraph{Précisions de la v1.2.} Deux règles que l'audit a rendues nécessaires. (i) Le seigneuriage $\Sigma_t$ porte sur le \emph{profit de la période} $\Pi^{CB}_t$, jamais sur le stock de fonds propres : le prototype, qui distribuait chaque semaine une fraction du stock, a créé $+50\,\%$ de dépôts la première année d'une économie à actifs (les fonds propres apparents incluaient une créance de refinancement sur la banque) et l'a conduite à l'effondrement en quarante-cinq ans. (ii) Seules les réserves \emph{adossées} à des actifs porteurs d'intérêt sont rémunérées, et la banque centrale mène des opérations structurelles pour que son portefeuille suive ses réserves ; rémunérer des réserves sans contrepartie est une perte structurelle, c'est-à-dire une monétisation cachée (2,5\,\% du PIB par an dans le prototype avant correction). Quand la dette publique disparaît (excédent structurel), la banque centrale n'a plus de titres à détenir : c'est l'un des motifs pour lesquels l'état de référence du jeu doit comporter une règle budgétaire à cible de dette.
\end{lecture}

Règle de Taylor pour les IA et le mode automatique, avec un taux naturel \emph{estimé} par apprentissage à partir du taux de Ramsey :
\begin{equation}
\begin{split}
i^{CB}_t &= \min\!\Big(\rho_i\, i^{CB}_{t-1} + (1-\rho_i)\,i^{\dagger}_t,\ \bar\imath\Big),\\
i^{\dagger}_t &= \pos{ \hat r^\ast_t + \min(\pi^e_t,\ \pi^s_t+0{,}5) + a_\pi(\pi^s_t - \pi^\ast) + a_y\, \hat y_t + a_e\,(\ln e_t - \ln \bar e)\,\ind_{\text{fixe}} },\\
\hat r^\ast_{t+1} &= \clip{\hat r^\ast_t + \lambda_{r}\,(\pi^s_t - \pi^\ast)}{-0{,}01}{0{,}06}, \qquad \hat r^\ast_0 = \rho+\sigma g_0, \qquad\\
\pi^s_t &= (1-\lambda_s)\pi^s_{t-1} + \lambda_s\cdot 12\,\clip{\ln\frac{P_t}{P_{t-1}}}{-\tfrac12}{\tfrac12},
\end{split}
\label{eq:taylor}
\end{equation}
\begin{lecture}[règle de Taylor et taux naturel]
\variables $i^{CB}_t$ : taux directeur ; $r^\ast_t$ : taux d'intérêt réel \emph{naturel} ou neutre, celui qui ne stimule ni ne freine l'économie ; $\pi^\ast$ : cible d'inflation ; $a_\pi>1$ : réaction à l'écart d'inflation ; $\hat y_t=\ln(Y_t/Y^{pot}_t)$ : écart de production, écart en pourcentage entre la production réelle et la production \emph{potentielle} (celle obtenue à plein emploi des capacités, $u=u^n$) ; $a_y$ : réaction à cet écart ; $a_e$ : réaction au change en régime fixe ; $\rho$ : taux de préférence pour le présent ; $\sigma$ : aversion au déplacement de la consommation ; $g^A$ : croissance de la productivité ; $\pos{\cdot}$ : le taux ne peut pas être négatif (borne zéro).
\sens La règle de Taylor (1993) décrit le comportement d'une banque centrale raisonnable : partir d'un taux « neutre » ($r^\ast+\pi^e$), le relever quand l'inflation dépasse la cible et quand l'économie surchauffe, l'abaisser dans le cas inverse. Le coefficient $a_\pi>1$ est le \emph{principe de Taylor} : il faut relever le taux nominal \emph{plus} que l'inflation pour que le taux réel monte et freine effectivement l'économie. Le taux naturel vient du modèle de Ramsey : à l'équilibre de long terme, le taux réel égale l'impatience des ménages plus la croissance pondérée ($r^\ast=\rho+\sigma g$) ; une économie qui croît vite a un taux naturel plus élevé.
\hyp La banque centrale connaît l'écart de production (en réalité mal mesuré en temps réel) ; la règle est utilisée par les IA et proposée en pilote automatique au joueur.
\limites La règle ne gère ni les bulles ni la stabilité financière ; c'est justement l'espace de décision laissé au joueur. La borne zéro empêche la règle de fonctionner en déflation profonde (\S\ref{sec:zlb}).
\paragraph{Ajouts de la v0.9.} $\hat r^\ast$ : taux naturel \emph{estimé} par la banque centrale, révisé chaque mois dans le sens de l'écart d'inflation ($\lambda_r=0{,}02$) : c'est un terme intégral, l'équivalent de ce que font les banques centrales réelles quand elles révisent leur estimation de $r^\ast$, et c'est ce qui garantit que l'inflation atteint la cible quel que soit le taux naturel effectif de l'économie. $\pi^s$ : inflation récente lissée, mesurée par la variation logarithmique mensuelle bornée puis annualisée (la mesure $(P_t/P_{t-1})^{12}-1$ transforme un mois à $+20\,\%$ en 790\,\%/an et bloquait le taux au plafond après une réforme). Le terme d'anticipations est borné par $\pi^s+50$ pts pour qu'une anticipation héritée d'une hyperinflation ne dicte pas le taux après la stabilisation. $\rho_i=0{,}95$ (lissage mensuel) et $\bar\imath=100\,\%$ (plafond). Ce que le prototype a montré : ancrée sur $r^\ast=\rho+\sigma g=3{,}5\,\%$, la règle maintenait un chômage de 7\,\% et une inflation de 1\,\% pendant trente ans, parce que le taux naturel \emph{effectif} de cette économie fermée est de 1 à 2\,\% (épargne élevée, demande autonome faible) ; avec l'apprentissage, l'inflation revient à 2,1\,\% et la banque centrale estime $\hat r^\ast\approx0{,}7\,\%$. Le lissage à 0,95 et la mesure en log ont fait disparaître un cycle limite d'un an (écart-type du chômage : 1,8 pt $\to$ 0,05 pt).
\paragraph{Précision de la v1.5.} L'initialisation $\hat r^\ast_0=\rho+\sigma g_0$ utilise la \emph{tendance} de productivité $g_0=1{,}5\,\%$, non la croissance réalisée $g$ : avec $\rho=0{,}01$ et $\sigma=1$, elle vaut 2,5\,\%, ce qui est le « taux théorique » de la \S\ref{sec:proto}. La croissance réalisée (1,8 à 1,9\,\%) donnerait 2,9\,\% ; la revue du cours a relevé que l'équation ne disait pas lequel des deux $g$ elle employait. Le choix de $g_0$ est celui du prototype et il est le bon : une banque centrale n'observe pas le taux naturel, elle l'estime à partir de la tendance, et l'apprentissage $\lambda_r$ fait le reste. La réforme \eqref{eq:reform} réinitialise $\hat r^\ast$ de la même façon.
\end{lecture}
où $\hat y = \ln(Y/Y^{pot})$ est l'écart de production, $Y^{pot}$ étant la production à $u = u^n$ et pleine utilisation des intrants.


\begin{figure}[htbp]
\centering
\resizebox{0.98\textwidth}{!}{%
\begin{tikzpicture}[x=1cm,y=1cm]
% colonne 1 : instruments
\node[pol] (fxp) at (0,9)   {Régime de change, $\chi^K$};
\node[pol] (i)   at (0,6.6) {Taux directeur $i^{CB}$\\ (borne $i^{CB}\ge 0$)};
\node[pol] (qe)  at (0,4.2) {QE : achats $\Delta B^{CB}$};
\node[pol] (mf)  at (0,1.6) {Financement monétaire $\Delta A^G$};
% colonne 2 : canaux
\node[var] (e)   at (5.8,9)   {change $e$};
\node[var] (tb)  at (5.8,7.4) {taux bancaires $i^L,i^D$};
\node[var] (ap)  at (5.8,5.8) {prix d'actifs $p^E,p^Z$,\\ rendement $i^B$};
\node[var] (res) at (5.8,4.0) {réserves excéd. $Res^{ex}$\\ (hors $M$)};
\node[var] (M)   at (5.8,1.6) {masse $M$, vitesse $V(\pi^e)$};
% colonne 3 : variables intermédiaires
\node[var] (imp)  at (11.4,9)   {prix importés,\\ exportations nettes};
\node[var] (I)    at (11.4,7.4) {investissement $q_j$,\\ consommation $E_H$};
\node[var] (V)    at (11.4,5.8) {richesse $V_h$, collatéral $\Lambda_j$};
\node[var] (cred) at (11.4,3.0) {crédibilité $c$,\\ anticipations $\pi^e$};
% colonne 4 : résultats
\node[var] (y)  at (16.8,6.6) {écart de production $\hat y$,\\ chômage $u$};
\node[var] (pi) at (16.8,3.0) {prix $p_j$ \eqref{eq:price}, salaires $W$ \eqref{eq:wage}\\ $\Rightarrow$ inflation $\pi$};

% instruments -> canaux
\draw[flux] (fxp) -- (e);
\draw[flux] (i.north east) -- node[lab,pos=0.45,above,sloped]{UIP} (e.south west);
\draw[flux] (i) -- (tb);
\draw[flux] (i.south east) -- (ap.west);
\draw[flux] (qe.north east) -- node[lab,pos=0.62,above,sloped]{$i^B\downarrow$ ($s_7$)} (ap.south west);
\draw[flux] (qe) -- (res);
\draw[flux] (mf) -- (M);
\draw[flux,dashed] (res) -- node[lab,pos=0.5,right]{si crédit : $\Delta L=\Delta D$} (M);
% canaux -> intermédiaires
\draw[flux] (e) -- (imp);
\draw[flux] (tb) -- (I);
\draw[flux] (ap.north east) -- (I.south west);
\draw[flux] (ap) -- (V);
\draw[flux] (M) -- node[lab,pos=0.66,above,sloped]{ancre $\upsilon$} (cred);
\draw[flux] (mf.north east) -- node[lab,pos=0.12,above,sloped]{$-\nu_3$} (cred.west);
% intermédiaires -> résultats
\draw[flux] (imp.east) to[out=0,in=100] (y.north);
\draw[flux] (I) -- (y);
\draw[flux] (V.east) -- (y.west);
\draw[flux] (y) -- (pi);
\draw[flux] (cred) -- (pi);
\node[lab] at (14.1,4.0) {$\pi^e$ dans $W$, $p$};
\draw[flux] (M.south east) .. controls (10.5,0.3) and (14,0.6) .. node[lab,pos=0.62,below]{terme $\xi_M$ \eqref{eq:monterm}} (pi.south west);
% v0.6 : bilan de banque centrale et dette consolidée
\node[var,fill=orange!10] (bcb) at (12.0,-2.8) {bilan BC \eqref{eq:ecb} : $E^{CB}$,\\ coupon figé $\bar c^{CB}$, $i^{res}$, $\Sigma$};
\node[var,fill=orange!10] (dette) at (18.6,-1.4) {charge de la dette $i^{app}B$\\ $\omega=B^{CB}/B\Rightarrow\theta^{cons}$ \eqref{eq:iapp}};
\draw[flux] (qe.west) .. controls (-3.8,2.0) and (-1.0,-5.2) .. node[lab,pos=0.72,below]{achats : coupon figé $\bar c^{CB}$} (bcb.south west);
\draw[flux] (res.east) .. controls (8.7,3.8) and (8.7,-2.0) .. node[lab,pos=0.78,below]{charge $i^{res}\,Res$} (bcb.north);
\draw[flux] (bcb.east) -- node[lab,pos=0.5,below]{$\Sigma$, $\omega$} (dette.south west);
\draw[flux] (pi.south east) -- node[lab,pos=0.5,right]{boule de neige $i^{app}-\pi-g$} (dette.north);
\draw[flux] (cred.south) .. controls (11.4,-6.0) and (-6.0,-6.0) .. node[lab,pos=0.55,below]{dollarisation $\theta_5$} (-6.0,4.0) .. controls (-6.0,12.0) and (0.5,11.8) .. (e.north);
% boucle de réaction
\draw[flux,gray] (y.north) .. controls (16.8,13.0) and (-4.2,13.0) .. node[lab,pos=0.5,above]{règle de Taylor \eqref{eq:taylor} ou décision du joueur} (i.west);
\end{tikzpicture}}
\caption{Transmission de la politique monétaire (v0.6 : ajout du bilan de banque centrale et de la charge de la dette consolidée, en orange). Le taux directeur agit par le crédit, les prix d'actifs et le change ; le QE agit d'abord sur les rendements obligataires et les prix d'actifs et ne devient inflationniste que si les réserves excédentaires se transforment en crédit ou financent des dépenses ; le financement monétaire direct alimente $M$ et détruit la crédibilité, ce qui active la boucle anticipations--vitesse de circulation--change (hyperinflation). En bas : le QE constitue un portefeuille à coupon figé financé par des réserves rémunérées au taux courant, de sorte qu'une hausse du taux directeur fait basculer le résultat de la banque centrale et se transmet immédiatement à la charge de la dette consolidée.}
\label{fig:mon}
\end{figure}

\subsection{Banque commerciale et monnaie endogène}
\label{sec:banques}
Les taux bancaires sont indexés sur le taux directeur, avec une prime de risque propre à l'emprunteur :
\begin{equation}
i^L_j = i^{CB} + m^L + \varrho_1\, \pos{\ell_j - \bar\ell} + \varrho_2\, \hat\rho^{NPL}_t, \qquad
i^D = \pos{i^{CB} - m^D},
\end{equation}
\begin{lecture}[taux bancaires]
\variables $i^L_j$ : taux du crédit au secteur $j$ ; $i^{CB}$ : taux directeur, coût auquel la banque se refinance ; $m^L$ : marge d'intermédiation ; $\varrho_1$ : prime de risque par unité de levier excédentaire $\pos{\ell_j-\bar\ell}$ ; $\hat\rho^{NPL}_t$ : taux récent de prêts non performants (\emph{non-performing loans}, crédits en défaut) et $\varrho_2$ son effet ; $i^D$ : taux des dépôts ; $m^D$ : marge sur les dépôts.
\sens La banque emprunte au taux directeur (ou rémunère les dépôts un peu moins) et prête plus cher, la différence couvrant ses coûts, ses pertes attendues et son profit. Elle facture plus cher un emprunteur très endetté et, après une vague de défauts, tout le monde : la mémoire des pertes rend le crédit cher longtemps après une crise. C'est le \emph{canal du crédit} de la politique monétaire : le taux directeur se transmet aux taux effectifs avec une marge variable.
\hyp Transmission complète et immédiate du taux directeur ; banque représentative unique.
\limites En réalité la transmission est partielle et retardée (taux fixes sur les prêts existants) ; en crise, la marge peut s'ouvrir violemment quel que soit le taux directeur (2008).
\end{lecture}
où $\hat\rho^{NPL}$ est le taux récent de créances douteuses (mémoire du cycle). L'offre de crédit est contrainte par le ratio de fonds propres réglementaire $\kappa^{CAR}$ (levier bancaire maximal) et par la liquidité :
\begin{equation}
L^{max}_t = \frac{E^{Bk}_t}{\kappa^{CAR}}, \qquad
\Lambda^{Bk}_t = \min\!\Big(1,\ \frac{L^{max}_t - L_t}{\sum_j \Delta L^{dem}_{j,t}}\Big)^{+},
\end{equation}
\begin{lecture}[contrainte de fonds propres et rationnement du crédit]
\variables $L^{max}_t$ : encours maximal de crédit ; $E^{Bk}_t$ : fonds propres de la banque (ce qu'elle possède en propre, différence entre ses actifs et ses dettes) ; $\kappa^{CAR}$ : ratio de fonds propres minimal exigé par la réglementation (\emph{capital adequacy ratio}, 8\,\% dans les accords de Bâle) ; $L_t$ : encours actuel ; $\Delta L^{dem}_{j,t}$ : nouveaux crédits demandés ; $\Lambda^{Bk}_t\in[0,1]$ : fraction des demandes que la banque peut satisfaire.
\sens Une banque ne peut prêter qu'un multiple de ses fonds propres ($1/\kappa^{CAR}$, soit 12,5 fois pour 8\,\%). Quand elle atteint ce plafond, elle rationne : elle refuse une partie des demandes même solvables. Ses fonds propres fondent avec les pertes sur crédits ; une crise réduit donc mécaniquement l'offre de crédit, ce qui aggrave la crise (\emph{credit crunch}). Inversement, un relâchement de $\kappa^{CAR}$ par le joueur libère du crédit et gonfle la masse monétaire : c'est le mécanisme de la \emph{monnaie endogène} (Moore 1988) : ce sont les prêts qui créent les dépôts, et non l'inverse.
\hyp Une banque unique ; pas de marché interbancaire ni de titrisation.
\limites La titrisation (revente des crédits) permettait avant 2008 de contourner cette contrainte, ce qui est précisément ce qui a fait déraper le levier ; le modèle peut le représenter par une baisse temporaire de $\kappa^{CAR}$.
\end{lecture}
et $\Lambda_j$ dans \eqref{eq:invest} est multiplié par $\Lambda^{Bk}$. Les crédits accordés créent des dépôts ($\Delta L = \Delta D$) : la masse monétaire est endogène ; le taux directeur agit sur la \emph{demande} de crédit (via $q_j$ et \eqref{eq:cH}) et sur le \emph{coût} du refinancement. Compte de résultat bancaire, pertes sur faillites et effet de la faillite bancaire :
\begin{equation}
\begin{split}
\Delta E^{Bk} &= i^L L - i^D D - i^{CB} L^{CB} + i^B B^{Bk} - \text{Pertes}^{NPL} - Div^{Bk} + \Delta E^{\'em},\\
Div^{Bk} &= \varpi^{Bk}\,\clip{\frac{CAR_t}{1{,}3\,\kappa^{CAR}}}{1}{3}\cdot\pos{E^{Bk} - \bar E^{Bk}}, \qquad CAR_t = \frac{E^{Bk}}{L+L^Z+L^E},\\
\Delta E^{\'em} &= \min\!\Big(1{,}3\,\kappa^{CAR}(L+L^Z+L^E) - E^{Bk},\ \ \varsigma^{\'em}\textstyle\sum_h D_h\Big)\cdot\ind_{\{0<CAR_t<1{,}1\kappa^{CAR}\}} .
\end{split}
\label{eq:bankeq}
\end{equation}
\begin{lecture}[compte de résultat de la banque]
\variables $\Delta E^{Bk}$ : variation des fonds propres ; $i^L L$ : intérêts reçus sur les crédits ; $i^D D$ : intérêts versés sur les dépôts ; $i^{CB}L^{CB}$ : intérêts payés à la banque centrale sur le refinancement ; $i^B B^{Bk}$ : intérêts reçus sur les obligations d'État détenues ; Pertes$^{NPL}$ : crédits passés en perte à la suite de faillites ; $Div^{Bk}$ : dividendes versés aux actionnaires de la banque.
\sens La banque gagne la différence entre ce qu'elle reçoit sur ses actifs et ce qu'elle paie sur ses ressources, moins ses pertes. Ses fonds propres sont le coussin qui absorbe les pertes ; s'ils deviennent négatifs, elle est insolvable et quelqu'un doit payer : l'État (renflouement) ou les déposants (faillite). L'enchaînement pertes $\to$ fonds propres $\to$ rationnement $\to$ récession $\to$ nouvelles pertes est le cycle du crédit de Minsky (1986), qui apparaît ici sans être programmé.
\hyp Pertes reconnues immédiatement (voir la forbearance, \S\ref{sec:zlb}, pour l'exception).
\limites Pas de coûts d'exploitation ni de revenus de commissions ; on les néglige car ils ne changent pas la dynamique.
\paragraph{Ajouts de la v1.1 : le capital bancaire est une variable à deux sens.} $\Delta E^{\text{ém}}$ : \emph{émission d'actions}. Quand le ratio de fonds propres passe sous $1{,}1\,\kappa^{CAR}$, la banque lève des fonds propres auprès des ménages, à hauteur d'au plus $\varsigma^{\text{ém}}=2\,\%$ de leurs dépôts par semaine, jusqu'à retrouver $1{,}3\,\kappa^{CAR}$ ; les ménages échangent des dépôts contre des actions bancaires (la masse monétaire baisse, la capacité de prêt monte). $Div^{Bk}$ : la distribution du capital excédentaire est \emph{accélérée} quand le ratio de couverture dépasse la cible, c'est-à-dire quand la demande de crédit est insuffisante pour employer le capital ; le facteur est borné à 3. Enfin, la recapitalisation publique se déclenche dès la \emph{sous-capitalisation} ($E^{Bk}<\tfrac12\kappa^{CAR}L$) et non seulement l'insolvabilité.
Ce que le prototype a montré : sans émission d'actions, le capital ne peut croître que par rétention de profits (quelques dixièmes de point de PIB par an) et la contrainte réglementaire devient une falaise. Une économie sortant d'une hyperinflation avec 3\,\% du PIB de fonds propres (contre 12\,\% en régime normal) voyait la reprise du crédit heurter le plafond en six mois : rationnement total, épuisement de la trésorerie du secteur des biens d'équipement, effondrement de son emploi (de 109 à 15), pénurie et inflation de 25\,\% sur les biens d'équipement pendant que le chômage atteignait 34\,\% ailleurs. Avec l'émission, la stabilisation post-hyperinflation est propre pour des monétisations de 25 à 40\,\% du PIB (chômage maximal 14 à 17\,\%, crédibilité reconstruite à son plafond, PIB $+33\,\%$ à vingt ans). \textbf{Propriété à retenir} : la réussite d'une stabilisation dépend de la capacité bancaire résiduelle autant que de l'ancre nominale ; c'est un objet de politique (recapitalisation, garantie, réglementation) et non seulement de politique monétaire.
\end{lecture}
Si $E^{Bk} < 0$ : renflouement public (coût $-E^{Bk}$ plus recapitalisation, inscrit au budget) ou pertes imputées aux déposants (baisse de $V_h$ et de $SoL_h$), au choix du joueur.
Une vague de faillites d'entreprises érode $E^{Bk}$, réduit $L^{max}$, donc $\Lambda_j$, donc l'investissement : c'est le cycle du crédit à la Minsky, qui émerge sans être scripté.

\paragraph{Ajouts de la v0.9 : banque insolvable.} Quand $E^{Bk}<0$, la banque ne prête plus, ne rémunère plus les dépôts et n'absorbe plus de dette publique ; l'État la recapitalise automatiquement au niveau $1{,}2\,\kappa^{CAR}$ de ses crédits par émission d'obligations souscrites par la banque centrale (au plus une fois par an), ce qui accroît la dette publique. Le prototype a montré que sans recapitalisation, une banque insolvable après une crise bloque le crédit pour toujours (chômage à 68\,\% quinze ans après une hyperinflation stabilisée) et que, sans le délai d'un an, la recapitalisation continue d'une banque qui perd chaque semaine est elle-même une monétisation. Le choix entre recapitalisation publique et pertes des déposants reste un levier du joueur (\S\ref{sec:run}) ; la recapitalisation automatique est le comportement par défaut de l'IA.

\subsection{Anticipations, crédibilité et inflation}
\label{sec:inflation}
L'inflation mesurée $\pi_t$ (éq.~\eqref{eq:cpi}) est \emph{émergente} : elle résulte des prix sectoriels \eqref{eq:price}, eux-mêmes poussés par les coûts (salaires \eqref{eq:wage}, intrants, importations) et par les déséquilibres offre/demande. Les anticipations qui bouclent le système sont :
\begin{align}
\pi^e_{t+1} &= \pi^e_t + \lambda_\pi\,(\pi_t - \pi^e_t) + c_t\,(\pi^\ast - \pi^e_t) + \upsilon\,\big(g^M_t - g^{Y}_t - \pi^e_t\big), \label{eq:expect}\\
c_{t+1} &= \operatorname{clip}\Big(c_t + \nu_1\,\ind_{\{|\pi_t - \pi^\ast| < \bar\epsilon\}} - \nu_2\,|\pi_t - \pi^\ast| - \nu_3\,\ind_{\{\Delta A^G_t > 0\}} - \nu_4\,\ind_{\text{BC dép.}} \nonumber\\
&\qquad\qquad ;\ 0,\ 1 - \nu_9 M^{hist}_t\Big), \label{eq:cred}\\
\lambda_\pi &= \lambda_0 + \lambda_1\, \ind_{\{\pi_t > \bar\pi^{hyper}\}}, \nonumber\\
M^{hist}_{t+1} &= (1-\delta_M)\, M^{hist}_t + \ind_{\text{abandon de la monnaie}} . \label{eq:hyper}
\end{align}
\begin{lecture}[anticipations d'inflation et crédibilité]
\variables $\pi^e_t$ : inflation anticipée ; $\pi_t$ : inflation observée ; $\lambda_\pi$ : vitesse de révision (part de l'erreur passée corrigée chaque mois) ; $c_t\in[0,1]$ : crédibilité de la banque centrale ; $\pi^\ast$ : cible annoncée ; $\upsilon$ : poids de l'ancre monétaire ; $g^M_t, g^Y_t$ : croissance de la masse monétaire et de la production réelle ; $\nu_1..\nu_4$ : vitesses de gain et de perte de crédibilité ; $\bar\epsilon$ : tolérance autour de la cible ; $\Delta A^G_t$ : financement monétaire du Trésor; $\bar\pi^{hyper}$ : seuil au-delà duquel les anticipations deviennent très réactives ($\lambda_1$ s'ajoute à $\lambda_0$) ; $\bar\pi_{24}$, $\bar\pi_{60}$ : inflation moyenne des 24 et 60 derniers mois; $\lambda_2$ : surcroît de réactivité quand l'inflation réapparaît après cinq ans d'absence (la surprise est plus grande) ; $M^{hist}$ : mémoire des abandons de monnaie (réforme monétaire, redénomination, remplacement par une devise), décroissant au taux $\delta_M$ ; $\nu_9$ : plafond de crédibilité par unité de mémoire.
\sens Les agents forment leur idée de l'inflation future de trois façons : en regardant l'inflation passée (adaptatif), en croyant la cible de la banque centrale si celle-ci est crédible (ancrage), et en tenant compte du fait qu'à long terme la monnaie créée en excès de la croissance réelle finit dans les prix (théorie quantitative). La crédibilité se gagne lentement (en tenant la cible) et se perd vite (en la ratant dans un sens ou dans l'autre : une déflation persistante est aussi un échec, et surtout en finançant l'État par la création monétaire). Elle est en outre \emph{plafonnée} par l'histoire : un pays qui a abandonné six monnaies en quarante ans ne peut pas dépasser une crédibilité de l'ordre de 0,5 quelle que soit sa performance courante, parce que la crédibilité est un stock et non un flux. Enfin, des agents qui n'ont pas vu d'inflation depuis cinq ans réagissent plus fort à sa réapparition (2021--2022). Quand l'inflation dépasse un seuil, les gens cessent de croire quoi que ce soit et n'anticipent plus qu'à partir du présent : les salaires et prix s'indexent presque instantanément, et l'inflation s'auto-entretient (hyperinflation).
\hyp Anticipations hybrides (Galí et Gertler 1999) plutôt que rationnelles ; crédibilité mesurée par un indice unique.
\limites Les anticipations réelles diffèrent entre ménages, entreprises et marchés financiers ; un seul $\pi^e$ est une simplification forte mais elle suffit à reproduire les régimes observés (ancrage, désancrage, hyperinflation).
\end{lecture}
Le terme en $\upsilon$ est l'ancre quantitativiste de long terme ($g^M$ : croissance de $M$, $g^Y$ : croissance réelle) : une création monétaire durablement supérieure à la croissance réelle finit par se refléter dans les anticipations, même si le canal de court terme passe par la demande. La crédibilité $c$ est le capital réputationnel de la banque centrale ; elle s'effondre en cas de financement monétaire du déficit. Au-delà du seuil $\bar\pi^{hyper}$, les anticipations deviennent quasi adaptatives instantanées et la vitesse de circulation augmente (régime de Cagan) : la demande de monnaie chute, $E_h$ est indexé sur l'inflation courante, et la boucle prix--salaires s'emballe.

\begin{lit}[courbe de Phillips réduite pour l'IA et l'équilibrage]
Au niveau agrégé, le système ci-dessus se comporte comme une courbe de Phillips hybride
$\pi_t \simeq (1-\omega)\,\pi^e_t + \omega\,\pi_{t-1} + \kappa_y \hat y_t + \kappa_m\, \pi^{imp}_t + \varepsilon_t$,
et l'écart de production comme une courbe IS
$\hat y_t \simeq -\sigma^{-1}(i_t - \pi^e_t - r^\ast) + \psi_G\,\hat g_t + \psi_x\,\hat x_t$
(impulsion budgétaire $\hat g$, exportations nettes $\hat x$). Ces formes réduites ne pilotent pas la simulation ; elles servent (i) à calibrer $\kappa_j,\phi_u,\iota$ pour obtenir des pentes plausibles, (ii) aux IA pour anticiper l'effet de leurs décisions.
\end{lit}

\subsection{Le levier de réforme monétaire}
\label{sec:reforme}
Le prototype a établi qu'une hyperinflation ne se stabilise pas par le seul arrêt de la création monétaire : les anticipations, l'indexation et la vitesse de circulation entretiennent la hausse des prix sans monnaie, et l'économie meurt (chômage 90\,\%). Il faut une \emph{ancre nominale nouvelle}, ce que l'histoire confirme (Rentenmark 1923, plan Real 1994). Le joueur dispose donc d'un levier de réforme, qui est aussi déclenché automatiquement par la redénomination :
\begin{equation}
\text{Réforme} : \quad P > 10^6 \ \text{ou décision} \ \Rightarrow\ 
\begin{cases}
\text{grandeurs nominales } \div 10^6, \quad M^{hist} \leftarrow M^{hist}+1,\\
\Delta A^G \leftarrow 0, \quad \pi^e,\pi^s \leftarrow \pi^\ast + 0{,}05,\\
c \leftarrow 0{,}8\,\max(1-\nu_9 M^{hist},\ \underline c),\\
\hat r^\ast \leftarrow \rho+\sigma g_0, \quad i^{CB} \leftarrow \hat r^\ast + \pi^e .
\end{cases}
\label{eq:reform}
\end{equation}
\begin{lecture}[réforme monétaire]
\variables $M^{hist}$ : nombre de monnaies abandonnées (mémoire monétaire, qui plafonne la crédibilité à $1-\nu_9M^{hist}$) ; $\underline c=0{,}05$ : plancher de crédibilité ; $\Delta A^G$ : avances de la banque centrale au Trésor ; $\pi^e,\pi^s$ : anticipations et inflation lissée de la règle.
\sens La réforme change l'unité de compte (ce qui n'a aucun effet réel), arrête le financement monétaire (ce qui est la condition nécessaire), et surtout remet les anticipations près de la cible avec une crédibilité partielle : c'est l'annonce d'un nouveau régime. La crédibilité restaurée est d'autant plus faible que le pays a déjà abandonné des monnaies. Le taux directeur est ramené au neutre : le prototype a montré qu'un taux laissé au plafond après la réforme (39\,\% nominal, $-26\,\%$ d'inflation, soit $+65\,\%$ réel) fait replonger l'économie.
\hyp La remise à zéro des anticipations est instantanée et uniforme ; en pratique elle exige un programme budgétaire crédible, que le joueur doit fournir sous peine de rechute (le prototype rechute si les avances reprennent).
\paragraph{Choix de conception, pas résultat théorique (v1.5).} La mémoire monétaire $M^{hist}$ et le plafond de crédibilité $1-\nu_9M^{hist}$ ne s'adossent à aucun modèle de la littérature ; les épisodes historiques (Allemagne 1923--1948, Argentine, Zimbabwe) suggèrent qu'un pays qui a déjà abandonné une monnaie stabilise plus difficilement la suivante, mais aucune estimation ne donne la forme ni le coefficient. C'est une règle de jeu, choisie pour rendre coûteuse la réforme répétée (propriété 12 : pas de méta-stratégie dominante), et elle doit être lue comme telle.
\limites En économie fermée, l'ancre est purement institutionnelle ; en économie ouverte (étape~2), un change fixe ou une caisse d'émission joueront ce rôle et le levier prendra la forme d'un choix de régime de souveraineté (\S\ref{sec:souv}).
\end{lecture}
Résultat du test de stabilisation : réforme déclenchée à 43\,\% d'inflation, inflation ramenée à $-3\,\%$ un an après, $+4\,\%$ à trois ans, $+2\,\%$ ensuite ; chômage de 14\,\% au pic (un an) puis retour vers 5\,\% en cinq ans ; crédibilité reconstruite à 0,85 sous le plafond 0,88.

\subsection{La machine à états : économie effondrée}
\label{sec:collapse}
Quand une hyperinflation n'est pas stabilisée, le système de prix cesse de fonctionner : les redénominations s'enchaînent, les prix relatifs perdent toute cohérence, les ventes s'effondrent alors que l'emploi subsiste. Le prototype a montré que ce régime est \emph{dégénéré} au sens numérique — il n'est plus décrit par les équations du modèle — et qu'aucune correction de mécanisme ne le résout sans en casser un autre. Le modèle bascule alors explicitement dans un régime de survie :
\begin{equation}
\begin{split}
\text{Entrée} &: \quad \pi^e_t \ge \bar\pi^e \text{ pendant 26 semaines,}\ \text{ou deux réformes en cinq ans} ;\\
\text{Survie} &: \quad L_j = \upsilon^{surv}(L^S - L^G)\,\tfrac{L_j}{\sum L}, \quad Y_j = A_jK_j^\alpha (HL_j)^{1-\alpha}, \quad \Delta K = -\tfrac12\delta K,\\
&\quad\ \ p_j,\ W_j,\ D_h \text{ indexés au taux } \pi^{surv} ;\ \text{ni crédit ni investissement net} ;\\
\text{Sortie} &: \quad \text{réforme \eqref{eq:reform}} \Rightarrow M \leftarrow \varsigma^{conv} PY \text{ (conversion, contre } \Delta A^G),\ \text{trésoreries et stocks}\\
&\quad\ \ \text{de redémarrage, banque au minimum réglementaire.}
\end{split}
\label{eq:collapse}
\end{equation}
\begin{lecture}[économie effondrée]
\variables $\bar\pi^e=1\,200\,\%$ : plafond des anticipations (le régime commence quand elles y restent six mois) ; $\upsilon^{surv}=0{,}6$ : part du potentiel produite en survie ; $\pi^{surv}=50\,\%$ : inflation résiduelle en unité stable (troc, devise, indexation généralisée) ; $\varsigma^{conv}=0{,}6$ : encaisses de transaction rétablies à la conversion, en part du PIB.
\sens En survie, l'économie tourne au ralenti dans une unité de compte de fait (le dollar, le sac de farine) : les prix et les salaires y sont stables, la monnaie officielle n'est qu'un multiplicateur commun. Le capital s'use sans être remplacé, le chômage est de l'ordre de 35\,\%, le niveau de vie de 40\,\% de la référence. La sortie est une réforme réussie : nouvelle unité, arrêt des avances, et \emph{conversion} — la nouvelle monnaie est émise en quantité suffisante pour que les échanges reprennent (sinon la reprise avorte faute d'encaisses), les entreprises reçoivent huit semaines de trésorerie et d'intrants, la banque est ramenée au minimum réglementaire.
\hyp Le régime de survie est stylisé : pas de marché, pas de crédit, une répartition figée. C'est un état d'attente, pas un modèle de l'économie informelle.
\paragraph{Choix de conception, pas résultat théorique (v1.5).} La machine à états n'a pas de référence dans la littérature : les modèles d'hyperinflation (Cagan 1956, Sargent--Wallace 1981) décrivent la divergence, pas ce qui se passe quand le système de prix a cessé de fonctionner. Les seuils (1\,200\,\% pendant six mois, deux réformes en cinq ans), la part produite en survie (0,6) et la conversion à la sortie (0,6) sont des choix de conception motivés par la stabilité numérique et par les récits d'économies effondrées, non des estimations. Le document ne fabrique pas de fondement après coup ; il documente ce que le prototype a exigé pour que la sortie réussisse.
\limites Ce que le prototype a montré : la sortie a échoué quatre fois avant de réussir — encaisses insuffisantes (M/PIB à 0,03), trésoreries d'entreprises héritées de la survie versées en dividendes (boom instantané), revenu permanent hérité de l'hyperinflation (consommation de 90\,\% des dépôts par semaine), prix d'actifs dans l'ancienne unité. Chaque échec était une variable d'état héritée sans signification dans la nouvelle unité : \emph{une réforme monétaire doit réinitialiser toutes les grandeurs nominales de référence}, pas seulement les prix. Après une sortie réussie (monétisation de 5\,\% du PIB pendant dix-huit ans, effondrement à l'an 15,6, réforme à l'an 18), le chômage revient à 2\,\% en trois ans et le PIB dépasse son niveau initial de 18\,\% à l'an 25 ; mais la crédibilité ne se reconstruit pas (mémoire monétaire) et l'inflation dérive vers 23\,\% en vingt ans : l'économie sortie de l'effondrement a besoin d'une ancre externe, ce que les régimes B ou E de la \S\ref{sec:souv} fournissent.
\end{lecture}

\subsection{Taux d'intérêt souverain et courbe des taux}
\label{sec:spread}
\begin{equation}
\begin{split}
i^B_t = i^{CB}_t + s_0\,\frac{\bar T}{\bar T_0}\Big(1 - \psi_{QE}\frac{B^{CB}_t}{B_t}\Big) + s_1\,\max\!\big(\varkappa^{for}_t,\ \underline\varkappa\,(1-\chi^K_i\varrho^B_i)\big)\,\pos{\frac{b^{mkt}_t}{\bar b_{i}} - 1}^2\\
+ s_2\, \pos{\Delta b^{proj}_t} + s_3\, D^{hist}_t + s_4\,(1-c_t) - s_6\,\ind_{\text{monnaie de réserve}} - s_8\,\ind_{\text{backstop}} - s_9\,\varrho^B,\\
 b^{mkt}_t = \frac{B_t - B^{CB}_t}{P_tY_t}, \qquad \varkappa^{for}_t = \frac{B^\ast_t}{B_t-B^{CB}_t},\\
\bar b_i = \bar b^{\,r\acute{e}gime}\cdot\min\!\Big(1{,}5,\ \frac{T_i/Y_i}{(T/Y)^{ref}}\Big), \qquad \Delta b^{proj}_t = \frac{i^{app}_t-\pi_t-g_t}{(1+\pi_t)(1+g_t)}\,b_t - pb_t,
\end{split}
\label{eq:spread}
\end{equation}
\begin{lecture}[taux d'intérêt de la dette publique]
\variables $i^B_t$ : taux exigé par les acheteurs d'obligations d'État ; $s_0\,\bar T/\bar T_0$ : prime de terme, proportionnelle à la maturité d'émission $\bar T$ (curseur du joueur, \S\ref{sec:iapp}) rapportée à une référence $\bar T_0$ ; $b_t=B_t/(P_tY_t)$ : ratio dette/PIB ; $\bar b_i$ : seuil effectif du pays, égal au seuil du régime multiplié par sa capacité fiscale relative $T_i/Y_i$ rapportée à une référence $(T/Y)^{ref}$ de 35\,\% du PIB (plafonné à 1,5) ; $\varkappa^{for}_t$ : part des non-résidents dans la dette \emph{de marché} $b^{mkt}$ (hors banque centrale), de sorte que la convexité ne porte que sur la fraction exposée à une fuite ; $\underline\varkappa=0{,}3$ : plancher de cette fraction, conditionné (v1.5) par la répression $\varrho^B$ et le contrôle des capitaux $\chi^K$ ; $\Delta b^{proj}_t$ : dérive projetée du ratio de dette (équation \eqref{eq:debtdyn} aux valeurs courantes), une dérive positive ajoutant $s_2$ par point de PIB ; $D^{hist}_t$ : mémoire des défauts passés ; $c_t$ : crédibilité monétaire ; $s_6$ : rabais de l'émetteur de la monnaie de réserve ; $\psi_{QE}B^{CB}/B$ : réduction de la prime de terme par la duration retirée du marché (le rabais $s_7$ des versions antérieures n'existe plus) ; $s_8$ : rabais du backstop d'union ; $s_9\varrho^B$ : répression financière (épargnants captifs).
\sens Le taux auquel un État emprunte est le taux sans risque plus une prime de terme ($s_0$), plus une prime qui monte avec la dette, les déficits, l'historique de défaut, la défiance monétaire et les entraves aux capitaux, et qui baisse quand la banque centrale achète la dette, quand la monnaie est une monnaie de réserve ou quand un backstop existe. La convexité en $b$ est ce qui rend la dette « soutenable jusqu'au jour où elle ne l'est plus » (Reinhart et Rogoff 2009). Trois précisions de la v0.7 en changent la portée. (i) L'excès est mesuré en \emph{relatif} au seuil : ce qui compte n'est pas d'être à 35\,\% du PIB mais d'être au double de ce que sa capacité fiscale permet. (ii) Le seuil dépend de la capacité fiscale : un État qui collecte 10\,\% du PIB a un seuil de 17\,\%, pas de 60\,\% ; c'est le dénominateur (les recettes), et non le numérateur (la dette), qui distingue un émergent d'un pays avancé. (iii) La convexité ne porte que sur la part de la dette détenue par des non-résidents : un pays dont la dette est détenue à 95\,\% par ses résidents (épargne captive, assureurs-vie, banque centrale) n'a pas de créanciers qui puissent partir ; il subit la prime de terme, l'inflation et la répression, pas la panique. C'est ce qui rend supportable un ratio de 230\,\% du PIB, et c'est mécanistique : $c_t$ et $s_6$ n'y sont pour rien. (iv) Le terme de déficit porte sur la \emph{dérive projetée} $\Delta b^{proj}$ et non sur le déficit primaire : un pays dont la croissance nominale excède son taux apparent peut financer un déficit primaire sans que sa dette dérive, et les créanciers le savent ; à l'inverse un excédent primaire ne rassure pas si le taux dépasse la croissance. C'est ce qui explique qu'un émetteur à déficit de 3,5\,\% et croissance nominale de 4,7\,\% paie moins qu'un membre d'union à déficit de 3\,\% et croissance nominale de 2,5\,\%.
\hyp Une prime réduite à des fondamentaux, complétée depuis la v0.6 par le saut discret de l'équation \eqref{eq:jump} pour les emprunteurs en monnaie qu'ils n'émettent pas.
\limites Le seuil $\bar b$ est un paramètre \emph{par régime} (cinq valeurs), non par pays : l'hétérogénéité entre pays passe par $T/Y$ et $\varkappa^{dom}$, qui sont des variables du modèle. Depuis la v0.8, il n'y a plus de prime d'illiquidité pour le contrôle des capitaux : un pays fermé paie par la répression ($\varrho^B$) et par $\varkappa^{dom}$, pas par un terme dédié. La part domestique est traitée comme exogène à court terme ; en réalité elle réagit aux taux (les résidents achètent quand la prime monte), ce qui est un stabilisateur supplémentaire non modélisé.
\paragraph{Réécriture de la v1.1 : hors banque centrale.} La convexité porte désormais sur la dette \emph{de marché} $b^{mkt}=(B-B^{CB})/PY$, c'est-à-dire celle qui doit être refinancée auprès d'investisseurs, et elle est pondérée par la part des non-résidents \emph{dans cette dette}, $\varkappa^{for}=B^\ast/(B-B^{CB})$, avec un plancher de 0,3 : même une dette entièrement domestique garde un coût convexe (éviction de l'épargne privée, risque de refinancement). Le rabais $-s_7B^{CB}/B$ de la v0.6 disparaît : il faisait un double comptage, la dette détenue par la banque centrale étant à la fois comptée dans le risque de refinancement et rabattue ensuite. Le canal de portefeuille de l'assouplissement quantitatif est conservé mais placé où il agit réellement, sur la \emph{prime de terme} : en retirant de la duration du marché, la banque centrale réduit $s_0$ d'une fraction $\psi_{QE}=0{,}5$ de la part de dette qu'elle détient. Effets mesurés : aucun changement en régime normal (la dette est sous le seuil) ; $-8$ pb sur le taux souverain lors d'un QE de 20\,\% du PIB ; dérive budgétaire inchangée (dette 1,12, taux maximal 12,0\,\% contre 11,7 auparavant). La lecture économique change en revanche : ce qui compte n'est pas « qui détient la dette » au sens large mais « quelle part doit être replacée sur un marché », la banque centrale étant une contrepartie de l'État consolidé et non un créancier.
\paragraph{Révision de la v1.5 : plancher conditionnel.} Le plancher $\underline\varkappa=0{,}3$ de la v1.1 représentait un risque de refinancement : même domestique, une dette doit être replacée, et l'épargnant peut refuser. Cet argument suppose que l'épargne ait une alternative. Il ne tient pas pour une économie à détention obligatoire ($\varrho^B=1$) et à sortie fermée ($\chi^K=1$), où le créancier captif ne peut ni partir ni refuser : la revue du cours a montré qu'il faisait payer 218 pb à la configuration 5 (dette détenue à 95\,\% par les résidents) pour une trentaine observée, ce qui contredisait la phrase (iii) ci-dessus. Le plancher devient $\underline\varkappa\,(1-\chi^K\varrho^B)$ : inchangé partout ailleurs, nul pour l'économie fermée à épargne captive. La forme multiplicative est vérifiée par le prototype (test H6) : le contrôle des capitaux seul ne change rien à la prime (écart nul), la répression seule agit par $-s_9$, les deux ensemble annulent le plancher. Ce qui ne bouge pas : la dérive budgétaire T9 est identique sous plancher 0,3, 0,15 ou 0 (taux maximal 29,5\,\% dans les trois cas), parce qu'en dynamique la prime est portée par la dérive projetée $s_2$ et la crédibilité $s_4$, non par la convexité ; le plancher est un enjeu de niveau statique. Ce que le prototype montre en plus, et qui est une inférence du modèle : sous répression, le même dérapage budgétaire coûte 5,4\,\% de taux souverain au lieu de 8,8 et une dette finale de 1,11 au lieu de 1,58, mais une inflation maximale de 7,7\,\% au lieu de 5,2 — la répression déplace le coût de la dette du taux vers l'inflation. Enfin, la répression $-s_9\varrho^B$, présente dans l'équation depuis la v0.8, n'existait pas dans le prototype avant la v1.5. Une limite assumée : l'équation construit $i^B$ sur le taux directeur \emph{courant}, donc sur une trajectoire anticipée plate ; en période de baisse anticipée (configuration 5), la prime de terme calculée (97 pb) excède l'écart observé, et cet écart d'anticipation est assumé, non corrigé.
\end{lecture}
\paragraph{Saut de prime et panique auto-réalisatrice.} Pour un pays qui emprunte dans une monnaie qu'il n'émet pas (régime B, ou dette en devise), la prime peut sauter sans changement des fondamentaux, comme la probabilité de dévaluation en change fixe (\S\ref{sec:change}) transposée du change à la dette :
\begin{equation}
\begin{split}
\Pr(\text{saut})_t &= \Phi\!\Big(\varsigma_0 + \varsigma_6\,(b_t - \bar b_i) + \varsigma_7\,\pos{-pb_t} - \varsigma_8\,\ind_{\text{backstop}} - \varsigma_9\, c_t + \varsigma_{11}\sum_{k\in\text{union}}\ind_{saut,k,t-1}\Big),\\
\rho^{pays}_t &\leftarrow \rho^{pays}_t + \Delta\rho^{saut} \ \text{si saut},
\end{split}
\label{eq:jump}
\end{equation}
le saut se résorbant ensuite au rythme $\lambda_\rho$ par mois si les fondamentaux le permettent.
\begin{lecture}[saut de prime souveraine]
\variables $\Pr(\text{saut})$ : probabilité mensuelle d'un décrochage de la prime ; $b-\bar b$ : excès de dette ; $\pos{-pb}$ : déficit primaire ; $\ind_{\text{backstop}}$ : 1 si le pays a accès au backstop de son union (ce qui suppose le respect des règles budgétaires) ; $c$ : crédibilité monétaire ; $\Delta\rho^{saut}$ : ampleur du saut (200 à 500 pb) ; $\lambda_\rho$ : vitesse de retour ; $\varsigma_0$ : constante fixant le niveau de base de la probabilité (v0.7) ; $\varsigma_{11}\sum_k \ind_{saut,k}$ : contagion, chaque autre membre de l'union en saut le mois précédent élevant la probabilité (v0.7).
\sens La dynamique de la dette \eqref{eq:debtdyn} admet deux équilibres quand la dette est élevée : un \og bon\fg{} où les taux sont bas parce que personne ne craint le défaut, un \og mauvais\fg{} où les taux sont hauts parce que tout le monde le craint, et où cette crainte rend le défaut probable (Calvo 1988, De Grauwe 2012). Le saut est le passage de l'un à l'autre. Un émetteur souverain ne le subit pas : il peut toujours monétiser (au prix de l'inflation). Un membre d'union sans backstop le subit ; avec backstop, la probabilité tombe presque à zéro. Pour le joueur, respecter la règle budgétaire supranationale est le prix de l'assurance qui empêche \eqref{eq:spread} de devenir explosive. La contagion (v0.7) est ce qui fait du régime B un jeu coopératif : le saut d'un membre triple la probabilité de saut des autres, y compris de ceux qui respectent les règles ; c'est la raison d'être du backstop, et ce qui donne aux membres vertueux un intérêt à ce que les autres soient secourus.
\hyp Saut de taille fixe ; retour graduel.
\limites Le saut est de taille fixe et la contagion est mécanique (nombre de membres en saut), sans pondération par la taille ou l'exposition croisée des banques ; le retour graduel suppose que les fondamentaux le permettent. Vérification numérique : avec $\varsigma_0=-2{,}55$, $\varsigma_6=3$, $\varsigma_7=20$, $\varsigma_8=1{,}05$, $\varsigma_9=1$, un membre à $b-\bar b_i=0{,}25$, $pb=-3\,\%$, $c=0{,}85$ a $z=-2{,}05$ sans backstop ($\Pr=2{,}0\,\%$ par mois) et $z=-3{,}10$ avec ($0{,}10\,\%$) ; un voisin en saut ($\varsigma_{11}=0{,}5$) porte ces valeurs à 6,1\,\% et 0,47\,\%.
\end{lecture}

avec $b = B/(PY)$ le ratio dette/PIB, $pb$ le solde primaire en \% du PIB, $D^{hist}$ une mémoire décroissante des défauts passés. La prime est convexe dans la dette (Reinhart--Rogoff) et pénalise le contrôle des capitaux (prime d'illiquidité). Le taux long privé suit $i^{long} = \bar\imath^{B} + \text{prime de terme}$.

\subsection{Taux de change}
\label{sec:change}
\paragraph{Change flottant (v1.0).} Le change est un prix d'actif : niveau cible (parité de pouvoir d'achat corrigée du différentiel de taux réel), rappel hebdomadaire vers cette cible, pression des flux nets de devises :
\begin{equation}
\begin{split}
e^{\,cible}_t &= e^{PPA}_t\,\exp\!\Big(-\theta_L\,(1-\chi^K)\big[(i_t-\pi^e_t) - (i^\ast_t-\pi^{e\ast}_t)\big] - \theta_{dev}\,\pi^{dev}_t\Big),\\
\Delta\ln e_t &= \lambda_e\,\big(\ln e^{\,cible}_t - \ln e_t\big) - \kappa_e\,\frac{NF_t}{M_t} + \varepsilon^e_t, \qquad NF_t = EX_t - IM_t + KA_t,\\
KA_t &= \Big[\kappa_K\,(1-\chi^K)\big[(i_t-\pi^e_t) - (i^\ast_t-\pi^{e\ast}_t)\big] - \kappa_{spec}\,\pi^{dev}_t\Big]\frac{M_t}{52},
\end{split}
\label{eq:fx}
\end{equation}
\begin{lecture}[taux de change flottant]
\variables $e_t$ : taux de change (unités de monnaie nationale par unité de monnaie étrangère : une hausse est une dépréciation) ; $i^\ast_t-i_t$ : écart entre le taux d'intérêt étranger et national ; $\rho^{pays}_t$ : prime de risque pays ; $\chi^K\in[0,1]$ : intensité du contrôle des capitaux (à 1, les capitaux ne circulent plus) ; $CA_t/(P_tY_t)$ : compte courant en \% du PIB ; $e^{PPA}_t=P_t/P^\ast_t$ : taux de change de parité de pouvoir d'achat (celui qui égaliserait les prix entre pays) ; $\Delta R^{FX}_t/H_t$ : interventions de la banque centrale (réserves utilisables) en proportion de la base monétaire ; $\theta_6(R^{FX}/H-\bar r^{FX})$ : effet de \emph{stock} des réserves, un stock supérieur à la normale $\bar r^{FX}$ signalant des ventes futures et exerçant une pression à l'appréciation (v0.6) ; $\theta_5(\pi^e-\pi^{e\ast})$ : différentiel d'inflation anticipée ; $\varepsilon^e_t$ : bruit.
\sens Quatre forces déplacent le change. (1) Les capitaux vont vers les taux d'intérêt élevés : un pays qui relève son taux voit sa monnaie s'apprécier (parité des taux d'intérêt non couverte, UIP), sauf si les capitaux sont bloqués. (2) Un déficit courant signifie que le pays vend plus de sa monnaie qu'il n'en achète : elle se déprécie. (3) À long terme, le change tend vers le niveau qui égalise les prix (parité de pouvoir d'achat, Cassel 1918) : un pays plus inflationniste voit sa monnaie se déprécier d'autant. (4) La banque centrale peut vendre ses réserves pour soutenir la monnaie ; mais un stock de réserves anormalement gros pèse en permanence sur le change, car les marchés anticipent qu'il sera un jour revendu (effet de stock). Le modèle de Dornbusch (1976) explique pourquoi le change \emph{surréagit} : les prix des biens sont lents, le change est instantané, donc il fait tout l'ajustement à court terme.
\hyp Coefficients constants ; anticipations de change implicites.
\limites Le change réel est notoirement imprévisible à court terme (Meese et Rogoff 1983) ; le modèle capture les tendances de moyen terme et les crises, pas le bruit quotidien, ce qui est suffisant pour un jeu à pas mensuel.
\paragraph{Réécriture de la v1.0.} Le change est traité comme un \emph{prix d'actif} : il a un niveau cible $e^{cible}$, la parité de pouvoir d'achat corrigée du différentiel de taux réel (un pays qui relève son taux voit sa monnaie s'apprécier \emph{immédiatement} d'environ $\theta_L=2$ fois le différentiel, puis revenir vers la PPA à mesure que les prix s'ajustent : c'est la surréaction de Dornbusch), et il y est rappelé chaque semaine ($\lambda_e=0{,}3$) ; s'y ajoute la pression des flux nets de devises $NF$ (balance commerciale plus flux de capitaux) rapportés à la masse monétaire ($\kappa_e=0{,}3$). Les flux de capitaux $KA$ suivent le différentiel de taux réel ($\kappa_K=2$ par an et par unité de différentiel), freinés par le contrôle des capitaux, moins la spéculation en régime fixe ($\kappa_{spec}=0{,}5$). Ce que le prototype a montré : la parité des taux écrite en \emph{flux} ($\theta_1(i^\ast-i)$ par semaine) donnait une appréciation de 2\,\% atteinte en 19 mois, sans surréaction ; la forme en niveau donne 4,4\,\% en 12 mois puis un retour à 1,5\,\% à cinq ans. La PPA n'est qu'un rappel lent et n'est pas imposée : sur un épisode de monétisation, la dépréciation cumulée suit le différentiel de prix cumulé à $\pm35\,\%$. Les termes $\theta_2\,CA$, $\theta_4\,\Delta R^{FX}$ et $\theta_6$ de la v0.9 sont absorbés dans $NF$ ; le terme $\theta_5(\pi^e-\pi^{e\ast})$ est absorbé dans $e^{PPA}$.
\end{lecture}
avec $e^{PPA} = P/P^\ast$, $\rho^{pays}$ une prime de risque (défaut, instabilité politique). Une hausse du taux domestique apprécie la monnaie (Dornbusch) ; un déficit courant la déprécie ; une vente de réserves ($\Delta R^{FX}<0$) la soutient. Le contrôle des capitaux ($\chi^K \to 1$) coupe le premier terme : le taux directeur retrouve son autonomie au prix de la prime dans \eqref{eq:spread} (triangle d'incompatibilité de Mundell).

\paragraph{Change fixe / étalon-or (v1.0).} $e_t=\bar e$ ; le taux directeur est \emph{importé} : $i_t = i^\ast_t + \tfrac12\pi^{dev}_t$ (régime E ; en union, taux commun) ; la banque centrale absorbe $\Delta R^{FX}=EX-IM+KA$. La probabilité perçue de dévaluation combine l'épuisement des réserves et l'incohérence des fondamentaux :
\begin{equation}
\pi^{dev}_t = \tfrac12\,\clip{\frac{R^{FX}_0 - R^{FX}_t}{R^{FX}_0 - \underline R}}{0}{1}\cdot F_t, \qquad
F_t = \begin{cases} \clip{(\pi_t-\pi^\ast_t)/0{,}05}{0}{1} & \text{1\up{re} génération}\\ 1 & \text{2\up{e} génération (option)} \end{cases}
\label{eq:pdev}
\end{equation}
La dévaluation (30\,\%, puis flottement, perte de crédibilité de 0,3) survient quand $R^{FX}<\underline R$ (deux mois d'importations), ou, en seconde génération, quand le gouvernement \emph{choisit} d'abandonner : $\pi^{dev}>5\,\%$ et $u>u^n+4$ pts, la défense par les taux coûtant plus qu'elle ne rapporte. Ce que le prototype a montré : (i) avec sa propre règle de Taylor, le pays ancré défendait la parité par les taux et l'attaque de première génération ne venait qu'à l'an 17 au lieu de 12 ; la perte d'autonomie doit être écrite ; (ii) une probabilité de dévaluation fondée sur les seules réserves produit des attaques auto-réalisatrices dès la troisième année sans aucune faute de politique, d'où le facteur de fondamentaux en première génération ; (iii) en seconde génération sans clause d'abandon, la hausse des taux importée attire des capitaux qui compensent la spéculation et l'attaque vient \emph{plus tard} qu'en première ; avec la clause, une récession chez le partenaire suffit à faire tomber une parité que rien ne menaçait (abandon à l'an 5,3, réserves encore à 3,6\,\% du PIB, chômage 12\,\%) ; (iv) le contrôle des capitaux retarde l'attaque (an 8,6 au lieu de 6,3) sans l'empêcher quand la parité est incohérente : la surchauffe passe par le déficit commercial. Sous étalon-or, la parité fixe la quantité de monnaie émissible : $H\le G/\bar g$, et la dévaluation à la manière d'E\&F consiste à abaisser $\bar g$.

\paragraph{Défense d'un plafond de change.} Le cas symétrique est celui d'un pays qui \emph{refuse l'appréciation} de sa monnaie (monnaie refuge, excédent courant structurel). Ses munitions sont illimitées : il crée sa propre monnaie pour acheter des devises, il ne peut pas manquer de réserves, il ne peut que grossir. La contrainte n'est donc pas $R^{FX}\ge\underline R$ mais la contrainte de bilan de \eqref{eq:ecb} : $E^{CB}/H \ge \underline\kappa^{CB}$. Chaque intervention réussie augmente $R^{FX}$, donc l'exposition au change qu'elle cherche à contenir ; une appréciation ultérieure de 10\,\% sur des réserves valant 120\,\% du PIB coûte 12\,\% du PIB de fonds propres. Le terme de stock $-\theta_6$ dans \eqref{eq:fx} rend la sortie coûteuse : vendre des réserves, c'est acheter sa propre monnaie, donc déclencher l'appréciation que l'on combat, et les marchés l'anticipent tant que le stock est gros. Le modèle produit ainsi l'abandon soudain d'un plancher par coût de bilan et non par épuisement des munitions ; l'abandon lui-même (Suisse, janvier 2015) provoque une appréciation immédiate, une perte de valorisation massive et une déflation importée.

\paragraph{Monnaie non convertible.} $\chi^K = 1$, change officiel $\bar e$ administré, marché noir des devises au taux $e^{bm} = \bar e\,(1+\mu_e\,(D^{FX}/S^{FX} - 1))$ ; les importations sont rationnées par les licences (décidées par le joueur) ; les exportateurs remettent leurs devises à la banque centrale.


\begin{figure}[htbp]
\centering
\resizebox{0.98\textwidth}{!}{%
\begin{tikzpicture}[x=1cm,y=1cm]
% ---- détenteurs ----
\node[agent] (men) at (2,8)     {Ménages $M,H$\\ épargne $S_h$, richesse $V_h$};
\node[agent] (bk)  at (9.6,8)   {Banque commerciale\\ $E^{Bk}$, ratio $\kappa^{CAR}$};
\node[agent] (bc)  at (16,8)    {Banque centrale\\ $H = C + Res$};
% ---- actifs ----
\node[actif] (dep)  at (0,4.5)    {Dépôts $D$\\ taux $i^D$};
\node[actif] (bond) at (4,4.5)    {Obligations d'État\\ $p^B$, taux $i^B$};
\node[actif] (eq)   at (8,4.5)    {Actions\\ $p^E$, $r^E$};
\node[actif] (im)   at (12,4.5)   {Immobilier\\ $p^Z$, stock $Z$};
\node[actif] (fx)   at (16,4.5)   {Devises\\ $e$, $R^{FX}$};
% ---- émetteurs ----
\node[agent] (tres)   at (4,1)    {Trésor\\ $\Delta B$, défaut $\hbar$};
\node[agent] (ent)    at (8,1)    {Entreprises\\ $q_j$, levier $\ell_j$};
\node[agent] (constr) at (12,1)   {Construction\\ $\Delta Z = I_Z$};
\node[agent,fill=gray!12] (rdm) at (16,1) {Reste du monde\\ $B^\ast$, $KA$};

% ---- ménages -> actifs ----
\draw[fluxf] (men.south west) -- node[lab,pos=0.55]{$s^D_h$} (dep.north);
\draw[fluxf] (men.south) -- node[lab,pos=0.55]{$s^B_h$} (bond.north west);
\draw[fluxf] (men.south east) -- node[lab,pos=0.35,above,sloped]{$s^E_h$} (eq.north west);
\draw[fluxf] (men.north) to[out=90,in=100] node[lab,pos=0.5,above]{logement (LTV)} (im.north);

% ---- banque -> actifs ----
\draw[fluxf] (bk.south west) -- node[lab,pos=0.42,above,sloped]{$B^{Bk}$} (bond.north east);
\draw[fluxf] (bk.south) -- node[lab,pos=0.55,left]{marge $m$} (eq.north east);
\draw[fluxf] (bk.south east) -- node[lab,pos=0.5,right]{$L^Z$} (im.north west);
\draw[fluxf] (bk.south west) .. controls (7.6,6.0) and (6.1,3.2) .. node[lab,pos=0.2,anchor=east]{crédits $L_j$,\\ intérêts $i^L$} (ent.north west);

% ---- banque centrale ----
\draw[fluxp] (bc.west) -- node[lab,pos=0.5,above]{refi. $L^{CB}$ à $i^{CB}$} (bk.east);
\draw[fluxp] (bc.north) to[out=90,in=80] node[lab,pos=0.5,above]{QE : $\Delta B^{CB}$} (bond.north east);
\draw[fluxp] (bc.south) -- node[lab,pos=0.5,right]{interventions} (fx.north);

% ---- émetteurs -> actifs ----
\draw[flux] (tres.north) -- (bond.south);
\draw[flux] (ent.north) -- node[lab,pos=0.5,right]{dividendes} (eq.south);
\draw[flux] (constr.north) -- (im.south);
\draw[fluxf] (rdm.north) -- node[lab,pos=0.5,right]{capitaux $KA$} (fx.south);
\draw[fluxf] (rdm.south) to[out=-100,in=-80] node[lab,pos=0.5,below]{dette externe $B^\ast$} (tres.south);

% ---- rétroactions ----
\draw[fluxf,red!70!black,dashed] (im.east) to[out=0,in=-20] node[lab,pos=0.6,right,text=red!70!black]{collatéral :\\ $L \le \mathrm{LTV}\,p^Z Z$} (bk.east);
\draw[fluxf,red!70!black,dashed] (eq.south west) to[out=-160,in=-90] node[lab,pos=0.28,below,text=red!70!black]{effet richesse $c^V_h$} (men.west);
\end{tikzpicture}}
\caption{Système financier. Quatre classes d'actifs avec un prix chacune (obligations, actions, immobilier, devises) plus les dépôts. Les banques prêtent contre collatéral (pointillés rouges : boucles prix d'actifs $\to$ crédit $\to$ prix d'actifs, et prix d'actifs $\to$ richesse $\to$ consommation) ; la banque centrale intervient par le taux de refinancement, le QE sur les obligations et le marché des changes.}
\label{fig:fin}
\end{figure}

\subsection{Actifs, prix d'actifs et bulles}
\label{sec:actifs}
Quatre actifs ont un prix de marché : obligations d'État (taux $i^B$, \S\ref{sec:spread}), actions ($p^E$), immobilier ($p^Z$) et devises ($e$, \S\ref{sec:change}). Ce sont les supports indispensables des crises de 1929 (actions, prêts sur marge), 1990 (foncier et actions au Japon) et 2008 (immobilier et crédit hypothécaire).

\paragraph{Immobilier.} Le stock de logements $Z$ est produit par le secteur construction ($\Delta Z = I_Z - \delta_Z Z$). La demande de logement combine un besoin de service (dans le LES) et une demande d'actif. Le coût d'usage et le prix évoluent selon
\begin{align}
uc_t &= i^{Z}_t - \pi^e_t + \delta_Z - \E_t g^Z, \qquad
D^Z_t = \sum_h N_h\,\zeta_h \frac{Y^d_h}{p^Z_t\, uc_t} \cdot \Lambda^{Z}_t, \qquad \Lambda^Z_t = \min\Big(1, \frac{\mathrm{LTV}_t\, p^Z_t}{p^Z_t - \text{apport}}\Big), \\
\frac{p^Z_{t+1}}{p^Z_t} &= 1 + \kappa_Z \frac{D^Z_t - S^Z_t}{D^Z_t + S^Z_t}, \qquad
\E_t g^Z = (1-\chi_Z)\, \bar g^Z + \chi_Z\, \overline{\Delta \ln p^Z}_{t-1,\,12},
\label{eq:pz}\\
I_{Z,t} &= Z_t\,\big(\delta_Z + \iota_Z \pos{p^Z_t / c^u_{constr,t} - 1}\big) \quad \text{(livraison différée de } T_Z \text{ mois)}.
\end{align}
\begin{lecture}[prix de l'immobilier et bulles]
\variables $uc_t$ : coût d'usage du logement, ce que coûte réellement le fait de détenir un logement pendant un an : intérêt hypothécaire réel ($i^Z-\pi^e$) plus dépréciation $\delta_Z$ moins plus-value attendue $\E_t g^Z$ ; $D^Z_t$ : demande de logements ; $\zeta_h$ : part du revenu que la strate consacre au logement ; $\Lambda^Z_t$ : accès au crédit hypothécaire, déterminé par la quotité $\mathrm{LTV}$ (\emph{loan-to-value}, part du prix que la banque accepte de financer) ; $S^Z_t$ : offre (stock existant plus constructions livrées) ; $\kappa_Z$ : vitesse d'ajustement du prix ; $\chi_Z\in[0,1]$ : poids des plus-values passées dans les anticipations ; $\overline{\Delta\ln p^Z}_{t-1,12}$ : hausse moyenne du prix sur les douze derniers mois ; $\bar g^Z$ : croissance « normale » ; $I_{Z,t}$ : mises en chantier ; $c^u_{constr}$ : coût de construction ; $T_Z$ : délai de construction.
\sens Le coût d'usage est la clef : si les ménages s'attendent à ce que les prix montent de 8\,\% par an alors que l'emprunt coûte 4\,\% réel, posséder un logement rapporte de l'argent, la demande explose et le prix monte, ce qui confirme l'anticipation. C'est une bulle auto-réalisatrice (Shiller 2000, Kindleberger 1978), rendue possible par $\chi_Z$ élevé (on extrapole le passé) et par un crédit facile (LTV haut, taux bas). Elle se retourne quand le coût d'usage remonte (hausse des taux) ou quand l'offre nouvelle arrive (avec le délai $T_Z$, d'où la surproduction de logements typique des fins de bulle : Espagne, Irlande, 2008).
\hyp Un marché du logement national ; anticipations extrapolatives ; offre élastique avec délai.
\limites Le foncier est en réalité très hétérogène (Tokyo vs campagne) ; un marché unique exagère la synchronisation. Les impôts fonciers et le marché locatif sont omis.
\end{lecture}
Le paramètre $\chi_Z$ (poids des plus-values passées dans les anticipations) est le moteur des bulles : pour $\chi_Z$ élevé et un crédit accommodant ($\mathrm{LTV}$ haut, $i^Z$ bas), la hausse s'auto-alimente jusqu'à ce que le coût d'usage ou l'offre nouvelle (délai $T_Z$) la retourne. L'immobilier sert de \textbf{collatéral} aux ménages ($L^Z_h \le \mathrm{LTV}\, p^Z Z_h$) et aux entreprises ($L_j \le \bar\ell\, p_K K_j + \mathrm{LTV}^{ent} p^Z Z_j$, ce qui est la clé du Japon de 1990) ; une baisse de $p^Z$ met des emprunteurs en fonds propres négatifs : $NPL^Z_t = \phi^Z \sum_h \pos{L^Z_h - p^Z_t Z_h}$.

\paragraph{Actions.} Le prix des actions équilibre l'offre (émissions nettes $\Delta E$) et la demande issue du portefeuille \eqref{eq:portfolio}, où le rendement attendu comprend les plus-values extrapolées :
\begin{equation}
\begin{split}
r^E_t = \frac{\sum_j (1-\theta_j)Div_j}{p^E_t E} + \E_t g^E, \qquad
\E_t g^E = (1-\chi_E)\, g^e + \chi_E\, \overline{\Delta \ln p^E}_{t-1,\,12},\\[2pt]
\frac{p^E_{t+1}}{p^E_t} = 1 + \kappa_E \frac{D^E_t - S^E_t}{D^E_t + S^E_t}.
\end{split}
\label{eq:pe}
\end{equation}
\begin{lecture}[prix des actions et prêts sur marge]
\variables $r^E_t$ : rendement attendu des actions ; $E$ : nombre d'actions (capital coté) ; $p^E_t$ : cours ; $\E_t g^E$ : plus-value attendue, mélange de la croissance fondamentale $g^e$ et de l'extrapolation des douze derniers mois (poids $\chi_E$) ; $D^E_t, S^E_t$ : demande et offre d'actions issues des portefeuilles ; $\kappa_E$ : vitesse d'ajustement du cours ; $p^{E,f}$ : valeur fondamentale (dividendes actualisés) ; $L^{marg}_h$ : prêts sur marge, crédit accordé pour acheter des actions ; $m_t$ : marge maximale, part du portefeuille finançable à crédit.
\sens Le rendement d'une action est son dividende rapporté au cours, plus la hausse de cours espérée. Quand les cours ont monté, les investisseurs extrapolent, achètent, et font monter les cours : même mécanique que l'immobilier. Le rapport cours/valeur fondamentale est l'indicateur de bulle affiché au joueur. Les prêts sur marge amplifient tout : quand le cours baisse, la valeur du collatéral baisse, la banque exige un remboursement (\emph{appel de marge}), l'investisseur doit vendre, ce qui fait baisser le cours : c'est la spirale d'octobre 1929, quand les marges autorisées atteignaient 90\,\%.
\hyp Marché national unique ; pas de nouvelles émissions en réaction à la bulle (elles existeraient et l'amortiraient un peu).
\limites Le modèle de dividendes actualisés est une référence, pas une prédiction ; les vrais fondamentaux sont incertains, ce qui laisse place à des bulles « rationnelles » que ce modèle ne distingue pas des irrationnelles.
\end{lecture}
La valeur fondamentale $p^{E,f} = Div/(r^L + \rho^E - g^e)$ est affichée au joueur ; l'écart $p^E/p^{E,f}$ est l'indicateur de bulle. Les ménages $H$ peuvent acheter à crédit (\textbf{prêts sur marge}) : $L^{marg}_h \le m_t\, p^E E_h$, la marge $m$ étant un instrument de régulation du joueur (Fed 1929 : $m \approx 0{,}9$). Un appel de marge force la vente : $S^E \uparrow$, $p^E \downarrow$, nouveaux appels --- la spirale de 1929.

\paragraph{Effet richesse et collatéral dans les équations existantes.} La richesse $V_h = D_h + p^B B_h + p^E E_h + p^Z Z_h - L_h$ entre dans la consommation via $c^V_h$ (\eqref{eq:cB}--\eqref{eq:cH}) et le levier $\ell_j$ de \eqref{eq:invest} est calculé sur la valeur de marché du collatéral : c'est par ces deux canaux que les prix d'actifs deviennent macroéconomiques.

\paragraph{Hystérèse de la dollarisation.} La part en devises désirée est celle du portefeuille ci-dessus ; la part effective ne redescend que lentement, car contrats, prix affichés, habitudes et infrastructure de paiement ont un coût de retour :
\begin{equation}
\begin{split}
s^{FX}_t &= \max\!\Big(s^{FX,\,d\acute{e}sir\acute{e}}_t,\ \phi_t\, s^{FX}_{t-1}\Big), \qquad \phi_t = 0{,}9995\,(1 - \kappa_d\,\varrho^{dedol}_t),\\
s^{FX,\,d\acute{e}sir\acute{e}}_t &= \clip{2\,\big(\max(\pi^e_t,\pi_t) - 0{,}05\big)}{0}{0{,}9},
\end{split}
\label{eq:sfx}
\end{equation}
avec $\varrho^{dedol}\in[0,1]$ la contrainte légale à l'usage de la monnaie nationale (\S\ref{sec:systemes}).
\begin{lecture}[hystérèse de la dollarisation]
\variables $s^{FX}_t$ : part effective de l'économie en devise étrangère ; $s^{FX,\,d\acute{e}sir\acute{e}}$ : part que les agents choisiraient aujourd'hui d'après les rendements et l'inflation anticipée ; $\phi_0$ : persistance mensuelle (0,995 : la dollarisation ne recule que de 6\,\% par an quand rien ne la pousse) ; $\varrho^{dedol}$ : intensité des obligations légales d'usage de la monnaie nationale (cours légal, impôts payables en monnaie nationale, cession obligatoire des devises, affichage), et $\kappa_d$ leur efficacité.
\sens On se dollarise en quelques mois de forte inflation et on met des décennies à en sortir : la part en devises ne peut baisser que de $(1-\phi)$ par mois, quelle que soit la performance monétaire courante. C'est pourquoi une stabilisation réussie sur les prix peut échouer sur la reconquête monétaire. La contrainte légale accélère la sortie, au prix d'un marché noir des devises et d'une fuite vers l'informel (\S\ref{sec:systemes}).
\hyp Asymétrie pure (hausse instantanée, baisse lente) ; le plancher n'est jamais franchi par les seuls rendements.
\limites Le degré de dollarisation réelle diffère selon les usages (dépôts, crédits, prix, salaires) ; un seul $s^{FX}$ les agrège.
\paragraph{Ajout de la v1.0.} La part désirée dépend de l'inflation \emph{réalisée} autant qu'anticipée (les agents fuient une monnaie qui fond avant que leurs anticipations aient rattrapé), et l'hystérèse est calibrée à $\phi=0{,}9995$ par semaine (demi-vie 27 ans). Test : après une monétisation de 25\,\% du PIB pendant trois ans (inflation 15\,\%) et une réforme, la part dollarisée monte à 18\,\% puis ne retombe qu'à 16\,\% cinq ans plus tard et 13\,\% treize ans plus tard, avec une inflation revenue à 2\,\% : la dédollarisation est un stock. En régime D, le taux réel des décisions d'investissement est $(1-s^{FX})r^L+s^{FX}r^\ast$ \eqref{eq:reff} et les encaisses nationales $(1-s^{FX})M$.
\end{lecture}

\subsection{Ruées bancaires, garantie des dépôts et prêteur en dernier ressort}
\label{sec:run}
Sans garantie des dépôts, une fraction des déposants retire ses fonds lorsque la banque paraît fragile :
\begin{equation}
\omega^{run}_t = \Phi\!\Big(\varsigma_1\,(\hat\rho^{NPL}_t - \bar\rho) + \varsigma_2 \pos{-E^{Bk}_t}/D_t + \varsigma_3\, \omega^{run}_{t-1} - \varsigma_4\, \ind_{\text{garantie}}\Big),
\label{eq:run}
\end{equation}
\begin{lecture}[ruées bancaires]
\variables $\omega^{run}_t$ : fraction des dépôts retirée pendant le tick ; $\Phi$ : fonction de répartition de la loi normale (transforme un score en probabilité entre 0 et 1) ; $\hat\rho^{NPL}_t-\bar\rho$ : excès de prêts non performants par rapport à la normale ; $\pos{-E^{Bk}_t}/D_t$ : insuffisance de fonds propres rapportée aux dépôts ; $\omega^{run}_{t-1}$ : retraits de la période précédente (la panique est contagieuse) ; $\ind_{\text{garantie}}$ : vaut 1 si une garantie publique des dépôts existe ; $\varsigma_1..\varsigma_4$ : sensibilités ; $L^{CB}$ : refinancement d'urgence de la banque centrale (prêteur en dernier ressort).
\sens Un déposant qui doute de sa banque retire son argent ; s'il voit les autres le faire, il se précipite (Diamond et Dybvig 1983) : la ruée est auto-réalisatrice, une banque solvable mais illiquide peut être détruite. La banque puise dans ses réserves, puis emprunte à la banque centrale contre des titres, puis vend ses actifs à perte, puis fait faillite. Sous l'étalon-or, la banque centrale ne peut pas prêter au-delà de sa couverture-or : c'est ce qui a transformé les paniques de 1930--33 en effondrement du système bancaire américain (Friedman et Schwartz 1963). La garantie des dépôts (FDIC, 1933) coupe la ruée à la racine, au prix d'un risque pour le contribuable et d'un aléa moral pour la banque.
\hyp Déposants homogènes ; contagion représentée par l'autocorrélation.
\limites En 2008 la ruée a frappé le financement de marché (repo, papier commercial) plutôt que les déposants ; le modèle le représente par une ruée sur $L^{CB}$ et la dette bancaire de court terme.
\end{lecture}
La banque honore les retraits avec ses réserves, puis avec le refinancement de la banque centrale $L^{CB}$ (limité aux titres éligibles et, sous étalon-or, par $H \le G/\bar g$), puis par ventes forcées d'obligations et de crédits à prix décoté (pertes, $E^{Bk}\downarrow$, contagion sur $p^B$ et $p^E$). Si les réserves sont épuisées : faillite, pertes des déposants ($V_h \downarrow$, $M \downarrow$ : contraction monétaire de 1930--33). La \textbf{garantie des dépôts} est une loi du joueur : elle annule quasiment $\omega^{run}$ mais crée un passif contingent de l'État et un aléa moral ($\kappa^{CAR}$ effectif plus faible).

\subsection{Borne zéro, trappe à liquidité et déflation par la dette}
\label{sec:zlb}
Le taux directeur est borné : $i^{CB}_t \ge \underline{\imath}$ ($\underline\imath = 0$, ou légèrement négatif). Le taux réel $r = i - \pi^e$ ne peut donc pas descendre sous $-\pi^e$ : en déflation, il \emph{monte} quand la banque centrale voudrait le baisser. Comme les dettes $L_j$, $L_h$, $B$ sont nominales, une baisse de $P$ élève leur poids réel et le levier $\ell_j$ (dénominateur $p_K K_j$ en baisse), ce qui réduit $\Lambda_j$ et déclenche des défauts : c'est la déflation par la dette de Fisher (1933), obtenue sans équation supplémentaire. En trappe à liquidité, la demande de crédit est insensible au taux ($q_j < 1$ quel que soit $i^L$ car $\E\Pi$ est déprimé) et le QE se traduit en réserves excédentaires (\S\ref{sec:qe}) ; seule la politique budgétaire, une dévaluation ou un relèvement crédible des anticipations $\pi^e$ (cible d'inflation, financement monétaire annoncé) font sortir de la trappe.

\paragraph{Banques zombies (forbearance).} Une banque sous-capitalisée a intérêt à reconduire les prêts d'emprunteurs insolvables plutôt qu'à constater la perte : la probabilité de reconduction est $\Pr(\text{reconduction}) = \Phi(\varsigma_5(\kappa^{CAR} - E^{Bk}/L))$. Les secteurs \og zombies\fg{} ainsi maintenus pèsent sur la productivité agrégée ($-\eta_Z\cdot \text{part}_{zombie}$ dans \eqref{eq:tfp}, Caballero--Hoshi--Kashyap 2008) : c'est la décennie perdue japonaise. Le joueur peut l'éviter par une recapitalisation forcée, coûteuse pour le budget.

\subsection{Assouplissement quantitatif, financement monétaire et hyperinflation}
\label{sec:qe}
Il faut distinguer trois agrégats : la base $H = C + Res$, les réserves excédentaires $Res^{ex} = Res - \varrho^{res} D$, et la masse monétaire $M = C + D$. Les prix réagissent à $M$ et à sa vitesse, pas à $H$ :
\begin{equation}
\text{terme monétaire de l'équation de prix \eqref{eq:price}} : \quad + \xi_M \Big[\ln\!\Big(\frac{M_t\, V(\pi^e_t)}{Y_t}\Big) - \ln P_t\Big], \qquad V(\pi^e) = V_0\, e^{a\,\pi^e} .
\label{eq:monterm}
\end{equation}
\begin{lecture}[encaisses réelles et vitesse de circulation]
\variables $M_t$ : masse monétaire (billets plus dépôts) ; $V(\pi^e)$ : vitesse de circulation, nombre de fois qu'une unité de monnaie change de main par an pour payer la production ; $V_0$ : vitesse de référence ; $a$ : sensibilité de la vitesse à l'inflation anticipée ; $Y_t$ : production réelle ; $P_t$ : niveau des prix ; $\xi_M$ : vitesse à laquelle un déséquilibre monétaire se répercute sur les prix.
\sens L'équation quantitative $MV=PY$ (Fisher 1911) dit que la monnaie en circulation, multipliée par le nombre de fois qu'elle circule, égale la valeur nominale de ce qui est acheté. Le terme ajouté à l'équation de prix dit que si la monnaie en circulation dépasse ce que l'économie « absorbe » à ses prix courants, les prix finissent par monter (les gens ont plus de monnaie qu'ils ne veulent en détenir et la dépensent : effet d'encaisses réelles de Patinkin). Il est lent ($\xi_M$ petit) : à court terme, l'inflation vient de la demande et des coûts, pas de la monnaie. La vitesse dépend de l'inflation anticipée (Cagan 1956) : quand la monnaie perd sa valeur, on s'en débarrasse vite, ce qui accélère encore les prix.
\hyp Vitesse fonction de $\pi^e$ seulement (en réalité aussi des taux d'intérêt et des innovations de paiement) ; $M$ exclut les réserves excédentaires des banques.
\limites La théorie quantitative est une relation de long terme ; l'utiliser à court terme conduit à surestimer l'effet inflationniste du QE, erreur fréquente en 2010--2015 que le modèle évite précisément par le faible $\xi_M$ et l'exclusion des réserves.
\paragraph{Vérification (v1.4).} Les deux cas du QE sont reproduits par le prototype selon la structure de détention de la dette : avec un ratio de liquidité bancaire ($\phi_{liq}=0{,}3$ : la banque détient 10\,\% du PIB de titres), le QE s'achète à la banque et la monnaie ne bouge pas ($-0{,}4\,\%$), le souverain baisse de 8 pb, les actions gagnent 0,8\,\%, l'inflation 0,1 pt ; sans ratio, il s'achète aux ménages et vaut monétisation ($M+15\,\%$, inflation $+0{,}9$ pt). Dans les deux cas l'effet sur l'emploi est nul et la banque centrale porte des pertes de coupon quand les taux remontent.
\paragraph{Limite assumée (v1.5) : l'ampleur de l'effet sur les taux.} Le modèle donne $-8$ pb pour des achats de 20\,\% du PIB ; les estimations empiriques (Gagnon et al. 2011, Krishnamurthy et Vissing-Jorgensen 2011, Joyce et al. 2011) situent l'effet entre 50 et 100 pb pour 10\,\% du PIB. L'écart est d'un ordre de grandeur et il est conforme aux équations : $\psi_{QE}=0{,}5$ ne porte que sur la prime de terme $s_0\bar T/\bar T_0$, soit un effet maximal de $50\,\text{pb}\times\omega$. Il est assumé plutôt que corrigé, pour deux raisons. La première est que \eqref{eq:spread} construit $i^B$ sur le taux directeur \emph{courant} : une part substantielle de l'effet mesuré du QE passe par le signal sur la trajectoire future des taux, canal que le modèle n'a pas ; les 50 à 100 pb observés ne sont donc pas comparables aux 8 pb calculés, qui isolent le seul canal de portefeuille. La seconde est que relever $\psi_{QE}$ ou $s_0$ pour atteindre les chiffres empiriques toucherait A3 et T9 et ferait du QE un levier de taux puissant là où le prototype a établi qu'il est un instrument d'actifs (propriété 8). Le cours l'écrit dans « ce que ce bloc ne sait pas faire ».
\end{lecture}
C'est un terme de déséquilibre monétaire (effet d'encaisses réelles) : il est nul à l'équilibre $M V = P Y$ et ne joue qu'à moyen terme ($\xi_M$ petit), mais il rend impossible une expansion durable de $M$ sans inflation. La vitesse dépend de l'inflation anticipée (Cagan 1956) : quand $\pi^e$ monte, la monnaie brûle les doigts, $V$ monte, et les prix montent plus que $M$.

\paragraph{QE (achats d'obligations).} La banque centrale achète $\Delta B^{CB}$ :
\begin{itemize}
\item auprès des banques : $\Delta Res^{ex} = \Delta B^{CB}$, $M$ inchangé ; effet uniquement par les rendements ($-s_7\, B^{CB}/B$ dans \eqref{eq:spread}), donc par $i^{long}$, $q_j$, $p^E$, $p^Z$ et $e$ (dépréciation via UIP) ;
\item auprès des ménages : $\Delta D_h = \Delta B^{CB}$, $M$ augmente et les ménages réallouent vers actions et immobilier (\eqref{eq:portfolio}) : effet richesse.
\end{itemize}
Les réserves excédentaires ne deviennent de la monnaie que si les banques prêtent ($\Delta L = \Delta D$), ce qui exige une demande de crédit ($q_j > 1$, $\Lambda^{Bk} > 0$). \emph{En trappe à liquidité, le QE gonfle $H$, pas $M$, et n'est pas inflationniste} (Japon 2001--06, zone euro 2015--19) : le modèle le reproduit parce que \eqref{eq:monterm} porte sur $M$. Sortie de la trappe avec des réserves excédentaires massives : le crédit repart, $M$ accélère, la banque centrale doit relever $i^{CB}$ ou vendre ses titres (pertes en capital pour elle et pour l'État).

\paragraph{Financement monétaire du déficit.} L'avance directe $\Delta A^G$ (ou le QE dont le produit est dépensé par l'État, ce qui revient au même pour $M$) crée des dépôts au fur et à mesure des dépenses publiques : $\Delta M = \Delta A^G$. Son effet passe par trois canaux simultanés : la demande (multiplicateur), le terme \eqref{eq:monterm}, et la crédibilité $c$ via $-\nu_3$ dans \eqref{eq:cred}.

\paragraph{Chemin vers l'hyperinflation (vérification).} Le seigneuriage réel que l'État peut extraire est $\sigma^{mon} = \pi\, m(\pi)$ avec $m = M/(PY) = V(\pi)^{-1}$, fonction en cloche (courbe de Laffer du seigneuriage) maximale en $\pi^{max} = 1/a$. Si le déficit à financer excède $\sigma^{max} = (ae)^{-1}V_0^{-1}$, il n'existe pas d'inflation stationnaire : le taux de croissance de $M$ doit augmenter sans cesse et l'inflation diverge (Cagan, Sargent--Wallace 1981). Dans le modèle, la séquence est : (1) $\Delta A^G > 0$ persistant $\Rightarrow$ $c \to 0$ et $\pi^e \uparrow$ par \eqref{eq:cred} et le terme $\upsilon$ de \eqref{eq:expect} ; (2) $V(\pi^e)\uparrow$ $\Rightarrow$ terme \eqref{eq:monterm} $\uparrow$ ; (3) degré d'indexation des salaires $\varpi_w \to 1$ (le coefficient de vitesse $\varpi$ des prix dans \eqref{eq:price} est un autre objet, non borné par 1) et raccourcissement des révisions (mensuelles puis hebdomadaires) au-delà de $\bar\pi^{hyper}$ ; (4) fuite devant la monnaie : la part des devises dans les portefeuilles \eqref{eq:portfolio} croît avec $\pi^e$ (terme $+a_4 \pi^e$ sur $s^{FX}_h$), $e$ se déprécie via \eqref{eq:fx}, augmenté d'un terme $\theta_5(\pi^e - \pi^{e\ast})$, et les prix importés répercutent ; (5) la base fiscale réelle s'érode (effet Olivera--Tanzi : recettes perçues avec retard sur des montants nominaux), le déficit réel augmente, donc $\Delta A^G$ augmente. Le point de non-retour est atteint lorsque le déficit réel dépasse $\sigma^{max}$. Sortie : arrêt du financement monétaire, ancre nominale (change fixe, nouvelle monnaie), consolidation budgétaire, ce qui reconstruit $c$ lentement.

\begin{design}[garantir que le QE \emph{puisse} être inoffensif et \emph{puisse} être catastrophique]
Les deux issues sont dans le modèle et dépendent des décisions : (i) QE sur titres détenus par les banques, en trappe à liquidité, avec banque centrale indépendante $\Rightarrow$ réserves oisives, légère hausse des prix d'actifs, pas d'inflation ; (ii) QE massif finançant un déficit primaire persistant, banque centrale dépendante $\Rightarrow$ perte de $c$, hausse de $V$, dépréciation, spirale. Le test 4 de la \S\ref{sec:tests} vérifie les deux cas.
\end{design}

%=====================================================================
\section{L'État : finances publiques}
\label{sec:etat}
%=====================================================================

\begin{joueur}[politique budgétaire]
\textbf{Recettes} : taux d'imposition du travail par strate $\tau_{W,h}$, du capital $\tau_K$, des sociétés $\tau_\Pi$, TVA/taxes à la consommation par bien $\tau^C_j$ (comme les taxes sur biens de Victoria~3), tarifs $\tau^M_{ij}$, capitation $T^{poll}$, dividendes des participations $\theta_j$, privatisations, seigneuriage $\Sigma$.\\
\textbf{Dépenses} : administration et salaires publics $G^{adm}$, éducation $G^{edu}$, santé $G^{san}$, transferts $Tr$ (chômage, retraites, minima), investissement public $G^{inv}$, subventions sectorielles $Sub_j$, R\&D publique, défense, intérêts $i^B B$.\\
\textbf{Dette} : montant et devise d'émission, maturité, choix de restructuration ; \textbf{nationalisations / privatisations} (\S\ref{sec:systemes}).
\end{joueur}

\subsection{Recettes et capacité administrative}
\begin{equation}
\begin{split}
T_t = \Big[\sum_h \tau_{W,h} W L_h + \tau_K \Pi^{distr}_h + \tau_\Pi \sum_j \pos{\Pi^{avant}_j} + \sum_j \tau^C_j p_j Q^{cons}_j + \text{Tarifs} + T^{poll} N \Big]\\
\cdot \varepsilon^{adm}_t \cdot \big(1 - \varepsilon^{ev} \bar\tau_t^{\,2}\big) + \sum_j \theta_j Div_j + \Sigma_t .
\end{split}
\end{equation}
\begin{lecture}[recettes publiques et courbe de Laffer]
\variables $T_t$ : recettes totales ; $\tau_{W,h}WL_h$ : impôt sur les salaires ; $\tau_K\Pi^{distr}_h$ : impôt sur les revenus du capital distribués ; $\tau_\Pi$ : impôt sur les sociétés ; $\tau^C_j p_jQ^{cons}_j$ : taxes à la consommation (TVA) sur les ventes aux ménages ; Tarifs : droits de douane ; $T^{poll}N$ : capitation ; $\varepsilon^{adm}_t\in[0,1]$ : efficacité de la collecte, fonction croissante et saturante des dépenses d'administration $G^{adm}/Y$ ; $\varepsilon^{ev}\bar\tau^2_t$ : part des recettes perdue par évasion et désincitation, croissant avec le carré du taux moyen $\bar\tau$ ; $\theta_jDiv_j$ : dividendes des entreprises publiques ; $\Sigma_t$ : seigneuriage (profit de la banque centrale reversé).
\sens Les recettes sont l'assiette multipliée par le taux, corrigée de deux réalités : on ne collecte bien que si l'administration fiscale est financée (le rôle de la bureaucratie dans Victoria~3), et des taux trop élevés font fuir l'assiette (travail au noir, exil fiscal, moindre effort). La courbe de Laffer en découle : les recettes augmentent avec le taux, atteignent un maximum, puis diminuent. Où se situe ce maximum est un objet de débat empirique (souvent estimé entre 60 et 70\,\% pour le taux marginal supérieur) ; dans le jeu, c'est $\varepsilon^{ev}$ qui le fixe.
\hyp Évasion fonction du taux moyen agrégé ; pas de concurrence fiscale explicite entre pays (elle est implicite via la mobilité du capital $KA$).
\limites La forme quadratique est une commodité ; la vraie courbe dépend du type d'impôt (la TVA fuit moins que l'impôt sur le capital).
\end{lecture}
L'efficacité de collecte $\varepsilon^{adm} = 1 - \exp(-\upsilon^{adm} G^{adm}/Y)$ joue le rôle de la bureaucratie de Victoria~3 : des institutions sous-financées collectent mal. Le terme d'évasion, quadratique dans le taux moyen $\bar\tau$, produit une courbe de Laffer.

\subsection{Effets réels de la dépense publique}
\begin{itemize}
\item $G^{inv}$ accumule le capital public : $K^G_{t+1} = (1-\delta_G)K^G_t + G^{inv}_t/p_K$, qui élève la productivité de tous les secteurs via \eqref{eq:tfp}.
\item $G^{edu}$ élève le capital humain : $H_{t+1} = H_t\big(1 + \eta_H \ln(1 + G^{edu}_t/(P_t N_t \bar g^{edu}))\big)$ avec un long délai (cohortes). Paramètres (v1.5) : $\eta_H=0{,}01$ par an, $\bar g^{edu}=5\,\%$ du PIB par tête (Mankiw, Romer et Weil 1992 : élasticité du produit au capital humain de l'ordre d'un tiers, atteinte par une dépense d'éducation d'une génération). \textbf{Spécifié, non simulé} : le prototype garde $H=1$ ; l'effet n'est visible qu'à quinze ou vingt ans, horizon qu'aucun test actuel ne mesure, et le levier $G^{edu}$ sera implémenté avec l'étape~5.
\item $G^{san}$ réduit la mortalité et augmente $SoL$ directement.
\item $G^{adm}$ et $G^{def}$ emploient $L^G$ (demande sur le marché du travail, salaires publics indexés).
\item Les transferts $Tr$ à la strate $B$ ont le multiplicateur le plus élevé ($c_B \approx 0{,}95$).
\item Les achats publics $G_j$ sont des ordres d'achat sur les marchés des biens.
\end{itemize}

\subsection{Taux apparent, maturité et dette consolidée}
\label{sec:iapp}
La dette n'est pas renouvelée chaque période au taux courant : seule la part qui arrive à échéance ou qui est nouvellement émise l'est. Le taux \emph{apparent} $i^{app}$ (intérêts payés rapportés au stock) suit le taux de marché $i^B$ avec une mémoire d'autant plus longue que la maturité moyenne $\bar T$ est élevée :
\begin{equation}
i^{app}_{t+1} = (1-\theta_t)\, i^{app}_t + \theta_t\, i^B_t, \qquad
\theta_t = \frac{1}{\bar T} + \frac{\pos{\text{Déficit}_t}}{B_t}, \qquad
\theta^{cons}_t = (1-\omega_t)\,\theta_t + \omega_t, \quad \omega_t = \frac{B^{CB}_t}{B_t},
\label{eq:iapp}
\end{equation}
et la charge d'intérêt du Trésor est $i^{app}_t B_t$. La maturité moyenne d'émission $\bar T$ est un curseur du joueur ; l'allonger coûte une prime de terme proportionnelle, $s_0\bar T/\bar T_0$ dans \eqref{eq:spread}.
\begin{lecture}[taux apparent à mémoire et maturité consolidée]
\variables $i^{app}$ : taux moyen effectivement payé sur le stock de dette ; $i^B$ : taux de marché sur les émissions nouvelles ; $\theta_t$ : fraction du stock refinancée au taux courant pendant l'année, somme de la part qui arrive à échéance ($1/\bar T$) et de la part nouvellement émise pour couvrir le déficit ; $\bar T$ : maturité moyenne d'émission, en années ; $\omega$ : part de la dette détenue par la banque centrale ; $\theta^{cons}$ : fraction de la dette \emph{consolidée} (Trésor et banque centrale réunis) qui se reprice chaque année.
\sens Une hausse du taux de marché ne se transmet à la charge de la dette qu'au rythme du refinancement : avec $\bar T=8$ ans et un déficit de 4\,\% du stock, $\theta\approx0{,}16$, il faut six ans pour que la moitié de la hausse soit répercutée. C'est ce qui rend une dette élevée supportable à court terme et dangereuse à long terme : la \og boule de neige\fg{} de \eqref{eq:debtdyn} est nulle aujourd'hui et devient dominante à dix ans par pur refinancement. La seconde formule dit que la part de la dette rachetée par la banque centrale a \emph{perdu} sa maturité : la banque centrale l'a financée par des réserves rémunérées au jour le jour au taux $i^{res}$ ; pour l'État consolidé, cette part se reprice instantanément. Un pays dont la banque centrale détient la moitié de la dette a $\theta^{cons}\approx0{,}55$ même si son Trésor émet à neuf ans : le QE a raccourci la maturité effective de la dette publique, et chaque point de hausse du taux directeur coûte immédiatement $\omega\, b$ points de PIB.
\hyp Maturité uniforme (pas de courbe complète des échéances) ; la dette détenue par la banque centrale est intégralement financée par des réserves rémunérées.
\limites La formule ignore les rachats anticipés et l'indexation sur l'inflation ; une part de dette indexée se reprice elle aussi immédiatement et s'ajoute à $\omega$ dans $\theta^{cons}$.
\end{lecture}

\subsection{Contrainte budgétaire et dynamique de la dette}
\begin{align}
\text{Déficit}_t &= G_t + Tr_t + Sub_t + i^{app}_t B_t - T_t ,\\
B_{t+1} &= B_t + \text{Déficit}_t - \Delta A^G_t - \Delta M^G_t , \label{eq:gbc}\\
\Delta b &= \frac{i^{app} - \pi - g}{(1+\pi)(1+g)}\, b - pb - \sigma^{mon},\qquad \sigma^{mon} = \frac{\Delta A^G}{PY}. \label{eq:debtdyn}
\end{align}
\begin{lecture}[contrainte budgétaire et dynamique de la dette]
\variables Déficit$_t$ : dépenses totales (achats $G$, transferts $Tr$, subventions $Sub$, intérêts $i^{app}B$ au taux apparent de \eqref{eq:iapp}) moins recettes $T$ ; $B_t$ : dette publique ; $\Delta A^G_t$ : avances de la banque centrale (financement monétaire) ; $\Delta M^G_t$ : variation du compte du Trésor ; $b$ : ratio dette/PIB ; $g$ : croissance réelle ; $\pi$ : inflation ; $pb$ : solde primaire en \% du PIB (positif si excédent) ; $\sigma^{mon}$ : seigneuriage en \% du PIB.
\sens Ce qui n'est pas couvert par l'impôt est emprunté ou créé. L'équation \eqref{eq:debtdyn} (Domar 1944, Blanchard 1990) est la plus importante des finances publiques : le ratio dette/PIB augmente de lui-même au rythme de l'écart entre le taux d'intérêt réel \emph{apparent} ($i^{app}-\pi$, celui effectivement payé sur le stock) et la croissance ($g$) : si l'État emprunte à 5\,\% et que l'économie croît de 2\,\% avec 2\,\% d'inflation, la dette gonfle de 1\,\% de son montant chaque année sans aucun nouveau déficit (effet \emph{boule de neige}). Un excédent primaire ou la création monétaire la réduisent. Quand $i^B-\pi<g$ (comme dans les années 2010), la dette se résorbe toute seule et un déficit primaire est soutenable ; quand les taux montent au-dessus de la croissance (années 1980, 2022), il faut des excédents. Comme le taux $i^B$ dépend lui-même de $b$ (équation \eqref{eq:spread}), il existe un seuil au-delà duquel la boule de neige devient incontrôlable : c'est la limite de crédit endogène, plus réaliste que le plafond fixe de Victoria~3.
\hyp Depuis la v0.6, le taux qui entre ici est le taux apparent $i^{app}$, qui ne suit le taux de marché qu'au rythme du refinancement \eqref{eq:iapp} : une hausse de $i^B$ ne devient une hausse de la boule de neige qu'après plusieurs années, sauf pour la part de dette détenue par la banque centrale.
\limites L'inflation « allège » la dette dans cette équation ; en réalité les créanciers l'anticipent et exigent des taux plus élevés (c'est le rôle de $\pi^e$ dans $i^B$ via $i^{CB}$), de sorte que seule l'inflation \emph{surprise} allège vraiment.
\end{lecture}
L'équation \eqref{eq:debtdyn} est la dynamique de Domar--Blanchard : la dette est soutenable si $i^B - \pi < g$ ou si le solde primaire compense l'effet boule de neige. Comme $i^B$ dépend de $b$ (éq.~\eqref{eq:spread}), il existe un seuil au-delà duquel la dynamique devient explosive : c'est la limite de crédit \emph{endogène}, plus riche que le plafond fixe de Victoria~3.

\paragraph{Ajout de la v1.3 : règle budgétaire de l'état de référence.} Les achats publics, l'investissement public et les minima sont multipliés par un facteur qui suit l'écart à une cible de dette :
\begin{equation}
\varphi^{fisc}_t = \clip{1 + \phi_b\,(b^\ast - b_t)}{0{,}6}{1{,}6}, \qquad \phi_b = 0{,}3,\quad b^\ast = 0{,}6 .
\label{eq:fiscrule}
\end{equation}
C'est la règle de l'IA budgétaire et celle de l'état de référence ; le joueur peut la désactiver. Ce que le prototype a montré : sans elle, l'excédent structurel de la calibration fait fondre la dette (0,61 $\to$ 0,04 en soixante ans) et prive la banque centrale de titres ; avec, la dette se stabilise entre 0,40 et 0,45 pour une correction moyenne de 6\,\%. Une première version rapportait la dette à cinquante-deux fois le PIB annuel et poussait la dette à 1,1 : la première chose qu'une règle budgétaire doit vérifier est l'unité de sa cible.

\paragraph{Précision de la v1.2 : qui achète la dette.} La banque n'absorbe la dette émise que jusqu'à 40\,\% de son bilan (limite de concentration) ; au-delà, la demande d'obligations des ménages devient \emph{sensible au rendement} — ils arbitrent leurs dépôts vers les titres à hauteur de $\kappa^{hB}(i^B-i^D)$, $\kappa^{hB}=20$ par point d'écart, jusqu'à 5\,\% de leurs dépôts par semaine — puis, si cela ne suffit pas, les dépenses sont rationnées. L'audit a montré les deux erreurs symétriques : une banque qui absorbe sans limite finance toute dérive budgétaire en hyperinflation ; un plafond sans demande élastique transforme une dette de 0,9 PIB en dépression par rationnement (le cas $\alpha=0{,}23$ du test de robustesse passait de 5 à 26\,\% de chômage). La demande élastique est le marché : la prime souveraine \eqref{eq:spread} monte, les ménages achètent, l'État paie.

\paragraph{Ajouts de la v0.9 : contrainte de trésorerie de l'État.} Les dépenses de la semaine (achats, investissement, salaires publics, transferts) sont exécutées au taux $\varrho^G=\min(1,\,T^{disp}/D^{prev})$, où $T^{disp}$, la trésorerie disponible, est le dépôt du Trésor plus les recettes attendues plus les avances de la banque centrale ; les intérêts impayés sont capitalisés dans la dette. La banque n'absorbe la dette émise que si elle est solvable (pondération souveraine nulle : elle n'est pas contrainte par $\kappa^{CAR}$ sur ces titres) ; sinon l'État se finance directement sur l'épargne des ménages jusqu'à 5\,\% de leurs dépôts par semaine, au taux $i^B$, et réduit ses dépenses si cela ne suffit pas. C'est l'écriture opérationnelle de la limite de crédit endogène de \S\ref{sec:etat} ; le prototype a montré que sans elle, des transferts indexés sur les salaires étaient financés sans limite pendant une hyperinflation et que l'arrêt des avances ne stoppait rien.

\subsection{Placement de la dette et défaut}
La demande d'obligations à chaque adjudication est $B^{dem} = \sum_h s^B_h S_h + \Delta B^{Bk} + \Delta B^{CB} + (1-\chi^K)\,B^{\ast,dem}(i^B - i^\ast - \rho^{pays})$. Si $B^{dem} < $ émission, le taux monte jusqu'à équilibre ; si aucun taux borné ($i^B \le \bar\imath$) ne permet le placement et que le financement monétaire est interdit ou refusé, l'État est en \textbf{défaut} :
\begin{equation}
B \leftarrow (1-\hbar)\,B, \qquad D^{hist} \leftarrow D^{hist} + 1, \qquad \text{pertes } \hbar B^{Bk} \text{ (banque)}, \ \hbar B_h \text{ (ménages)},
\label{eq:default}
\end{equation}
\begin{lecture}[défaut souverain]
\variables $\hbar\in(0,1)$ : décote (\emph{haircut}), part de la dette non remboursée ; $B$ : dette ; $D^{hist}$ : compteur de défauts, qui alimente la prime de risque ; $B^{Bk}, B_h$ : dette détenue par la banque et par les ménages, qui subissent la perte ; $T^{exc}$ : durée d'exclusion des marchés externes.
\sens Quand plus personne ne veut prêter à un taux supportable et que la création monétaire est exclue, l'État ne rembourse pas tout. Il efface une partie de sa dette, mais paie trois fois : ses créanciers domestiques (banque et ménages) perdent de la richesse, sa réputation le fera emprunter plus cher pendant des années, et l'étranger lui ferme ses marchés (Eaton et Gersovitz 1981). Si la banque domestique détient beaucoup de dette publique, le défaut peut la ruiner et déclencher une crise bancaire : la « boucle diabolique » de la zone euro en 2011.
\hyp Défaut partiel unique ; négocié (décote et exclusion modérées) ou unilatéral (fortes).
\limites Les procédures réelles (clubs de Paris et de Londres, FMI) sont longues et conditionnelles ; le modèle les résume par la décote et la durée d'exclusion.
\end{lecture}
avec exclusion des marchés externes pendant $T^{exc}$ mois.
La décote $\hbar$ est choisie par le joueur (restructuration négociée) ou imposée (défaut unilatéral, décote plus forte et exclusion plus longue). Un défaut sur la dette détenue par la banque domestique peut la mettre en faillite (\S\ref{sec:banques}) : boucle dette souveraine--banques. La dette en devise étrangère $B^\ast$ (\og péché originel\fg) est réévaluée par $e$ : une dépréciation gonfle mécaniquement $b$.

\subsection{Le plafond de crédit \og à la Victoria 3\fg{} comme garde-fou}
Pour la lisibilité du joueur, on affiche un \emph{plafond indicatif} $\bar B = \varkappa_1 \sum_j M_j + \varkappa_2 \sum_h V_h + \varkappa_3 PY$ (épargne privée mobilisable). Il n'est pas contraignant en soi mais borne $B^{dem}$ domestique : au-delà, seul le placement extérieur ou la monétisation restent possibles.

%=====================================================================
\section{Systèmes économiques : du libéralisme à la planification}
\label{sec:systemes}
%=====================================================================

Le joueur n'a pas à choisir un \og camp\fg{} au départ : il déplace des curseurs institutionnels, chacun ayant un coût de transition et des effets sur les équations précédentes.

\begin{center}\small
\begin{tabular}{p{3.6cm}p{1.6cm}p{9.6cm}}
\toprule
\textbf{Curseur} & \textbf{Symbole} & \textbf{Effet mécanique}\\
\midrule
Propriété publique du capital, par secteur & $\theta_j \in[0,1]$ & Part des dividendes allant au budget ; part de l'investissement décidée par l'État ; au-delà de $\theta_j > 0{,}5$ le secteur est \og public\fg{} et obéit aux directives.\\
Prix et salaires administrés & $\varrho^p_j \in\{0,1\}$ & Le prix $p_j = \bar p_j$ (ou, pour le travail, la grille salariale) est fixé par le joueur ; l'éq.~\eqref{eq:price} (resp. \eqref{eq:wage}) est désactivée ; pénurie ou excédent physiques.\\
Planification quantitative & $\varrho^{plan}\in[0,1]$ & Part de la production allouée par plan (\S\ref{sec:plan}) ; contrainte budgétaire lâche pour les secteurs publics.\\
Ouverture commerciale & $\tau^M$, quotas & Tarifs, licences, embargos.\\
Contrôle des capitaux, convertibilité & $\chi^K$ & Voir \S\ref{sec:change}.\\
Indépendance de la banque centrale & binaire & Autorise/interdit $\Delta A^G$ ; effet sur $c$.\\
État-providence & $Tr$, $\varrho$ & Niveau des transferts et de l'assurance chômage.\\
Régulation bancaire & $\kappa^{CAR}$, $\varrho^{res}$ & Levier bancaire maximal.\\
Règle budgétaire constitutionnelle & $\varrho^{fisc}$ & Plafond de déficit structurel, frein à l'endettement ou prélèvement plafonné sur le fonds souverain ; engagement difficilement réversible (délai et coût politique de sortie $\lambda^{fisc}$) ; abaisse $\rho^{pays}$ et conditionne le backstop en régime B ; critiques : bride $G^{inv}$ et repose sur une estimation contestable du cycle ($Y^{pot}$).\\
Maturité d'émission & $\bar T$ & Mémoire du taux apparent \eqref{eq:iapp} ; allonger coûte une prime de terme proportionnelle ($s_0\bar T/\bar T_0$).\\
Rémunération des réserves & $\Delta^{res}$ & Écart entre taux directeur et taux des réserves ; transfert aux banques contre transmission du taux directeur ; pèse sur $\Sigma$ \eqref{eq:ecb}.\\
Contrainte d'usage de la monnaie nationale & $\varrho^{dedol}$ & Cours légal, impôts en monnaie nationale, cession des devises, affichage ; abaisse le plancher de $s^{FX}$ \eqref{eq:sfx} ; coût : marché noir des devises ($\varsigma^{bm}$), perte d'activité formelle ($-\varepsilon^{dedol}\varrho^{dedol}$ sur $A$).\\
Répression financière & $m^D$, $\varrho^B$ & Marge sur dépôts administrée (taux des dépôts maintenu sous le taux directeur) et détention obligatoire de dette publique par les banques ($\varrho^B$ de leurs actifs) ; transfert explicite des épargnants vers les emprunteurs $\mathcal{T}^{rep}_t = (i^{CB}_t - i^D_t)\,D_t$, dont une part $\varrho^B$ va au Trésor via $-s_9\varrho^B$ dans \eqref{eq:spread} ; c'est un impôt invisible.\\

\bottomrule
\end{tabular}
\end{center}

\subsection{Nationalisation et privatisation}
\begin{equation}
\text{Nationaliser } \Delta\theta_j : \quad \text{coût} = \Delta\theta_j\, V^E_j\,(1 - \varsigma^{exp}) , \qquad
\varsigma^{exp} \in [0,1] \text{ : taux d'expropriation}.
\end{equation}
\begin{lecture}[nationalisation et privatisation]
\variables $\Delta\theta_j$ : part du secteur transférée à l'État ; $V^E_j$ : valeur de marché du secteur ; $\varsigma^{exp}\in[0,1]$ : taux d'expropriation (0 : indemnisation intégrale au prix du marché ; 1 : confiscation) ; $\rho^{pays}$ : prime de risque pays ; $\Lambda_j$ : accès au crédit des secteurs privés restants.
\sens Nationaliser en indemnisant coûte le prix du marché ; exproprier est gratuit sur le moment mais fait fuir les capitaux, renchérit la dette publique, dissuade l'investissement privé restant (chacun craint d'être le prochain) et peut provoquer des représailles. Privatiser rapporte une recette exceptionnelle mais l'État perd les dividendes futurs : c'est un arbitrage entre trésorerie immédiate et revenu de long terme, souvent mal fait sous la pression budgétaire.
\hyp La valeur de marché est celle de l'équation \eqref{eq:portfolio} ; les effets réputationnels décroissent avec le temps.
\limites Les nationalisations réelles changent aussi la gestion (objectifs d'emploi plutôt que de profit), ce que le modèle capture par la pénalité de planification quand $\theta_j$ devient majoritaire.
\end{lecture}
Une expropriation ($\varsigma^{exp} > 0$) économise le budget mais déclenche : fuite des capitaux ($\rho^{pays} \uparrow$, $KA<0$), hausse durable de $\rho^{pays}$ et du \emph{spread} \eqref{eq:spread}, baisse de $\Lambda_j$ dans les secteurs privés restants (peur de l'expropriation réduit l'investissement), représailles commerciales possibles des autres joueurs. Une privatisation est une recette exceptionnelle égale à $\Delta\theta_j V^E_j$ (décotée si le marché est peu profond).

\subsection{Planification input--output}
\label{sec:plan}
Le joueur planificateur fixe des cibles de \emph{demande finale} $d^{plan}$ (consommation par tête par strate, investissement par secteur, dépenses de l'État, exportations). Les productions brutes nécessaires sont données par le système de Leontief :
\begin{equation}
x^{plan} = (I - \mathbf{A})^{-1}\, d^{plan}, \qquad \mathbf{A} = [a_{kj}],
\label{eq:leontief}
\end{equation}
\begin{lecture}[planification par le tableau de Leontief]
\variables $d^{plan}$ : vecteur des demandes finales visées par le plan (consommation, investissement, État, exportations, pour chaque bien) ; $\mathbf{A}=[a_{kj}]$ : matrice des coefficients techniques (ce qu'il faut de $k$ pour produire une unité de $j$) ; $x^{plan}$ : vecteur des productions brutes nécessaires ; $(I-\mathbf{A})^{-1}$ : inverse de Leontief, qui additionne les besoins directs et tous les besoins indirects en cascade ; $Y^{cap}_j$ : capacité de chaque secteur.
\sens Pour livrer 100 voitures, il faut de l'acier ; pour produire cet acier il faut du charbon ; pour extraire ce charbon il faut des machines, donc encore de l'acier... L'inverse de Leontief (1941) résout d'un coup cette cascade infinie et donne la production brute que chaque secteur doit atteindre. C'est l'outil même du Gosplan soviétique et des planificateurs français de l'après-guerre. Le rapport entre production requise et capacité mesure la \emph{tension du plan} : au-delà de 1, le plan est physiquement infaisable et il faut choisir des priorités, ce que le joueur fait en ordonnant les secteurs.
\hyp Coefficients fixes ; information parfaite du planificateur sur les capacités (c'est cette hypothèse que la pénalité d'inefficience remet en cause).
\limites Le calcul suppose que les secteurs exécutent le plan ; le modèle introduit ensuite les écarts (pénuries, files, marché noir).
\end{lecture}
puis confrontées aux capacités \eqref{eq:prod}. L'écart $\max_j\, x^{plan}_j / Y^{cap}_j$ mesure la \emph{tension du plan}. Le plan alloue la main-d'œuvre ($L_j$ par directive, plein emploi mécanique), les intrants (priorité aux secteurs stratégiques) et l'investissement (le taux d'investissement $I/Y$ est une variable de décision, ce qui permet la croissance \og extensive\fg{} par accumulation forcée).

\paragraph{Contrainte budgétaire lâche (Kornai).} Les secteurs publics ne font pas faillite : $Sub_j = \pos{-\Pi_j}$ est automatiquement inscrit au budget. En contrepartie, les pertes cumulées ne disciplinent plus la production, et l'incitation à l'efficacité disparaît.

\paragraph{Pénalités d'inefficience.} 
\begin{equation}
\varepsilon^{plan}_{j,t} = \varepsilon_1\, \varrho^{plan}\,\theta_j + \varepsilon_2\, \varrho^{plan}\, \ln(|\mathcal{J}|\cdot \text{complexité}_t) + \varepsilon_3\, \max(0, \text{tension}_t - 1),
\end{equation}
\begin{lecture}[pénalité d'inefficience de la planification]
\variables $\varepsilon^{plan}_{j,t}$ : perte annuelle de croissance de la productivité ; $\varrho^{plan}\in[0,1]$ : degré de planification quantitative ; $\theta_j$ : part publique du secteur ; $|\mathcal{J}|\cdot$complexité : nombre de biens multiplié par un indice de sophistication de l'économie ; tension : rapport entre plan et capacité ; $\varepsilon_1,\varepsilon_2,\varepsilon_3$ : poids ; $\varepsilon_4$ : perte de taux d'utilisation des capacités.
\sens Trois sources d'inefficience, chacune documentée. (1) Sans profit ni faillite, l'incitation à améliorer disparaît : c'est la \emph{contrainte budgétaire lâche} de Kornai (1980). (2) Le planificateur ne peut pas connaître les millions d'informations locales que les prix agrègent spontanément ; plus l'économie est complexe, plus il se trompe : c'est l'argument de Hayek (1945), qui explique pourquoi l'URSS a bien réussi l'acier et mal l'électronique. (3) Un plan trop tendu crée des goulots d'étranglement en cascade. Le message pour le joueur : la planification est un outil puissant pour une économie simple en rattrapage, et de plus en plus coûteux à mesure que l'économie se diversifie.
\hyp Pénalité additive sur la croissance de $A$, calibrée sur l'écart de croissance observé entre les deux blocs pendant la guerre froide (0,5 à 1,5 point par an).
\limites Ce sont des ordres de grandeur historiques, pas des lois ; la littérature (Allen 2003, Harrison) débat de la part respective des incitations, de l'information et de l'effort militaire dans le ralentissement soviétique.
\end{lecture}
qui s'applique à la croissance de $A_j$ dans \eqref{eq:tfp}, et un facteur $(1 - \varepsilon_4 \varrho^{plan})$ sur le niveau d'utilisation des capacités. Le second terme formalise le problème de l'information de Hayek : plus l'économie est complexe (nombre de biens, sophistication), plus la planification centralisée perd d'efficacité relative --- la planification est donc plus performante pour une économie simple en rattrapage (industrialisation lourde) que pour une économie diversifiée à la frontière.

\paragraph{Prix administrés, pénuries et marché noir.} Quand $\varrho^p_j = 1$ et $D_j > S_j$ : rationnement (files d'attente $q^{file}$ dans \eqref{eq:sol}), marché noir au prix $p^{bm}_j = \bar p_j\,(1 + \mu_{bm}\, (D_j/S_j - 1))$ absorbant une fraction $\varsigma^{bm}$ des quantités, manque à gagner fiscal, et \textbf{inflation réprimée} : l'écart $p^{bm}_j/\bar p_j$ est mémorisé et se libère brutalement si le joueur libéralise les prix (choc à la 1992). Quand $D_j < S_j$ : stocks invendus, gaspillage, subventions.

\paragraph{Monnaie passive.} En économie planifiée intégrale, le crédit est alloué par la banque d'État, les taux ne jouent plus de rôle allocatif ; la monnaie des ménages excédentaire (revenu non dépensable faute d'offre) s'accumule en \og surplomb monétaire\fg{} $\Omega = \sum_h D_h - \bar\upsilon\, P\, Q^{cons}$, qui alimente marché noir et inflation réprimée.

\subsection{Résumé des arbitrages}
\begin{center}\small
\begin{tabular}{p{4.2cm}p{5.2cm}p{5.2cm}}
\toprule
& \textbf{Économie de marché} & \textbf{Économie planifiée}\\
\midrule
Croissance de la PGF & Élevée (concurrence, rattrapage) & Réduite par $\varepsilon^{plan}$\\
Accumulation du capital & Décidée par $q$ ; sensible aux taux et au crédit & Décidée par le plan ; taux d'investissement libre\\
Emploi & Chômage cyclique ($u$ autour de $u^n$) & Plein emploi mécanique, sous-emploi caché\\
Prix & Inflation ouverte, ajustement continu & À prix administrés ($\varrho^p_j=1$) : prix stables, pénuries, files, marché noir, inflation réprimée. À prix libres ($\varrho^p_j=0$) : surcapacité, déflation, excédents exportés (v1.5)\\
Inégalités & Élevées (dividendes concentrés) & Faibles ($\theta_j$ élevé, grille salariale)\\
Stabilité financière & Cycles du crédit, crises bancaires et de change & Pas de crise financière ; crise de balance des paiements possible\\
Recettes de l'État & Impôts ; sensibles à l'évasion & Excédents des entreprises publiques ; pertes à couvrir\\
Résilience aux chocs & Ajustement par les prix & Ajustement par les quantités (files)\\
\bottomrule
\end{tabular}
\end{center}
Un point contre-intuitif (v1.5) : la ligne « Prix » dépend de $\varrho^p$, pas de $\varrho^{plan}$. La même contrainte budgétaire lâche (Kornai) produit la pénurie quand les prix sont administrés — l'excès de demande passe par les quantités — et la surcapacité déflationniste quand ils sont libres — l'excès d'offre subventionnée passe par les prix, et par les exportations en économie ouverte (configuration 5 de la \S\ref{sec:configs}). Le mécanisme est déjà dans le modèle ; c'est le tableau de la v1.4 qui ne décrivait que le premier cas. Aucune des deux colonnes ne domine ; les régimes hybrides (économie mixte, socialisme de marché : $\theta_j$ élevé mais $\varrho^p_j = 0$, ou dirigisme sectoriel) sont possibles et ont leurs propres combinaisons d'effets.


%=====================================================================
\section{Crises économiques et financières : mécanismes et vérification}
\label{sec:crises}
%=====================================================================
Les crises doivent émerger. Cette section vérifie, scénario par scénario, que les équations contiennent la chaîne causale complète et identifie les conditions (paramètres, institutions, décisions) sous lesquelles elle se déclenche. La figure~\ref{fig:crise} résume la structure commune.


\begin{figure}[htbp]
\centering
\resizebox{0.98\textwidth}{!}{%
\begin{tikzpicture}[x=1cm,y=1cm,every node/.style={align=center}]
% --- phase 1 : euphorie (rangée haute, gauche -> droite) ---
\node[crise] (c1) at (0,6)    {crédit abondant\\ ($i^{CB}$ bas, $\kappa^{CAR}$ faible)};
\node[crise] (c2) at (4.6,6)  {prix d'actifs $\uparrow$\\ \eqref{eq:pz}--\eqref{eq:pe}};
\node[crise] (c3) at (9.2,6)  {collatéral et richesse $\uparrow$\\ $\Lambda_j\uparrow$, $c^V V_h\uparrow$};
\node[crise] (c4) at (13.8,6) {levier $\ell\uparrow$,\\ extrapolation $\chi$};
\node[crise] (c5) at (18.0,6) {choc : $i^{CB}\uparrow$,\\ offre nouvelle, matières premières};
% --- phase 2 : retournement (rangée basse, droite -> gauche) ---
\node[crise] (c6) at (18.0,2.4) {prix d'actifs $\downarrow$\\ marges, LTV violés};
\node[crise] (c7) at (13.8,2.4) {défauts, $NPL\uparrow$\\ $E^{Bk}\downarrow$, ruées \eqref{eq:run}};
\node[crise] (c8) at (9.2,2.4)  {crédit rationné\\ $\Lambda^{Bk}\downarrow$, $I\downarrow$};
\node[crise] (c9) at (4.6,2.4)  {$\hat y\downarrow$, $u\uparrow$,\\ déflation $\pi<0$};
\node[crise] (c10) at (0,2.4)   {dette réelle $L/P\uparrow$,\\ taux réel $\uparrow$ à la ZLB};
% --- contraintes et sorties ---
\node[crise] (c11) at (13.8,-0.8) {étalon-or : $H\le G/\bar g$,\\ pas de prêteur en dernier ressort};
\node[crise,fill=green!10] (c12) at (4.6,-0.8) {sorties : QE, garantie des dépôts,\\ relance, dévaluation, défaut};
% --- flèches phase 1 ---
\draw[flux] (c1) -- (c2); \draw[flux] (c2) -- (c3); \draw[flux] (c3) -- (c4); \draw[flux] (c4) -- (c5);
\draw[flux] (c3.north) -- ++(0,1.3) -| node[lab,pos=0.25,above]{boucle Minsky (euphorie)} (c1.north);
% --- descente ---
\draw[flux] (c5) -- (c6);
% --- flèches phase 2 ---
\draw[flux] (c6) -- (c7); \draw[flux] (c7) -- (c8); \draw[flux] (c8) -- (c9); \draw[flux] (c9) -- (c10);
\draw[flux] (c10.north) -- ++(0,1.0) -| node[lab,pos=0.25,above]{boucle Fisher (déflation par la dette)} (c6.north);
% --- contraintes ---
\draw[flux] (c11) -- node[lab,pos=0.5,right]{aggrave} (c7);
\draw[flux,dashed] (c12) -- node[lab,pos=0.5,right]{atténue} (c9);
\end{tikzpicture}}
\caption{Anatomie d'une crise financière émergente : phase d'euphorie (Minsky, rangée haute), retournement puis déflation par la dette (Fisher, rangée basse). Chaque boîte renvoie à une équation du modèle ; aucun élément n'est scripté. Les sorties de crise sont des décisions des joueurs.}
\label{fig:crise}
\end{figure}


\subsection{Les ingrédients communs}
\begin{enumerate}
\item \textbf{Un actif à prix extrapolatif} (\eqref{eq:pz}, \eqref{eq:pe}) avec $\chi_Z, \chi_E > \chi^{crit}$ : la hausse passée alimente la demande présente. Le seuil critique est atteint lorsque le gain en capital extrapolé dépasse le coût d'usage, $\chi\, g^{passé} > i - \pi^e + \delta$ : bulle si le taux réel est bas et l'extrapolation forte.
\item \textbf{Un lien crédit--collatéral} : $\mathrm{LTV}$, marges $m$, levier $\bar\ell$ évalués aux prix de marché, ce qui rend $\Lambda_j$ et $L^{max}$ procycliques.
\item \textbf{Un système bancaire à levier} ($\kappa^{CAR}$) exposé aux mêmes actifs (hypothèques, prêts sur marge, obligations).
\item \textbf{Des dettes nominales} et une politique monétaire bornée (\S\ref{sec:zlb}).
\item \textbf{Un déclencheur} : retournement endogène (coût d'usage, offre nouvelle de logements après $T_Z$) ou décision (hausse de $i^{CB}$, durcissement de $\mathrm{LTV}$ ou de $m$).
\end{enumerate}

\subsection{Vérification des scénarios historiques}
\begin{center}\small
\begin{longtable}{p{2.2cm}p{4.0cm}p{5.7cm}p{3.0cm}}
\toprule
\textbf{Scénario} & \textbf{Conditions initiales et décisions} & \textbf{Chaîne d'équations} & \textbf{Signature attendue}\\
\midrule
\endfirsthead
\toprule
\textbf{Scénario} & \textbf{Conditions initiales et décisions} & \textbf{Chaîne d'équations} & \textbf{Signature attendue}\\
\midrule
\endhead
\textbf{1929 et Grande Dépression} & Étalon-or ($H \le G/\bar g$), pas de garantie des dépôts, marge $m \approx 0{,}9$, $\chi_E$ élevé, taux bas puis hausse de $i^{CB}$ (1928--29), pays créancier accumulant l'or sans le monétiser (stérilisation) & Bulle actions \eqref{eq:pe} à crédit sur marge $\Rightarrow$ hausse de $i^{CB}$ $\Rightarrow$ $r^E$ attendu $<$ coût $\Rightarrow$ ventes, appels de marge $\Rightarrow$ $p^E \downarrow$ $\Rightarrow$ $V_h \downarrow$, $E_H \downarrow$ \eqref{eq:cH} ; pertes bancaires $\Rightarrow$ ruées \eqref{eq:run} ; refinancement limité par l'or $\Rightarrow$ faillites, $M \downarrow$ $\Rightarrow$ déflation via \eqref{eq:price} et \eqref{eq:monterm} $\Rightarrow$ dette réelle $\uparrow$, $\Lambda_j \downarrow$ \eqref{eq:invest} $\Rightarrow$ $I \downarrow$, $u \uparrow$ ; pays en change fixe importent la déflation via \eqref{eq:armington} & Chute de $p^E$ de 60--80\,\%, de $M$ de 25--30\,\%, déflation cumulée de 20--25\,\%, $u > 20\,\%$ ; sortie plus rapide pour les pays qui dévaluent (abandon de l'or : $\bar g \uparrow$, $H$ libéré)\\
\midrule
\textbf{Japon 1986--2000} & Industriel mûr, $\chi_Z$ et $\chi_E$ élevés, crédit garanti par le foncier ($\mathrm{LTV}^{ent}$ élevé), $i^{CB}$ bas après appréciation du yen (accord du Plaza : $e \downarrow$ imposé), puis hausse rapide de $i^{CB}$ en 1989--90 & Boucle \eqref{eq:pz} $\to$ collatéral $\to$ $\Lambda_j$ $\to$ $I$ et \eqref{eq:pe} $\Rightarrow$ retournement par $i^{CB}$ $\Rightarrow$ $p^Z \downarrow$ sur dix ans (délai $T_Z$ et offre excédentaire) $\Rightarrow$ $NPL^Z$, $E^{Bk} \downarrow$ ; forbearance (\S\ref{sec:zlb}) $\Rightarrow$ zombies, $g^A \downarrow$ ; ZLB atteinte, déflation douce ; QE en réserves oisives (\S\ref{sec:qe}) & Prix du foncier $\div 2$ à $\div 3$, croissance $\approx 1\,\%$, déflation de 0 à $-1\,\%$, dette publique $b \to 2$ sans crise (épargne domestique, $c$ élevée, monnaie de réserve secondaire)\\
\midrule
\textbf{2008} & $\mathrm{LTV} \to 1$, $\kappa^{CAR}$ faible, garantie des dépôts (pas de ruée des déposants mais ruée sur le financement de marché : $L^{CB}$ et dette bancaire de court terme), $i^{CB}$ bas 2001--04 puis hausse 2004--06, $\chi_Z$ élevé & \eqref{eq:pz} $\Rightarrow$ offre de logements $I_Z$ après $T_Z$ $+$ hausse de $i^Z$ $\Rightarrow$ $uc \uparrow$, $p^Z \downarrow$ $\Rightarrow$ $NPL^Z$ $\Rightarrow$ $E^{Bk} < \kappa^{CAR} L$ $\Rightarrow$ $\Lambda^{Bk} \to 0$ (crise de crédit) $\Rightarrow$ $I \downarrow$, $E \downarrow$ (richesse) $\Rightarrow$ $\hat y \downarrow$ ; contagion externe par $B^\ast$ et Armington ; réponses : ZLB, QE (réserves oisives, $i^B \downarrow$), renflouement ($b \uparrow$) & Récession de 4--6\,\% du PIB, $u$ doublé, pas de déflation durable grâce au QE et à la relance ; $b$ $+30$ pts ; reprise lente\\
\midrule
\textbf{Crise de change / dette externe (1997, 2001)} & Petit pays ouvert, change fixe, $B^\ast$ élevé en devise étrangère, $R^{FX}$ faible, entrées de capitaux ($\chi^K = 0$) & Boom du crédit financé par $KA$ $\Rightarrow$ $CA < 0$ $\Rightarrow$ $R^{FX} \downarrow$ $\Rightarrow$ $\pi^{dev} \uparrow$ (\S\ref{sec:change}) $\Rightarrow$ sorties, attaque $\Rightarrow$ dévaluation forcée $\Rightarrow$ $B^\ast$ réévalué, bilans privés en devises $\Rightarrow$ défauts, $E^{Bk} \downarrow$ $\Rightarrow$ défaut souverain possible \S\ref{sec:etat} & Dévaluation de 40--70\,\%, récession de 5--10\,\%, inflation importée, puis reprise par les exportations\\
\midrule
\textbf{Hyperinflation (1923, 1989, 2008 Zimbabwe)} & Déficit primaire $> \sigma^{max}$, banque centrale dépendante, $\Delta A^G$ persistant, souvent dette en devises ou réparations & Séquence (1)--(5) de \S\ref{sec:qe} & $\pi$ mensuelle $> 50\,\%$, $V \times 10$, dollarisation, effondrement des recettes réelles ; stabilisation brutale possible par ancre nominale\\
\midrule
\textbf{Choc pétrolier (1973)} & Rentiers pétroliers formant un cartel (décision de joueurs : quota $\Rightarrow$ $S^w \downarrow$) & $p^w_{res} \uparrow$ $\Rightarrow$ coûts $c^u_j$ $\uparrow$ via intrants Leontief $\Rightarrow$ $p_j \uparrow$ \eqref{eq:price} et $\hat y \downarrow$ (stagflation) ; indexation $\varpi$ $\Rightarrow$ persistance ; dilemme de Taylor & Inflation $+5$ à $+10$ pts, récession, transferts de richesse vers les rentiers, recyclage des pétrodollars ($KA$)\\
\bottomrule
\end{longtable}
\end{center}

\subsection{Grille de couverture : dix configurations contemporaines}
\label{sec:configs}
Les dix configurations ci-dessous sont des archétypes de jeu anonymisés, pas des pays. Chacune est reproductible depuis ses conditions initiales, sans script, et \textbf{active au moins un mécanisme qu'aucune autre n'active} (colonne \og mécanisme couvert en propre\fg{} de la table de la \S\ref{sec:archetypes} et première ligne de chaque case). La colonne \og valeurs calculées\fg{} donne, avec la calibration de la \S\ref{sec:calib} : la fraction de refinancement $\theta$ et sa version consolidée $\theta^{cons}$ \eqref{eq:iapp}, l'exposition $\omega b$, la prime de \eqref{eq:spread} terme par terme (en points de base au-dessus de $i^{CB}$), la probabilité mensuelle de saut \eqref{eq:jump} lorsqu'elle s'applique, la boule de neige $(i^{app}-\pi-g)\,b/[(1+\pi)(1+g)]$ en points de PIB par an, et l'état stationnaire du chômage naturel \eqref{eq:un}, $u^{n\ast}=0{,}75\,u+1{,}25$ (\%). Toutes ces valeurs sont recalculables à la main : c'est leur fonction.

\begin{center}\footnotesize
\begin{longtable}{p{1.6cm}p{3.1cm}p{3.4cm}p{2.9cm}p{3.5cm}}
\toprule
\textbf{Config.} & \textbf{Conditions initiales} & \textbf{Chaîne d'équations activée} & \textbf{Signature attendue} & \textbf{Valeurs calculées}\\
\midrule
\endfirsthead
\toprule
\textbf{Config.} & \textbf{Conditions initiales} & \textbf{Chaîne d'équations activée} & \textbf{Signature attendue} & \textbf{Valeurs calculées}\\
\midrule
\endhead
\textbf{1. Union endettée} (B) & $b=1{,}18$, $T/Y=0{,}44$, $\varkappa^{dom}=0{,}5$, $pb=-3\,\%$, déficit total 5\,\%, $c=0{,}9$, $\omega=0{,}25$, $\bar T=8{,}5$, backstop actif, $i^{U}=2$, $\pi=1{,}8$, $g=0{,}7$, $u=7{,}3$ & Boule de neige \eqref{eq:debtdyn} via $i^{app}$ \eqref{eq:iapp} ; backstop conditionnel $-s_8$ et saut \eqref{eq:jump} ; blocage politique (soutien \S\ref{sec:politique} bloquant la consolidation) ; contagion & Boule de neige $\approx0$ aujourd'hui, positive et dominante à dix ans par pur refinancement ; effort de consolidation de 2,7 pts de PIB aujourd'hui à 5,0 pts à dix ans, soit \textbf{0,25 pt par année de report} ; pas de crise tant que le backstop tient ; saut si perte du backstop & $\bar b_i=1{,}19$ ($b<\bar b_i$) ; $\theta=0{,}160$, $\theta^{cons}=0{,}37$ ; $\omega b=0{,}29$ ; $\Delta b^{proj}=-0{,}23+3{,}0=2{,}8$ ; prime $=+106$ (terme, $\bar T=8{,}5$, réduite de $\psi_{QE}\omega=12{,}5\,\%$) $+0$ (conv. : $b^{mkt}=0{,}89<\bar b_i$) $+83$ (dérive) $+30$ ($c$) $-50$ (backstop) $=\mathbf{169}$ pb, $i^B=3{,}7$ ; $\Pr(\text{saut})<0{,}01\,\%$ avec backstop, 0,22\,\% sans ; boule de neige $-0{,}23$ (à $i^{app}=2{,}3$), $+0{,}6$ à $i^{app}\to3{,}0$ ; $u^{n\ast}=6{,}7$\\
\midrule
\textbf{2. Union vertueuse} (B) & $b=0{,}63$, $T/Y=0{,}42$, $pb=-1\,\%$, $c=0{,}9$, $\omega=0{,}25$, $\bar T=7$, $\varrho^{fisc}$ constitutionnel actionné (sortie partielle sectorielle), $u=3{,}5$ & Règle budgétaire \S\ref{sec:systemes} et son coût de sortie $\lambda^{fisc}$ ; $\omega$ maximal ; contrepartie de la contagion dans \eqref{eq:jump} (subit le saut des autres) ; $G^{inv}$ libéré & Prime minimale de l'union ; croissance $\approx g_0$ ; l'assouplissement de la règle relève $G^{inv}$ et $b$ sans saut ; intérêt objectif au backstop des autres & $\bar b_i=1{,}14$ ; $\theta=0{,}175$, $\theta^{cons}=0{,}38$ ; $\omega b=0{,}16$ ; $\Delta b^{proj}=-0{,}73+1{,}0=0{,}3$ ; prime $=+87$ (terme, réduite de $12{,}5\,\%$) $+0$ (conv.) $+8+30-50=\mathbf{76}$ pb, $i^B=2{,}76$ (contre $\approx60$ pb observés : l'écart de la v0.8 est résolu par la réécriture de \eqref{eq:spread}) ; $\Pr<0{,}01\,\%$ (0,5\,\% si un voisin saute) ; boule de neige $-0{,}73$ ; $u^{n\ast}=3{,}9$\\
\midrule
\textbf{3. Émetteur de réserve} (A) & $b=1{,}0$, $T/Y=0{,}27$, $\varkappa^{dom}=0{,}7$, $pb=-3{,}5\,\%$, déficit 6,5\,\%, $c=0{,}9$, $\omega=0{,}15$, $\bar T=6$, $\chi_E=0{,}7$, $\tau^M$ actifs, $i^{CB}=3$, $\pi=2{,}7$, $g=2$, $u=4{,}2$ & $-s_6$ et sa décroissance graduelle avec $b$ et $c$ ; bulle d'actifs extrapolative \eqref{eq:pe} ; tarifs \eqref{eq:armington} et représailles ; monnaie de réserve dans $p^w$ & Boule de neige négative qui \emph{réduit} la prime via $\Delta b^{proj}$, aucune contrainte immédiate ; érosion graduelle et non discrète du rabais ; $p^E/p^{E,f}$ au-delà de 1,5 ; inflation importée par les tarifs & $\bar b_i=0{,}62$ ; $\theta=0{,}23$, $\theta^{cons}=0{,}35$ ; $\omega b=0{,}15$ ; $\Delta b^{proj}=-1{,}24+3{,}5=2{,}3$ ; prime $=+79$ (terme, $\bar T=6$, réduite de $7{,}5\,\%$) $+29$ (conv. sur $b^{mkt}=0{,}85$, $\varkappa^{for}=0{,}35$) $+68$ (dérive) $-40$ ($s_6$) $+30=\mathbf{166}$ pb, $i^B=4{,}7$ ; saut n.a. (émetteur) ; boule de neige $-1{,}24$ ; $u^{n\ast}=4{,}4$\\
\midrule
\textbf{4. Dette monétisée en sortie de trappe} (A) & $b=2{,}3$, $\varkappa^{dom}=0{,}95$, $pb=-1{,}5\,\%$, déficit 6\,\%, $c=0{,}95$, $\omega=0{,}5$, $\bar T=9$, $i^{CB}=0{,}5$, $\pi=2{,}5$, $g=0{,}5$, $u=2{,}5$ & $\theta^{cons}$ \eqref{eq:iapp} ; bilan de banque centrale \eqref{eq:ecb} (coupon figé $\approx0{,}3\,\%$) ; crédibilité symétrique \eqref{eq:cred} (undershoot passé) ; détention domestique dans \eqref{eq:spread} & Boule de neige fortement négative (l'inflation efface la dette) ; chaque point de $i^{CB}$ coûte immédiatement $\omega b\approx1{,}15$ pt de PIB ; dividende $\Sigma$ nul dès la première hausse ; arbitrage : tolérer l'inflation ou préserver le pouvoir d'achat & $\bar b_i=0{,}80$ ; $\theta=0{,}137$, $\theta^{cons}=\mathbf{0{,}57}$ ; $\omega b=\mathbf{1{,}15}$ ; $\Delta b^{proj}=-4{,}9+1{,}5<0$ ; prime $=+96$ (terme, $\bar T=9$, réduite de $25\,\%$) $+34$ (conv. sur $b^{mkt}=1{,}15$ ; $\varkappa^{for}=0{,}10$, plancher 0,3 actif car $\chi^K\varrho^B=0$) $+0+15=\mathbf{145}$ pb, $i^B=1{,}95$ ; saut n.a. ; boule de neige $\mathbf{-4{,}9}$ ; $u^{n\ast}=3{,}1$\\
\midrule
\textbf{5. Planifié en déflation} (A, $\chi^K=1$) & $b=1{,}2$ (dont 0,6 d'entités publiques locales garanties), $T/Y=0{,}26$, $\varkappa^{dom}=0{,}95$, $pb=-4\,\%$, $c=0{,}85$, $\varrho^B=1$, $\varrho^{plan}$ et $\theta_j$ élevés, $\chi_Z=0{,}7$ puis retournement, $L_{rur}/N<\bar\ell_{rur}$, $\pi=-0{,}5$, $g=4{,}5$, $u=5{,}2$ & Kornai (subventions automatiques, \S\ref{sec:plan}) ; Fisher (\S\ref{sec:zlb}) ; $\varepsilon^{plan}$ dans \eqref{eq:tfp} ; répression financière $-s_9$ et $m^D$ administré ; post-Lewis dans \eqref{eq:wage} ; forbearance & Déflation par la dette ; taux réel qui monte pendant que la banque centrale assouplit ; croissance qui ralentit par pénalité d'inefficience ; pas de crise financière ouverte (répression) mais surplomb & $\bar b_i=0{,}59$ ; $\theta=0{,}20$, $\theta^{cons}=0{,}24$ ; $\omega b=0{,}06$ ; $\Delta b^{proj}=-1{,}7+4{,}0=2{,}3$ ; prime $=+97$ (terme) $+27$ (conv. sur $b^{mkt}=1{,}14$ ; $\varkappa^{for}=0{,}05$, plancher nul car $\chi^K\varrho^B=1$) $+69$ (dérive) $+45$ ($c$) $-150$ (répression) $=\mathbf{88}$ pb, $i^B=2{,}4$ (218 pb sous le plancher inconditionnel de la v1.1 ; $\approx30$ pb observés, l'écart résiduel étant la prime de terme calculée sur une trajectoire plate du taux directeur) ; saut n.a. ; boule de neige $-1{,}7$ ; $u^{n\ast}=5{,}15$\\
\midrule
\textbf{6. Refuge à la borne zéro} (C) & $b=0{,}38$, $T/Y=0{,}28$, $pb=+0{,}5\,\%$, $c=0{,}95$, $\bar T=10$, $i^{CB}=0$, $H\approx PY$, $R^{FX}\approx PY$, $E^{CB}/H=0{,}2$, $\underline\kappa^{CB}=-0{,}10$, $\pi=0{,}2$, $g=1{,}2$, $u=2{,}4$ & Défense d'un plafond (\S\ref{sec:change}) ; contrainte $E^{CB}/H$ \eqref{eq:ecb} ; effet de stock $\theta_6$ dans \eqref{eq:fx} ; borne zéro & Déflation importée ; interventions sans effet durable ; une appréciation de 30\,\% détruit 30 pts de PIB de $E^{CB}$ et fait franchir $\underline\kappa^{CB}$ : abandon du plafond par coût de bilan, pas par manque de munitions & $\bar b_i=0{,}48$ ; $\theta=\theta^{cons}=0{,}10$ ; $\omega b=0$ ; prime $=+143$ (terme, $\bar T=10$) $+15=\mathbf{158}$ pb, $i^B=1{,}6$ ; saut n.a. ; boule de neige $-0{,}3$ ; $E^{CB}/H$ après $+30\,\%$ : $-0{,}10$ (seuil atteint) ; $u^{n\ast}=3{,}05$\\
\midrule
\textbf{7. Rentier à fonds externe} (A) & $b=0{,}40$, $T/Y=0{,}42$, $Y_{res}/Y=0{,}2$, $\varsigma^F=1$, $\beta^{dom}=0$, $F=3{,}2\,PY$, règle de prélèvement 3\,\% de $F$, $c=0{,}95$, $i^{CB}=4$, $\pi=3$, $g=1{,}5$, $u=4$ & \eqref{eq:fonds} avec $\beta_F=0{,}7$ ; règle de prélèvement \S\ref{sec:systemes} ; stérilisation du change mais pas de la demande ; \eqref{eq:lbd} & Maladie hollandaise neutralisée sur $e$ ; surchauffe domestique persistante ($\hat y>0$) par le prélèvement de 9,6\,\% du PIB ; une année boursière à $-20\,\%$ coûte $0{,}7\times20\times3{,}2=45$ pts de PIB au fonds, quatre ans de prélèvements & $\bar b_i=0{,}96$ ; $\theta=0{,}143$ ; $\omega b=0$ ; prime $=+100+15=\mathbf{115}$ pb, $i^B=5{,}15$ ; saut n.a. ; boule de neige $-0{,}4$ ; écart-type annuel du rendement du fonds $\approx0{,}7\sigma(p^{E,w})+8\,\%\approx14\,\%$, soit $\pm45$ pts de PIB ; $u^{n\ast}=4{,}25$\\
\midrule
\textbf{8. Rentier à ancrage rigide} (E) & $b=0{,}30$, $T/Y=0{,}30$ (dont rente), $\varkappa^{dom}=0{,}6$, $pb=-2\,\%$ au prix courant de la ressource, $c=0{,}9$ (importée), $\bar T=6$, $\beta^{dom}=0{,}6$, $\varpi_E=0{,}3$, $R^{FX}/M=0{,}6$, $i^{CB}=i^{ancre}=4$, $\pi=2$, $g=3$, $u=3$ & $i^{CB}$ importé ; $\beta^{dom}$ élevé dans \eqref{eq:fonds} (bulle d'actifs locale) ; point mort de la rente ($p^w$ tel que $pb=0$) ; rupture d'ancrage par $\pi^{dev}$ (\S\ref{sec:change}) quand $R^{FX}/M\to\varpi_E$ & Cycle calé sur $p^w$ ; aucune arme monétaire ; le point mort monte avec les dépenses domestiques ; sous le point mort trois ans, les réserves passent sous $\varpi_E M$ et la parité tombe & $\bar b_i=0{,}43$ ; $\theta=\theta^{cons}=0{,}23$ ; $\Delta b^{proj}=-0{,}14+2{,}0=1{,}9$ ; prime $=+86+56+30=\mathbf{172}$ pb, $i^B=5{,}7$ ; $\Pr(\text{saut})=0{,}03\,\%$/mois au prix courant, 2\,\% si $pb=-8\,\%$ ; boule de neige $-0{,}14$ ; $u^{n\ast}=3{,}5$\\
\midrule
\textbf{9. Émergent à faible capacité fiscale} (C/D) & $b=0{,}35$ dont moitié en devises, $T/Y=0{,}10$, $\varkappa^{dom}=0{,}5$, $pb=-3\,\%$, $c=0{,}6$, $D^{hist}=0{,}5$, $\chi^K=0{,}5$, $\bar T=5$, remises $Rem/Y=8\,\%$, $L_{rur}/N=0{,}6$, $i^{CB}=9$, $\pi=7$, $g=4$, $u=6$ & $\bar b_i$ modulé par $T/Y$ dans \eqref{eq:spread} ; pré-Lewis (Phillips plate, \eqref{eq:migr}) ; remises dans $Tr_h$ et $CA$ ; dette en devises réévaluée par $e$ ; saut \eqref{eq:jump} & Prime de plusieurs centaines de pb à 35\,\% de dette ; une dépréciation de 10\,\% ajoute 1,75 pt à $b$ ; les remises soutiennent $SoL$ et $s^{FX}$ ; un saut tous les quatre ans en moyenne & $\bar b_i=\mathbf{0{,}17}$ ; $\theta=0{,}34$ ; $\Delta b^{proj}=+0{,}3+3{,}0=3{,}3$ ; prime $=+71$ (terme, $\bar T=5$) $+326$ (conv.) $+99+75+120=\mathbf{691}$ pb, $i^B=15{,}9$ ; $\Pr(\text{saut})=\mathbf{2{,}2\,\%}$/mois ; boule de neige $+0{,}3$ ; $u^{n\ast}=5{,}75$\\
\midrule
\textbf{10. Post-hyper\-inflation dollarisé} (D) & $b=0{,}9$ (arriérés), $T/Y=0{,}15$, $\varkappa^{dom}=0{,}3$, $c=0{,}2$, $M^{hist}=6$, $D^{hist}=2$, $s^{FX}=0{,}8$, nouvelle monnaie à couverture métallique, $\pi=5$, $g=3$, $u=9$ & $M^{hist}$ plafonnant $c$ à $1-\nu_9M^{hist}=0{,}28$ \eqref{eq:cred} ; hystérèse \eqref{eq:sfx} ; taux effectif \eqref{eq:reff} ; seigneuriage réduit à $(1-s^{FX})$ & Prix stabilisés à un chiffre ; reconquête monétaire impossible ($s^{FX}$ plancher $\varphi$) ; taux directeur sans effet réel ; exclusion des marchés ; seigneuriage détruit & $\bar b_i=0{,}17$ ; $\theta=0{,}22$ ; prime $>7\,000$ pb $\Rightarrow$ \textbf{exclusion} (financement concessionnel à $i^{app}=3$) ; $\Pr(\text{saut})=43\,\%$/mois si accès au marché ; $c\le0{,}28$ quelle que soit la performance ; $u^{n\ast}=8{,}0$\\
\bottomrule
\end{longtable}
\end{center}

\paragraph{Couverture et redondances.} Chaque configuration active un mécanisme propre : 1 le backstop et le blocage politique, 2 la règle constitutionnelle et la contrepartie de la contagion, 3 le rabais de réserve et les tarifs, 4 la maturité consolidée et la crédibilité symétrique, 5 Kornai et la répression, 6 la contrainte de bilan, 7 le placement externe et la règle de prélèvement, 8 le taux importé sans backstop et le point mort, 9 la capacité fiscale et les remises, 10 la mémoire monétaire et l'hystérèse. Deux recouvrements partiels subsistent et sont assumés : 1 et 2 partagent le régime B (ils sont les deux faces d'un même jeu coopératif), 7 et 8 partagent la rente (ils s'opposent précisément sur $\beta^{dom}$ et le régime). Aucune configuration n'est déductible d'une autre.

\subsection{Indicateurs d'alerte fournis au joueur}
Ratio crédit/PIB et son écart à la tendance (BRI), $p^Z/(\text{loyer})$ et $p^Z/Y^d$, $p^E/p^{E,f}$, levier bancaire $L/E^{Bk}$, part des prêts non performants, $CA/Y$, $R^{FX}$ en mois d'importations, $b$ et part de la dette en devises, écart $\pi - \pi^\ast$, $B^{CB}/B$, et pour l'investissement le rapport $q/q^\ast$ avec $q^\ast=1+g_0/(\iota\Gamma)$ (v1.5 : un $q$ de 1,1 sur un sentier à 2\,\% n'est pas un signal d'investissement). Ces indicateurs sont aussi les entrées des règles macroprudentielles que le joueur peut adopter ($\mathrm{LTV}$ contracyclique, coussin de fonds propres).

%=====================================================================
\section{Soutien politique et conditions de fin}
\label{sec:politique}
%=====================================================================
Chaque strate évalue le gouvernement :
\begin{equation}
\Lambda^{pol}_h = \lambda_0 + \lambda_1\, \Delta_{12}\ln SoL_h - \lambda_2\, u_h - \lambda_3\, |\pi_t - \pi^{ref}| - \lambda_4\, \tau_h - \lambda_5\, (\text{Gini} - \text{Gini}^{ref}_h)\cdot \sgn_h ,
\end{equation}
\begin{lecture}[soutien politique]
\variables $\Lambda^{pol}_h$ : soutien de la strate $h$ au gouvernement ; $\Delta_{12}\ln SoL_h$ : croissance du niveau de vie sur douze mois ; $u_h$ : chômage de la strate ; $|\pi_t-\pi^{ref}|$ : écart de l'inflation à un niveau jugé normal ; $\tau_h$ : pression fiscale sur la strate ; Gini$-$Gini$^{ref}_h$ : écart des inégalités à ce que la strate juge acceptable, avec un signe $\sgn_h$ qui dépend de la strate (la strate basse pénalise les inégalités, la strate haute pénalise leur réduction forcée) ; $\lambda_0..\lambda_5$ : poids ; $\varpi_h$ : poids politique de chaque strate selon le régime (suffrage universel ou censitaire).
\sens Les citoyens jugent le gouvernement sur ce qu'ils ressentent : leur niveau de vie s'améliore-t-il, ont-ils un emploi, les prix sont-ils stables, paient-ils trop d'impôts, la société est-elle juste selon leurs critères. Chaque strate pondère différemment. Le soutien agrégé conditionne la capacité à réformer (une réforme impopulaire prend plus de temps) et, sous un seuil, déclenche des troubles puis un changement de gouvernement.
\hyp Fonction linéaire ; mémoire courte (douze mois), conforme aux résultats de la science politique sur le vote économique.
\limites C'est une version très allégée des groupes d'intérêt de Victoria~3 ; on l'enrichira si le jeu politique doit devenir un objet en soi.
\end{lecture}
la strate haute pénalisant l'imposition du capital et l'inflation (érosion de la richesse), la strate basse le chômage et l'inégalité. Le soutien agrégé $\Lambda^{pol} = \sum_h \varpi_h \Lambda^{pol}_h$ (pondérations dépendant du régime politique) module la capacité à faire passer les réformes (délai $\propto 1/\Lambda^{pol}$) et, sous un seuil, déclenche troubles (baisse de $A$, grèves : $L_j$ réduit), puis changement de gouvernement (fin de partie ou passage de la main). Il s'agit d'une version allégée des groupes d'intérêt et des radicaux de Victoria~3.

\textbf{Score} : combinaison pondérée du niveau de vie moyen, de sa croissance, de la stabilité (variance de $\pi$ et $u$), de la position extérieure nette et du soutien politique ; des objectifs asymétriques par pays (rattrapage, hégémonie monétaire, égalité) permettent des parties non symétriques.

%=====================================================================
\section{Chocs et événements}
%=====================================================================
Les chocs exogènes sont peu nombreux et paramétrés ; l'essentiel de la dynamique est endogène.
\begin{itemize}
\item Productivité : $\varepsilon^A_{j,t} = \rho_A \varepsilon^A_{j,t-1} + \eta_t$, écarts-types faibles ; récoltes (alimentation) plus volatiles.
\item Matières premières : découverte / épuisement modifiant $A_{\text{énergie}}$ d'un pays.
\item Confiance : choc sur $\rho^{pays}$ ou sur $\Lambda^{Bk}$ (panique), déclenché avec une probabilité croissante en $\ell$ moyen et $b$.
\item Politique : changement de régime chez une IA, guerre commerciale, embargo.
\end{itemize}

%=====================================================================
\section{Analyse d'impact : que se passe-t-il quand une variable bouge ?}
\label{sec:impacts}
%=====================================================================
Cette section est le mode d'emploi du modèle et la base des tests de stabilité. Pour chaque variable importante, on décrit l'effet d'une \textbf{hausse} (\up) \emph{toutes choses égales par ailleurs} : d'abord les \textbf{effets directs} (les équations où la variable apparaît), puis les \textbf{effets de second tour} (les réactions des autres agents), et enfin ce qui change quand la variable \textbf{baisse} (\dn) lorsque l'effet n'est pas simplement symétrique. Une flèche entre parenthèses renvoie à l'équation qui porte l'effet. Les délais indiqués sont ceux attendus avec la calibration de la section~\ref{sec:calib} ; ils sont à vérifier au prototype.

\newcommand{\impact}[1]{\subsubsection*{#1}}
\newcommand{\direct}{\par\noindent\textbf{Effets directs (\up).} }
\newcommand{\second}{\par\noindent\textbf{Second tour.} }
\newcommand{\baisse}{\par\noindent\textbf{Baisse (\dn).} }

\subsection{Leviers de politique monétaire}

\impact{Taux directeur $i^{CB}$}
\direct Les taux bancaires $i^L, i^D$ montent d'autant \eqref{eq:spread} ; le taux réel $r^L$ monte ; $q_j$ baisse et l'investissement recule \eqref{eq:invest} ; la strate haute reporte sa consommation \eqref{eq:cH} ; le coût d'usage du logement monte, $D^Z$ et $p^Z$ baissent \eqref{eq:pz} ; la valeur des actions baisse (actualisation) \eqref{eq:pe} ; le change s'apprécie par l'afflux de capitaux \eqref{eq:fx} ; le taux des obligations monte et le service de la dette publique aussi.
\second La demande recule, $\hat y$ baisse, le chômage monte avec 6 à 12 mois de retard \eqref{eq:labour} ; la courbe de Phillips ralentit les salaires \eqref{eq:wage}, l'appréciation du change fait baisser les prix importés : l'inflation recule avec 12 à 18 mois de retard \eqref{eq:price}. La baisse des prix d'actifs réduit le collatéral ($\Lambda_j$) et la richesse ($c^V V_h$), ce qui amplifie la baisse de la demande : c'est pourquoi une hausse de taux après une bulle peut la faire éclater. Les emprunteurs à taux variable voient leur revenu disponible baisser \eqref{eq:yd}, les épargnants le voient monter : transfert des jeunes endettés vers les retraités épargnants. La dette publique se dégrade doublement (intérêts plus élevés, croissance plus faible) \eqref{eq:debtdyn}. Si la hausse est jugée crédible, $\pi^e$ baisse et la crédibilité $c$ remonte \eqref{eq:cred}, ce qui réduit à terme le coût de la désinflation.
\baisse Symétrique avec deux asymétries. À la borne zéro, on ne peut plus baisser : la trappe à liquidité (\S\ref{sec:zlb}). Et une baisse en présence d'anticipations extrapolatives ($\chi_Z,\chi_E$ élevés) nourrit une bulle plus sûrement qu'une hausse ne la dégonfle : l'effet est asymétrique dans le temps (la bulle met des années à se former, des mois à éclater).

\impact{Assouplissement quantitatif $\Delta B^{CB}$}
\direct Le taux des obligations baisse ($-s_7$) \eqref{eq:spread}, donc le taux long et le coût de la dette publique ; les réserves excédentaires $Res^{ex}$ gonflent (achats aux banques) ou les dépôts et $M$ (achats aux ménages) ; la banque centrale gagne le coupon, versé au Trésor en seigneuriage $\Sigma$.
\second Les épargnants privés de rendement obligataire se reportent sur les actions et l'immobilier \eqref{eq:portfolio} : $p^E, p^Z$ montent, effet richesse positif, hausse des inégalités (les actifs sont détenus par la strate $H$). Le change se déprécie (rendements plus bas) et soutient les exportations. Si les banques trouvent des emprunteurs ($q_j>1$), les réserves se transforment en crédit et $M$ accélère ; sinon elles restent oisives et rien n'arrive aux prix : c'est le cas japonais et européen (\S\ref{sec:qe}). Le QE ne devient inflationniste que par la monnaie, jamais par un seuil de défiance automatique (v0.8) : la configuration 4, dont la banque centrale détient la moitié de la dette sans perte de crédibilité, l'exige.
\baisse Le resserrement quantitatif (ventes ou non-renouvellement) fait remonter $i^B$, exerce une pression baissière sur les actifs et peut révéler des pertes en capital au bilan de la banque centrale (achats à taux bas, ventes à taux hauts), qui deviennent une charge pour le Trésor. La sortie d'un QE massif est plus délicate que l'entrée. Depuis la v0.6, le QE a un second coût : il raccourcit la maturité consolidée de la dette publique ($\theta^{cons}$ dans \eqref{eq:iapp}), de sorte que toute hausse future du taux directeur se transmet immédiatement à la charge de la dette via le taux des réserves, et fait basculer le résultat de la banque centrale \eqref{eq:ecb} en perte.

\impact{Financement monétaire du déficit $\Delta A^G$}
\direct La base et la masse monétaires augmentent d'autant ; la dette publique n'augmente pas \eqref{eq:gbc} ; la crédibilité chute ($-\nu_3$) \eqref{eq:cred} et $\pi^e$ monte.
\second La dépense publique financée soutient la demande (multiplicateur) : à court terme, croissance et emploi. Puis le terme d'encaisses \eqref{eq:monterm} et la vitesse $V(\pi^e)$ poussent les prix ; les salaires s'indexent \eqref{eq:wage} ; la fuite vers les devises déprécie le change et renchérit les importations \eqref{eq:fx}. Le taux des obligations monte ($s_4(1-c)$), ce qui rend le financement de marché plus cher et pousse à monétiser davantage : cercle vicieux. Les recettes fiscales réelles s'érodent (Olivera--Tanzi). Si le déficit réel dépasse le seigneuriage maximal, l'hyperinflation est inévitable (\S\ref{sec:qe}). Un usage ponctuel et annoncé comme tel (guerre, catastrophe) avec une banque centrale par ailleurs indépendante coûte peu de crédibilité ; c'est la persistance qui tue.
\baisse L'arrêt du financement monétaire est la condition nécessaire de toute stabilisation ; il ne rétablit pas la crédibilité à lui seul (elle remonte lentement, $\nu_1$) et provoque une récession de stabilisation le temps que $\pi^e$ redescende.

\impact{Réserves obligatoires $\varrho^{res}$ et ratio de fonds propres $\kappa^{CAR}$}
\direct Une hausse réduit le multiplicateur de crédit : $L^{max}$ baisse \eqref{eq:spread}, $\Lambda^{Bk}$ peut passer sous~1, le crédit est rationné.
\second Investissement et achats immobiliers freinés ; $p^Z$ et $p^E$ baissent ; masse monétaire ralentie ; banques plus solides (coussin de pertes plus épais), donc moins de ruées \eqref{eq:run} et moins de risque de renflouement public. Instrument \emph{macroprudentiel} : il freine le crédit sans relever les taux pour toute l'économie, ce qui en fait l'outil naturel contre une bulle immobilière lorsque l'inflation est basse.
\baisse Libère du crédit, gonfle $M$, stimule les actifs ; historiquement, l'assouplissement réglementaire précède les crises (1920s, 2000s). Une baisse en pleine crise (forbearance réglementaire) peut au contraire éviter un credit crunch, au prix d'un système fragile.

\impact{Quotité hypothécaire LTV et marge boursière $m$}
\direct L'accès au crédit immobilier $\Lambda^Z$ et aux prêts sur marge augmente ; $D^Z$ et $D^E$ montent, donc $p^Z$ et $p^E$ \eqref{eq:pz}--\eqref{eq:pe}.
\second Les plus-values nourrissent l'extrapolation ($\chi$) ; le collatéral gonfle et permet plus de crédit : boucle de Minsky. Le levier des ménages monte ; à la première baisse de prix, une part des emprunteurs passe en fonds propres négatifs ($NPL^Z$), les appels de marge forcent les ventes. Le LTV est la variable qui, à elle seule, détermine si un retournement de prix est une correction ou une crise systémique.
\baisse Freine les prix d'actifs, réduit le levier ; si la baisse est brutale au sommet d'une bulle, elle la crève (Japon 1990 : plafonnement du crédit immobilier).

\impact{Contrôle des capitaux $\chi^K$}
\direct Le canal UIP du change se ferme \eqref{eq:fx} : le taux directeur peut diverger du taux mondial sans mouvement de capitaux ; la prime souveraine monte ($s_5$) \eqref{eq:spread} ; le placement externe de la dette se tarit.
\second Autonomie monétaire retrouvée (le triangle de Mundell : on ne peut avoir à la fois change fixe, capitaux libres et politique monétaire autonome) ; protection contre les attaques spéculatives ; mais coût du capital plus élevé, marché noir des devises si la monnaie est surévaluée, et moindre discipline budgétaire (la prime de marché ne joue plus son rôle d'alerte).
\baisse L'ouverture attire les capitaux si les taux et la croissance sont attractifs (boom du crédit, appréciation, déficit courant) et expose à leur retrait brutal (\S\ref{sec:crises}, crises de 1997). Ouvrir avant d'avoir un système bancaire solide et des réserves est l'erreur classique.

\impact{Régime de change : passage au change fixe (ou à l'or)}
\direct Le change cesse de bouger ; la banque centrale absorbe tout déséquilibre par ses réserves $\Delta R^{FX}=CA+KA$ ; le taux directeur doit suivre celui du pays d'ancrage (sauf contrôle des capitaux).
\second Discipline importée : l'inflation converge vers celle de l'ancre, la crédibilité peut monter ; mais perte de l'amortisseur : un choc (baisse du prix du pétrole, récession chez le partenaire) ne peut plus être absorbé par une dépréciation et passe entièrement par la baisse des salaires et de l'emploi. Si l'inflation reste supérieure à celle de l'ancre, la surévaluation s'accumule ($e$ s'écarte de $e^{PPA}$), le déficit courant se creuse, les réserves fondent et la probabilité de dévaluation perçue monte : attaque spéculative (\S\ref{sec:change}). Sous l'or, la contrainte s'étend au prêteur en dernier ressort, ce qui rend les paniques bancaires létales (\S\ref{sec:crises}).
\baisse Défendre un \emph{plafond} (refuser l'appréciation) est le cas miroir : munitions illimitées, mais bilan qui grossit et contrainte $E^{CB}/H$ qui finit par forcer l'abandon (\S\ref{sec:change}). Le passage au flottant restaure l'autonomie ; l'abandon de l'or ou d'une parité surévaluée produit une dépréciation immédiate, une poussée d'inflation importée et un rebond des exportations ; historiquement, les pays qui ont quitté l'or les premiers (1931) sont sortis les premiers de la Dépression.

\impact{Cible d'inflation $\pi^\ast$ et indépendance de la banque centrale}
\direct La cible entre dans la règle de Taylor et dans l'ancrage des anticipations ($c(\pi^\ast-\pi^e)$) \eqref{eq:expect} ; l'indépendance interdit $\Delta A^G$ et retire la pénalité $\nu_4$ de crédibilité.
\second Relever la cible (de 2 à 4\,\%) donne plus de marge avant la borne zéro et allège la dette réelle, mais coûte de la crédibilité si le relèvement est perçu comme un renoncement. L'indépendance élève $c$ à terme, abaisse la prime souveraine ($s_4$) et rend chaque désinflation moins coûteuse en chômage ; elle retire au joueur l'arme de la monétisation en cas d'urgence.
\baisse Une cible plus basse renforce l'ancrage mais rapproche de la déflation ; la dépendance de la banque centrale offre une flexibilité de court terme payée par une prime de risque permanente.

\subsection{Leviers de politique budgétaire}

\impact{Impôt sur le travail $\tau_{W,h}$}
\direct Le revenu disponible de la strate baisse \eqref{eq:yd} ; sa consommation baisse proportionnellement à $c_h$ \eqref{eq:cB} : l'effet est maximal si l'impôt frappe la strate $B$ ($c_B\approx0{,}95$), minimal sur la strate $H$ ($c_H\approx0{,}5$, qui épargne la moitié) ; les recettes montent moins que proportionnellement (évasion, $\varepsilon^{ev}$).
\second Baisse de la demande, de $\hat y$, de l'emploi ; pression à la hausse des salaires bruts (les salariés cherchent à préserver leur net) qui renchérit le coût du travail ; le déficit se réduit et $i^B$ baisse un peu. Progressivité : taxer $H$ plutôt que $B$ réduit les inégalités et coûte moins en demande, mais réduit l'épargne disponible pour le crédit et la dette publique \eqref{eq:portfolio}.
\baisse Relance par la consommation, d'autant plus efficace qu'elle cible les bas revenus ; creuse le déficit ; en surchauffe, nourrit l'inflation. Une baisse d'impôt sur $H$ se transforme surtout en épargne (effet ricardien partiel).

\impact{Impôt sur le capital $\tau_K$ et sur les sociétés $\tau_\Pi$}
\direct $\tau_\Pi$ réduit le profit net \eqref{eq:profit}, donc $q_j$ et l'investissement \eqref{eq:invest} ; $\tau_K$ réduit le revenu de la strate $H$ et les rendements après impôt de tous les placements, ce qui abaisse la valeur des actions.
\second Moins d'investissement aujourd'hui, moins de capital et de productivité demain : l'effet sur la croissance est lent mais durable. Les inégalités baissent. La fuite des capitaux (si $\chi^K<1$) réduit l'assiette : la courbe de Laffer est plus pentue pour le capital que pour le travail. Les recettes financent des dépenses dont l'effet dépend de leur nature (voir ci-dessous) : un transfert de $\tau_\Pi$ vers $G^{inv}$ peut être neutre ou positif pour la croissance.
\baisse Stimule l'investissement et les prix d'actifs, creuse les inégalités et le déficit ; concurrence fiscale : si un voisin baisse $\tau_\Pi$, les capitaux migrent ($KA$).

\impact{Taxes à la consommation $\tau^C_j$ et tarifs $\tau^M$}
\direct La TVA élève le prix payé par les ménages sans changer le prix reçu par le producteur ; les tarifs élèvent le prix rendu des importations \eqref{eq:armington}, déplacent la demande vers les producteurs domestiques et rapportent des recettes.
\second La TVA sur un bien de subsistance ($\gamma_{hj}>0$) frappe les pauvres (régressive) et comprime toute leur autre consommation \eqref{eq:les} ; sur un bien de luxe, elle est progressive. Un tarif protège un secteur (emploi, apprentissage par la pratique \eqref{eq:lbd}) au prix d'intrants plus chers pour les secteurs aval, d'une inflation importée et de représailles des partenaires (autres joueurs) ; le change s'apprécie un peu (moins d'importations), ce qui rogne l'avantage. Protéger une industrie naissante (Hamilton, List) peut être gagnant si l'apprentissage est fort et le tarif temporaire ; protéger une industrie mûre ne l'est pas.
\baisse La libéralisation commerciale baisse les prix, élargit le choix, expose les secteurs non compétitifs ; le bilan agrégé est positif mais les perdants sont concentrés (soutien politique).

\impact{Investissement public $G^{inv}$}
\direct Ordres d'achat sur la construction et l'équipement ; $K^G$ s'accumule et élève la productivité de tous les secteurs \eqref{eq:tfp} avec délai ; le déficit se creuse.
\second Multiplicateur de demande à court terme (surtout en récession, où il ne fait pas monter les prix) ; hausse de $Y^{pot}$ à long terme, donc allègement du ratio dette/PIB par le dénominateur \eqref{eq:debtdyn}. Rendements décroissants ($\zeta<1$) : les premières routes valent plus que les dernières. En surchauffe, éviction du privé (hausse des prix de la construction et des taux).
\baisse Économie budgétaire immédiate, dégradation lente et peu visible de la productivité : la coupe la plus tentante et la plus coûteuse à long terme.

\impact{Éducation et santé $G^{edu}, G^{san}$}
\direct $H$ croît avec un délai d'une génération ; la mortalité baisse et $SoL$ monte directement.
\second Productivité et salaires plus élevés à terme, mobilité sociale (passage de $B$ à $M$), fécondité plus basse (transition démographique) ; aucun rendement visible avant dix à vingt ans : le joueur qui les finance en paie le coût et rarement le bénéfice, comme les gouvernants réels.
\baisse Aucun effet visible à court terme ; décrochage durable de $H$ et de la frontière technologique.

\impact{Transferts $Tr$ (chômage, retraites, minima)}
\direct Revenu disponible de la strate $B$ \up, consommation \up{} à $c_B$ ; déficit \up.
\second Multiplicateur le plus élevé du modèle ; stabilisateur automatique (les allocations chômage montent en récession sans décision) ; réduction des inégalités et hausse du soutien politique de $B$ ; effet désincitatif sur l'offre de travail si le remplacement $\varrho$ est proche de~1 ($u^n$ monte) ; hausse du salaire de réservation, donc des salaires \eqref{eq:wage}. Les retraites, indexées sur la population âgée, deviennent un coût croissant dans une économie vieillissante.
\baisse Austérité avec le multiplicateur le plus fort : récession immédiate, hausse des inégalités et de la contestation ; réduction du déficit inférieure à la coupe (recettes perdues).

\impact{Subventions $Sub_j$}
\direct Maintiennent en vie des secteurs non rentables ; déficit \up.
\second Emploi préservé à court terme ; prix du bien maintenu bas si la subvention est répercutée ; contrainte budgétaire lâche : le secteur cesse d'améliorer sa productivité (Kornai) ; les ressources restent piégées dans un secteur déclinant au lieu de migrer. Une subvention à un secteur exportateur peut être contestée par les autres joueurs.
\baisse Faillites, chômage sectoriel concentré, puis réallocation ; effet positif sur la productivité agrégée après deux à trois ans.

\impact{Dette publique $b$ et déficit primaire $pb$}
\direct Un déficit primaire plus élevé soutient la demande et accroît $b$ \eqref{eq:debtdyn} ; la prime de risque monte de façon convexe \eqref{eq:spread}.
\second À faible dette, l'effet dominant est keynésien (relance). À dette élevée, la hausse de $i^B$ se transmet aux taux privés (éviction), les intérêts absorbent une part croissante des recettes, la dynamique boule de neige s'enclenche, et les épargnants exigent une prime de défaut $\rho^{def}$ qui réduit la demande d'obligations \eqref{eq:portfolio}. La dette détenue par la banque domestique lie le sort de l'État et des banques. La dette en devises est réévaluée à chaque dépréciation : un pays qui emprunte en monnaie étrangère perd le contrôle de son ratio.
\baisse La consolidation réduit la prime mais coûte de la croissance à court terme ; sa composition compte : couper $G^{inv}$ est le plus coûteux à long terme, couper $Tr$ le plus coûteux à court terme, relever $\tau_{W,H}$ ou $\tau_K$ le moins coûteux en demande mais le plus fuyant.

\subsection{Leviers institutionnels}

\impact{Propriété publique $\theta_j$}
\direct Les dividendes du secteur vont au budget \eqref{eq:div} ; l'investissement du secteur est décidé par l'État ; au-delà de 50\,\%, le secteur obéit aux directives et bénéficie de la contrainte budgétaire lâche.
\second Recettes hors impôt (utile pour un pays à faible capacité fiscale) ; réduction des inégalités (la strate $H$ perd ses dividendes) ; pénalité de productivité croissante avec $\varrho^{plan}$ \eqref{eq:tfp} ; possibilité d'investir contre la rentabilité (industrialisation lourde, infrastructures) ; risque de sur-emploi et de pertes chroniques. Le mode d'acquisition (indemnisation ou expropriation) détermine la réaction des capitaux privés.
\baisse La privatisation rapporte une recette exceptionnelle, restaure les incitations, concentre la richesse ; mal séquencée (avant la mise en place d'institutions de marché), elle produit des oligarques et une chute de production (Russie 1992--98).

\impact{Prix administrés $\varrho^p_j$}
\direct Le prix cesse de réagir au déséquilibre \eqref{eq:price} ; si le prix fixé est inférieur à l'équilibre, pénurie, rationnement, files ($q^{file}$), marché noir ; s'il est supérieur, invendus et subventions.
\second Inflation \emph{réprimée} : l'écart entre prix officiel et prix de marché noir s'accumule et se libère brutalement à la libéralisation (Russie 1992, Vénézuéla). Le signal-prix disparu, les producteurs n'ont plus de raison d'accroître l'offre du bien en pénurie ; le surplomb monétaire des ménages gonfle. Utile à court terme contre un choc (prix du pain en guerre), destructeur à long terme.
\baisse La libération des prix provoque une hausse ponctuelle (rattrapage de l'inflation réprimée) puis résorbe les pénuries en quelques mois ; elle est politiquement coûteuse et doit être accompagnée de transferts.

\impact{Degré de planification $\varrho^{plan}$}
\direct Allocation par le plan (Leontief) plutôt que par les prix ; plein emploi mécanique ; taux d'investissement choisi ; pénalité de productivité \eqref{eq:tfp} croissante avec la complexité de l'économie.
\second Phase 1 (dix à vingt ans) : croissance extensive rapide par accumulation forcée du capital et transfert de main-d'œuvre rurale ; inégalités faibles ; pas de cycle financier. Phase 2 : rendements décroissants du capital, pénalité d'inefficience qui croît avec la diversification, pénuries chroniques, retard technologique cumulé ; stagnation. Le plan est un outil de rattrapage, pas de frontière.
\baisse La transition vers le marché est un choc : chute de production le temps que les prix et les entreprises s'ajustent (courbe en J), chômage, inégalités ; réussie si elle est graduelle (Chine) ou accompagnée d'institutions solides (Pologne), catastrophique sinon (Russie).

\impact{Ouverture commerciale (biais domestique $\varphi_{ii}$, quotas, blocs)}
\direct Plus d'ouverture réduit $\varphi_{ii}$ : les acheteurs comparent davantage les prix étrangers \eqref{eq:armington}.
\second Baisse des prix, spécialisation selon les avantages comparatifs, croissance de la productivité par la concurrence et l'accès à la frontière ($\chi$) ; secteurs perdants concentrés ; plus grande transmission des chocs extérieurs et du change à l'inflation domestique.
\baisse Fermeture : protection, autonomie, prix plus élevés, perte de productivité relative à long terme, marché intérieur seul moteur (viable pour un pays très peuplé, pas pour un petit pays).

\subsection{Nouveaux leviers de la v0.6}

\impact{Maturité d'émission $\bar T$}
\direct $\theta$ baisse \eqref{eq:iapp} : le taux apparent réagit plus lentement aux taux de marché ; prime de terme $s_0\bar T/\bar T_0$ \eqref{eq:spread}.
\second Assurance contre une hausse future des taux, payée d'avance ; en période de taux bas, allonger fige un coût faible pour longtemps (bon) ; en période de taux hauts, allonger fige un coût élevé (mauvais). Le QE défait en partie l'allongement : $\theta^{cons}$ dépend de $\omega$, non de $\bar T$ seul.
\baisse Raccourcir économise la prime de terme et expose au risque de refinancement : chaque hausse de taux se transmet vite, et un saut de prime \eqref{eq:jump} devient immédiatement coûteux.

\impact{Rémunération des réserves $\Delta^{res}$}
\direct Un écart plus grand entre taux directeur et taux des réserves réduit la charge $i^{res}Res$ de la banque centrale \eqref{eq:ecb} : le seigneuriage $\Sigma$ remonte.
\second Les banques, moins rémunérées sur leurs réserves, répercutent sur les dépôts ($i^D$ baisse) : transfert des déposants vers le budget ; si l'écart est trop grand, les banques fuient les réserves et la transmission du taux directeur se dégrade.
\baisse Rémunération pleine : transmission parfaite, transfert maximal du budget vers les banques en période de taux hauts.
\impact{Règle budgétaire constitutionnelle $\varrho^{fisc}$}
\direct Le déficit structurel est plafonné : le joueur perd la main sur $pb$ en surchauffe et la retrouve en récession (clause de cycle) ; $\rho^{pays}$ baisse ; en régime B, le backstop est acquis.
\second Prime souveraine durablement plus basse, donc boule de neige plus faible et espace budgétaire retrouvé à long terme ; mais l'investissement public est la variable d'ajustement la plus facile à couper (\S\ref{sec:impacts}), et la règle repose sur une estimation de $Y^{pot}$ que le modèle lui-même ne connaît qu'imparfaitement : une règle mal calibrée impose l'austérité en récession. Deux rentiers identiques, l'un avec règle de prélèvement sur son fonds, l'autre sans, divergent en dix ans : le premier accumule et surchauffe peu, le second dépense la rente et subit la maladie hollandaise.
\baisse Sortir de la règle prend $\lambda^{fisc}$ mois et coûte du soutien politique et de la prime ; la flexibilité retrouvée sert rarement à investir.

\impact{Contrainte d'usage de la monnaie nationale $\varrho^{dedol}$}
\direct Le plancher de $s^{FX}$ baisse plus vite \eqref{eq:sfx} ; part de l'activité en marché noir des devises \up{} ; productivité formelle \dn.
\second Reconquête progressive du seigneuriage et de la transmission monétaire \eqref{eq:reff} ; mais si la crédibilité ne suit pas ($c$ plafonnée par $M^{hist}$), les agents contournent (marché noir), l'assiette fiscale se réduit et la contrainte nourrit la méfiance qu'elle voulait combattre. Efficace seulement après stabilisation crédible.
\baisse Laisser faire : dollarisation durable, politique monétaire importée, pas de prêteur en dernier ressort ; c'est un choix défendable pour un petit pays qui a renoncé.

\impact{Répression financière $m^D$, $\varrho^B$}
\direct Taux des dépôts \dn{} ; transfert $\mathcal{T}^{rep}$ des épargnants vers les emprunteurs ; taux souverain \dn{} ($-s_9\varrho^B$) \eqref{eq:spread}.
\second Désendettement public sans impôt visible (les années 1945--1980 des économies avancées) ; épargne des ménages découragée ou déplacée vers les actifs réels et les devises ($s^{FX}$ \up), bulles immobilières ; crédit mal alloué (les banques prêtent à l'État plutôt qu'aux entreprises, $\Lambda_j$ \dn) ; inflation tolérée pour que le taux réel des dépôts soit négatif. Fonctionne avec contrôle des capitaux ($\chi^K$), s'effondre sans.
\baisse La libéralisation financière remonte les taux réels, révèle la charge de la dette et relance l'épargne ; mal séquencée, elle produit une crise bancaire (Amérique latine 1980).

\impact{Capacité fiscale $T/Y$ et part domestique de la dette $\varkappa^{dom}$}
\direct Une hausse de $T/Y$ relève le seuil $\bar b_i$ \eqref{eq:spread} : la même dette devient moins chère, et le saut \eqref{eq:jump} moins probable ($b-\bar b_i$ baisse). Une hausse de $\varkappa^{dom}$ réduit la convexité au prorata.
\second $T/Y$ monte avec $G^{adm}$ (efficacité de collecte) et avec le développement (formalisation) : la marge d'endettement d'un émergent croît avec sa capacité à lever l'impôt, pas avec son PIB. $\varkappa^{dom}$ monte avec l'épargne domestique, la répression financière ($\varrho^B$ force les banques à détenir la dette) et le QE ($\omega$) ; il baisse avec l'ouverture du compte de capital. Le prix de la détention domestique est ailleurs : répression des épargnants, boucle banques--souverain, transmission accélérée des hausses de taux via $\omega$.
\baisse Une baisse de $T/Y$ (récession, évasion) abaisse le seuil et peut faire basculer un pays de la zone sûre à la zone convexe sans qu'il ait emprunté davantage ; une baisse de $\varkappa^{dom}$ (fuite des résidents) ouvre la convexité d'un coup.

\impact{Mutualisation $\varsigma^{mut}$}
\direct Une part plus grande des pertes de la banque centrale de l'union et des défauts est répartie selon la clé : le coût d'un défaut pour le pays défaillant baisse, celui des autres monte.
\second Aléa moral : à $\varsigma^{mut}$ élevé, la discipline ne peut venir que de la règle budgétaire supranationale ; à $\varsigma^{mut}$ faible, le backstop n'écarte que le mauvais équilibre de \eqref{eq:jump} et laisse la solvabilité au pays. Pour le membre vertueux, $\varsigma^{mut}$ est le prix de l'assurance contre la contagion.
\baisse À $\varsigma^{mut}=0$, l'union est une zone de change fixe irrévocable sans solidarité : c'est le régime E collectif.

\impact{Placement du fonds souverain $\beta^{dom}$}
\direct La part domestique alimente $D^E$ et l'investissement national ; la part externe alimente $KA$ et neutralise l'appréciation \eqref{eq:fx}.
\second À $\beta^{dom}$ élevé : bulle d'actifs locale, surchauffe, salaires et prix domestiques en hausse, maladie hollandaise malgré le fonds, et rendement du fonds corrélé au cycle national plutôt qu'au cycle mondial (perte de diversification). À $\beta^{dom}=0$ : le pays exporte sa rente et importe le risque boursier mondial ; la pression politique pour \og investir chez nous\fg{} est le mécanisme par lequel $\beta^{dom}$ dérive dans le temps.
\baisse Ramener $\beta^{dom}$ vers zéro après une phase de placement domestique force la vente d'actifs nationaux : baisse de $p^E$ et effet richesse négatif.

\impact{Contagion $\varsigma_{11}$ et régime E}
\direct Un saut chez un membre relève la probabilité de saut des autres de $\varsigma_{11}$ dans l'argument de \eqref{eq:jump} ; en régime E, le taux directeur est importé et la contrainte de réserves porte sur $R^{FX}\ge\varpi_E M$.
\second La contagion donne aux membres vertueux un intérêt au sauvetage et rend le backstop rationnel ; elle rend aussi possible une crise systémique de l'union à partir d'un seul membre si $\varsigma^{mut}$ et le backstop sont faibles. En régime E, un choc sur la ressource se transmet intégralement au budget puis aux réserves, sans amortisseur monétaire ; la seule défense est le fonds souverain externe ($\beta^{dom}$ bas).
\baisse Sortir du régime E (flotter) rend l'autonomie au prix d'une dévaluation immédiate et d'une perte de la crédibilité importée : $c$ retombe à sa valeur propre.

\subsection{Variables endogènes clés}

\impact{Inflation anticipée $\pi^e$}
\direct Salaires \up{} \eqref{eq:wage} ; indexation des prix \up{} \eqref{eq:price} ; taux réel \dn{} pour un taux nominal donné (stimule $q_j$ et la consommation de $H$) ; vitesse de circulation \up{} \eqref{eq:monterm} ; part en devises \up{}, change \dn{} \eqref{eq:fx} ; coût d'usage du logement \dn{} (attrait des actifs réels).
\second L'inflation réalisée monte et valide l'anticipation : auto-réalisation. La banque centrale doit relever $i^{CB}$ plus que la hausse de $\pi^e$ (principe de Taylor) pour rétablir un taux réel restrictif ; sinon la spirale s'installe. La crédibilité baisse si l'inflation dépasse la cible. La variable la plus dangereuse du modèle parce qu'elle se nourrit d'elle-même.
\baisse Une baisse de $\pi^e$ vers zéro ou en dessous (déflation anticipée) élève le taux réel à la borne zéro, retarde les achats (pourquoi acheter aujourd'hui ce qui sera moins cher demain), alourdit les dettes réelles : la spirale déflationniste est le miroir de la spirale inflationniste et plus difficile à combattre.

\impact{Crédibilité de la banque centrale $c$}
\direct Ancrage plus fort de $\pi^e$ sur la cible \eqref{eq:expect} ; prime souveraine \dn{} ($s_4$) \eqref{eq:spread}.
\second Chaque choc inflationniste (pétrole, dépréciation) est absorbé sans spirale, donc avec moins de hausse de taux et moins de chômage : la crédibilité est un capital qui abaisse le coût de toutes les politiques futures. Elle se construit lentement (années de cible tenue) et se détruit en quelques mois (monétisation, cible ratée durablement).
\baisse Chaque désinflation devient coûteuse (sacrifice ratio élevé) ; les marchés exigent une prime ; la marge de manœuvre budgétaire se réduit.

\impact{Chômage $u$}
\direct Salaires \dn{} \eqref{eq:wage} ; transferts \up{} (stabilisateur) ; niveau de vie \dn{} \eqref{eq:sol} ; soutien politique \dn.
\second Baisse de la consommation de $B$ (multiplicateur négatif), baisse des recettes fiscales, hausse du déficit ; désinflation avec retard ; chômage de longue durée : perte de qualification ($H$ \dn), mobilité descendante ($B$ gonfle), hausse de $u^n$ (hystérèse). La banque centrale baisse $i^{CB}$ (Taylor) ; à la borne zéro, seul le budget peut agir.
\baisse Sous $u^n$ : tensions salariales, inflation, puis resserrement monétaire ; au point de Lewis pour un pays dual, décollage des salaires.

\impact{Prix d'un bien $p_j$}
\direct Rentabilité du secteur \up{} \eqref{eq:profit}, donc $q_j$, investissement et embauches \up{} \eqref{eq:invest}--\eqref{eq:labour} ; consommation du bien \dn{} et, s'il est de subsistance, toute autre consommation \dn{} \eqref{eq:les} ; coût unitaire des secteurs aval \up{} (intrants) \eqref{eq:price} ; indice des prix \up{} au poids $w_j$ ; parts d'importation \up{} \eqref{eq:armington}.
\second L'offre augmente et le prix revient : c'est le mécanisme d'ajustement du marché ; sa vitesse dépend du délai d'investissement. Un choc sur un intrant universel (énergie) se propage à tous les prix (stagflation), tandis qu'un choc sur un bien final reste local. En prix administré, rien de ceci ne se produit : la pénurie persiste.
\baisse Symétrique ; un prix durablement sous le coût conduit à la faillite du secteur et à la dépendance aux importations.

\impact{Productivité $A_j$}
\direct Production par travailleur \up{} ; coût unitaire \dn{}, donc prix \dn{} ou marge \up{} ; salaires \up{} par la part $\gamma_A$ \eqref{eq:wage}.
\second Le PIB potentiel monte, ce qui ouvre un écart de production négatif si la demande ne suit pas (chômage technologique transitoire) ; le taux naturel $r^\ast$ monte ; la compétitivité extérieure s'améliore (exportations \up) ; la dette publique s'allège par le dénominateur. C'est la seule source de croissance durable du niveau de vie.
\baisse Une stagnation de $A$ avec des salaires qui continuent de suivre $\pi^e$ produit une inflation par les coûts et une perte de compétitivité (Italie 1995--2015).

\impact{Prix des actifs $p^Z$, $p^E$}
\direct Richesse $V_h$ \up{} donc consommation \up{} ($c^V$) ; collatéral \up{} donc crédit \up{} ($\Lambda_j$, $\Lambda^Z$) ; mises en chantier \up{} ; inégalités \up{} (les actifs sont concentrés).
\second Boucle de Minsky si $\chi$ élevé et crédit facile ; hausse du levier ; la banque centrale voit une inflation modérée des biens et ne réagit pas (années 1920, 2000) ; l'offre nouvelle de logements arrive avec retard et fait basculer le marché.
\baisse Effet richesse négatif, appels de marge, fonds propres négatifs des emprunteurs, $NPL$ \up, $E^{Bk}$ \dn, credit crunch, récession ; déflation par la dette si les prix des biens suivent. L'asymétrie est totale : la baisse détruit plus vite que la hausse ne construit.

\impact{Fonds propres bancaires $E^{Bk}$}
\direct Capacité de crédit $L^{max}$ \up{} ; probabilité de ruée \dn{} \eqref{eq:run}.
\second Plus de crédit, plus d'investissement, plus de $M$ ; en excès, financement de projets marginaux et de bulles.
\baisse Rationnement du crédit indépendant de la demande ; forbearance et zombies si la banque est sous-capitalisée mais non fermée (\S\ref{sec:zlb}) ; ruées sans garantie des dépôts ; renflouement public qui transfère la perte au ratio dette/PIB (Irlande 2010 : $+40$ points).

\impact{Compte courant $CA$ et réserves de change $R^{FX}$}
\direct Un excédent apprécie le change (flottant) ou accumule des réserves (fixe) \eqref{eq:fx} ; un déficit fait l'inverse.
\second Déficit persistant : dette externe croissante, intérêts versés à l'étranger qui aggravent le déficit, vulnérabilité à un retrait des capitaux ; réserves qui fondent en change fixe jusqu'à l'attaque. Excédent persistant : appréciation qui érode la compétitivité (ou réserves qui gonflent et créent de la monnaie si elles ne sont pas stérilisées), tensions avec les partenaires. Un pays rentier a un excédent volatil qui suit le prix mondial.
\baisse Voir ci-dessus ; la variable à surveiller en mois d'importations pour un pays en change fixe.

\impact{Niveau de vie $SoL$ et population $N$}
\direct $SoL$ \up{} : soutien politique \up, fécondité \up{} puis \dn{} (transition), migration entrante ; $N$ \up{} : offre de travail \up{}, demande de subsistance \up{}, pression sur la terre \eqref{eq:agr}, marché intérieur plus grand ($\varphi_{ii}$).
\second Une population qui croît plus vite que le capital et la terre fait baisser le salaire vers la subsistance (Malthus) ; une population qui croît avec le capital fournit un dividende démographique ; une population qui vieillit alourdit les retraites et réduit la participation.
\baisse Dépopulation : rareté du travail, salaires \up, terre abondante, marché intérieur qui se contracte.

\impact{Prix mondial des ressources $p^w_{res}$}
\direct Pour le rentier : rente \up, recettes \up, compte courant \up, change \up{} (appréciation) ; pour l'importateur : coûts \up{} sur tous les secteurs via Leontief, inflation par les coûts, compte courant \dn.
\second Rentier : maladie hollandaise si la rente est dépensée ; investissements surdimensionnés calés sur un prix élevé ; effondrement budgétaire à la baisse suivante. Importateur : stagflation, dilemme de la banque centrale (relever les taux contre l'inflation ou les baisser contre la récession) ; le pétrodollar revient sous forme de capitaux ($KA$).
\baisse Pour le rentier, crise budgétaire et de change (Russie 1998, Venezuela 2014) d'autant plus violente que la dépendance est forte ; pour l'importateur, un stimulus gratuit.

\impact{Terre $T$ et réforme agraire}
\direct Plus de terre par travailleur : salaire rural \up, rente par hectare \dn, migration rurale \dn{} \eqref{eq:migr}.
\second La redistribution de la terre (réforme agraire) ne change pas $T$ mais transfère la rente des propriétaires ($H$) aux paysans ($B$) : inégalités \dn, consommation \up{} (propension plus forte), soutien de $B$ \up{} et de $H$ \dn{} ; effets sur la productivité ambigus (petites exploitations mieux entretenues, mais moins de capital) ; c'est historiquement l'acte fondateur des décollages asiatiques (Japon 1946, Corée 1950, Taïwan 1953).
\baisse Une terre plus rare (croissance démographique, dégradation) enclenche la pression malthusienne, sauf mécanisation ou importation d'alimentation.

%=====================================================================
\section{Prototype de l'étape 1 : état stationnaire, tests, propriétés}
\label{sec:proto}
%=====================================================================
L'étape~1 de la feuille de route a été programmée (Python, ~600 lignes : quatre biens, trois strates, une banque, une banque centrale, un État, pas hebdomadaire, grand livre comptable) et soumise à quinze tests. Cette section en donne l'état stationnaire, les résultats et les propriétés qui en découlent ; les fichiers (\texttt{model.py}, \texttt{tests.py}, \texttt{results.json}) accompagnent le document.

\subsection{État stationnaire obtenu (calibration de la \S\ref{sec:calib})}
\label{sec:ssref}
\paragraph{Deux états de référence, deux rôles (v1.5).} Le document contient deux états stationnaires : celui-ci (économie fermée, sans marchés d'actifs) et celui de la \S\ref{sec:proto3} (avec actifs : $K/Y$ 2,2, $I/Y$ 11\,\% $+$ 3\,\% de logement, part salariale 0,72). Ils diffèrent d'un quart sur le stock de capital et ne décrivent pas la même économie. La v1.5 déclare leurs rôles : \textbf{l'économie à actifs est la référence de jeu} (les parties démarrent avec actifs) ; \textbf{l'économie fermée est la référence de calibration} (elle est complète, ses grandeurs se recalculent depuis les équations, et c'est elle que le cours utilise pour ses applications chiffrées). L'écart entre les deux est le chantier ouvert du partage de la valeur ajoutée (\S\ref{sec:chantiers}) ; il n'est pas résolu et la question de la référence sera reprise quand il le sera. Les valeurs ci-dessous ont été re-mesurées sur le prototype de la v1.5 et recalculées depuis les équations là où c'est possible : $I/Y=(\delta+g)\,K/Y=(0{,}041+0{,}018)\times2{,}85=0{,}17$ (et non 0,18), $q^\ast=1+g/(\iota\Gamma)=1{,}11$.
\begin{center}\small
\begin{tabular}{p{7.4cm}p{4.2cm}p{3.6cm}}
\toprule
\textbf{Variable} & \textbf{Valeur (années 20--30)} & \textbf{Repère observé}\\
\midrule
Croissance réelle & 1,9\,\%/an & 1,5--2,5\,\%\\
Inflation / anticipée / cible & 2,1\,\% / 2,1\,\% / 2,0\,\% & cible atteinte\\
Chômage / chômage naturel (hystérèse) & 3,9\,\% / 4,5\,\% & 4--8\,\%\\
Taux directeur & 3,4\,\% & ---\\
Taux réel directeur / naturel estimé $\hat r^\ast$ / théorique $\rho+\sigma g$ & 1,3\,\% / 0,7\,\% / 2,5\,\% & voir propriété 2\\
Taux des dépôts / du crédit / souverain / apparent & 2,4 / 4,6 / 4,4 / 4,3\,\% & ---\\
$C/Y$ (TVA incluse) $\cdot$ $I/Y$ $\cdot$ achats publics $\cdot$ inv. public $\cdot$ salaires publics & 0,74 $\cdot$ 0,17 (recalculé : $(\delta+g)K/Y$) $\cdot$ 0,085 $\cdot$ 0,030 $\cdot$ 0,064 & 0,60--0,70 $\cdot$ 0,16--0,18 (hors logement) $\cdot$ $\approx$0,20\\
Part salariale (charges incl.) / dividendes / épargne des ménages & 0,65 / 0,087 / 0,050 & 0,60--0,68 / 0,03--0,08 / 0,05--0,10\\
$K/Y$ (nominal, $p_KK/PY$) $\cdot$ $K^G/Y$ $\cdot$ $q$ ($=q^\ast$) & 2,85 (2,93 en v1.4) $\cdot$ 0,57 $\cdot$ 1,11 & 2,5--3,5 $\cdot$ 0,5--0,7 $\cdot$ $\approx$1\\
$M/Y$ $\cdot$ dépôts des ménages $\cdot$ crédits aux entreprises $\cdot$ levier $\ell$ & 0,84 $\cdot$ 0,79 $\cdot$ 0,89 $\cdot$ 0,32 (0,08 dans l'économie à actifs) & 0,8--1,2 $\cdot$ --- $\cdot$ 0,6--1,0 $\cdot$ 0,3--0,5\\
Fonds propres bancaires / ratio CAR effectif & 0,122 / 12,4\,\% & exigence 9\,\%\\
Dette publique (ménages $\cdot$ banque $\cdot$ BC) & 0,37 (0,31 $\cdot$ 0,03 $\cdot$ 0,03) & stable\\
Emploi privé / public / actifs / population & 517 / 60 / 600 / 1\,000 & ---\\
Parts sectorielles de l'emploi (alim., énergie, conso., équip.) & 21,9\,\% / 17,4\,\% / 40,2\,\% / 20,4\,\% & ---\\
Revenu par tête B : M : H $\cdot$ Gini (3 classes) & 1 : 1,09 : 2,04 $\cdot$ 0,09 & trop plat (pas d'actifs)\\
\bottomrule
\end{tabular}
\end{center}

\subsection{Résultats des quinze tests}
\begin{center}\small
\begin{longtable}{p{4.6cm}p{7.6cm}p{1.6cm}}
\toprule
\textbf{Test} & \textbf{Résultat} & \textbf{Verdict}\\
\midrule
\endfirsthead
\toprule
\textbf{Test} & \textbf{Résultat} & \textbf{Verdict}\\
\midrule
\endhead
T1 État stationnaire, 30 ans & u 4,0\,\%, $\pi$ 2,1\,\%, croissance 1,8\,\%, dette 0,61 $\to$ 0,36, 0 faillite, identités comptables à $10^{-12}$ & PASS\\
T2 Choc monétaire $+1$ pt un an & v1.0 : $\Delta u$ $+3{,}4$ pt (15 mois), $\Delta$PIB $-2{,}9\,\%$, $\Delta\pi$ $-1{,}7$ pt (13 mois) ; v0.9 : $-6{,}4\,\%$ & PASS ; amplitude $\approx2\times$ l'empirique (propriété 5)\\
T3 Relance $+1\,\%$ PIB, deux ans & multiplicateur à taux fixe 0,79 (an 1) / 1,57 (an 2) ; sous règle de Taylor 0,24 / $-0{,}07$ & PASS\\
T4 Monétisation de dépenses $+5\,\%$ PIB, trois ans & crédibilité $-0{,}46$, $\Delta\pi$ $+4{,}7$ pts, pic à cinq ans & PASS\\
T5 Hyperinflation ($+30\,\%$ PIB monétisés, permanent) & $\pi$ : 26\,\% (an 7), 110\,\% (an 9), 58\,700\,\% (an 11) ; 4 réformes ; fini sur 20 ans & PASS\\
T5b Stabilisation (réforme à l'an 8, $\pi=43\,\%$) & $\pi$ $-3\,\%$ (an 9), $+4\,\%$ (an 11), $+2\,\%$ ensuite ; u 14\,\% au pic puis 5\,\% ; crédibilité 0,85 & PASS\\
T6 Choc sectoriel ($-50\,\%$ du capital énergie) & retour à 90\,\% en 8 ans, 0 zombie & PASS\\
T7 Stabilité stochastique (30 graines, chocs de PGF 2\,\%/an) & u max 7,4\,\%, $|\pi|$ max 7,8\,\%, aucun plancher, identités exactes & PASS\\
T8 Effondrement de la demande ($-8\,\%$ PIB, deux ans) & u 17,6\,\% au pic, demi-vie de reprise 1 an, aucun effet permanent & PASS\\
T9 Dérive budgétaire ($+6\,\%$ PIB, douze ans) & dette 0,50 $\to$ 1,12 puis 1,09 ; taux apparent max 9\,\%, prime max 11,5 pts & PASS\\
T10 Cycle résiduel (40 ans) & écart-type du chômage 0,08 pt puis 0,01 pt & PASS\\
T11 Identités comptables & bilan BC à $10^{-12}$ ; aucun dépôt ni réserve négatifs ; dette $=$ somme des détenteurs & PASS\\
T12 Expériences budgétaires & $+2$ pts d'impôt sur la strate basse : u $+0{,}8$ pt, niveau de vie $-3\,\%$ ; sur la haute : u $+0{,}2$ pt, $-1\,\%$ ; inv. public $+2$ pts : PIB $+13\,\%$ à 25 ans & PASS\\
T13 Robustesse $\pm30\,\%$ sur 8 paramètres & 16 économies stables, u moyen de 3,7 à 5,5\,\% & PASS\\
S1 Sensibilité à l'épargne (information) & propensions $\times0{,}8$ : u 20\,\%, taux zéro ; $\times0{,}9$ : u 8,8\,\% & ---\\
\bottomrule
\end{longtable}
\end{center}

\subsection{Propriétés identifiées de l'économie fermée}
\begin{enumerate}
\item \textbf{La cohérence stock-flux est exacte dans tous les régimes}, hyperinflation et réformes comprises, parce que le grand livre interdit les découverts : toute création monétaire passe par quatre portes identifiables (crédit, absorption de dette publique par une banque solvable, avances, recapitalisation).
\item \textbf{Le taux naturel effectif est inférieur au taux théorique.} L'économie a un taux wicksellien élevé ($MPK-\delta\approx6\,\%$) et un taux keynésien bas (1 à 2\,\%) : à $\rho+\sigma g$, l'épargne excède la demande autonome et l'économie s'installe dans un sous-emploi stable. Deux réponses coexistent en v0.9 : la banque centrale \emph{apprend} le taux naturel (l'inflation atteint la cible), et la prime de risque $\rho_E$ relève la demande d'investissement. L'écart résiduel entre taux réel effectif (1,3\,\%) et théorique (2,5\,\%) est la prime de risque non recyclée en richesse ; il se refermera avec les marchés d'actifs (étape~3). Le paramètre décisif est la propension à consommer : à $-20\,\%$, trappe à liquidité.
\item \textbf{L'inflation est déterminée par les coûts, la courbe de Phillips par le marché du travail.} Les entreprises servent les commandes et ajustent leurs capacités ; le tâtonnement ne voit presque jamais d'excès de demande ; les prix suivent le coût unitaire majoré de la marge normale. Conséquences : une relance n'est inflationniste que si elle réduit le chômage ; la monétisation d'un déficit \emph{existant} ne l'est pas, c'est la dépense monétisée qui l'est ; l'hyperinflation exige plein emploi (surenchère), perte de crédibilité (indexation sur l'inflation réalisée) et fuite devant la monnaie (vitesse), et se développe alors en quelques années.
\item \textbf{Le partage de la valeur ajoutée doit être écrit} : dans une économie à production déterminée par la demande, rien ne rappelle le salaire vers la productivité marginale ; la marge normale et l'indexation sur la productivité le font. Sans eux, régime d'accumulation sans salaires.
\item \textbf{La transmission monétaire est correcte en signe et en délais, trop forte en amplitude} ($-6\,\%$ de PIB pour $+1$ pt, contre $-0{,}5$ à $-1{,}5$ observé). Cause : l'élasticité de $q$ au taux vaut $1/(r+\rho_E+\delta)\approx13$ et l'investissement est un quart du PIB. La prime de risque l'a réduite d'un quart ; les retards de transmission (taux long lissé, ajustement partiel) et l'accélérateur ont été essayés et rejetés parce qu'ils engendrent des cycles de trois ans. Chantier ouvert pour l'étape~2, où le change offrira un second canal et où une règle prospective pourra être testée.
\item \textbf{Le multiplicateur budgétaire dépend du régime monétaire} (0,8 à 1,6 à taux fixe, nul sous règle de Taylor) et de la strate (taxer la strate basse coûte quatre fois plus de chômage). L'investissement public a un rendement élevé parce qu'il agit sur la productivité de tous les secteurs.
\item \textbf{Les planchers irréversibles sont fermés.} Quatre mécanismes produisaient des économies mortes : banque insolvable jamais recapitalisée, faillites destructrices, État en découvert, indexation sans monnaie. Tous quatre sont neutralisés et testés (T5b, T6, T8).
\item \textbf{Deux instabilités connues} : l'accélérateur (cycles de Samuelson-Hicks) et les retards de transmission ; le modèle est stable parce que l'investissement répond au $q$ courant et que la politique lisse fortement.
\item \textbf{Transitoire initial} : trois à cinq ans d'ajustement (chômage jusqu'à 7\,\%) ; les parties devront démarrer sur un état sauvegardé après dix ans de simulation.
\end{enumerate}

\subsection{Étape 2 : monde à $N$ pays}
\label{sec:proto2}
Le monde relie $N$ économies fermées identiques par le commerce (Armington sur les biens finals, parts graduelles, énergie échangeable en option), les comptes de change et les réserves (\S\ref{sec:sfc}), les flux de capitaux et les cinq régimes de souveraineté (\S\ref{sec:souv}). Le pays~0 émet la monnaie de référence. Treize tests, tous passés ; identités de bilan des banques centrales à $10^{-12}$ dans tous les régimes.

\begin{center}\small
\begin{longtable}{p{4.4cm}p{8.0cm}p{1.4cm}}
\toprule
\textbf{Test} & \textbf{Résultat} & \textbf{Verdict}\\
\midrule
\endfirsthead
\toprule
\textbf{Test} & \textbf{Résultat} & \textbf{Verdict}\\
\midrule
\endhead
E1 Monde symétrique, flottant & $e=0{,}996\pm0{,}0001$, balance $-0{,}2\,\%$, inflation 1,8\,\% partout, part importée 16\,\% & PASS\\
E2 PPA de long terme (A monétise $+4\,\%$ PIB, 6 ans) & inflation max 8,7\,\% ; dépréciation cumulée / différentiel de prix $=1{,}35$ & PASS\\
E3 Dornbusch (A $+1$ pt, un an) & appréciation 4,4\,\% à 12 mois, retour à 1,5\,\% à cinq ans : surréaction & PASS\\
E4 Productivité (A $+10\,\%$) & PIB relatif $+14\,\%$, appréciation réelle 2\,\%, déficit commercial $-2{,}2\,\%$ & PASS\\
E5 Krugman (A ancré, monétise $+3\,\%$) & réserves de 8 à 2 mois d'importations par déficit commercial, attaque à l'an 11,8, sorties de 5\,\% du PIB, dévaluation 30\,\%, inflation $1{,}8\to7{,}8\,\%$, déflation chez le partenaire & PASS\\
E6 Union (B dérive de $+5\,\%$, 12 ans) & change constant, déficit de B $-2{,}5\,\%$, solde de règlement croissant, dette 0,9, prime souveraine 5,1\,\% ; A subit l'inflation de B & PASS\\
E7 Stabilité (15 graines, chocs de PGF) & $e\in[0{,}95;1{,}05]$, chômage max 10,9\,\%, identités exactes & PASS\\
E8 Trois pays, énergie échangeable & $e$ à $\pm0{,}5\,\%$, balances $\pm0{,}2\,\%$, inflation 1,8\,\%, part importée 20\,\% & PASS\\
E9 Choc pétrolier importé (exportateur $-25\,\%$ en six mois) & stagflation : inflation $1{,}5\to2{,}4\,\%$, chômage $4{,}6\to6{,}0\,\%$ & PASS\\
E10 Dollarisation (monétisation 25\,\%, réforme) & $s^{FX}$ 18\,\% à la réforme, 13\,\% treize ans plus tard, dépréciation totale 22\,\% & PASS\\
E11 Contrôle des capitaux (taux 1 pt sous l'ancre, deux ans) & sans contrôle : attaque à l'an 6,3 ; avec : an 8,6 ; le déficit commercial contourne le contrôle & PASS\\
E12 Attaque auto-réalisatrice (récession du partenaire, aucune faute) & première génération : épuisement à l'an 8,5 ; seconde : abandon à l'an 5,3, réserves 3,6\,\% du PIB, chômage 12\,\% & PASS\\
E13 Transmission en économie ouverte ($+1$ pt, un an) & ouverte : PIB $-2{,}9\,\%$, $\pi$ $-2{,}6$ pt, appréciation 4,8\,\% ; fermée : $-2{,}5\,\%$, $-1{,}9$ pt : le change accélère la désinflation sans alléger le coût réel & PASS\\
\bottomrule
\end{longtable}
\end{center}

\paragraph{Enseignements de l'étape 2.}
\begin{enumerate}
\item Le règlement international doit porter sur les montants payés, et la conservation des réserves mondiales est un test indispensable en plus des identités locales (\S\ref{sec:sfc}).
\item Les parts de marché doivent s'ajuster graduellement ; l'élasticité seule ne rend pas le commerce stable (\S\ref{sec:marches}).
\item Un choc de dotation brutal casse une économie : un doublement de la productivité d'un secteur en une semaine produit une déflation sectorielle de 57\,\%, l'épuisement des trésoreries sous salaires rigides, puis une inflation de pénurie ; le même choc étalé sur deux ans est absorbé. Les hétérogénéités de dotation des dix configurations (\S\ref{sec:configs}) seront introduites à la calibration initiale ou graduellement, jamais par saut.
\item La perte d'autonomie monétaire d'un pays ancré doit être écrite (taux importé), sinon il défend la parité par les taux et les crises de change disparaissent.
\item Les attaques de seconde génération n'existent que si le gouvernement peut choisir d'abandonner ; sans clause de coût, la défense par les taux gagne toujours.
\item Le contrôle des capitaux achète du temps, pas l'autonomie : la surchauffe passe par le déficit commercial.
\item La dédollarisation est un stock (hystérèse de 27 ans de demi-vie).
\item Le change accélère la désinflation après une hausse de taux mais n'en allège pas le coût réel : l'appréciation déplace la demande vers les importations. L'excès de sensibilité de l'investissement (facteur $\approx2$) reste le chantier de la calibration.
\end{enumerate}

\subsection{Étape 3 : actifs, bulles, ruées, assouplissement quantitatif}
\label{sec:proto3}
Le prototype de l'étape~1 a été étendu aux actifs : immobilier (stock par strate, prix, crédit hypothécaire, construction avec délai), actions (valeur de marché, prêts sur marge), effet richesse sur la consommation, ruées bancaires, assouplissement quantitatif avec coupon figé. Cinq tests, tous passés.

\paragraph{Forme opérationnelle retenue.} Les anticipations de plus-value sont \emph{réelles} (extrapolation de la hausse observée nette de l'inflation, v1.2) et la valeur fondamentale actualise les loyers au coût d'usage réel net de la croissance ; sans quoi l'inflation était comptée deux fois et la valeur fondamentale devenait négative dès $\pi^\ast=8\,\%$. Le prix immobilier suit l'excès du rendement attendu (loyer implicite plus plus-value extrapolée) sur le coût d'usage, avec un rappel vers la valeur fondamentale (loyers actualisés, poids $\psi_{fund}=0{,}5$) et une pression de l'offre nouvelle ; la construction vaut le remplacement plus une réponse des promoteurs au prix relatif, plafonnée à trois fois le remplacement et livrée après 26 semaines. Le crédit hypothécaire est borné par trois contraintes simultanées : le capital bancaire, un service de la dette inférieur à 30\,\% du revenu, et la solvabilité de la banque ; l'encours sur le parc existant tend vers $0{,}35\,(\mathrm{LTV}/0{,}8)\,p^ZZ$. Les défauts portent sur les prêts récents (30\,\% de l'encours, environ trois ans) originés à la quotité sur le prix moyen des trois dernières années : ils apparaissent quand le prix courant passe sous $\mathrm{LTV}$ fois ce prix d'origine. Le cours des actions égalise la part de richesse désirée à la capitalisation, rappelé vers la valeur fondamentale (dividendes actualisés au rendement requis net de la croissance, \emph{taux d'actualisation plancher 3\,\%} : sans plancher la valeur fondamentale divergeait à taux zéro et transformait la trappe en hyperinflation, v1.3), avec prêts sur marge (marge $\times5\,\%$ de la capitalisation), appels de marge et pertes bancaires. L'effet richesse ($1{,}5\,\%$ par an sur chaque classe) s'ajoute à la consommation \emph{sans} entrer dans la cible d'Euler, qui ne porte que sur la richesse liquide.

\begin{center}\small
\begin{longtable}{p{4.0cm}p{8.4cm}p{1.4cm}}
\toprule
\textbf{Test} & \textbf{Résultat} & \textbf{Verdict}\\
\midrule
\endfirsthead
\toprule
\textbf{Test} & \textbf{Résultat} & \textbf{Verdict}\\
\midrule
\endhead
A1 État stationnaire avec actifs & u 3,3\,\%, $\pi$ 2,3\,\%, \textbf{taux réel 3,0\,\% contre 1,1\,\% sans actifs}, taux naturel estimé 2,6\,\% $=$ théorique 2,5\,\% ; richesse immobilière 1,1 PIB, actions 1,6 PIB, crédit hypothécaire 0,40 PIB & PASS\\
A2 Bulle immobilière ($\chi_Z=0{,}7$, LTV 0,95) & pic $\times2{,}3$ à dix ans, creux $-73\,\%$ en deux ans, créances douteuses 2,7\,\%, chômage 6,2\,\% au pic puis retour ; à $\chi_Z=0{,}3$ : $\times1{,}7$ sans krach & PASS\\
A3 QE à la borne zéro (choc $-8\,\%$ PIB, QE 10\,\% PIB/an) & $-8$ pb sur le souverain, actions $+1{,}3\,\%$, immobilier $+2\,\%$, $+1$ pt d'inflation, $M+22\,\%$, \textbf{aucun effet mesurable sur l'emploi} ; pertes de la banque centrale supérieures de 16 pts de PIB & PASS\\
A4 1929 synthétique (sans garantie des dépôts, prêteur en dernier ressort plafonné par l'or, marge 0,9, $\chi_E=0{,}7$, hausse de 3 pts) & bourse $\times2{,}2$ puis $-98\,\%$ ; 94 ruées, faillite bancaire, déflation $-23\,\%$, chômage 25\,\% ; \emph{avec} garantie des dépôts : aucune ruée, déflation $-5{,}5\,\%$, chômage 4,5\,\% & PASS\\
A5 Stabilité stochastique (12 graines) & chômage max 6,1\,\%, $|\pi|$ max 6,4\,\%, $p^Z\in[0{,}68\,;1{,}89]$, identités exactes & PASS\\
\bottomrule
\end{longtable}
\end{center}

\paragraph{Enseignements de l'étape 3.}
\begin{enumerate}
\item \textbf{L'écart de taux réel se ferme.} La prime de risque $\rho_E$, capitalisée dans les prix d'actifs et rendue à la consommation par l'effet richesse, porte le taux réel d'équilibre de 1,1 à 3,0\,\% : le taux naturel estimé par la banque centrale rejoint $\rho+\sigma g$. Le chantier ouvert depuis la v0.9 est clos, et la règle de Taylor peut s'ancrer sur le taux théorique sans entretenir de sous-emploi.
\item \textbf{Trois bornes rendent les bulles survivables} : le capital bancaire, la capacité de remboursement des ménages et la capacité de construction. Sans elles, toute bulle finit en hyperinflation par création monétaire hypothécaire ; avec elles, une bulle à $\chi_Z=0{,}7$ fait $\times2{,}3$ puis $-73\,\%$ et l'économie se relève en quelques années.
\item \textbf{Les institutions décident de l'amplitude} : mêmes chocs et mêmes anticipations, une dépression à 25\,\% de chômage sans garantie des dépôts ni prêteur en dernier ressort, une récession à 4,5\,\% avec. C'est la propriété annoncée au \S\ref{sec:crises}, vérifiée.
\item \textbf{Le QE est un instrument d'actifs, pas d'emploi} : dans cette économie il abaisse les rendements, soutient les prix d'actifs et coûte cher au bilan de la banque centrale, sans effet mesurable sur le chômage. Le cas « réserves oisives » du \S\ref{sec:qe} exige des banques détentrices de titres publics, ce que la structure de détention du prototype (épargne absorbée par les ménages) ne produit pas encore.
\item \textbf{Le bloc d'actifs est fragile à la calibration} : sept versions ont divergé avant d'être stables. Les fautes, dans l'ordre : prix sans valeur fondamentale (bulle spontanée à $\chi=0{,}3$) ; crédit hypothécaire non borné (3\,\% du PIB de monnaie créée par semaine) ; richesse illiquide entrée dans la cible d'Euler ; prêts sur marge créés instantanément ; cible de fonds propres bancaires ignorant les hypothèques ; pression d'offre trop forte (cycle du porc à six mois) ; banque centrale sans portefeuille en contrepartie de ses réserves. Chaque correction est une règle de bord ou une calibration cohérente, aucune n'ajoute de mécanisme.
\end{enumerate}

\subsection{Audit du prototype avant l'étape 4}
\label{sec:audit}
\paragraph{Méthode.} Trois séries : (a) soixante ans sans choc ; (b) quarante valeurs de leviers modifiées \emph{à l'an 5} d'une économie calibrée, effets mesurés aux ans 15--25 ; (c) huit politiques extrêmes maintenues vingt ans. Critères d'alerte : chômage $>25\,\%$, inflation $>30\,\%$, PIB $\pm15\,\%$, niveau de vie $+10\,\%$ (méta-stratégie candidate), réforme monétaire déclenchée.

\paragraph{État stationnaire à soixante ans.} u 3,0\,\% (max 4,5), inflation 2,2\,\% (max 3,5), taux réel 2,8\,\%, croissance 1,8\,\%, aucune faillite, identités à $10^{-12}$. Dérives lentes : la dette fond (0,61 $\to$ 0,04, excédent structurel), le prix immobilier nominal fait $\times4{,}3$ (prix réel stable), les actions restent 15 à 35\,\% sous la valeur comptable. Aucune ne diverge ; la première appelle une règle budgétaire pour l'état de référence.

\paragraph{Leviers (base : u 3,0\,\%, $\pi$ 2,0\,\%, $r$ 2,2\,\%).}
\begin{center}\small
\begin{tabular}{p{4.6cm}p{6.2cm}p{4.2cm}}
\toprule
\textbf{Levier} & \textbf{Effet aux ans 15--25} & \textbf{Lecture}\\
\midrule
Expansion budgétaire (G 8,5 $\to$ 16\,\% ; minima $\times2$ ; impôts sur le travail, TVA ou cotisations à zéro) & u au plancher (1,8\,\%), inflation 5--6\,\%, taux réel 10--12\,\%, dette 1,0--1,1 puis rationnement des dépenses (54--82\,\%), crédibilité 0,5--0,7, \textbf{PIB $-10\,\%$} (I/Y 11 $\to$ 3\,\%), niveau de vie 0 à $+2{,}5\,\%$ & Pas de repas gratuit : l'éviction de l'investissement domine\\
Austérité (G 2\,\%, TVA 30\,\%, cotisations 40\,\%, minima 0) & u 6--12\,\%, niveau de vie $-10$ à $-18\,\%$ & Coût réel immédiat\\
Investissement public 3 $\to$ 8\,\% & u 1,8\,\%, inflation 5\,\%, dette 0,9, niveau de vie $+2\,\%$ & Rendement long, surchauffe courte\\
Quotité 0,5 / 1,0 & $\pm2\,\%$ de PIB & Pas de bulle sans extrapolation\\
Marge boursière 1,5 & u 9,6\,\% (max 20), krach, niveau de vie $-8\,\%$ & Le plus dangereux des leviers financiers\\
Extrapolation $\chi=0{,}9$ & cycles, chômage max 13\,\% & Attendu\\
Exigence de capital 0,02 / 0,15 & quasi nul sans choc & Ne compte qu'en crise (test avec choc à faire)\\
Cible d'inflation 0 / 6\,\% & atteintes (0,0 / 6,2\,\%) & La règle avec apprentissage tient\\
Marge bancaire 0 / 5\,\% ; prime de risque 0 / 5\,\% & taux réel 3,1 / $-1{,}0$ ; 3,8 / 0,5\,\% & Canaux cohérents\\
Garantie des dépôts, étalon-or & nul hors crise & Institutions dormantes\\
\bottomrule
\end{tabular}
\end{center}
Politiques extrêmes : taux à 8\,\% permanent, coût net sans effondrement ; QE 5\,\% du PIB par an permanent, sans effet réel ; monétisation 2\,\% par an, inflation 3\,\% et crédibilité nulle (c'est la capacité de seigneuriage naturelle) ; réformes monétaires tous les deux ans, crédibilité 0,02 sans divergence (levier non exploitable) ; \textbf{taux à zéro permanent} et \textbf{monétisation à 5\,\% permanente} : hyperinflation, réformes en série, mort économique (chômage 75--90\,\%, PIB $-72\,\%$). Ces deux cas sont les seuls effondrements ; l'état de mort reste plus violent que l'histoire (Zimbabwe : $-50\,\%$) et devra être adouci par une économie de survie (troc, dollarisation spontanée) à l'étape~4.

\paragraph{Méta-stratégies.} Aucune stratégie dominante après corrections : impôts à zéro, dépenses maximales, QE permanent, monétisation modérée, réformes répétées, taux bas, quotité maximale sont tous payés (inflation, dette et rationnement, crédibilité, éviction, krach) ou neutres. Vigilance : « impôts sur le travail à zéro » garde $+2{,}5\,\%$ de niveau de vie à vingt ans pour 5,7\,\% d'inflation et une dette de 1,1 — la sanction (prime, saut, défaut) est au-delà de l'horizon testé.

\paragraph{Faille méthodologique.} Un premier balayage modifiait les paramètres \emph{avant} la calibration ; chaque valeur produisait une économie initiale différente (un G à 16\,\% calibrait un PIB réel 44\,\% plus grand) et les comparaisons de niveau de vie étaient sans objet ($+62\,\%$ apparents). Règle, pour l'audit comme pour le jeu : les configurations de la \S\ref{sec:configs} se calibrent à l'état initial ; les leviers agissent ensuite.

\subsection{Étape 4 : hétérogénéité, et les cinq chantiers de l'audit}
\label{sec:proto4}
Dix tests, tous passés ; les trente-et-un tests antérieurs restent verts.
\begin{center}\small
\begin{longtable}{p{4.4cm}p{8.0cm}p{1.4cm}}
\toprule
\textbf{Test} & \textbf{Résultat} & \textbf{Verdict}\\
\midrule
\endfirsthead
\toprule
\textbf{Test} & \textbf{Résultat} & \textbf{Verdict}\\
\midrule
\endhead
C1 Règle budgétaire (60 ans) & dette stabilisée entre 0,40 et 0,45 (0,61 $\to$ 0,04 sans règle), u 2,7\,\%, $\pi$ 2,3\,\%, correction moyenne 6\,\% & PASS\\
C2 Machine à états (monétisation 5\,\% PIB dix-huit ans, réforme à l'an 18) & effondrement à l'an 15,6 ; survie : u 36\,\%, PIB 97\,\% du niveau initial ; après réforme : u 1,8\,\% à l'an 21, 2,1\,\% à l'an 25, PIB $+18\,\%$ ; pas de rechute & PASS\\
C3 Bulle sous exigence de capital 0,02 vs 0,09 & fonds propres au creux 2,0 contre 8,1\,\% du PIB, chômage 12,5 contre 12,1\,\% ; l'émission d'actions évite la faillite dans les deux cas & PASS\\
C4 Horizon des sanctions (impôts sur le travail à zéro, 40 ans) & niveau de vie $+2{,}4\,\%$ (ans 15--25) puis $-10\,\%$ (ans 30--40), PIB $-18\,\%$, dette 1,26, taux souverain max 28\,\%, dépenses rationnées à 66\,\% : la sanction arrive entre l'an 25 et l'an 30 & PASS\\
C5 Actions : valeur fondamentale & cours moyen 0,75 (ans 20--60) ; la sous-valorisation est un trait de calibration (dividendes/valeur comptable), pas un défaut ; plancher d'actualisation indispensable & PASS\\
H1 Cohortes (40 ans) & âgés 15 $\to$ 25\,\%, dette 1,30 contre 0,38, chômage 2,3 contre 3,7\,\%, croissance 1,0 contre 1,9\,\% & PASS\\
H2 Lewis (réservoir 50\,\%) & ruraux 299 $\to$ 119 en quarante ans, chômage urbain $+0{,}8$ pt, PIB $+37\,\%$ & PASS\\
H3 Climat (récolte $-30\,\%$) & inflation $2{,}1\to6{,}4\,\%$, niveau de vie $-12{,}5\,\%$, retour en trois ans & PASS\\
H4 Ressource (120 ans de réserves, 80 ans) & plateau quarante ans, prix relatif de l'énergie $\times1{,}9$, PIB $-18\,\%$ à quatre-vingts ans, chômage 5\,\% & PASS\\
H5 Archétype agraire & niveau de vie $-8{,}7\,\%$, volatilité de l'inflation $\times50$, chômage 7,9\,\%, croissance 2,4\,\% (absorption du réservoir) & PASS\\
\bottomrule
\end{longtable}
\end{center}
\paragraph{Enseignements.} (1) Le Lewis, le vieillissement et le canal alimentaire émergent des blocs sans paramètre ad hoc. (2) L'épuisement d'une ressource est absorbable si des substituts existent : le plancher de productivité décide. (3) Toute actualisation du modèle doit avoir un taux plancher (deux valeurs fondamentales ont failli dans le même sens). (4) La mort économique est un régime, pas un mécanisme : elle se traite par une machine à états, et sa sortie exige de réinitialiser toutes les grandeurs nominales de référence. (5) Le premier balayage d'audit avait comparé des économies calibrées différemment ; les leviers s'évaluent en cours de partie.

\subsection{Variantes de calibration testées pour l'investissement (à essayer en économie ouverte)}
Le prototype retient l'option B. Les autres sont disponibles par simple changement de paramètres et documentées pour l'étape~2.
\begin{center}\small
\begin{tabular}{p{5.2cm}cccccp{4.2cm}}
\toprule
\textbf{Option} & $I/Y$ & $K/Y$ & part sal. & $r$ réel & $u$ & \textbf{Commentaire}\\
\midrule
B. $\rho_E=2\,\%$ (retenue) & 0,18 & 2,9 & 0,65 & 1,3\,\% & 3,9\,\% & part salariale $\approx1-\alpha$, stabilité excellente\\
C. $\rho_E=3\,\%$ & 0,22 & 3,1 & 0,60 & 0,7\,\% & 4,6\,\% & $K/Y$ au repère, taux réel trop bas\\
D. $\rho_E=5\,\%$ & 0,17 & 3,1 & 0,53 & $\approx0$ & 12\,\% & trappe à liquidité : la demande ne suit plus\\
E. $\rho_E=2\,\%$, achats publics 9\,\% & 0,17 & 2,7 & 0,66 & 1,7\,\% & 3,7\,\% & compromis ; dette stable à 0,40\\
F. $\rho_E=3\,\%$, effet richesse doublé ($c^V_H=0{,}12$) & 0,21 & 3,0 & 0,61 & 0,8\,\% & 4,3\,\% & anticipe l'étape~3 (richesse en actifs)\\
G. $\rho_E=3\,\%$, baisse d'impôts & 0,13 & 2,3 & 0,64 & 2,4\,\% & 3,6\,\% & inefficace : le déficit fait monter le taux et évince l'investissement\\
\bottomrule
\end{tabular}
\end{center}
La lecture d'ensemble : relever $K/Y$ exige des profits plus élevés, donc des salaires réels plus bas, donc une demande et un taux naturel plus bas. Cet arbitrage est intrinsèque à une économie fermée sans actifs ; il se dénoue quand la prime de risque est capitalisée dans des prix d'actifs qui soutiennent la consommation, ce que l'étape~3 apportera.

%=====================================================================
\section{Calibration initiale proposée}
\label{sec:calib}
%=====================================================================
Valeurs de départ (à ajuster lors du prototypage ; base annuelle sauf mention) :
\begin{center}\small
\begin{longtable}{p{4.2cm}p{3.0cm}p{7.6cm}}
\toprule
\textbf{Paramètre} & \textbf{Valeur} & \textbf{Source / justification}\\
\midrule
\endfirsthead
\toprule
\textbf{Paramètre} & \textbf{Valeur} & \textbf{Source / justification}\\
\midrule
\endhead
$\alpha_j$ (part du capital) & 0,30--0,40 & Comptabilité nationale ; plus élevé pour énergie et intermédiaires\\
$\sigma$ & 1 (log) & Élasticité de substitution intertemporelle unitaire\\
$\beta$ & $1/(1+\rho)$ & \\
$c_B, c_M, c_H$ & 0,95 ; 0,80 ; 0,50 & Propensions à consommer décroissantes (Kaldor)\\
$c^V_h$ & 0,03 & Effet richesse\\
$\gamma_{hj}$ & alimentation 40\,\% du panier bas & Subsistance ; nul pour les services haut de gamme\\
$\kappa_j$, $\bar\kappa$ & 0,10 ; 0,05/semaine & Vitesse de tâtonnement ; borne d'ajustement\\
$\varpi$ (indexation des prix) & 0,5 par tick, sur $\pi^e/52$ & Coefficient de vitesse d'indexation des prix \eqref{eq:price} (v1.5 : la ligne « 0,02/mois » de la v1.4 était mal libellée)\\
$\phi_u$ (Phillips) & 0,3 & Pente salaire--chômage\\
$\lambda_L$ & 0,25/semaine & Ajustement de l'emploi\\
$\bar\ell$, $\eta_\ell$ & 0,6 ; 2 & Levier de référence ; rationnement\\
$\kappa^{CAR}$ & 8--10\,\% & Bâle\\
$m^L$, $m^D$ & 2\,\% ; 1\,\% & Marges bancaires\\
$a_\pi$, $a_y$ (Taylor) & 1,5 ; 0,5 & Taylor (1993)\\
$\lambda_0$ (anticipations) & 0,2/mois & Adaptatif lent ; $\lambda_1 = 0{,}6$ au-delà de $\bar\pi^{hyper}=50\,\%$/an\\
$\upsilon$ (ancre quantitativiste) & 0,05/mois & \\
$\nu_1,\nu_2,\nu_3$ (crédibilité) & 0,01 ; 0,5 ; 0,2 & \\
$\epsilon_j$ (Armington) & 2--4 & Élasticités de substitution ; faibles pour l'énergie\\
$\varphi_{ii}/\varphi_{ik}$ & 3 & Biais domestique\\
$\hbar$ (décote) & 0,3--0,6 & Restructurations historiques\\
$\varepsilon_1,\varepsilon_2$ (planification) & 0,005 ; 0,003 & Perte de croissance de la PGF de 0,5 à 1,5 pt/an\\
$\eta_G$, $\zeta$ & 0,02 ; 0,5 & Rendement du capital public\\
$\alpha_T$ (terre) & 0,25 & Part de la rente foncière dans l'agriculture\\
$\lambda_{mig}$ & 0,02/an par unité d'écart & Harris--Todaro\\
$\eta_{LBD}$ & 0,02 & Apprentissage par la pratique dans l'industrie\\
$\kappa^w$, $\sigma(\varepsilon^w)$ & 0,15 ; 0,25/an & Volatilité des matières premières\\
$\chi_Z, \chi_E$ & 0,2 (calme) -- 0,7 (bulle) & Poids de l'extrapolation ; seuil de bulle $\approx 0{,}5$\\
$\mathrm{LTV}$, $m$ & 0,7--0,95 ; 0,5--0,9 & Instruments macroprudentiels du joueur\\
$\delta_Z$, $T_Z$ & 0,015 ; 24 mois & Dépréciation et délai de construction\\
$\xi_M$, $a$, $V_0$ & 0,02/mois ; 2 ; 1,5 & Encaisses réelles ; Cagan : $\pi^{max} = 1/a$\\
$\varsigma_1..\varsigma_4$ (ruée) & 5 ; 10 ; 0,8 ; 3 & \\
$\eta_Z$ (zombies) & 0,01 & Caballero--Hoshi--Kashyap\\
$s_8$, $s_9$ & 50 pb ; 150 pb & Backstop ; répression financière ($s_7$ retiré en v1.1 : voir $\psi_{QE}$)\\
$\bar T$, $\bar T_0$ & 4--12 ans ; 7 ans & $\theta=1/\bar T+$déficit/stock : 0,12 à $\bar T=8$--9 ans et déficit nul ; 0,17 à 6 ans\\
$\varsigma_6$, $\varsigma_7$, $\varsigma_9$, $\Delta\rho^{saut}$, $\lambda_\rho$ & 3 ; 20 ; 1 ; 300 pb ; 0,05/mois & Saut de prime (voir $\varsigma_0$, $\varsigma_8$, $\varsigma_{11}$ ci-dessous)\\
$\Delta^{res}$, $\underline\kappa^{CB}$, $\bar E^{CB}$ & 0--25 pb ; $-0{,}10$ ; 2\,\% de $H$ & Rémunération des réserves ; tolérance politique aux fonds propres négatifs (10\,\% de $H$) ; fonds propres statutaires\\
$\theta_6$ & 0,3 & Effet de stock des réserves sur le change\\
$\nu_9$, $\delta_M$ & 0,12 ; 0,02/an & Crédibilité : plafond par abandon de monnaie (six abandons plafonnent $c$ à 0,28) ; décroissance très lente de la mémoire\\
$\phi_0$, $\kappa_d$, $\varepsilon^{dedol}$ & 0,995/mois ; 0,4 ; 0,01 & Hystérèse de la dollarisation ; efficacité et coût de la contrainte légale\\
$\bar r^F$, $\beta_F$, $\sigma(\varepsilon^F)$ & 4\,\% ; 0,7 ; 8\,\%/an & Fonds souverain ; règle de prélèvement 3\,\%\\
$\psi_B, \psi_M, \psi_H$, $\overline{cov}$ & 0,15 ; 0,30 ; 0,20 ; 0,25 & Épargne de précaution : à couverture nulle, $c_B=0{,}81$, $c_M=0{,}56$, $c_H=0{,}40$ ; couverture pleine à 25\,\% du PIB de dépenses sociales\\
$\eta_u$, $\eta'_u$, $u^n_0$ & 0,03/mois ; 0,01/mois ; 5\,\% & Hystérèse du chômage (Blanchard--Summers) ; deux ans à $u=u^n+4$ pts relèvent $u^n$ d'environ 1,5 pt\\
$s_1$ & 0,06 & Convexité en excès relatif $(b^{mkt}/\bar b_i-1)^2$ sur la dette de marché, portée par $\max(\varkappa^{for},\underline\varkappa(1-\chi^K\varrho^B))$ (réécritures v1.1 et v1.5)\\
$s_0$, $\bar T_0$ ; $s_2$, $s_3$, $s_4$, $s_6$ & 100 pb à $\bar T_0=7$ ans ; 30 pb/pt ; 150 pb par défaut (décroissance 5\,\%/an) ; 300 pb ; 40 pb & Prime de terme proportionnelle à la maturité (v0.8 : fusion de $s_0$ et $s_T$) ; dérive projetée $\Delta b^{proj}$ ; historique de défaut ; passage de $c=1$ à $c=0$ ; privilège de l'émetteur de réserve\\
$\bar b^{A..E}$, $(T/Y)^{ref}$ & 0,8 ; 0,95 ; 0,6 ; 0,4 ; 0,5 ; 0,35 & Seuils par régime, multipliés par la capacité fiscale relative $\min(1{,}5,\,(T/Y)/0{,}35)$ (v0.7) : 0,17 pour un pays collectant 10\,\% du PIB\\
$\varsigma_0$, $\varsigma_8$, $\varsigma_{11}$ & $-2{,}55$ ; 1,05 ; 0,5 & Saut : $\Phi^{-1}(0{,}02)=-2{,}05$ et $\Phi^{-1}(0{,}001)=-3{,}09$ fixent la constante et l'écart de backstop ; contagion (v0.7)\\
$\varsigma^{mut}$ & 0,2 (régime B) & Part mutualisée des pertes de \eqref{eq:ecb} et des défauts (v0.7)\\
$\beta^{dom}$ & 0 (fonds externe) à 0,6 (rentier ancré) & Part du fonds souverain placée dans le pays (v0.7)\\
$\varpi_E$ & 0,3 (1,0 en \emph{currency board}) & Seuil de réserves du régime E en proportion de $M$ (v0.7)\\
$Rem/Y$ & 0--10\,\% & Remises des migrants, dans $Tr_h$ et $CA$ (v0.7)\\
$\lambda^{fisc}$ & 24 mois & Délai de sortie d'une règle constitutionnelle\\

\midrule
\multicolumn{3}{l}{\emph{Révisions issues du prototype (v0.9)}}\\
$\rho_E$ & 0,02 & Prime de risque sur le capital (option B ; variantes 0,03--0,05 en \S\ref{sec:proto})\\
$\iota$, plafond & 0,20 ; $2\delta$ & Sensibilité de l'investissement à $q$ ; investissement brut $\le3\delta K$\\
$\underline u$, $\bar u$ & 0,80 ; 0,95 & Fenêtre d'utilisation des capacités conditionnant l'investissement net\\
$\delta_j$ & 0,035 ; 0,045 ; 0,04 ; 0,05 & Dépréciation par secteur (l'ancienne valeur moyenne 6,4\,\% donnait un remplacement de 19\,\% du PIB)\\
$\psi$, $\bar m$ & 0,9 ; 8 semaines & Dividendes prioritaires ; trésorerie cible\\
$c_h$, $c^V_h$ & (0,95 ; 0,90 ; 0,80) ; (0 ; 0,04 ; 0,08) & Propensions révisées : à $-20\,\%$ l'économie tombe en trappe à liquidité\\
Répartition des salaires & 55 / 30 / 15\,\% & Masse salariale par strate ; revenu par tête B:M:H = 1 : 1,1 : 2,0\\
$\gamma_A$, $\phi_x$ & 1,0 ; 0,3 & Pass-through complet de la productivité ; surenchère en pénurie de main-d'œuvre (Phillips convexe)\\
$\mu_c$, marge normale $m_j$ & 0,05/semaine ; $(r^\ast+\rho_E)p_KK_j/(c^u_jY_j)$ & Rappel symétrique vers le coût majoré\\
$\lambda_r$, $\lambda_s$, $\rho_i$, $\bar\imath$ & 0,02 ; 0,10 ; 0,95 ; 1,0 & Apprentissage de $\hat r^\ast$ ; lissage de l'inflation de la règle ; lissage du taux ; plafond\\
$\rho$, $m^L$, $m^D$ & 0,01 ; 0,01 ; 0,01 & Préférence pour le présent ($r^\ast_{th}=\rho+\sigma g_0=2{,}5\,\%$) ; marges bancaires\\
$\tau_{W,h}$, $\tau_\Pi$, $G/Y$, $G^{inv}/Y$ & (0,05 ; 0,12 ; 0,25) ; 0,20 ; 0,085 ; 0,03 & Budget quasi équilibré à l'état stationnaire (dette stable à 0,4--0,8)\\
$\eta_G$, $\zeta$, $\delta_G$, $K^G_0/Y$ & 0,02 ; 0,5 ; 0,03 ; 0,5 & Capital public dans la PGF (neutre au ratio initial)\\
$\phi_F$, réorganisation & 0,05 ; 13 semaines & Faillite non destructrice\\
Recapitalisation & $1{,}2\,\kappa^{CAR}L$, $\le$ 1/an & Sauvetage public automatique d'une banque insolvable\\
Réforme & $P>10^6$ ; $\pi^e\leftarrow\pi^\ast+0{,}05$ ; $c\leftarrow0{,}8\times$plafond & Redénomination et ancre nouvelle \eqref{eq:reform}\\
Bornes & prix $[0{,}7;1{,}19]$/sem., salaires $[0{,}9;2]$/mois, $\pi^e\le12$, $v$ saturée à $\pi^e=1$, prime $\le0{,}3$ & Régimes extrêmes seulement\\
Commandes suivies, stock cible & $\le1{,}1\times$offre ; 4 semaines & Amortisseurs du cycle des stocks\\
\midrule
\multicolumn{3}{l}{\emph{Révisions et ajouts de la v1.0}}\\
$s_q$, $\lambda_q$ & 0,03 ; 1/52 & Réponse saturante de l'investissement autour du $q$ de long terme\\
$\lambda_L$ (équipement), $\lambda_{div}$ & $\lambda_L/2$ ; 0,1 & Rétention de main-d'œuvre dans les biens d'équipement ; dividendes sur profit lissé\\
$\epsilon$, $\varphi_{ii}$, $\tau^M$, $\lambda_\varsigma$ & 3 ; 0,8 ; 2\,\% ; 0,1/sem. & Armington : élasticité, biais domestique, droits de douane, ajustement des parts (part importée 16--20\,\% du PIB)\\
$\epsilon_{\acute{e}n}$, $\varphi_{\acute{e}n}$ & 2 ; 0,6 & Énergie échangeable (option)\\
$\theta_L$, $\lambda_e$, $\kappa_e$, $\kappa_K$, $\kappa_{spec}$ & 2 ; 0,3/sem. ; 0,3 ; 2/an ; 0,5 & Change : niveau cible, rappel, pression des flux ; flux de capitaux ; spéculation\\
$R^{FX}_0$, $\underline R$, dévaluation & 8 mois ; 2 mois d'importations ; 30\,\%, perte de crédibilité 0,3 & Régime E\\
Seconde génération (option) & $\pi^{dev}>0{,}05$ et $u>u^n+0{,}04$ & Clause d'abandon d'Obstfeld\\
$s^{FX}$ : pente, seuil, $\phi$ & 2 ; 5\,\% ; 0,9995/sem. & Dollarisation, hystérèse\\
\midrule
\multicolumn{3}{l}{\emph{Ajouts de la v1.1}}\\
$\varsigma^{\text{ém}}$, seuils d'émission & 2\,\% des dépôts/sem. ; déclenchement à $1{,}1\,\kappa^{CAR}$, cible $1{,}3\,\kappa^{CAR}$ & Émission d'actions bancaires \eqref{eq:bankeq}\\
$\varpi^{Bk}$, borne & 5\,\%/sem. du capital excédentaire ; facteur borné à $[1;3]$ & Dividendes bancaires accélérés quand le capital est sans emploi\\
Seuil de recapitalisation publique & $E^{Bk}<\tfrac12\kappa^{CAR}L$, au plus une fois par an & Sous-capitalisation et non seulement insolvabilité\\
$\psi_{QE}$, $\underline\varkappa$ & 0,5 ; 0,3 & Canal de portefeuille sur la prime de terme ; plancher de la part exposée, conditionnel depuis la v1.5 \eqref{eq:spread}\\
$\lambda_{lop}$ & 0,05/sem. (option) & Rappel vers le prix mondial (énergie, matières premières)\\
$\rho_Z$, $\delta_Z$, $T_Z$, $\psi_{fund}$, $\nu_{build}$ & 2\,\% ; 2\,\%/an ; 26 sem. ; 0,5 ; plafond 3$\times$ le remplacement & Immobilier : prime de risque, dépréciation, délai de construction, rappel fondamental, promoteurs\\
LTV, marge, $\chi_Z$, $\chi_E$ & 0,8 (0,95 en bulle) ; 0,5 (0,9 en 1929) ; 0,3 (0,7 en bulle) ; 0,3 (0,7) & Curseurs macroprudentiels et extrapolation\\
$c^V_Z$, $c^V_E$, $L^Z/Y$, $L^E$ & 1,5\,\%/an chacun ; 0,40 ; marge $\times5\,\%$ de la capitalisation & Effet richesse (hors cible d'Euler) ; encours hypothécaire ; prêts sur marge\\
Ruées & $\Phi(40(\rho^{NPL}-1\,\%) + 5\pos{-E^{Bk}}/D - 2\,\ind_{\text{garantie}} - 1{,}5)$, mensuel & Probabilité de retrait \eqref{eq:run}\\
\midrule
\multicolumn{3}{l}{\emph{Ajouts de la v1.2}}\\
$\underline u^n$ & 3\,\% & Plancher frictionnel du chômage naturel \eqref{eq:un} (distinct de $u^n_0=5\,\%$, niveau structurel initial ; v1.5 : les deux symboles sont séparés dans le glossaire)\\
$\omega_G$ & 0,15 & Poids des services publics dans le niveau de vie \eqref{eq:sol}\\
Plafond d'absorption bancaire, $\kappa^{hB}$ & 40\,\% du bilan ; 20/pt, $\le5\,\%$ des dépôts/sem. & Placement de la dette publique\\
Réserves rémunérées & $\min(Res, B^{CB}+A^G+L^{CB})$ & Portage de la banque centrale \eqref{eq:ecb}\\
\midrule
\multicolumn{3}{l}{\emph{Ajouts de la v1.3}}\\
$\phi_b$, $b^\ast$ & 0,3 ; 0,6 & Règle budgétaire \eqref{eq:fiscrule}\\
Effondrement : $\bar\pi^e$, durée, réformes ; $\upsilon^{surv}$, $\pi^{surv}$, $\varsigma^{conv}$ & 1 200\,\%, 26 sem., 2 en 5 ans ; 0,6 ; 50\,\% ; 0,6 & Machine à états \eqref{eq:collapse}\\
$\mathcal R_0$, $\varrho^{plat}$, $\underline a$ & 150 ans ; 0,5 ; 0,4 & Ressource épuisable\\
$\alpha_T$, $\sigma_\omega$ & 0,25 ; 0 (0,05--0,08 en agraire) & Terre et climat\\
Cohortes : naissances, activité, retraite, pension & 1,2\,\%/an ; 20 ans ; 45 ans, 18 ans ; 50\,\% & Démographie\\
$w^{rur}$, $\mu$, $c^{mig}$ & 0,5 ; 0,05/an ; 0,2 & Migration rurale (Lewis)\\
Plancher d'actualisation & 3\,\% & Valeurs fondamentales (actions, immobilier)\\
\midrule
\multicolumn{3}{l}{\emph{Ajouts de la v1.4}}\\
$\phi_{liq}$ & 0 (0,3 en option) & Ratio de liquidité : titres publics détenus par la banque en part des dépôts\\
Natalité, $\eta_f$, $\eta_m$ & 1,4\,\% des actifs ; 0,5 ; 0,3 & Cohortes : population stable ; fécondité et mortalité endogènes (option)\\
Bornes de jeu & $\chi\in[0;0{,}9]$, LTV $[0{,}3;1]$, marge $[0;1{,}5]$, $\kappa^{CAR}$ $[0{,}02;0{,}3]$, $\pi^\ast$ $[0;0{,}10]$ & Hors bornes : cycles d'actifs extrêmes\\
$\psi^{E}_{fund}$ & 2 & Rappel des actions vers la valeur fondamentale\\
\midrule
\multicolumn{3}{l}{\emph{Ajouts et révisions de la v1.5}}\\
$\varrho^B$, $\chi^K$ (dans \eqref{eq:spread}) & 0 par défaut ; 1 pour la configuration 5 & Répression financière et contrôle des capitaux, désormais dans le prototype ; leur produit conditionne le plancher $\underline\varkappa$\\
$\varpi_w$ & 1,0, affaibli par $\omega_u$ & Degré d'indexation des salaires \eqref{eq:wage}, borné par 1 ; c'est lui qui « tend vers 1 » en hyperinflation, non $\varpi$\\
$\eta_H$, $\bar g^{edu}$ & 0,01/an ; 5\,\% du PIB par tête & Capital humain (\S\ref{sec:etat}) : spécifié, non simulé\\
$\zeta_q$ & 0,1 & Poids des files d'attente dans \eqref{eq:sol} : spécifié, non simulé avant l'étape~5\\
$q^\ast$ & $1+g_0/(\iota\Gamma)$ & Référence de l'indicateur d'investissement du joueur\\
\bottomrule
\end{longtable}
\end{center}
\paragraph{Paramètres fixés en dur.} Les paramètres qui n'influencent aucune configuration de la \S\ref{sec:configs} sont hors de l'espace d'équilibrage : $\mu$ et $\lambda_{inv}$ (moyenne mobile et stocks), $s^\ast_j$, $\lambda_S$, $T_K$, $\epsilon$, $\lambda_\rho$, $\delta_M$. Les paramètres retirés en v0.8 ($s_5$, $s_T$, $s_{10}$, $\nu_5$, $\bar\varphi$, $\nu_6$, $\lambda_2$, $\vartheta$, $\varrho^{tier}$, $D^{CB}$, $\mu^{FX}$, $\delta_i$), en v1.0 ($\theta_1$, $\theta_2$, $\theta_3$, remplacés par $\theta_L$, $\lambda_e$, $\kappa_e$, $\kappa_K$, $\kappa_{spec}$), en v1.1 ($s_7$, remplacé par $\psi_{QE}$) et en v1.5 ($\eta_{agg}$) n'existent plus. Les lignes initiales de la table qui donnaient $\iota=0{,}10$, $\delta_j\in[0{,}05;0{,}08]$, $\rho=0{,}02$ et $u^n=5\,\%$ ont été supprimées en v1.5 : les valeurs en vigueur sont celles des blocs de révision ($\iota=0{,}20$, $\delta_j$ sectoriels, $\rho=0{,}01$, $u^n_0=5\,\%$ et $\underline u^n=3\,\%$). Règle de lecture : quand une équation et une table diffèrent, l'équation fait foi ; quand un bloc de révision et la table initiale diffèrent, la révision fait foi.

\subsection{Tests de cohérence à faire passer au prototype}
\label{sec:tests}
\begin{enumerate}
\item \textbf{État stationnaire} : sans chocs ni décision, un pays doit converger vers $u \approx u^n$, $\pi \approx \pi^\ast$, $g \approx g_0$, $b$ stable.
\item \textbf{Choc monétaire} : $+1$ pt de $i^{CB}$ pendant un an doit produire une baisse de l'écart de production de l'ordre de 0,5--1 pt et de l'inflation de 0,3--0,5 pt avec un pic à 12--18 mois (bosse VAR standard), une appréciation du change.
\item \textbf{Choc budgétaire} : $+1$ pt de PIB de transferts à la strate $B$ doit donner un multiplicateur de 0,8--1,5 selon le régime de change et l'écart de production.
\item \textbf{Monétisation} : un financement monétaire de 5\,\% du PIB par an pendant trois ans doit dégrader la crédibilité et faire décoller l'inflation avec retard.
\item \textbf{Planification} : le passage à $\varrho^{plan}=1$ avec $I/Y$ forcé à 35\,\% doit produire une décennie de croissance élevée puis un plafonnement (rendements décroissants, $\varepsilon^{plan}$).
\item \textbf{Cohérence stock-flux} : la somme de tous les bilans doit être nulle à chaque tick, à la précision numérique.
\item \textbf{Bulle et krach} : avec $\chi_Z = 0{,}7$, $\mathrm{LTV} = 0{,}95$ et un taux réel de 0\,\% pendant cinq ans, $p^Z$ doit s'écarter d'au moins 50\,\% de sa valeur d'usage, puis chuter après une hausse de 300 pb de $i^{CB}$, avec récession et hausse des $NPL$ ; avec $\chi_Z = 0{,}2$ et $\mathrm{LTV} = 0{,}7$, pas de bulle.
\item \textbf{QE bénin} : à la ZLB avec $q_j < 1$, un QE de 20\,\% du PIB auprès des banques doit laisser $M$ et $\pi$ presque inchangés et baisser $i^B$ de 50--100 pb.
\item \textbf{1929 synthétique} : étalon-or, pas de garantie des dépôts, $m = 0{,}9$ : une hausse de 200 pb doit produire une déflation cumulée $> 15\,\%$ et une chute de $M$ ; la même hausse avec garantie des dépôts et change flottant ne doit produire qu'une récession ordinaire.
\item \textbf{Maladie hollandaise} : un rentier dont $p^w_{res}$ double doit voir $e$ s'apprécier, $Y_{man}/Y$ baisser et $A_{man}$ décrocher de la frontière en dix ans, sauf si 80\,\% de la rente vont au fonds souverain.
\item \textbf{Point de Lewis} : un pays agraire peuplé doit connaître vingt ans de croissance industrielle à salaires stables, puis une accélération des salaires quand $L_{rur}/N < \bar\ell_{rur}$.
\item \textbf{Transmission du taux marginal au taux apparent} : avec $\bar T=8$ ans et un déficit de 4\,\% du stock, une hausse permanente de 200 pb de $i^B$ doit relever $i^{app}$ de 100 pb à six ans et de 160 pb à dix ans ; avec $\bar T=4$ ans, de 100 pb à trois ans. Avec $\omega=0{,}5$, la moitié de la hausse doit être répercutée dès le premier mois via $i^{res}$.
\item \textbf{Saut de prime avec et sans backstop} : un membre d'union à $b=1{,}2$ et $pb=-3\,\%$ doit subir un saut de prime avec une probabilité supérieure à 50\,\% dans les trois ans sans backstop, et inférieure à 5\,\% avec ; après le saut, la dynamique \eqref{eq:debtdyn} doit devenir explosive sans consolidation d'au moins 2 pts de PIB.
\item \textbf{Abandon de plancher par coût de bilan} : un pays de régime C avec $R^{FX}=1{,}2\,PY$ et $\underline\kappa^{CB}=-0{,}10$ doit abandonner son plancher après une pression d'appréciation cumulée de 10 à 15\,\% ; l'abandon doit produire une appréciation immédiate de 10 à 20\,\% et une déflation dans les douze mois.
\item \textbf{Stabilisation post-hyperinflation} : un pays de régime D avec $M^{hist}=5$ et $s^{FX}=0{,}8$ qui tient sa cible pendant cinq ans doit voir $c$ plafonner sous $1-\nu_9M^{hist}\approx0{,}4$ et $s^{FX}$ rester au-dessus de 0,55 ; une hausse de 500 pb de $i^{CB}$ doit réduire $\hat y$ cinq fois moins qu'en régime A.
\item \textbf{Seigneuriage et charge de la dette} : avec $\omega=0{,}5$ et $\bar c^{CB}=0{,}3\,\%$, une hausse de $i^{CB}$ de 0 à 2\,\% doit annuler le dividende $\Sigma$ dès la première année et le suspendre tant que $E^{CB}<\bar E^{CB}$, tandis que la charge d'intérêt consolidée monte de $\omega b\times2$ points de PIB.
\item \textbf{Chômage structurel émergent} (v0.7) : partant de $u^n_0=5\,\%$, un chômage maintenu à 8\,\% doit faire converger $u^n$ vers $7{,}3\,\%$ ($0{,}75\times8+1{,}25=7{,}25$) et un chômage maintenu à 2,5\,\% vers $3{,}1\,\%$, avec une demi-vie de convergence de $\ln 2/0{,}04\approx17$ mois.
\item \textbf{Probabilité de saut} (v0.7) : à $b-\bar b_i=0{,}25$, $pb=-3\,\%$, $c=0{,}85$, la probabilité mensuelle de \eqref{eq:jump} doit valoir $2{,}0\,\%\pm0{,}2$ sans backstop et $0{,}10\,\%\pm0{,}02$ avec ; un voisin en saut doit les porter à 6\,\% et 0,5\,\%.
\item \textbf{Capacité fiscale} (v0.7) : deux pays à $b=0{,}35$, $\varkappa^{dom}=0{,}5$, $pb=-2\,\%$, $c=0{,}8$, l'un collectant 40\,\% du PIB et l'autre 10\,\%, à dérive projetée $\Delta b^{proj}=+2$ pts identique, doivent avoir des primes de \eqref{eq:spread} séparées d'au moins 300 pb (calcul : 220 et 546 pb, écart 326).
\item \textbf{Détention domestique} (v0.7) : un pays à $b=2{,}3$, $\bar b_i=0{,}8$, doit payer une convexité de $105\pm10$ pb à $\varkappa^{dom}=0{,}95$ et de $1\,050\pm50$ pb à $\varkappa^{dom}=0{,}5$ ; à taux directeur constant, la boule de neige doit rester négative dans le premier cas tant que $\pi+g>i^{app}$.
\end{enumerate}

%=====================================================================
\section{Feuille de route}
%=====================================================================
\begin{enumerate}
\item \textbf{Prototype fermé} (\emph{réalisé}, \S\ref{sec:proto}) : un pays, quatre biens, trois strates, banque, banque centrale, État ; quinze tests passés ; état stationnaire à la cible d'inflation ; propriétés et chantiers ouverts identifiés (amplitude de la transmission monétaire, transitoire initial).
\item \textbf{Multi-pays} (\emph{réalisé}, \S\ref{sec:proto2}) : monde à $N$ pays, Armington à parts graduelles, change prix d'actif, réserves et règlement exact, cinq régimes de souveraineté, contrôle des capitaux, dollarisation, attaques de première et seconde génération, énergie échangeable ; treize tests passés.
\item \textbf{Finance avancée} : marché obligataire avec adjudications, prime de risque, défaut, boucle banques--souverain, contrôle des capitaux ; actifs (actions, immobilier), collatéral, ruées, QE ; reproduction des scénarios de la \S\ref{sec:crises}.
\item \textbf{Hétérogénéité} (\emph{réalisé}, \S\ref{sec:proto4}) : terre et climat, ressource épuisable, cohortes, migration rurale, archétype agraire ; règle budgétaire et machine à états ; dix tests passés.
\item \textbf{Systèmes économiques} : curseurs institutionnels, planification Leontief, prix administrés, marché noir, transition.
\item \textbf{Souveraineté monétaire et bilans (v0.6--0.8)} : régimes A--D, taux apparent et maturité consolidée, bilan de banque centrale, saut de prime, hystérèse de la dollarisation ; reproduction des sept configurations contemporaines de la \S\ref{sec:configs}.
\item \textbf{Jouabilité} : soutien politique, scores, événements, tutoriel, tableau de bord (indicateurs macro, bilans, graphiques).
\end{enumerate}

%=====================================================================
\section{Pistes écartées}
\label{sec:ecartees}
%=====================================================================
Ces éléments sont réels mais hors budget de complexité ; ils sont notés pour une version ultérieure.
\begin{itemize}
\item \textbf{Mécanique de guerre, mobilisation, sanctions.} La v0.6 avait introduit une décote pays sur le prix des ressources et une fraction utilisable des réserves pour un producteur sous sanctions ; ce cas n'est dans aucune des dix configurations et les deux paramètres ont été retirés en v0.8. Ils se réintroduisent en deux lignes si un joueur \og sanctionneur\fg{} est ajouté.
\item \textbf{Secteur promoteur immobilier avec passif de préventes.} La chaîne collatéral $\to$ crédit $\to$ prix couvre la dynamique agrégée ; le passif de préventes n'ajoute qu'un canal de contagion vers les ménages acheteurs, que l'on représente déjà par $NPL^Z$.
\item \textbf{Intermédiation non bancaire détaillée.} Une baisse temporaire de $\kappa^{CAR}$ représente le contournement réglementaire ; un secteur séparé doublerait le module bancaire.
\item \textbf{Chaînes de valeur multi-frontières.} L'Armington à une frontière suffit pour les mécanismes de change et de tarifs ; les chaînes longues ne changent que l'ampleur de la transmission, pas sa direction.
\item \textbf{Démographie par cohortes fines.} Trois cohortes (jeunes, actifs, âgés) donnent le dividende démographique et la charge des retraites ; le détail par âge n'ajoute que de la précision.
\item \textbf{Sous-secteurs des administrations publiques.} Les entités publiques locales endettées de la configuration~5 sont traitées comme une dette garantie implicite ajoutée à $B$ lors de la crise ; une comptabilité séparée serait plus fine sans changer le mécanisme.
\item \textbf{Divergence entre statistiques affichées et variables réelles.} Elle mériterait mieux qu'un renvoi : un joueur qui ment sur son déficit (Grèce 2009) ou sur sa croissance est un cas de jeu intéressant, et la mécanique serait légère (un bruit de mesure sur les indicateurs affichés aux autres joueurs, révélé avec délai, avec saut de $\rho^{pays}$ à la révélation). Nous ne l'implémentons pas ici parce qu'elle touche à l'interface multijoueur plus qu'au modèle.
\item \textbf{Compétitivité hors-prix endogène} (v0.7). Faire de $\varphi_{ik}$ une variable répondant à la R\&D et à la qualité serait juste sur le fond mais coûterait un sous-système ; l'Armington à préférences fixes reste le bon compromis, et \eqref{eq:lbd} capture déjà l'effet de la taille de l'industrie sur sa productivité.
\item \textbf{Épargne de précaution par incertitude politique} (v0.7). L'équation \eqref{eq:cB} lie la propension à consommer à la couverture sociale ; un pays à couverture maximale dont le taux d'épargne monte quand même relève d'un mécanisme mal identifié empiriquement (incertitude sur la fiscalité future, épargne de retraite privée, effet richesse négatif anticipé). Il est mentionné dans les limites de \eqref{eq:cB} plutôt que modélisé.
\item \textbf{Effet d'agglomération $\eta_{agg}$} (retiré en v1.5). Une phrase de la \S\ref{sec:hetero} faisait croître $A_j$ avec la part urbaine ; le terme n'avait ni équation dans la table de calibration ni dérivation dans le cours, et il comptait deux fois ce que le rattrapage $\chi$ de \eqref{eq:tfp} et l'épuisement du réservoir de Lewis produisent déjà. Krugman (1991) reste une piste si l'étape~5 fait apparaître une géographie des secteurs.
\item \textbf{Tiering, fonds pour risques généraux, dominance budgétaire, surprise inflationniste} (retirés en v0.8). Quatre raffinements de la v0.6 qui décrivaient correctement des débats réels de banque centrale mais ne changeaient la trajectoire d'aucune configuration ; le report à nouveau des pertes et la suspension du dividende suffisent comme amortisseur budgétaire.
\end{itemize}

%=====================================================================
\section{Simplifications}
\label{sec:simplif}
%=====================================================================
La v0.6 a ajouté ; la v0.7 retire ce qui ne servait pas. Trois retraits sont appliqués, deux sont proposés.
\begin{itemize}
\item \textbf{Archétypes et régimes fusionnés} (appliqué). La table des six archétypes du \S\ref{sec:hetero} et la table des régimes A--E se recouvraient (le petit pays ouvert \emph{est} le régime C, l'émetteur de réserve \emph{est} un cas de A). Elles deviennent une seule grille dotations $\times$ souveraineté (\S\ref{sec:archetypes}) dont les cases sont les configurations de la \S\ref{sec:configs}. Une table remplace deux, et la grille de validation devient l'index du document.
\item \textbf{Curseur $\varrho^w$ fusionné dans $\varrho^p$} (appliqué). Les salaires administrés n'apparaissaient qu'une fois hors table ; un seul curseur \og prix et salaires administrés\fg{} suffit, le salaire étant le prix du travail.
\item \textbf{Paramètres fixés en dur ou retirés} (appliqué, \S\ref{sec:calib}) : ceux qui n'influencent aucune configuration de la grille sortent de l'espace d'équilibrage (v0.7) ; onze ont été retirés du modèle en v0.8 (table en tête de document).
\item \textbf{Proposé : fusionner les tests 1929 et 6 (refuge)} de la \S\ref{sec:tests} en un seul test paramétré par le régime de change, les deux ne différant que par le signe de la contrainte de bilan (or : $H\le G/\bar g$ ; plafond : $E^{CB}/H\ge\underline\kappa^{CB}$).
\item \textbf{Proposé : retirer la strate moyenne $M$} et la remplacer par une pondération continue entre les comportements $B$ et $H$ ($\lambda_M$ devient une distribution). Cela supprime une strate et ses paramètres sans perdre les inégalités ; nous ne l'appliquons pas parce que la strate $M$ porte le soutien politique de la \S\ref{sec:politique}.
\end{itemize}
La v0.8 applique le principe jusqu'au bout : quatre équations allégées, onze paramètres retirés, document plus court que la v0.7.

%=====================================================================
\section{Arbitrages et questions ouvertes}
\label{sec:arbitrages}
%=====================================================================
\subsection*{État des chantiers ouverts (v1.5)}
\label{sec:chantiers}
Chaque chantier ouvert depuis la v0.9 a été vérifié sur le prototype. La v1.5 ajoute les points du relevé d'incohérences établi en préparant le cours (défauts du document plutôt que du modèle) et leur état. \textbf{Clos} : résolu et vérifié par un test ; \textbf{partiel} : amélioré, résidu documenté ; \textbf{ouvert} : non résolu, avec ce qui a été essayé.
\begin{center}\small
\begin{longtable}{p{0.9cm}p{4.6cm}p{1.3cm}p{7.6cm}}
\toprule
\textbf{Ouvert} & \textbf{Chantier} & \textbf{État} & \textbf{Vérification / résidu}\\
\midrule
\endfirsthead
\toprule
\textbf{Ouvert} & \textbf{Chantier} & \textbf{État} & \textbf{Vérification / résidu}\\
\midrule
\endhead
v0.9 & Amplitude de la transmission monétaire ($-6\,\%$ de PIB pour $+1$ pt un an) & clos & Avec actifs : $-1{,}3\,\%$ de PIB, $+1{,}8$ pt de chômage, immobilier $-12\,\%$ — dans la fourchette empirique. La richesse absorbe une part de l'ajustement ; la saturation autour de $\bar q$ (v1.0) fait le reste\\
v0.9 & Taux réel effectif sous le théorique & clos & 3,0\,\% contre 2,5\,\% théorique avec actifs (\S\ref{sec:proto3})\\
v0.9 & Investissement et $K/Y$ en bas de fourchette & ouvert & Avec actifs : $I/Y$ 11\,\% (entreprises) $+$ 3\,\% (logement), $K/Y$ 2,2 et en baisse lente (1,9 à quarante ans) ; part salariale 0,72 $>1-\alpha$, profits comprimés. Essayés sans effet : $\iota$ 0,3, $\rho_E$ 0,01, fenêtre d'utilisation, $\delta$ plus bas. Le problème est le partage de la valeur ajoutée avec indexation complète sur la productivité\\
v0.9 & Transitoire initial (trois à cinq ans) & clos & Chômage maximal 4,5\,\% sur les cinq premières années de l'économie à actifs\\
v0.9 & Distribution trop plate & partiel & Richesse par tête B:M:H $=1:3:14$ (réaliste) ; revenu par tête $1:1{,}1:2{,}0$, Gini 0,09 : la strate haute ne tire pas assez de revenu de son capital (dividendes 9\,\% du PIB, rendement 5\,\%)\\
v1.0 & Dotations hétérogènes par saut & clos & Règle adoptée et appliquée : les archétypes se calibrent à l'état initial (\S\ref{sec:proto4}, H5), les chocs de dotation sont graduels\\
v1.0 & Seconde génération : curseur & clos & Option \texttt{second\_gen} et seuils d'abandon exposés comme paramètres\\
v1.0 & Exportateur d'énergie en déficit & ouvert & Avec prix mondial : la balance de l'exportateur passe de $-1{,}7$ à $+0{,}6\,\%$ après un choc d'offre (la rente apparaît), mais le déficit structurel subsiste. Une élasticité d'Armington $<1$ (0,7) est \emph{instable} (les parts croissent avec le prix : effondrement). Un bien homogène exige un marché commun sans différenciation, hors du cadre Armington\\
v1.1 & Structure de détention de la dette (QE à réserves oisives) & clos & Ratio de liquidité $\phi_{liq}=0{,}3$ : la banque détient 10\,\% du PIB de titres, le QE lui achète, monnaie inchangée, souverain $-8$ pb, inflation $+0{,}1$ pt (test A3)\\
v1.1 & Forbearance et zombies & ouvert & Non implémenté (scénario japonais)\\
v1.1 & Amplitude des cycles d'actifs en configuration extrême & clos & Bornes de jeu : $\chi\le0{,}9$, LTV $\in[0{,}3;1]$, marge $\le1{,}5$, $\kappa^{CAR}\in[0{,}02;0{,}3]$, $\pi^\ast\le10\,\%$\\
v1.1 & Portage des actifs sur le monde à $N$ pays & partiel & Deux économies à actifs dans le monde : stable, identités exactes ; pas encore de détention d'actifs étrangers ni de contagion par les portefeuilles\\
v1.1 & Prix mondial généralisé aux ressources & ouvert & Voir « exportateur d'énergie » : la structure de marché commun manque\\
v1.2 & Règle budgétaire & clos & Test C1\\
v1.2 & État de mort économique & clos & Machine à états, test C2 (\S\ref{sec:collapse})\\
v1.2 & Exigence de capital et crises & clos & Test C3 : effet réel mais faible ; l'émission d'actions le précède\\
v1.2 & Horizon des sanctions & clos & Test C4 : sanction entre l'an 25 et l'an 30\\
v1.2 & Actions : part désirée fixe & clos & Rappel fondamental, plancher d'actualisation (C5)\\
v1.3 & Inflation chronique post-effondrement & clos & C'est un artefact d'économie \emph{fermée} : sans prix étranger, rien n'ancre les anticipations après la réforme. Dans le monde à deux pays, la même économie relevée converge à 2\,\% d'inflation et reconstruit sa crédibilité (0,88) en flottant comme en ancrage ; la taille de la conversion n'y change rien\\
v1.3 & Ressources et rente sur $N$ pays & ouvert & Voir « exportateur d'énergie »\\
v1.3 & Actions sous la valeur comptable & partiel & 0,77 avec un rappel fondamental renforcé ; un poids de 5 donne 0,85--0,98 mais double la volatilité du chômage\\
v1.3 & Cohortes et fécondité endogène & clos & Option \texttt{fert\_endo} : transition démographique reproduite\\
v1.3 & Étape 5 (systèmes économiques) & ouvert & Non commencée ; $q^{file}$, $\zeta_q$ et $G^{edu}$ y sont rattachés\\
v1.5 & Grille \S\ref{sec:configs} calculée avec l'ancienne forme de \eqref{eq:spread} & clos & Recalculée terme par terme ; cinq configurations changent (169, 76, 166, 145, 88 pb). L'écart « union vertueuse à 13 pb contre 60 observés » de la v0.8 est résolu par la réécriture (76 pb)\\
v1.5 & Plancher $\underline\varkappa$ inconditionnel & clos & Conditionné par $\chi^K\varrho^B$ ; répression $\varrho^B$ et $s_9$ ajoutés au prototype ; test H6\\
v1.5 & Deux états de référence & partiel & Rôles déclarés (\S\ref{sec:ssref}) : actifs $=$ jeu, fermée $=$ calibration ; l'écart de $K/Y$ reste le chantier du partage de la valeur ajoutée\\
v1.5 & Table \S\ref{sec:calib} : lignes périmées, $u^n$ ambigu, $\varpi$ à trois échelles & clos & Six lignes supprimées, liste des paramètres retirés complétée, $\underline u^n$ et $u^n_0$ séparés, $\varpi$ (prix, vitesse) et $\varpi_w$ (salaires, degré) distingués\\
v1.5 & (35) : $g$ ou $g_0$ dans $\hat r^\ast_0$ & clos & $g_0$, comme le prototype ; 2,5\,\%\\
v1.5 & $q=1$ présenté comme seuil de rentabilité & clos & $q^\ast=1+g/(\iota\Gamma\Lambda\Lambda^{Bk})=1{,}11$ ; indicateur joueur $q/q^\ast$\\
v1.5 & Effet du QE ($-8$ pb pour 20\,\% du PIB) & clos (écart assumé) & Conforme aux équations ($\psi_{QE}$ sur la seule prime de terme) ; la littérature (50--100 pb pour 10\,\%) inclut un effet d'anticipation de trajectoire que \eqref{eq:spread}, construite sur le taux courant, n'a pas. Écrit dans les limites de \S\ref{sec:qe}\\
v1.5 & Accélérateur financier inactif ($\ell<\bar\ell$) & ouvert & Levier 0,32 en économie fermée, \textbf{0,08} avec actifs : aucun seuil ne l'active (balayage de $\bar\ell$ 0,6 $\to$ 0,35 sans effet). Forme continue à concevoir avec le chantier du partage de la valeur ajoutée, dont la cause est la même\\
v1.5 & $\zeta_q$, $q^{file}$, $\eta_{agg}$, $\eta_H$ sans spécification & clos & $\eta_{agg}$ retiré ; $\eta_H$ et $q^{file}$ spécifiés avec valeur, non simulés (étape~5)\\
v1.5 & (41), (43) sans référence & clos & Encadrés « choix de conception, pas résultat théorique »\\
v1.5 & Colonne « Prix » de \S\ref{sec:systemes} & clos & Dédoublée selon $\varrho^p$ : pénurie à prix administrés, déflation à prix libres\\
\bottomrule
\end{longtable}
\end{center}
\paragraph{Ce qui reste à faire, par ordre d'importance.} (1) Le partage de la valeur ajoutée dans l'économie à actifs (part salariale 0,72, capital en baisse lente, levier des entreprises 0,08) : c'est le seul résidu qui touche l'état de référence, et il porte désormais aussi l'accélérateur financier et le choix entre les deux états de référence. Piste : l'indexation complète des salaires sur la productivité ($\gamma_A=1$) est peut-être trop forte quand une part du revenu passe par la richesse ; et le financement par actions évince le crédit bancaire plus qu'il ne le complète. (2) Un marché mondial de biens homogènes hors Armington, pour les ressources et leur rente. (3) La forbearance (scénario japonais). (4) La détention d'actifs étrangers (contagion). (5) L'étape~5.

\paragraph{Écarts acceptés (v0.8, révisés en v1.5).} Le recentrage laissait des écarts connus entre primes calculées et observées. Deux sont résolus par les réécritures de \eqref{eq:spread} : l'union vertueuse passe de 13 à 76 pb ($\approx60$ observés) et le pays planifié de 79 à 88 pb, après être monté à 218 sous le plancher inconditionnel de la v1.1 ($\approx30$ observés au-dessus du taux directeur, 2026). Restent le refuge (158 contre $\approx50$) et l'émetteur de réserve (166 contre 50--100) ; dans les deux cas l'excès est la prime de terme de 100 pb à sept ans, calibrée sur l'écart historique moyen dix ans/taux directeur, que des courbes aplaties par les anticipations de baisse ne reproduisent pas. Aucun ne change la trajectoire : dans les quatre cas la boule de neige reste négative, aucun saut ne se déclenche, et le classement des dix configurations par prime est le bon. Un écart serait inacceptable s'il faisait sauter la configuration 2, exploser la dette de la 4 ou disparaître l'inflation de la 10 ; ce n'est le cas d'aucun. La correction tentante de la v0.8 (un rabais de QE $-s_7\omega$ décroissant) est devenue sans objet : $s_7$ n'existe plus depuis la v1.1.

\paragraph{Choix faits en v0.7.}
\begin{itemize}
\item \textbf{Convexité en excès relatif plutôt qu'absolu.} La correction proposée (terme additif $-s_{10}\varkappa^{dom}$) produisait des primes négatives pour tout pays à dette moyenne et détention domestique élevée ; nous avons préféré pondérer la convexité par $(1-\varkappa^{dom})$, ce qui ne peut pas devenir négatif. Et la convexité mesurée en $(b-\bar b)^2$ ne pouvait pas séparer de 300 pb deux pays à $b=0{,}35$ : à ce niveau de dette, $s_1(0{,}18)^2$ vaut 16 pb quelle que soit la capacité fiscale. La forme relative $(b/\bar b_i-1)^2$ le fait (écart de 326 pb) et donne le bon ordre de grandeur pour la configuration 4 (105 pb au lieu de 1\,125). C'est la seule modification de forme que nous ayons imposée par rapport au cahier des charges.
\item \textbf{$s_0$ à 100 pb et non 50.} Avec $s_2$ à 30 pb par point et un backstop à $-50$ pb, une union endettée à $pb=-3\,\%$ ne pouvait atteindre les 100 à 130 pb d'écart observés au-dessus du taux de l'union qu'avec une prime de terme de 100 pb ; c'est aussi l'ordre de grandeur historique de l'écart dix ans/taux directeur en régime normal.
\item \textbf{Prime d'illiquidité supprimée} (v0.8). Elle décrivait la même chose que la répression financière (des épargnants captifs) ; un seul terme suffit.
\item \textbf{Remises sans équation.} Elles sont un transfert reçu par les ménages et un revenu du compte courant ; les deux variables existent ($Tr_h$, revenus nets de $CA$). Les redéfinir coûte deux phrases et zéro équation, ce qui a permis de retenir les trois ajouts de priorité 3 dans le budget de trois équations.
\item \textbf{Régime E sans équation}, par réutilisation de $\pi^{dev}$ avec un seuil $\varpi_E M$.
\item \textbf{Fusion des archétypes et des régimes} plutôt que leur maintien côte à côte : c'est le retrait le plus visible et il rend la grille de validation lisible comme un index.
\end{itemize}
\paragraph{Calibrations contestées.}
\begin{itemize}
\item $s_6=40$ pb \og décroissant\fg{} : la décroissance devrait être explicite (par exemple $s_6\cdot c\cdot\ind_{b<1{,}5}$) ; nous la laissons en prose faute d'un mécanisme documenté de perte du statut de réserve, qui est un événement rare et politique.
\item $\varsigma_8=1{,}05$ retrouve les deux cibles, mais l'hypothèse que le backstop réduit la probabilité d'un facteur 20 et non 200 est une opinion : les épisodes observés (2012) sont compatibles avec les deux.
\item \textbf{Dérive projetée plutôt que déficit primaire} (v0.7.1). Le terme $s_2$ appliqué au seul déficit primaire faisait payer 214 pb à l'émetteur de réserve alors que sa croissance nominale (4,7\,\%) excède son taux apparent (3,4\,\%) et annule sa boule de neige. Faire porter $s_2$ sur $\Delta b^{proj}$, l'équation \eqref{eq:debtdyn} évaluée aux valeurs courantes, prend la croissance en compte sans variable nouvelle : la configuration 3 tombe à 177 pb, la 4 à 99, la 5 à 79. Le revers est la configuration 2 (union vertueuse), qui tombe à 13 pb au-dessus du taux de l'union contre 60 observés : le rabais de QE $-s_7\omega$ ($-75$ pb pour $\omega=0{,}25$) y est trop généreux, et devrait probablement décroître avec le resserrement quantitatif. Nous laissons ce point ouvert.
\end{itemize}
\paragraph{Questions laissées au commanditaire.}
\begin{itemize}
\item La part domestique $\varkappa^{dom}$ doit-elle rester exogène ou réagir à la prime (les résidents rachètent quand les étrangers vendent) ? C'est un stabilisateur réel et un terme de plus.
\item Le seuil $(T/Y)^{ref}=0{,}35$ est-il universel ou doit-il dépendre du niveau de développement (un pays riche à faible État peut emprunter davantage qu'un pays pauvre à même $T/Y$) ?
\item Faut-il appliquer la fusion proposée de la strate $M$ dans la v0.8, au prix d'une refonte du soutien politique ?
\item Les questions de la v0.6 restantes : backstop binaire ou gradué ; $\underline\kappa^{CB}$ fixe ou fonction du soutien politique ; incrémentation de $M^{hist}$ lors d'une hyperinflation sans changement de monnaie. La contagion, elle, est tranchée (activée).
\end{itemize}

\newpage
\appendix
\section{Glossaire de la notation}
\label{sec:glossaire}
\small
\begin{longtable}{p{4.3cm}p{7.8cm}p{3.3cm}}
\toprule
\textbf{Symbole} & \textbf{Signification} & \textbf{Où}\\
\midrule
\endfirsthead
\toprule
\textbf{Symbole} & \textbf{Signification} & \textbf{Où}\\
\midrule
\endhead
\multicolumn{3}{l}{\emph{Indices et opérateurs}}\\
$t$ ; $j,k$ ; $h$ ; $i$ & temps ; secteurs/biens ; strates ($B$ basse, $M$ moyenne, $H$ haute) ; pays & partout\\
$\bar x$, $\hat x$, $x^\ast$, $\Delta x$, $\E_t$ & référence/cible ; écart ou visé ; étranger/équilibre ; variation ; anticipation & \S\ref{sec:rappels}\\
$\pos{x}$, $\ind_{\{\cdot\}}$, clip & partie positive ; indicatrice ; bornage & \S\ref{sec:rappels}\\
\midrule
\multicolumn{3}{l}{\emph{Production et entreprises}}\\
$Y_j$, $Y^{cap}_j$, $\hat Y_j$, $Q_j$ & production, capacité, production visée, ventes du secteur $j$ & \eqref{eq:prod}\\
$A_j$, $A^{front}$ & productivité globale des facteurs ; frontière mondiale & \eqref{eq:tfp}\\
$K_j$, $I_j$, $\delta_j$, $p_K$ & capital, investissement, taux de dépréciation, prix du capital & \eqref{eq:invest}\\
$L_j$, $L^\ast_j$, $W_j$, $\tau_S$ & emploi, emploi souhaité, salaire, cotisations employeur & \eqref{eq:labour}\\
$H$ & capital humain moyen & \eqref{eq:prod}\\
$X_{kj}$, $a_{kj}$, $\mathbf{A}$ & intrants, coefficients techniques, matrice de Leontief & \eqref{eq:prod}, \eqref{eq:leontief}\\
$S^{inv}_j$, $s^\ast_j$ & stocks, stock cible & \S\ref{sec:stocks}\\
$\Pi_j$, $M_j$, $Div_j$, $\psi_j$, $\bar m$ & profit, trésorerie, dividendes, taux de distribution, trésorerie cible & \eqref{eq:profit}--\eqref{eq:div}\\
$q_j$, $\iota$, $\Lambda_j$, $\ell_j$, $\bar\ell$ & $q$ de Tobin, sensibilité, rationnement, levier, levier de référence & \eqref{eq:invest}\\
$\theta_j$ & part du capital du secteur détenue par l'État & \eqref{eq:div}, \S\ref{sec:systemes}\\
$c^u_j$ & coût unitaire de production & \eqref{eq:price}\\
$T$, $\alpha_T$ & terre agricole, part de la terre & \eqref{eq:agr}\\
$\mathcal{R}$, $Y_{res}$, $p^w_{res}$, $\mathcal{S}_{res}$, $\tau^{roy}$, $F$ & réserves, extraction, prix mondial, rente, redevance, fonds souverain & \S\ref{sec:hetero}\\
\midrule
\multicolumn{3}{l}{\emph{Ménages}}\\
$N_h$, $N$, $L^S$, $\varsigma$, $u$, $u^n_t$ & population de la strate, totale, active, participation, chômage, chômage naturel courant (hystérèse) & \S\ref{sec:demog}\\
$Y^d_h$, $E_h$, $S_h$, $V_h$ & revenu disponible, dépense, épargne, richesse nette & \eqref{eq:yd}\\
$c_h$, $c^V_h$, $\beta$, $\sigma$, $\rho$ & propensions à consommer, escompte, aversion intertemporelle, préférence pour le présent & \eqref{eq:cB}--\eqref{eq:cH}\\
$x_{hj}$, $\gamma_{hj}$, $\beta_{hj}$ & consommation, subsistance, parts de dépense & \eqref{eq:les}\\
$SoL_h$, $P$, $\pi$, $w_j$, Gini & niveau de vie, indice des prix, inflation, poids du panier, inégalités & \eqref{eq:sol}--\eqref{eq:cpi}\\
$\omega_h$ & part du capital privé détenue par la strate & \eqref{eq:yd}\\
$s^D_h$, $s^B_h$, $s^E_h$, $s^{FX}_h$ & parts de portefeuille (dépôts, obligations, actions, devises) & \eqref{eq:portfolio}\\
$L_{rur}$, $L_{urb}$, $W_{rur}$, $\lambda_{mig}$ & main-d'œuvre rurale/urbaine, revenu rural, vitesse de migration & \eqref{eq:migr}\\
\midrule
\multicolumn{3}{l}{\emph{Marchés et commerce}}\\
$D_j$, $S_j$, $z_j$, $\kappa_j$, $\bar\kappa$, $\varpi$ & demande, offre, déséquilibre, vitesses d'ajustement, coefficient de vitesse d'indexation des prix (par tick, sur $\pi^e/52$) & \eqref{eq:price}\\
$\varpi_w$, $\omega_u$ & degré d'indexation des salaires (borné par 1 ; $\to1$ en hyperinflation), affaiblissement par le chômage & \eqref{eq:wage}\\
$p_j$, $\bar p_j$, $p^{bm}_j$ & prix de marché, prix administré, prix de marché noir & \eqref{eq:price}, \S\ref{sec:plan}\\
$e$, $e^{PPA}$, $\bar e$ & taux de change, parité de pouvoir d'achat, parité fixe & \eqref{eq:fx}\\
$\tau^M$, $\tau^{tr}$, $\varphi_{ik}$, $\epsilon_j$ & tarif, coût de transport, préférence, élasticité de substitution & \eqref{eq:armington}\\
$EX$, $IM$, $TB$, $CA$, $KA$, $R^{FX}$ & exportations, importations, balance commerciale, compte courant, compte de capital, réserves & \S\ref{sec:bop}\\
$\chi^K$ & contrôle des capitaux (0 libre, 1 fermé) & \eqref{eq:fx}\\
\midrule
\multicolumn{3}{l}{\emph{Monnaie et finance}}\\
$i^{CB}$, $i^L$, $i^D$, $i^B$, $i^Z$, $i^{long}$ & taux directeur, du crédit, des dépôts, des obligations, hypothécaire, long & \S\ref{sec:monnaie}\\
$r$, $r^\ast$, $r^L$ & taux réel, taux naturel, taux réel du crédit & \eqref{eq:taylor}, \eqref{eq:invest}\\
$\pi^e$, $\pi^\ast$, $c$, $\lambda_\pi$, $\upsilon$ & inflation anticipée, cible, crédibilité, vitesse de révision, ancre monétaire & \eqref{eq:expect}--\eqref{eq:cred}\\
$H$, $M$, $C$, $D$, $Res$, $Res^{ex}$ & base monétaire, masse, billets, dépôts, réserves bancaires, réserves excédentaires & \S\ref{sec:qe}\\
$V(\pi^e)$, $V_0$, $a$, $\xi_M$ & vitesse de circulation et paramètres, poids du terme d'encaisses & \eqref{eq:monterm}\\
$B$, $B^{Bk}$, $B^{CB}$, $B_h$, $B^\ast$, $b$ & dette publique et ses détenteurs, dette externe, ratio dette/PIB & \S\ref{sec:etat}\\
$\Delta B^{CB}$, $\Delta A^G$, $\Sigma$ & QE, avances au Trésor, seigneuriage & \S\ref{sec:qe}\\
$E^{Bk}$, $\kappa^{CAR}$, $L^{max}$, $\Lambda^{Bk}$, $NPL$ & fonds propres bancaires, ratio réglementaire, plafond de crédit, rationnement, créances douteuses & \S\ref{sec:banques}\\
$\omega^{run}$ & fraction des dépôts retirée (ruée) & \eqref{eq:run}\\
$p^E$, $p^{E,f}$, $r^E$, $\chi_E$, $m$ & cours des actions, valeur fondamentale, rendement attendu, extrapolation, marge & \eqref{eq:pe}\\
$p^Z$, $Z$, $uc$, $\mathrm{LTV}$, $\chi_Z$, $T_Z$ & prix immobilier, stock, coût d'usage, quotité, extrapolation, délai & \eqref{eq:pz}\\
$\rho^{pays}$, $\rho^{def}$, $\pi^{dev}$, $D^{hist}$, $\hbar$ & prime pays, défaut perçu, probabilité de dévaluation, historique de défaut, décote & \S\ref{sec:change}, \S\ref{sec:etat}\\
\midrule
\multicolumn{3}{l}{\emph{État et institutions}}\\
$\tau_{W,h}$, $\tau_K$, $\tau_\Pi$, $\tau^C_j$, $T^{poll}$ & impôts sur le travail, le capital, les sociétés, la consommation, capitation & \S\ref{sec:etat}\\
$G$, $G^{adm}$, $G^{edu}$, $G^{san}$, $G^{inv}$, $Tr$, $Sub$ & dépenses publiques par nature, transferts, subventions & \S\ref{sec:etat}\\
$K^G$, $\varepsilon^{adm}$, $\varepsilon^{ev}$ & capital public, efficacité de collecte, évasion & \S\ref{sec:etat}\\
$pb$, Déficit & solde primaire, déficit total & \eqref{eq:gbc}\\
$\varrho^p_j$, $\varrho^{plan}$, $\varepsilon^{plan}$ & prix et salaires administrés, degré de planification, pénalité & \S\ref{sec:systemes}\\
$\Lambda^{pol}_h$ & soutien politique de la strate & \S\ref{sec:politique}\\
\midrule
\multicolumn{3}{l}{\emph{Ajouts v0.6}}\\
$i^{app}$, $\theta$, $\theta^{cons}$, $\bar T$, $\omega$ & taux apparent, fraction refinancée, fraction consolidée, maturité d'émission, part de la dette à la banque centrale & \eqref{eq:iapp}\\
$E^{CB}$, $\bar c^{CB}$, $i^{res}$, $\Delta^{res}$, $\bar E^{CB}$, $\underline\kappa^{CB}$ & fonds propres, coupon figé, taux des réserves et écart au taux directeur, fonds propres statutaires, seuil de bilan & \eqref{eq:ecb}\\
$\theta_6$ & effet de stock des réserves sur le change & \eqref{eq:fx}\\
$s_8$, $s_9$, $\bar b^{A..E}$ & backstop, répression, seuil de dette par régime & \eqref{eq:spread}\\
$\Pr(\text{saut})$, $\Delta\rho^{saut}$, $\lambda_\rho$, $\varsigma_6..\varsigma_9$ & saut de prime et paramètres & \eqref{eq:jump}\\
$M^{hist}$, $\nu_9$, $\delta_M$ & mémoire monétaire et paramètres de crédibilité & \eqref{eq:cred}\\
$s^{FX}$, $\phi_0$, $\varrho^{dedol}$, $\kappa_d$, $r^{eff}$, $M^{nat}$ & dollarisation, persistance, contrainte légale, taux effectif, monnaie nationale & \eqref{eq:sfx}, \eqref{eq:reff}\\
$r^F$, $\bar r^F$, $\beta_F$, $p^{E,w}$ & rendement du fonds, marché mondial des actions & \eqref{eq:fonds}\\
$cov$, $\psi_h$, $\overline{cov}$ & couverture sociale et épargne de précaution & \eqref{eq:cB}\\
$u^n_0$, $\underline u^n$, $\eta_u$, $\eta'_u$ & chômage naturel structurel initial (5\,\%), plancher frictionnel (3\,\%), hystérèse & \eqref{eq:un}\\
$\varrho^{fisc}$, $\lambda^{fisc}$, $\varrho^B$, $\mathcal{T}^{rep}$ & règle budgétaire, délai de sortie, détention obligatoire, transfert de répression & \S\ref{sec:systemes}\\
\midrule
\multicolumn{3}{l}{\emph{Ajouts de la v0.7}}\\
$\varkappa^{dom}$, $(T/Y)^{ref}$, $\bar b_i$, $\Delta b^{proj}$ & part domestique de la dette, capacité fiscale de référence, seuil effectif du pays, dérive projetée du ratio de dette & \eqref{eq:spread}\\
$\varsigma_0$, $\varsigma_{11}$ & constante et contagion du saut de prime & \eqref{eq:jump}\\
$\varsigma^{mut}$ & part mutualisée des pertes dans une union & \S\ref{sec:souv}\\
$\beta^{dom}$ & part du fonds souverain placée dans le pays & \eqref{eq:fonds}\\
$\varpi_E$ & seuil de réserves du régime E, en proportion de $M$ & \S\ref{sec:souv}\\
$Rem$ & remises des migrants (dans $Tr_h$ et $CA$) & \S\ref{sec:bop}\\
\midrule
\multicolumn{3}{l}{\emph{Ajouts de la v0.9}}\\
$\rho_E$, $\Gamma_j$, $\underline u$, $\bar u$ & prime de risque sur le capital ; taux d'utilisation et sa fenêtre & \eqref{eq:invest}\\
$m_j$, $f_t$ & marge normale ; facteur de flexibilité des prix & \eqref{eq:price}\\
$\pi^{ref}$, $\omega_u$, $x^L$, $\phi_x$, $\hat g^{prod}$ & inflation de référence de l'indexation ; poids d'indexation ; excès de demande de travail et surenchère ; productivité apparente lissée & \eqref{eq:wage}\\
$\hat r^\ast$, $\lambda_r$, $\pi^s$, $\lambda_s$, $\rho_i$, $\bar\imath$ & taux naturel estimé et son apprentissage ; inflation lissée de la règle ; lissage et plafond du taux & \eqref{eq:taylor}\\
$\varrho^G$ & taux d'exécution des dépenses publiques (contrainte de trésorerie) & \S\ref{sec:etat}\\
$v(\pi^e)$ & facteur de vitesse relâchant les contraintes de caisse & \S\ref{sec:sfc}\\
$R_j$ & remboursement de dette sur trésorerie excédentaire & \eqref{eq:div}\\
\midrule
\multicolumn{3}{l}{\emph{Ajouts de la v1.0}}\\
$\bar q_j$, $s_q$, $\lambda_q$ & $q$ de long terme, largeur et vitesse de la réponse saturante & \eqref{eq:invest}\\
$e^{cible}$, $\theta_L$, $\lambda_e$, $NF$, $\kappa_e$, $\kappa_K$, $\kappa_{spec}$ & cible de change, sensibilité au différentiel de taux, rappel, flux nets de devises, pressions & \eqref{eq:fx}\\
$\pi^{dev}$, $R^{FX}_0$ & probabilité de dévaluation, réserves initiales & \eqref{eq:pdev}\\
$\varsigma^{eff}$, $\lambda_\varsigma$ & parts Armington effectives et leur vitesse d'ajustement & \eqref{eq:armington}\\
\midrule
\multicolumn{3}{l}{\emph{Ajouts de la v1.1}}\\
$\Delta E^{\text{ém}}$, $\varsigma^{\text{ém}}$, $CAR$ & émission d'actions bancaires, sa borne, ratio de couverture & \eqref{eq:bankeq}\\
$b^{mkt}$, $\varkappa^{for}$, $\psi_{QE}$ & dette de marché (hors banque centrale), part non-résidente de cette dette, canal de portefeuille du QE & \eqref{eq:spread}\\
$p^w_j$, $\lambda_{lop}$ & prix mondial d'un bien homogène et vitesse de rappel & \S\ref{sec:marches}\\
$\psi_{fund}$, $\nu_{build}$, $L^E$ & rappel vers la valeur fondamentale, réponse des promoteurs, prêts sur marge & \S\ref{sec:proto3}\\
\midrule
\multicolumn{3}{l}{\emph{Ajouts de la v1.2}}\\
$\underline u^n$, $g^{obt}$, $\omega_G$, $\kappa^{hB}$ & plancher du chômage naturel, consommation publique effective par tête, son poids, élasticité de la demande d'obligations & \eqref{eq:un}, \eqref{eq:sol}\\
\midrule
\multicolumn{3}{l}{\emph{Ajouts de la v1.3}}\\
$\varphi^{fisc}$, $\phi_b$, $b^\ast$ & facteur de la règle budgétaire, sa pente, cible de dette & \eqref{eq:fiscrule}\\
$\upsilon^{surv}$, $\pi^{surv}$, $\varsigma^{conv}$ & production, inflation résiduelle et conversion en régime de survie & \eqref{eq:collapse}\\
$\varrho^{plat}$, $\underline a$, $\omega_t$, $L^{rur}$, $w^{rur}$, $\mu$, $c^{mig}$ & plateau et plancher de la ressource ; aléa climatique ; réservoir rural, revenu rural, vitesse et coût de migration & \S\ref{sec:hetero}\\
\bottomrule
\end{longtable}
\normalsize
\newpage
\section*{Références}
\begin{itemize}\small
\item Armington, P. (1969), \emph{A Theory of Demand for Products Distinguished by Place of Production}, IMF Staff Papers.
\item Bernanke, B., Gertler, M., Gilchrist, S. (1999), \emph{The Financial Accelerator in a Quantitative Business Cycle Framework}, Handbook of Macroeconomics.
\item Blanchard, O. (1990), \emph{Suggestions for a New Set of Fiscal Indicators}, OECD.
\item Cagan, P. (1956), \emph{The Monetary Dynamics of Hyperinflation}, in Friedman (ed.), Studies in the Quantity Theory of Money.
\item Cass, D. (1965) ; Koopmans, T. (1965) ; Ramsey, F. (1928), \emph{A Mathematical Theory of Saving}, Economic Journal.
\item Dornbusch, R. (1976), \emph{Expectations and Exchange Rate Dynamics}, JPE.
\item Eaton, J., Gersovitz, M. (1981), \emph{Debt with Potential Repudiation}, REStud.
\item Galí, J. (2008), \emph{Monetary Policy, Inflation, and the Business Cycle}, Princeton UP ; Galí, Gertler (1999), \emph{Inflation Dynamics: A Structural Econometric Analysis}, JME ; Galí, López-Salido, Vallés (2007), JEEA.
\item Godley, W., Lavoie, M. (2007), \emph{Monetary Economics: An Integrated Approach to Credit, Money, Income, Production and Wealth}, Palgrave.
\item Hayek, F. (1945), \emph{The Use of Knowledge in Society}, AER.
\item Hicks, J. (1937), \emph{Mr. Keynes and the Classics}, Econometrica.
\item Kornai, J. (1980), \emph{Economics of Shortage} ; (1992), \emph{The Socialist System}.
\item Krugman, P. (1979), \emph{A Model of Balance-of-Payments Crises}, JMCB.
\item Lange, O. (1936), \emph{On the Economic Theory of Socialism}, REStud.
\item Leontief, W. (1941), \emph{The Structure of American Economy}.
\item Minsky, H. (1986), \emph{Stabilizing an Unstable Economy}.
\item Mundell, R. (1963) ; Fleming, J. M. (1962).
\item Phillips, A. W. (1958) ; Friedman, M. (1968), \emph{The Role of Monetary Policy}, AER.
\item Reinhart, C., Rogoff, K. (2009), \emph{This Time Is Different}.
\item Solow, R. (1956), \emph{A Contribution to the Theory of Economic Growth}, QJE.
\item Stone, R. (1954), \emph{Linear Expenditure Systems and Demand Analysis}, Economic Journal ; Geary, R. (1950).
\item Taylor, J. (1993), \emph{Discretion versus Policy Rules in Practice}, Carnegie-Rochester.
\item Tobin, J. (1969), \emph{A General Equilibrium Approach to Monetary Theory}, JMCB ; Hayashi, F. (1982), Econometrica.
\item Woodford, M. (2003), \emph{Interest and Prices}.
\item Caballero, R., Hoshi, T., Kashyap, A. (2008), \emph{Zombie Lending and Depressed Restructuring in Japan}, AER.
\item Fisher, I. (1933), \emph{The Debt-Deflation Theory of Great Depressions}, Econometrica.
\item Harris, J., Todaro, M. (1970), \emph{Migration, Unemployment and Development}, AER ; Lewis, W. A. (1954), \emph{Economic Development with Unlimited Supplies of Labour}.
\item Kindleberger, C. (1978), \emph{Manias, Panics, and Crashes} ; Shiller, R. (2000), \emph{Irrational Exuberance}.
\item Krugman, P. (1987), \emph{The Narrow Moving Band, the Dutch Disease\ldots}, JDE ; Matsuyama, K. (1992), JET ; Corden, Neary (1982).
\item Sargent, T., Wallace, N. (1981), \emph{Some Unpleasant Monetarist Arithmetic} ; Olivera (1967), Tanzi (1977).
\item Friedman, M., Schwartz, A. (1963), \emph{A Monetary History of the United States} ; Eichengreen, B. (1992), \emph{Golden Fetters}.
\item Koo, R. (2008), \emph{The Holy Grail of Macroeconomics: Lessons from Japan's Great Recession}.
\item Paradox Interactive, \emph{Victoria 3 Wiki} : pages Market, Goods, Building, Treasury, Infrastructure (vic3.paradoxwikis.com) ; Dev Diary \#12 \emph{Treasury}.
\item EBTX\_, \emph{Economic and Financial Mod (E\&F) --- V4}, Steam Workshop (descriptions des versions V2 à V4, notes de mise à jour).
\end{itemize}

\end{document}
